\documentclass[11pt,twocolumn,twoside]{book}

\usepackage[margin=0.75in,top=1in,bottom=1in,headheight=14pt]{geometry}
\usepackage{titlesec}
\usepackage{amsmath,amssymb,amsfonts,amsthm}
\usepackage{multicol}
\usepackage{enumitem}
\usepackage{hyperref}
\usepackage{fancyhdr}
\usepackage{booktabs}
\usepackage{microtype}
\usepackage[utf8]{inputenc}
\usepackage[T1]{fontenc}
\usepackage{lmodern}
\usepackage{xcolor}
\usepackage{textcomp}
\usepackage{sectsty}
\usepackage{parskip}
\usepackage{makeidx}
\usepackage{graphicx}
\makeindex

\pagestyle{fancy}
\fancyhf{}
\fancyhead[LE]{\textit{\leftmark}}
\fancyhead[RO]{\textit{\leftmark}}
\fancyfoot[LE,RO]{\thepage}
\renewcommand{\headrulewidth}{0.4pt}
\renewcommand{\footrulewidth}{0.4pt}

\hypersetup{
    colorlinks=true,
    linkcolor=blue!60!black,
    urlcolor=blue!60!black,
    citecolor=blue!60!black,
}

\chapterfont{\LARGE\bfseries\color{blue!70!black}}
\sectionfont{\Large\bfseries\color{blue!70!black}}
\subsectionfont{\normalsize\bfseries\color{blue!50!black}}

\newcommand{\term}[1]{\hypertarget{#1}{}{\raggedright\noindent\textbf{\large\textcolor{MidnightBlue}{#1}}\index{#1}\par}\nobreak\vspace{1pt}\noindent\ignorespaces}
\newcommand{\pronun}[1]{}
\newcommand{\abbrev}[1]{\textsc{#1}}

\definecolor{MidnightBlue}{RGB}{25,25,112}

\raggedbottom
\sloppy
\setlength{\emergencystretch}{3em}

\setlength{\columnsep}{0.3in}
\setlength{\parindent}{0pt}
\setlength{\parskip}{2pt}

\title{\textbf{\Huge Mathematical Dictionary}\\[0.5em]
\Large Terms, Theorems, and Concepts --- A--Z Entries}
\author{Chaman Singh Verma}
\date{\today}

\begin{document}

\onecolumn
\frontmatter
\maketitle
\newpage
\thispagestyle{empty}
\vspace*{\fill}

%\begin{center}
%\includegraphics[width=0.6\textwidth]{frontpage.png}
%\end{center}

\vspace*{\fill}
\newpage

\thispagestyle{fancy}

\chapter*{Preface}

\noindent Mathematical dictionaries have traditionally been arranged alphabetically, an arrangement well suited for looking up an unfamiliar term when one already knows its name. But mathematics is not learned alphabetically, nor is it practised that way. A topologist thinking through a homotopy argument reaches for topological terms, not a mixture of topology wedged between tensor analysis and Thom's first isotopy lemma.

\medskip

\noindent This dictionary is organised alphabetically within chapters covering the broad areas of mathematics: Algebra, Analysis, Geometry, Topology, Number Theory, and Applied Mathematics. Each chapter contains terms listed alphabetically. The result is a reference that mirrors the structure of mathematical knowledge itself, with cross-references redirecting the reader from variant names and alternative spellings to the main entry.

\medskip

\noindent The entries are drawn from classical and modern mathematical literature. Each entry is written in a concise single-line format, with the term displayed in bold blue and automatically indexed. The dictionary runs to over 1600 entries across A--Z chapters, covering foundational concepts through advanced theory.

\medskip

\noindent This work is a living document. The LaTeX source is designed to be easily extended, corrected, and reorganised. Contributions, corrections, and suggestions are welcome.

\medskip

\noindent\textit{Chaman Singh Verma}


\chapter*{Mathematical Disclaimer}

\noindent \textbf{Educational and Reference Purpose Only}: This publication, \textit{Mathematical Dictionary}, is prepared and compiled strictly for educational, academic, historical, and reference purposes. It is intended for mathematics students, educators, and researchers.

\medskip

\noindent \textbf{No Pedagogical Claims}: The definitions and theorems presented herein are standard mathematical facts drawn from established literature. They do not constitute pedagogical recommendations or curriculum guidance.

\medskip

\noindent \textbf{Limitation of Liability}: Neither the author nor the publishers accept any legal responsibility or liability for any errors, omissions, or interpretive outcomes resulting from the use or application of any information presented in this work.

\newpage
\twocolumn
\mainmatter

\chapter*{A}
\label{chap:a}

\term{Abelian Category}
A category in which every hom-set is an abelian group and composition is bilinear, and in which every morphism has a kernel and a cokernel. In such categories the isomorphism theorems hold and every morphism factors through its image. The categories of abelian groups, of modules over a ring, and of sheaves of abelian groups on a space are all abelian, while general preadditive categories need not be. Abelian categories are the natural setting for homological algebra, where exact sequences, kernels, and cokernels behave as in the module case.

\term{Abelian Group}
A group in which the group operation is commutative, so $ab = ba$ for all elements $a, b$. Abelian groups are the additive groups of rings, modules, and vector spaces.

\term{Abel's Theorem}
A theorem asserting that if a power series converges at a point on its circle of convergence, then its sum has a limit at that point equal to the value of the series there. It also refers to Abel's result that the general quintic equation cannot be solved by radicals.

\term{Abilene Paradox}
A social phenomenon in which a group makes a decision contrary to the private preferences of every member, because each person assumes the others want it. Not a mathematical paradox, but a well-known example of pluralistic ignorance in group decision making.

\term{Absolute Convergence}
A property of a series $\sum a_n$ for which $\sum |a_n|$ converges. An absolutely convergent series is unconditionally convergent, meaning its sum is independent of the order of its terms.

\term{Absolutely Continuous}
A function $f$ on $[a,b]$ is absolutely continuous if for every $\varepsilon > 0$ there is $\delta > 0$ such that the total variation of $f$ over finitely many disjoint subintervals of total length less than $\delta$ is less than $\varepsilon$. By the fundamental theorem of calculus, such functions are exactly those with $f(x) = f(a) + \int_a^x g(t)\,dt$ for a Lebesgue integrable $g$. A measure $\mu$ is absolutely continuous with respect to $\nu$, written $\mu \ll \nu$, if every $\nu$-null set is $\mu$-null; then the Radon--Nikodym theorem yields a density $g$ with $\mu = g\nu$.

\term{Absorbing Element}
An element $z$ of a set with a binary operation such that $z \ast x = x \ast z = z$ for every element $x$. Zero is the absorbing element for multiplication.

\term{Abstract Algebra}
The branch of mathematics studying algebraic structures such as groups, rings, fields, modules, and their homomorphisms, abstracting away from the concrete nature of the elements.

\term{Accumulation Point}
A point $p$ of a topological space such that every neighbourhood of $p$ contains a point of a set $S$ different from $p$. Also called a limit point or cluster point.

\term{Adaptive Mesh Refinement}
A numerical technique that locally refines or coarsens a computational mesh based on error estimates, improving accuracy where needed.

\term{Additive Inverse}
For an element $a$ of a group written additively, the element $-a$ such that $a + (-a) = 0$, the identity element.

\term{Adjoint}
The operator $A^*$ associated with a linear map $A$ between inner product spaces satisfying $\langle Ax, y \rangle = \langle x, A^* y \rangle$ for all vectors $x, y$. Also called the conjugate transpose in the matrix setting.

\term{Adjunction}
A pair of functors $F \dashv G$ with a natural bijection $\operatorname{Hom}(FX, Y) \cong \operatorname{Hom}(X, GY)$, a fundamental concept in category theory.

\term{Albanese}
The Albanese variety of a complex algebraic variety $X$ is the universal abelian variety arising as the target of the Albanese morphism $\mathrm{alb}_X : X \to \mathrm{Alb}(X)$, a covariant analogue of the Jacobian and Picard constructions built from the space of global holomorphic $1$-forms. It is characterized by the universal property that every morphism from $X$ to an abelian variety factors through it. For a smooth projective curve the Albanese variety is canonically isomorphic to its Jacobian, and the Albanese morphism recovers the Abel--Jacobi map.

\term{Algebra}
A vector space equipped with a bilinear multiplication. Also the branch of mathematics studying symbolic manipulation and the structures built from sets and operations.

\term{Algebraic Closure}
A field extension $\bar{k}$ of a field $k$ that is algebraically closed and algebraic over $k$. Every field has an algebraic closure, unique up to isomorphism.

\term{Algebraic Extension}
A field extension $K/F$ in which every element of $K$ is algebraic over $F$, i.e. a root of some non-zero polynomial over $F$.

\term{Algebraic Number}
A complex number that is a root of some non-zero polynomial with rational coefficients. Numbers that are not algebraic are called transcendental.

\term{Algebraic Structure}
A set equipped with one or more operations satisfying a specified collection of axioms, such as a group, ring, field, or lattice.

\term{Algorithm}
A finite, unambiguous procedure that, given an input, produces an output after finitely many elementary steps. Algorithms are the objects of study of computability theory and complexity theory.

\term{Allais Paradox}
An observed choice pattern in which people's decisions between risky gambles violate the independence axiom of expected utility theory, even for careful reasoners. In the classic pair of gambles, most prefer the sure million over a gamble with a small chance of nothing, yet prefer the gamble offering a $0.10$ chance at five million over the one offering a $0.11$ chance at one million, the reversal hinging on the $0.01$ difference in probability. Such behaviour prompted alternative theories of risk such as prospect theory.

\term{Almost Everywhere}
A property that holds on a set whose complement has measure zero. Functions that differ on a null set are considered equal almost everywhere.

\term{Almost Sure Convergence}
The convergence of a sequence of random variables $X_n \to X$ on an event of probability 1, also called convergence with probability one.

\term{Alternating Group}
The group $A_n$ of even permutations of a set of $n$ symbols. It is a normal subgroup of the symmetric group $S_n$ and is simple for $n \geq 5$.

\term{Analysis}
The branch of mathematics built on limits, continuity, differentiation, and integration, including real analysis, complex analysis, and functional analysis.

\term{Analytic Continuation}
The extension of an analytic function to a larger domain by power series or functional relations, unique when the larger domain is connected.

\term{Analytic Function}
A function that is locally representable by a convergent power series. A function is holomorphic if and only if it is analytic.

\term{Annihilator}
For a subspace $S$ of a vector space $V$ with dual $V^*$, the subspace $S^\circ$ of all linear functionals vanishing on $S$.

\term{Annulus}
The region between two concentric circles, $\{z : r < |z| < R\}$ in the complex plane. It is a standard domain for Laurent series expansions.

\term{Antiderivative}
A function $F$ whose derivative is a given function $f$, so that $F' = f$. The set of all antiderivatives of $f$ differs by constants.

\term{Antisymmetric Relation}
A relation $R$ on a set such that $xRy$ and $yRx$ together imply $x = y$. Partial orders are antisymmetric relations.

\term{Archimedean Property}
The property of an ordered field that for any positive elements $a, b$ there exists a natural number $n$ with $na > b$. The real numbers satisfy it; non-archimedean fields do not.

\term{Argument Principle}
The theorem that $\frac{1}{2\pi i}\int_\gamma \frac{f'(z)}{f(z)}\,dz = N - P$, where $N$ is the number of zeros and $P$ the number of poles of a meromorphic function $f$ inside the closed curve $\gamma$.

\term{Arrow's Impossibility Theorem}
A theorem stating that with at least three alternatives, no social welfare function can simultaneously satisfy unanimity, independence of irrelevant alternatives, and non-dictatorship. It shows that no fair, consistent voting method exists for aggregating individual preferences, and it is a central result of welfare economics.

\term{Artinian Ring}
A ring in which every descending chain of left ideals stabilizes, so the descending chain condition holds on left ideals. By Hopkins's theorem a left Artinian ring is left Noetherian and its Jacobson radical is nilpotent, a result with no analogue for general Noetherian rings. Finite rings, matrix rings over division rings, and finite products of these are Artinian, while polynomial rings and most integral domains are not. Artinian rings play a central role in noncommutative ring theory, where the Artinian simple rings are exactly the matrix rings over division rings.

\term{Arzelà--Ascoli Theorem}
A criterion for a family of continuous functions on a compact metric space to be relatively compact in the uniform topology: it suffices that the family be equicontinuous and pointwise bounded. Consequently, a sequence of continuous functions on a compact metric space that is equicontinuous and pointwise bounded has a uniformly convergent subsequence, and $C(X)$ is compact in the uniform topology under these hypotheses.

\term{Asymptote}
A line that a curve approaches arbitrarily closely but does not in general meet, describing the behaviour of the curve at infinity.

\term{Automorphic Form}
A smooth function on a symmetric space, or more generally on the adelic points of a reductive group, that is invariant up to a fixed factor under the left action of an arithmetic or congruence subgroup. Classical modular forms are the case of weight $k$ on the upper half-plane, transforming as $f\bigl(\frac{az+b}{cz+d}\bigr) = (cz+d)^k f(z)$. Automorphic forms generate automorphic representations by right translation, and their Hecke eigenvalues are tied to $L$-functions, making them central objects in the Langlands program.

\term{Automorphism}
An isomorphism from a mathematical structure to itself. The set of automorphisms of a structure forms a group under composition.

\term{Axiom}
A statement taken to be self-evidently true, serving as a starting point for further reasoning. Modern mathematics relies on explicit axiom systems such as the Zermelo--Fraenkel axioms.

\term{Axiom of Choice}
The set-theoretic axiom stating that for any collection of non-empty sets there is a choice function selecting one element from each. It is independent of the Zermelo--Fraenkel axioms.


\chapter*{B}
\label{chap:b}

\term{Backward Stability}
A numerical algorithm is backward stable if its computed answer is the exact answer to a nearby problem. In symbols, the output for input $x$ is the exact result for $x+\Delta x$, where $\|\Delta x\|$ is small compared with $\|x\|$. \par\noindent\textbf{Applications:} backward stability explains why matrix algorithms can be trusted even when round-off error is unavoidable.

\term{Baire Category Theorem}
The Baire category theorem says that a complete metric space cannot be built as a countable union of thin sets, called nowhere-dense sets. Equivalently, a countable intersection of open dense sets is dense. \par\noindent\textbf{Applications:} it proves that certain objects exist in analysis and supports the uniform boundedness and open mapping theorems.

\term{Banach Fixed-Point Theorem}
If a map on a complete metric space brings points strictly closer together, then it has exactly one fixed point, and repeated iteration converges to it. This is also called the contraction mapping principle. \par\noindent\textbf{Applications:} it proves existence and uniqueness of solutions to equations and gives a practical iteration method for finding them.

\term{Banach Space}
A Banach space is a normed vector space in which every Cauchy sequence converges to a point in the space. Completeness means that a calculation whose successive approximations become arbitrarily close does not leave the space. \par\noindent\textbf{Applications:} Banach spaces provide the setting for infinite-dimensional approximation, integral equations, optimization, and differential equations.

\term{Banach--Steinhaus Theorem}
The uniform boundedness principle says that a family of bounded linear operators on a Banach space is uniformly bounded on the unit ball if each individual vector is sent to a bounded family of values. Its proof uses the Baire category theorem. \par\noindent\textbf{Applications:} it prevents hidden blow-up in families of approximation operators and is useful in functional analysis and numerical methods.

\term{Banach--Tarski Paradox}
Using a finite number of non-measurable pieces and rigid motions, a solid ball in three-dimensional space can be reassembled into two balls of the same size. This does not violate conservation of volume because the pieces have no well-defined volume; the construction uses the axiom of choice. \par\noindent\textbf{Applications:} the paradox clarifies the limits of assigning volume to arbitrary sets and the role of symmetry in measure theory.

\term{Base Case}
The base case is the first statement checked in a proof by mathematical induction. After it is established, an induction step shows that truth for one integer implies truth for the next. \par\noindent\textbf{Applications:} base cases are essential in recursive algorithms, proofs about sequences, and definitions of data structures.

\term{Basis}
A basis of a vector space is a linearly independent set that spans the whole space. Every vector has one and only one finite linear combination of basis vectors. \par\noindent\textbf{Applications:} coordinates, Fourier expansions, finite-element methods, and computer graphics all represent objects using a suitable basis.

\term{Bayes' Theorem}
Bayes' theorem relates conditional probabilities: $P(A\mid B)=P(B\mid A)P(A)/P(B)$ when $P(B)>0$. It updates the probability of a hypothesis after new evidence is observed. \par\noindent\textbf{Applications:} Bayesian statistics, medical diagnosis, spam filtering, fault detection, and machine learning use this update rule.

\term{Bernoulli Numbers}
The Bernoulli numbers $B_n$ are defined by $x/(e^x-1)=\sum_{n\ge0}B_nx^n/n!$. They occur in formulas for sums of powers and in the Euler--Maclaurin formula. \par\noindent\textbf{Applications:} they improve numerical integration and connect combinatorics with special values of the Riemann zeta function.

\term{Bernoulli Trial}
A Bernoulli trial is an experiment with two possible outcomes, usually called success and failure, with probabilities $p$ and $1-p$. Independent repetitions form a Bernoulli process. \par\noindent\textbf{Applications:} Bernoulli trials model yes/no tests, defective/nondefective items, clicks, infections, and binary classification outcomes.

\term{Berry's Paradox}
Berry's paradox arises from the phrase “the smallest positive integer not definable in fewer than twenty words.” The phrase appears to define such an integer in fewer than twenty words, creating a contradiction between informal language and formal definitions. \par\noindent\textbf{Applications:} it motivates precise languages in logic, computer science, and the study of what can be described or computed.

\term{Bertrand's Paradox}
Bertrand's paradox shows that the probability of a randomly chosen chord of a circle being longer than a side of an inscribed equilateral triangle depends on how “random chord” is defined. \par\noindent\textbf{Applications:} it teaches that a probability model must specify its sampling procedure, a principle important in statistics, simulation, and scientific experiments.

\term{Besov Space}
Besov spaces $B^s_{p,q}$ measure the smoothness of functions using how their detail changes across scales. The parameter $s$ describes smoothness, while $p$ and $q$ describe size and how different scales are combined; Sobolev spaces are important special cases. \par\noindent\textbf{Applications:} Besov spaces are used in partial differential equations, wavelet compression, image denoising, and nonlinear approximation.

\term{Bessel Function}
Bessel functions solve the differential equation $x^2y''+xy'+(x^2-n^2)y=0$. The functions $J_n$ and $Y_n$ are its standard independent solutions. \par\noindent\textbf{Applications:} they describe vibrations of circular membranes, heat flow in cylinders, electromagnetic waves, and acoustic modes.

\term{Bijection}
A bijection is a function that is both one-to-one and onto: every input has a different output, and every target is reached exactly once. A bijection proves that two sets have the same number of elements, even when they are infinite. \par\noindent\textbf{Applications:} bijections are used in counting, data encoding, coordinate changes, and proofs about infinite sets.

\term{Bilinear Form}
A bilinear form is a function $B:V\times W\to k$ that is linear in each input when the other is held fixed. Inner products and matrix expressions $x^TAy$ are examples. \par\noindent\textbf{Applications:} bilinear forms describe energy, angles, quadratic optimization, finite-element models, and the geometry of conic sections.

\term{Binary Operation}
A binary operation combines two elements of a set and returns an element of the same set: $*:S\times S\to S$. Addition and multiplication are examples, but the operation need not be commutative or associative. \par\noindent\textbf{Applications:} binary operations define the basic structure of groups, rings, logic gates, and algebraic data types.

\term{Binary Relation}
A binary relation between sets $A$ and $B$ is a collection of ordered pairs $(a,b)$ describing which elements are related. A function is a special relation in which each $a$ is paired with exactly one $b$. \par\noindent\textbf{Applications:} relations represent networks, database links, ordering, logical statements, and allowable transitions in algorithms.

\term{Binomial Coefficient}
The binomial coefficient $\binom nk=n!/[k!(n-k)!]$ counts the ways to choose $k$ objects from $n$ without regard to order. It is also the coefficient of $x^ky^{n-k}$ in $(x+y)^n$. \par\noindent\textbf{Applications:} binomial coefficients occur in probability, combinatorics, polynomial expansion, and error-correcting codes.

\term{Binomial Distribution}
The binomial distribution gives the probability of $k$ successes in $n$ independent Bernoulli trials: $P(X=k)=\binom nkp^k(1-p)^{n-k}$. Its parameters are the number of trials $n$ and success probability $p$. \par\noindent\textbf{Applications:} it models counts such as successful tests, conversions, defects, and disease cases in a fixed sample.

\term{Binomial Theorem}
The binomial theorem expands $(x+y)^n$ as $\sum_{k=0}^n\binom nkx^{n-k}y^k$. It converts a power of a sum into a sum of simpler terms. \par\noindent\textbf{Applications:} it is used in algebra, probability, numerical approximations, and the derivation of many combinatorial identities.

\term{Bipartite Graph}
A bipartite graph has its vertices divided into two sets, with every edge joining the two different sets. It has no odd-length cycle. \par\noindent\textbf{Applications:} bipartite graphs model job assignments, matching students to projects, network flow, recommendation systems, and molecular structures.

\term{Birch--Swinnerton-Dyer Conjecture}
For an elliptic curve $E$ over $\mathbb Q$, this conjecture relates the number of independent rational points on $E$ to how many times its $L$-function vanishes at $s=1$. It also predicts the leading coefficient using arithmetic quantities such as the regulator and Tate--Shafarevich group. \par\noindent\textbf{Applications:} it links analytic formulas with rational solutions and guides research in arithmetic geometry.

\term{Birthday Paradox}
In a group of 23 people, the probability that two share a birthday is already greater than $1/2$, assuming 365 equally likely birthdays. The probability of no match among $n$ people is $\prod_{k=0}^{n-1}(365-k)/365$. \par\noindent\textbf{Applications:} the effect is important in collision probabilities, hash tables, computer security, and the design of identifiers.

\term{Borel--Cantelli Lemmas}
If events $E_n$ satisfy $\sum P(E_n)<\infty$, then with probability one only finitely many of them occur. If the events are independent and $\sum P(E_n)=\infty$, then infinitely many occur with probability one. \par\noindent\textbf{Applications:} these lemmas test whether rare events eventually stop occurring and help prove convergence in probability theory.

\term{Borel Measure}
A Borel measure assigns sizes or probabilities to Borel sets, the sets generated from open sets by countable unions, intersections, and complements. On familiar spaces such as $\mathbb R^n$, it includes the usual notions of length, area, and volume. \par\noindent\textbf{Applications:} Borel measures form the foundation of probability, integration, and geometric measurement.

\term{Borel Set}
A Borel set is built from open sets using countable unions, countable intersections, and complements. The Borel sets form the smallest $\sigma$-algebra containing all open sets. \par\noindent\textbf{Applications:} intervals, regions, and countable combinations of them can be assigned probabilities or measures in analysis and statistics.

\term{Borsuk--Ulam Theorem}
Every continuous map from the sphere $S^n$ to $\mathbb R^n$ gives the same value at some pair of opposite points. For $n=2$, one consequence is that at some pair of antipodal points on Earth, temperature and atmospheric pressure are simultaneously equal. \par\noindent\textbf{Applications:} the theorem supports fair-division results such as the ham-sandwich theorem and studies of continuous data on spheres.

\term{Bound}
A bound is a number that limits a set, sequence, or function from above or below. A set is bounded if it has both an upper and a lower bound. \par\noindent\textbf{Applications:} bounds control error, guarantee stability, establish convergence, and quantify worst-case performance in algorithms.

\term{Boundary}
The boundary of a set $S$ consists of points whose every neighborhood meets both $S$ and its complement. For an interval $[0,1]$, the boundary is $\{0,1\}$. \par\noindent\textbf{Applications:} boundaries specify interfaces, constraints, and edge conditions in geometry, image processing, and differential equations.

\term{Boundary-value Problem}
A boundary-value problem asks for a solution of a differential equation subject to conditions on the boundary of a domain, such as fixed values (Dirichlet conditions) or fixed normal derivatives (Neumann conditions). \par\noindent\textbf{Applications:} it models temperature in a solid, vibration of a membrane, fluid flow, and electrostatic fields.

\term{Bounded Function}
A function is bounded if all its values lie between two fixed numbers: $m\le f(x)\le M$. For example, sine is bounded between $-1$ and $1$. \par\noindent\textbf{Applications:} boundedness prevents uncontrolled output and is used in stability analysis, probability, optimization, and signal processing.

\term{Bounded Linear Operator}
A linear map $T:X\to Y$ between normed spaces is bounded if $\|Tx\|\le C\|x\|$ for some constant $C$. For linear maps, boundedness is equivalent to continuity, and $\|T\|=\sup_{\|x\|\le1}\|Tx\|$ measures the largest amplification. \par\noindent\textbf{Applications:} operator norms quantify sensitivity in differential equations, numerical linear algebra, and signal transformations.

\term{Boy or Girl Paradox}
The probability that two children are both girls depends on how the information “at least one is a girl” was obtained. Under the usual equally likely interpretation it is $1/3$, while choosing a random child and finding that child is a girl gives $1/2$. \par\noindent\textbf{Applications:} the paradox teaches careful conditioning, which is essential in medical testing, surveys, and statistical reporting.

\term{Braess's Paradox}
Adding a road to a congested network can make every traveller's journey longer because selfish route choices can create a worse equilibrium. \par\noindent\textbf{Applications:} the paradox informs traffic planning, communication networks, supply chains, and the design of systems where local decisions affect global performance.

\term{Branch Point}
A branch point is a point around which a multi-valued complex function, such as $\sqrt z$ or $\log z$, cannot be made single-valued and continuous without choosing a branch cut. Going once around the point may change the value. \par\noindent\textbf{Applications:} branch points occur in complex integration, wave propagation, fluid flow, and mathematical physics.

\term{Brouwer Fixed-Point Theorem}
Every continuous map from a closed ball in $\mathbb R^n$ to itself has at least one fixed point $f(x)=x$. Continuity and the fact that the map stays inside the ball are essential. \par\noindent\textbf{Applications:} the theorem proves equilibrium existence in economics and game theory and helps establish solutions of nonlinear equations.

\term{Brownian Motion}
Brownian motion is a continuous random process whose independent increments are normally distributed and whose variance grows with elapsed time. It is the mathematical model of random diffusion. \par\noindent\textbf{Applications:} it models molecular motion, particle diffusion, financial price models, heat transport, and stochastic differential equations.

\term{Burali--Forti Paradox}
The collection of all ordinal numbers cannot be a set: if it were, it would have an ordinal larger than every ordinal in the collection, including itself. This contradiction helped motivate modern axioms that distinguish sets from proper classes. \par\noindent\textbf{Applications:} it clarifies the foundations of set theory and prevents uncontrolled definitions of “the set of everything.”

\term{Burnside's Lemma}
If a finite group $G$ acts on a finite set $X$, the number of distinct orbits is $|X/G|=|G|^{-1}\sum_{g\in G}|X^g|$, where $X^g$ is the set fixed by $g$. It counts objects while treating symmetric copies as the same. \par\noindent\textbf{Applications:} it counts necklaces, colorings, chemical structures, and configurations under rotation or reflection.

\term{Busy Beaver}
The Busy Beaver function gives the largest number of steps, or ones written, by any halting Turing machine with a fixed number of states. It grows faster than every computable function and therefore cannot itself be computed by an algorithm. \par\noindent\textbf{Applications:} Busy Beaver illustrates the limits of computation, undecidability, and the difference between simple programs and predictable behavior.

\chapter*{C}
\label{chap:c}

\term{Calculus}
The study of limits, derivatives, integrals, and infinite series, divided into differential and integral calculus.

\term{Calculus of Variations}
The branch of analysis concerned with minimizing or maximizing functionals, maps from a space of functions to $\mathbb{R}$, such as $J(y) = \int_a^b F(x, y, y')\,dx$. Critical points satisfy the Euler--Lagrange equations $\frac{\partial F}{\partial y} - \frac{d}{dx}\frac{\partial F}{\partial y'} = 0$. The classical brachistochrone problem motivated the subject, and Noether's theorem ties symmetries to conservation laws. Extrema are classified as weak or strong. It underlies optimal control theory and provides variational formulations of elliptic PDEs.

\term{Coordinate Ring}
The ring of polynomial functions on an algebraic variety $V$, defined as the quotient $k[x_1, \ldots, x_n] / I(V)$, where $I(V)$ is the ideal of polynomials vanishing on $V$. The variety is recovered as $\operatorname{Spec} A$ for the coordinate ring $A$.

\term{C*-algebra}
A Banach algebra over the complex numbers with an involution $*$ satisfying $\|a^*a\| = \|a\|^2$; the mathematical framework for quantum mechanics.

\term{Cantor--Bernstein Theorem}
If there are injections $f: A \to B$ and $g: B \to A$, then there is a bijection between $A$ and $B$. Consequently, two sets have equal cardinality as soon as each injects into the other, so cardinalities can be compared using injections alone. The theorem is also called the Schröder--Bernstein theorem.

\term{Cantor's Paradox}
The assumption of a universal set $V$ of all sets yields a contradiction: by Cantor's theorem $|\mathcal{P}(V)| > |V|$, yet $\mathcal{P}(V)$ is a set and hence a member of $V$. It shows there is no set of all sets and motivates the ZF axioms.

\term{Cantor's Theorem}
For every set $X$, the power set $\mathcal{P}(X)$ has strictly greater cardinality than $X$: there is no surjection $X \to \mathcal{P}(X)$. The proof is Cantor's diagonal argument, which shows that no list $x_1, x_2, \ldots$ of elements can cover $\mathcal{P}(X)$. Since $\mathbb{R}$ has the same cardinality as $\mathcal{P}(\mathbb{N})$, the uncountability of the reals follows as a corollary.

\term{Canonical Bundle}
The canonical bundle of a smooth complex manifold $X$ of dimension $n$ is the determinant line bundle $K_X = \det T^*X = \bigwedge^n T^*X$, whose sections are holomorphic $n$-forms. Positivity or negativity of $K_X$ governs Kodaira vanishing and the minimal model program: if $K_X$ is ample, the canonical ring is finitely generated and $X$ admits a canonical model. $K_X$ is holomorphically trivial precisely when $X$ is Calabi--Yau, as for complex tori, and the Kodaira dimension measures the growth of $H^0(X, K_X^{\otimes m})$.

\term{Cardinality}
The size of a set, measured by the number of its elements; two sets have equal cardinality when a bijection exists between them.

\term{Cartesian Product}
The set $A \times B = \{(a,b) : a \in A, b \in B\}$ of all ordered pairs with first coordinate in $A$ and second in $B$.

\term{Casorati--Weierstrass Theorem}
If $f$ has an isolated essential singularity at $z_0$, then the image of every punctured neighbourhood of $z_0$ is dense in $\mathbb{C}$: near an essential singularity the function comes arbitrarily close to every complex value. This contrasts sharply with removable singularities, where $f$ has a finite limit, and with poles, where $|f(z)| \to \infty$.

\term{Category}
A mathematical structure consisting of objects and morphisms between them, together with an associative composition law and identity morphisms.

\term{Category Theory}
The abstract study of mathematical structures and the structure-preserving maps between them, via categories, functors, and natural transformations.

\term{Cauchy Problem}
An initial-value problem for a partial differential equation, typically specifying the value of the solution and its normal derivative on a hypersurface (the initial surface), along with the equation to be satisfied in the domain.

\term{Cauchy--Kovalevskaya Theorem}
The classical local existence and uniqueness theorem for analytic partial differential equations. A first-order system of PDEs with analytic coefficients and analytic initial data, written in the normal form $\partial_t u_i = F_i(x, t, u, \nabla u)$, has a unique analytic solution in a neighbourhood of the initial surface $t = 0$. It establishes that analytic Cauchy problems are locally well-posed, though it says nothing about the smoothing behaviour of non-analytic equations.

\term{Cauchy--Schwarz Inequality}
In any inner product space, $|\langle u, v \rangle|^2 \leq \langle u, u \rangle \langle v, v \rangle$ with equality iff $u$ and $v$ are linearly dependent. For vectors in $\mathbb{R}^n$: $(\sum a_i b_i)^2 \leq (\sum a_i^2)(\sum b_i^2)$. For integrals: $(\int fg)^2 \leq \int f^2 \cdot \int g^2$. Fundamental to analysis, it implies the triangle inequality and defines the notion of angle in abstract inner product spaces.

\term{Cauchy Integral Theorem}
If $f$ is holomorphic on a simply connected domain $D$ and $\gamma$ is any closed curve contained in $D$, then $\oint_\gamma f(z)\,dz = 0$. This is the foundation of complex analysis and implies the Cauchy integral formula and the residue theorem.

\term{Cauchy Sequence}
A sequence $(x_n)$ in a metric space is Cauchy if for every $\epsilon > 0$ there exists $N$ such that $d(x_n, x_m) < \epsilon$ for all $n, m \geq N$. A metric space is complete iff every Cauchy sequence converges. The real numbers are complete, but $\mathbb{Q}$ is not: the sequence $3, 3.1, 3.14, 3.141, \ldots$ approximating $\pi$ is Cauchy but has no limit in $\mathbb{Q}$.

\term{Cauchy's Integral Formula}
If $f$ is holomorphic on a simply connected domain $D$ and $\gamma$ is a positively oriented simple closed curve in $D$, then for any $a$ inside $\gamma$: $f(a) = \frac{1}{2\pi i} \oint_\gamma \frac{f(z)}{z-a}\,dz$. This formula shows that holomorphic functions are infinitely differentiable and completely determined by their values on any closed curve. All higher derivatives follow: $f^{(n)}(a) = \frac{n!}{2\pi i} \oint_\gamma \frac{f(z)}{(z-a)^{n+1}}\,dz$.

\term{Cauchy's Integral Theorem}
If $f$ is holomorphic on a simply connected domain $D$ and $\gamma$ is any closed curve contained in $D$, then $\oint_\gamma f(z)\,dz = 0$. By homotopy invariance the integral vanishes as long as $\gamma$ is null-homotopic in the domain of holomorphy, so holomorphic functions have path-independent contour integrals. As the foundation of contour integration, the theorem leads to Cauchy's integral formula and the residue theorem.

\term{Cauchy's Theorem}
If $f$ is holomorphic on a simply connected domain $D$ and $\gamma$ is a closed curve in $D$, then $\oint_\gamma f(z)\,dz = 0$. This is the foundation of complex analysis. It generalises the Fundamental Theorem of Calculus: path independence of integrals of analytic functions. Consequences include Liouville's theorem, the Maximum Modulus Principle, and the Identity Theorem.

\term{Cayley's Theorem}
Every group $G$ is isomorphic to a subgroup of the symmetric group $S_G$ acting on $G$ by left multiplication. Thus every abstract group is a permutation group. For a finite group of order $n$, this embeds $G$ into $S_n$. This result, due to Arthur Cayley (1854), justifies the study of permutation groups as the primary tool for understanding all groups.

\term{Cayley--Hamilton Theorem}
Every square matrix satisfies its own characteristic equation: if $p(\lambda) = \det(A - \lambda I)$, then $p(A) = 0$. For a $2 \times 2$ matrix $A = \begin{pmatrix} a & b \\ c & d \end{pmatrix}$, this gives $A^2 - (a+d)A + (ad-bc)I = 0$. The theorem allows matrix powers to be reduced and is used in computing matrix functions and solving systems of linear differential equations.

\term{Center (Group Theory)}
The center of a group $G$ is $Z(G) = \{g \in G : gh = hg \text{ for all } h \in G\}$. It is a normal subgroup of $G$. For a non-abelian group, the center may be trivial: $Z(S_n) = \{e\}$ for $n \geq 3$. The center measures how far a group is from being abelian. If $G/Z(G)$ is cyclic, then $G$ is abelian.

\term{Central Limit Theorem}
If $X_1, X_2, \ldots, X_n$ are independent identically distributed random variables with mean $\mu$ and variance $\sigma^2$, then $\frac{\bar{X}_n - \mu}{\sigma/\sqrt{n}}$ converges in distribution to the standard normal $\mathcal{N}(0,1)$ as $n \to \infty$. This explains why the normal distribution is ubiquitous: sums and averages of independent random variables tend toward normality regardless of the original distribution.

\term{Centralizer}
The subgroup of elements that commute with a given element of a group: $C_G(a) = \{g \in G : ga = ag\}$.

\term{Chain Complex}
A sequence of abelian groups or modules $A_* = \cdots \to A_{n+1} \xrightarrow{d_{n+1}} A_n \xrightarrow{d_n} A_{n-1} \to \cdots$ with $d_n \circ d_{n+1} = 0$, so that $\operatorname{im} d_{n+1} \subseteq \ker d_n$. The homology $H_n(A_*) = \ker d_n / \operatorname{im} d_{n+1}$ measures how far the sequence is from exact. Exact sequences have vanishing homology; quasi-isomorphisms are chain maps inducing isomorphisms on homology, and the mapping cone encodes the homological features of a map. Chain complexes are the central objects of homological algebra, with complexes of sheaves arising in algebraic geometry.

\term{Chain Rule}
If $f: \mathbb{R}^n \to \mathbb{R}^m$ and $g: \mathbb{R}^m \to \mathbb{R}^p$ are differentiable, then the derivative of $g \circ f$ is $D(g \circ f)(x) = Dg(f(x)) \cdot Df(x)$, the product of Jacobian matrices. In single-variable calculus: $\frac{d}{dx} g(f(x)) = g'(f(x)) \cdot f'(x)$. The chain rule is the fundamental differentiation rule for composite functions.

\term{Chaos}
The sensitive dependence on initial conditions exhibited by some deterministic dynamical systems, in which nearby trajectories separate exponentially.

\term{Characteristic Class}
A characteristic class assigns to each vector bundle $E \to X$ a cohomology class $c(E) \in H^*(X)$ that is natural: $c(f^*E) = f^*c(E)$. Important examples are the Chern classes $c_i$ for complex bundles, the Pontryagin and Stiefel--Whitney classes for real bundles, and the Euler class. The Whitney sum formula governs behaviour under direct sums. Characteristic classes obstruct triviality of bundles and, in obstruction theory, the existence of sections; they also enter index theory, as in the Atiyah--Singer index theorem.

\term{Characteristic Function}
(Probability) The function $\varphi_X(t) = E[e^{itX}]$ of a random variable $X$, which determines its distribution uniquely.

\term{Characteristic Polynomial}
The polynomial $\det(\lambda I - A)$ of a square matrix $A$, whose roots are the eigenvalues.

\term{Characteristic Variety}
For a linear partial differential operator $P$ of order $m$ on a manifold $X$ with principal symbol $p_m(x, \xi)$, the characteristic variety is the zero set $\{p_m(x, \xi) = 0\}$ in the cotangent bundle $T^*X$. Solvability and uniqueness hold away from it: Holmgren's theorem gives unique continuation of analytic Cauchy problems across non-characteristic hypersurfaces, while microlocal analysis shows that singularities of solutions propagate along the bicharacteristic flow within the characteristic variety. For a coherent $D_X$-module, the characteristic variety is the support of the associated graded module, and its involutivity (Cauchy--Kovalevskaya--Kashiwara theorem) forces dimension at least $\dim X$.

\term{Chebotarev Density Theorem}
For a finite Galois extension $L/K$ of number fields with Galois group $G$, the Chebotarev density theorem asserts that the primes of $K$ unramified in $L$ whose Frobenius automorphism lies in a fixed conjugacy class $C \subseteq G$ have natural density $|C|/|G|$. It generalises Dirichlet's theorem on primes in arithmetic progressions, which is recovered for cyclotomic extensions of $\mathbb{Q}$, and determines the distribution of primes that split completely in $L$, which have density $1/|G|$. The theorem is foundational in algebraic number theory and arithmetic geometry, underlying the Brauer--Siegel theorem and results on Galois representations.

\term{Chebyshev Polynomials}
The Chebyshev polynomials $T_n(x)$ are defined by $T_n(\cos \theta) = \cos(n\theta)$. They satisfy the recurrence $T_0(x) = 1$, $T_1(x) = x$, $T_{n+1}(x) = 2xT_n(x) - T_{n-1}(x)$. On $[-1,1]$, $|T_n(x)| \leq 1$. They are optimal for polynomial interpolation (minimising the Runge phenomenon) and appear in numerical analysis, approximation theory, and the solution of differential equations.

\term{Chebyshev's Inequality}
For a random variable $X$ with mean $\mu$ and finite variance $\sigma^2$, $P(|X - \mu| \geq k\sigma) \leq \frac{1}{k^2}$ for any $k > 0$. This gives a universal bound on tail probabilities. For any distribution, at least 75\% of values lie within 2 standard deviations of the mean, and at least 89\% within 3 standard deviations.

\term{Chern Classes}
Characteristic classes associated with complex vector bundles over a base space $X$. The total Chern class is $c(E) = 1 + c_1(E) + c_2(E) + \cdots \in H^*(X; \mathbb{Z})$. The first Chern class $c_1$ of a line bundle corresponds to the Euler class and measures topological twisting. Chern classes are fundamental in differential geometry, topology, and mathematical physics, particularly in the Atiyah--Singer Index Theorem.

\term{Chern--Weil Theory}
A method for constructing characteristic classes (Chern, Pontryagin, Euler) of vector bundles using curvature forms. Given a connection $\nabla$ on a bundle $E$ with curvature $\Omega$, one forms invariant polynomials in $\Omega$ whose cohomology classes are independent of the connection. This bridges differential geometry and topology, realising topological invariants via geometric data.

\term{Chi-squared (\ensuremath{\chi^2}) Distribution}
A continuous probability distribution that arises when summing the squares of $k$ independent standard normal random variables. It is widely used in hypothesis testing to determine "goodness of fit" between observed and expected data.

\term{Chi-squared (\ensuremath{\chi^2}) Test}
A statistical test used to determine if there is a significant difference between the expected frequencies and the observed frequencies in one or more categories of data.

\term{Chinese Remainder Theorem}
If the moduli $m_1, \ldots, m_k$ are pairwise coprime, then the system of congruences $x \equiv a_i \pmod{m_i}$ has a unique solution modulo $m_1 \cdots m_k$. Ring-theoretically this is the statement $R/(ab) \cong R/(a) \times R/(b)$ for coprime ideals $a$ and $b$ of a ring $R$; for $R = \mathbb{Z}$ it recovers the classical version. It is a workhorse of integer arithmetic, allowing computations modulo a product to be reduced to its smaller coprime factors.

\term{Christoffel Symbols}
Coefficients $\Gamma^k_{ij}$ describing the Levi--Civita connection on a Riemannian manifold with metric $g_{ij}$: $\Gamma^k_{ij} = \frac{1}{2} g^{kl}(\partial_i g_{jl} + \partial_j g_{il} - \partial_l g_{ij})$. They appear in the geodesic equation $\ddot{x}^k + \Gamma^k_{ij}\dot{x}^i\dot{x}^j = 0$ and in the covariant derivative $\nabla_i v^j = \partial_i v^j + \Gamma^j_{ik}v^k$. They are not tensors but transform in a controlled way under coordinate changes.

\term{Church--Turing Thesis}
The thesis that every function computable by an algorithm is computable by a Turing machine, and vice versa; equivalently, that effective computability coincides with computability by the lambda calculus or by general recursive functions. It pins down the informal notion of computability in a precise mathematical form. Since it identifies an intuitive concept with a formal one, it is a thesis rather than a proved theorem, although the equivalence among the various formal models is a theorem.

\term{Class Field Theory}
The classification of abelian extensions of number fields and local fields: for a number field $K$, finite abelian extensions $L/K$ correspond to open subgroups of the idele class group, and the Galois group $G(L/K)$ is isomorphic, via the reciprocity map, to the corresponding quotient. Artin reciprocity generalizes quadratic and higher reciprocity laws; the Hilbert class field, the maximal abelian unramified extension of $K$, has Galois group isomorphic to the ideal class group. The Kronecker--Weber theorem, that every abelian extension of $\mathbb{Q}$ is cyclotomic, is the base case.

\term{Classical Solution}
A solution of a differential equation that is differentiable as many times as required by the equation.

\term{Classification of Finite Simple Groups}
The complete classification, assembled over decades in tens of thousands of published pages, asserts that every finite simple group is one of: a cyclic group $\mathbb{Z}/p\mathbb{Z}$ of prime order, an alternating group $A_n$ ($n \geq 5$), a member of one of the infinite families of groups of Lie type, or one of the 26 sporadic groups culminating in the Monster. Its consequences are far-reaching: structural theorems for all finite groups, such as the Odd Order Theorem, follow from it, and the classification underpins much of modern group theory.

\term{Classification of Surfaces}
Every closed connected surface is homeomorphic to a sphere with $g$ handles (orientable case) or with $k$ crosscaps (non-orientable case); equivalently, to a connected sum of $g$ tori or of $k$ projective planes. Thus a closed surface is determined up to homeomorphism by whether it is orientable and by a single integer invariant, the genus.

\term{Clebsch--Gordan Coefficients}
Coefficients expressing the coupling of two angular momenta in quantum mechanics. The product of two irreducible representations of $SU(2)$ with spins $j_1$ and $j_2$ decomposes as $\bigoplus_{j=|j_1-j_2|}^{j_1+j_2} j$. The Clebsch--Gordan coefficients $\langle j_1 m_1 j_2 m_2 | jm \rangle$ give the expansion of coupled basis states in terms of product basis states. They are real numbers tabulated in standard references.

\term{Closed Graph Theorem}
The theorem that a linear map between Banach spaces with a closed graph is bounded (continuous).

\term{Closed Map}
A function $f: X \to Y$ between topological spaces is closed if the image of every closed set in $X$ is closed in $Y$. A continuous bijection from a compact space to a Hausdorff space is a closed map, hence a homeomorphism. The projection $\mathbb{R} \times [0,1] \to \mathbb{R}$ is closed, but $\mathbb{R} \times (0,1) \to \mathbb{R}$ is not: the hyperbola $\{(x, 1/x) : x > 0\}$ is closed in the domain but maps to $(0, \infty)$, which is not closed.

\term{Closed Set}
A subset $C$ of a topological space is closed if its complement is open. Equivalently, $C$ contains all its limit points. In $\mathbb{R}$ with the standard topology, closed sets include intervals $[a,b]$, finite sets, and the Cantor set. The closed sets are precisely the intersections of families of closed intervals, forming a topology dual to the open sets.

\term{Closure}
The closure of a set $S$ in a topological space is the smallest closed set containing $S$, equivalently the union of $S$ and its limit points.

\term{Closure Operator}
A closure operator on a set $X$ is a function $\operatorname{cl}: \mathcal{P}(X) \to \mathcal{P}(X)$ satisfying: $A \subseteq \operatorname{cl}(A)$, $A \subseteq B \implies \operatorname{cl}(A) \subseteq \operatorname{cl}(B)$, and $\operatorname{cl}(\operatorname{cl}(A)) = \operatorname{cl}(A)$. The topological closure, the algebraic closure in field theory, and the subgroup generated by a set are all closure operators.

\term{Cluster Point (Accumulation Point)}
A point $x$ is a cluster point of a sequence $(x_n)$ if every neighbourhood of $x$ contains infinitely many terms of the sequence. A point $p$ is a cluster point of a set $S$ if every neighbourhood of $p$ contains a point of $S$ distinct from $p$. For sequences in compact spaces, every sequence has a cluster point (Bolzano--Weierstrass).

\term{Coarea Formula}
A generalisation of Fubini's theorem for integrals over level sets. If $f: \mathbb{R}^n \to \mathbb{R}^m$ is smooth with $n \geq m$ and $g: \mathbb{R}^n \to [0,\infty)$ is integrable, then $\int_{\mathbb{R}^n} g(x) J_m f(x)\,dx = \int_{\mathbb{R}^m} \left(\int_{f^{-1}(y)} g\,d\mathcal{H}^{n-m}\right) dy$, where $J_m f$ is the Jacobian determinant of $f$. Used in geometric measure theory and the study of Sobolev functions.

\term{Coastline Paradox}
The measured length of a coastline, border, or other natural boundary increases without bound as the measuring unit shrinks, because the curve has detail at every scale. Coastlines are fractal-like, with length depending on the scale of measurement rather than being an intrinsic property.

\term{Cobordism}
Two closed oriented manifolds $M$ and $N$ of dimension $n$ are cobordant if there exists a compact oriented manifold $W$ of dimension $n+1$ whose boundary is $M \sqcup (-N)$. Cobordism is an equivalence relation, and the cobordism classes form the cobordism ring $\Omega_n^{SO}$. Thom computed this ring, showing $\Omega_*^{SO} \otimes \mathbb{Q}$ is a polynomial algebra generated by complex projective spaces $\mathbb{CP}^{2k}$.

\term{Code (Error-Correcting)}
A code is a subset $C \subseteq \Sigma^n$ of words of length $n$ over an alphabet $\Sigma$. The Hamming distance $d(u,v)$ counts positions where $u$ and $v$ differ. The minimum distance $d_{\min}$ determines error-correcting capability: a code with $d_{\min} = 2t+1$ can correct up to $t$ errors. The Hamming code $(7,4)$ has $d_{\min} = 3$ and corrects single-bit errors.

\term{Coding Theory}
The study of error-detecting and error-correcting codes, subsets $C$ of $\Sigma^n$ chosen so that distinct codewords stay far apart in Hamming distance. Linear codes over finite fields are vector subspaces with parameters $[n,k,d]$, and their minimum distance determines error-correcting capability. Bounds such as the Hamming, Singleton, and Gilbert--Varshamov bounds constrain the achievable parameters. Constructions over finite fields, including Reed--Solomon codes, connect coding theory to cryptography, design theory, and algebraic geometry.

\term{Codomain}
The set into which a function maps; the target set, which contains the range as a subset.

\term{Coefficient of Variation}
The coefficient of variation is $CV = \sigma/\mu$, the ratio of standard deviation to mean. It is a dimensionless measure of dispersion, useful for comparing variability across datasets with different units or scales. A $CV > 1$ indicates high relative variability; a $CV < 0.1$ indicates low relative variability.

\term{Combination}
An unordered selection of $k$ elements from a set of $n$, counted by the binomial coefficient $\binom{n}{k}$.

\term{Cofibration}
A continuous map $i: A \to X$ of topological spaces possessing the homotopy extension property: any homotopy of $i(a)$ extends to a homotopy of all of $X$ fixing $A$. Every map $f$ factors as a homotopy equivalence followed by a cofibration via the mapping cylinder. In model categories, cofibrations form one of the three classes of distinguished maps; cofibrant objects are those whose map from the initial object is a cofibration, and cofibrant replacement resolves objects up to weak equivalence. The notion is dual to fibrations.

\term{Coherence (Category Theory)}
Coherence theorems in category theory state that all diagrams built from the structural natural isomorphisms (associator, unitors, braiding) commute. Mac Lane's coherence theorem for monoidal categories asserts that any two parallel morphisms built from associators and unitors are equal. This justifies suppressing parentheses in monoidal categories without ambiguity.

\term{Coherent Sheaf}
A coherent sheaf on a scheme or complex space $X$ is a sheaf of $\mathcal{O}_X$-modules that is locally finitely presented: locally isomorphic to the cokernel of a map $\mathcal{O}_X^{\oplus p} \to \mathcal{O}_X^{\oplus q}$. The coherent sheaves form an abelian category, since kernels, cokernels, and extensions of coherent sheaves are coherent. On proper schemes, coherent sheaves have finite-dimensional cohomology groups; over $\mathbb{C}$, Serre's GAGA comparison identifies analytic and algebraic cohomology. For an ample line bundle $L$, Serre vanishing gives $H^i(X, \mathcal{F} \otimes L^{\otimes m}) = 0$ for $i > 0$ and $m$ sufficiently large.

\term{Cohomology}
Cohomology is the contravariant functor $H^n(X; G)$ associated to a topological space $X$ and coefficient group $G$, computed from the dual of the singular chain complex. Unlike homology, cohomology carries a natural ring structure (the cup product), making $H^*(X) = \bigoplus H^n(X)$ a graded ring. Cohomology is more computable than homology in many cases and carries richer structure.

\term{Colimit}
A categorical construction dual to a limit: the colimit of a diagram $D: J \to \mathcal{C}$ is the universal cocone, an object $\operatorname{colim} D$ with maps $D_j \to \operatorname{colim} D$ through which every other cocone factors uniquely. Direct limits, pushouts, and coproducts are special cases. Colimits of sets and modules exist and are built from generators and relations. The construction is a left adjoint to the diagonal functor, and filtered colimits commute with finite limits in categories such as modules.

\term{Commutative Algebra}
The study of commutative rings, their ideals, and modules, forming the algebraic foundation of algebraic geometry and number theory.

\term{Compact Operator}
A bounded linear operator $T: X \to Y$ between Banach spaces is compact if the image of the unit ball is relatively compact (its closure is compact). Equivalently, for every bounded sequence $(x_n)$, the sequence $(Tx_n)$ has a convergent subsequence. Compact operators generalise matrices to infinite dimensions; the spectral theory of compact operators closely mirrors finite-dimensional linear algebra.

\term{Compact Space}
A topological space is compact if every collection of open sets that covers it has a finite subcollection that still covers it. In $\mathbb{R}^n$, the Heine--Borel theorem says that compactness is equivalent to being closed and bounded. A continuous real-valued function on a compact space attains both a maximum and a minimum. Applications: compactness helps prove that optimization problems have solutions and that limits or subsequences exist.

\term{Compactness Theorem}
In first-order logic, a set of sentences has a model if and only if every finite subset of it has a model; the substantive direction is that finite satisfiability forces satisfiability of the whole set. It follows from Gödel's completeness theorem and implies the Löwenheim--Skolem theorems and the existence of nonstandard models of arithmetic. Compactness is widely used to transfer finite results to infinite settings, as when the four-colour theorem for finite planar graphs is extended to infinite graphs.

\term{Completeness (Metric Space)}
A metric space $(X,d)$ is complete if every Cauchy sequence in $X$ converges to a point in $X$. The real numbers are complete; the rationals are not. The completion of a metric space $X$ is a complete metric space $\hat{X}$ containing $X$ as a dense subspace, unique up to isometry. The reals are the completion of the rationals.

\term{Complexity Class}
A set of decision problems solvable within prescribed resource bounds, such as time or space as functions of input length. Major classes include $\mathsf{P}$, $\mathsf{NP}$, $\mathsf{PSPACE}$, and $\mathsf{EXP}$. Reductions order problems by difficulty; a problem is complete for a class if it lies in the class and every member reduces to it, as in NP-completeness. Hierarchy theorems show the corresponding time and space classes are strictly increasing. Whether $\mathsf{P} = \mathsf{NP}$ remains open.

\term{Complex Manifold}
A real manifold of even dimension with an atlas whose transition functions are biholomorphic, modeled on open subsets of $\mathbb{C}^n$ with holomorphic coordinate changes. Complex manifolds carry a holomorphic tangent bundle and Dolbeault cohomology groups $H^{p,q}$. Riemann surfaces, the complex dimension-one case, and complex projective space $\mathbb{CP}^n$ are the primary examples. An almost-complex structure is a pointwise condition that is integrable exactly when the Nijenhuis tensor vanishes, by the Newlander--Nirenberg theorem.

\term{Complex Number}
A complex number is $z = a + bi$ where $a, b \in \mathbb{R}$ and $i^2 = -1$. The complex plane $\mathbb{C}$ is identified with $\mathbb{R}^2$. The modulus is $|z| = \sqrt{a^2+b^2}$ and the argument is $\arg(z) = \arctan(b/a)$. Complex numbers form an algebraically closed field: every non-constant polynomial has a root in $\mathbb{C}$ (Fundamental Theorem of Algebra).

\term{Composition Series}
A composition series of a group $G$ is a finite chain $G = G_0 \triangleright G_1 \triangleright \cdots \triangleright G_n = \{e\}$ where each $G_{i+1}$ is a maximal normal subgroup of $G_i$, so each factor $G_i/G_{i+1}$ is simple. The Jordan--Hölder theorem states that any two composition series have the same length and the same multiset of simple factors (up to isomorphism).

\term{Concave Function}
A function $f$ satisfying $f(tx + (1-t)y) \geq tf(x) + (1-t)f(y)$ for $0 \leq t \leq 1$; its graph lies above its chords.

\term{Concentration of Measure}
In high-dimensional spaces, most of the volume (or measure) concentrates near the boundary or equator. For the unit sphere $S^{n-1} \subset \mathbb{R}^n$, the measure of an $\epsilon$-neighbourhood of the equator $\{x : |x_1| \leq \epsilon\}$ approaches 1 as $n \to \infty$. Lévy's lemma quantifies this: Lipschitz functions on $S^n$ are nearly constant. Concentration of measure is foundational in asymptotic geometric analysis and probability.

\term{Condition Number}
A measure of the sensitivity of the output of a function to small changes in the input. For a linear system $Ax = b$, the condition number of $A$ is $\kappa(A) = \|A\|\|A^{-1}\|$; a large condition number indicates an "ill-conditioned" problem where small perturbations in $b$ can lead to large errors in $x$.

\term{Conditional Expectation}
The expectation of a random variable given information about another, $E[X \mid \mathcal{F}]$, itself a random variable measurable with respect to the conditioning $\sigma$-algebra.

\term{Conditional Probability}
The conditional probability of $A$ given $B$ is $P(A|B) = \frac{P(A \cap B)}{P(B)}$ for $P(B) > 0$. Bayes' theorem relates conditional probabilities: $P(A|B) = \frac{P(B|A)P(A)}{P(B)}$. Conditional expectation $E[X|Y]$ is a random variable measurable with respect to $\sigma(Y)$ satisfying $\int_A X\,dP = \int_A E[X|Y]\,dP$ for all $A \in \sigma(Y)$.

\term{Conditioning}
The sensitivity of the output of a numerical problem to small changes in its input, measured by the condition number.

\term{Condorcet Paradox}
Even when each voter's preferences are transitive, majority voting can produce a cycle: A beats B, B beats C, C beats A. It shows majority rule need not yield a consistent social ordering and motivates Condorcet's impossibility results.

\term{Cone}
A solid with a circular base and a single vertex, or in convex geometry the set of all non-negative scalar multiples of a given set.

\term{Conformal Map}
A holomorphic function $f: D \to \mathbb{C}$ is conformal at $z_0$ if $f'(z_0) \neq 0$. Conformal maps preserve angles (in magnitude and orientation) between curves. By the Riemann Mapping Theorem, every simply connected proper open subset of $\mathbb{C}$ is conformally equivalent to the unit disk. Conformal maps are used in fluid dynamics, electrostatics, and cartography.

\term{Conjugate}
The complex conjugate $\bar{z}$ of $z = x + iy$ is $x - iy$. Conjugate elements of a group are related by conjugation $g \mapsto aga^{-1}$.

\term{Connected Space}
A topological space is connected if it cannot be written as the union of two non-empty disjoint open sets. A space is path-connected if any two points can be joined by a continuous path. Path-connected implies connected, but not conversely (the topologist's sine curve is connected but not path-connected). The interval $[0,1]$ is connected; the Cantor set is totally disconnected.

\term{Connectedness}
A topological space is connected if it cannot be written as the union of two disjoint non-empty open sets.

\term{Connection}
A structure on a manifold or vector bundle governing parallel transport and defining covariant differentiation.

\term{Consistency}
The property of a numerical scheme that its truncation error tends to zero as the mesh is refined.

\term{Contact Manifold}
A contact manifold is an odd-dimensional smooth manifold $M^{2n+1}$ with a contact structure: a maximally non-integrable hyperplane field locally the kernel of a 1-form $\alpha$ with $\alpha \wedge (d\alpha)^n$ nowhere vanishing. By Darboux's theorem every contact structure is locally equivalent to $\alpha = dz - \sum_{i=1}^n y_i\,dx_i$ in suitable coordinates. The Reeb vector field determined by $\alpha$ generates a flow preserving $\alpha$, the Reeb flow. Contact geometry is the odd-dimensional counterpart of symplectic geometry, with applications in Hamiltonian dynamics, low-dimensional topology, and geometric quantisation.

\term{Continuity}
See \term{Continuous Function}.

\term{Continuous Function}
A function $f: X \to Y$ between topological spaces is continuous if the preimage of every open set in $Y$ is open in $X$. Equivalently, for every $x \in X$ and every neighbourhood $V$ of $f(x)$, there exists a neighbourhood $U$ of $x$ with $f(U) \subseteq V$. In metric spaces, continuity is equivalent to the $\epsilon$--$\delta$ definition.

\term{Continuum}
A connected compact metric space; often used loosely for the real line as an uncountable continuum.

\term{Continuum Hypothesis}
The assertion that there is no set whose cardinality lies strictly between the countable and the continuum: no $X$ satisfies $|\mathbb{N}| < |X| < |\mathbb{R}|$. Gödel showed the hypothesis is consistent with ZFC by constructing the constructible universe $L$, and Cohen proved its negation consistent by the method of forcing. Hence the continuum hypothesis is independent of the ZFC axioms.

\term{Contour Integral}
The integral of a function along a curve in the complex plane, $\int_\gamma f(z)\,dz$.

\term{Contraction Mapping Principle}
Also called the Banach Fixed-Point Theorem. If $(X,d)$ is a complete metric space and $T: X \to X$ satisfies $d(Tx, Ty) \leq c\, d(x,y)$ for some $0 \leq c < 1$, then $T$ has a unique fixed point, and the iterates $x_{n+1} = Tx_n$ converge to it for any starting point $x_0$. This is the foundation for existence-uniqueness proofs in ODEs, integral equations, and numerical analysis.

\term{Contradiction}
A statement that is false under all interpretations. Proof by contradiction assumes the negation of a claim and derives an impossibility.

\term{Convergence in Distribution}
The convergence of a sequence of random variables such that their cumulative distribution functions converge at all continuity points; the weakest of the common modes of convergence.

\term{Convergence Almost Everywhere}
Also called almost sure convergence or convergence with probability one. A sequence of random variables $X_n$ converges almost everywhere to $X$ if $P(\{\omega : X_n(\omega) \to X(\omega)\}) = 1$. This is stronger than convergence in probability and implies it.

\term{Convergence in Probability}
The convergence of a sequence of random variables $X_n$ to $X$ such that $P(|X_n - X| > \epsilon) \to 0$ for every $\epsilon > 0$.

\term{Convergence (Mode of)}
Modes of convergence for sequences of random variables: (a) Almost sure: $P(\lim X_n = X) = 1$. (b) In probability: $P(|X_n - X| > \epsilon) \to 0$ for all $\epsilon > 0$. (c) In $L^p$: $E[|X_n - X|^p] \to 0$. (d) In distribution: $F_{X_n}(x) \to F_X(x)$ at continuity points. Implications: a.s. $\implies$ probability $\implies$ distribution; $L^p$ convergence $\implies$ probability convergence.

\term{Convergent Series}
An infinite series whose sequence of partial sums has a finite limit, called its sum.

\term{Convex Function}
A function $f: I \to \mathbb{R}$ on an interval $I$ is convex if $f(tx + (1-t)y) \leq tf(x) + (1-t)f(y)$ for all $x,y \in I$ and $t \in [0,1]$. Equivalently, the epigraph $\{(x,y) : y \geq f(x)\}$ is a convex set. A twice-differentiable function is convex iff $f''(x) \geq 0$. Jensen's inequality states $f(E[X]) \leq E[f(X)]$ for convex $f$.

\term{Convex Hull}
The convex hull of a set $S \subseteq \mathbb{R}^n$ is the smallest convex set containing $S$, equivalently the set of all convex combinations $\sum \lambda_i x_i$ with $x_i \in S$, $\lambda_i \geq 0$, $\sum \lambda_i = 1$. Carathéodory's theorem states that any point in the convex hull of $S \subset \mathbb{R}^n$ is a convex combination of at most $n+1$ points of $S$.

\term{Convex Optimization}
The minimization of a convex function over a convex set, for which every local optimum is global.

\term{Convex Polytope}
The convex hull of finitely many points in $\mathbb{R}^n$, equivalently a bounded intersection of finitely many closed half-spaces by the Minkowski--Weyl theorem. Its faces form a lattice: vertices, edges, and facets are the 0-, 1-, and $(n-1)$-dimensional faces, encoded by the $f$-vector. The polar, or dual, polytope reverses the face lattice. Euler's formula $V - E + F = 2$ holds in three dimensions. Polytopes are central to linear programming and combinatorial geometry.

\term{Convex Set}
A set $C$ in a vector space is convex if for any $x, y \in C$, the line segment $[x,y] = \{tx + (1-t)y : 0 \leq t \leq 1\}$ lies entirely in $C$. Examples: balls, half-spaces, simplices, the positive orthant. The intersection of convex sets is convex. A function is convex iff its epigraph is a convex set. Convex sets are central to optimisation and functional analysis.

\term{Coordinate Chart}
A coordinate chart on a manifold $M$ is a pair $(U, \varphi)$ where $U \subset M$ is open and $\varphi: U \to \mathbb{R}^n$ is a homeomorphism onto an open subset of $\mathbb{R}^n$. An atlas is a collection of charts whose domains cover $M$. The transition maps $\varphi_\beta \circ \varphi_\alpha^{-1}$ between overlapping charts must be smooth for a smooth manifold.

\term{Copula}
A copula is a multivariate distribution function whose one-dimensional margins are uniform on $[0,1]$. By Sklar's theorem, every joint distribution function $F$ with margins $F_1, \ldots, F_n$ admits the decomposition $F(x_1, \ldots, x_n) = C(F_1(x_1), \ldots, F_n(x_n))$ with $C$ unique when the margins are continuous, so the copula captures dependence independently of the marginal laws. Principal families include the Archimedean copulas (Clayton, Frank, Gumbel) and the Gaussian copula. Copulas are widely used in finance, risk management, and dependence modelling.

\term{Corollary}
A proposition that follows almost immediately from a previously proved theorem.

\term{Correlation Coefficient}
The Pearson correlation coefficient between random variables $X$ and $Y$ is $\rho_{XY} = \frac{\operatorname{Cov}(X,Y)}{\sigma_X \sigma_Y}$, measuring linear dependence with values in $[-1,1]$. $\rho = 1$ means perfect positive linear relationship; $\rho = -1$ means perfect negative; $\rho = 0$ means no linear relationship. Note: zero correlation does not imply independence.

\term{Correspondence Theorem}
For a group homomorphism $f: G \to H$, the lattice of subgroups of $G$ containing $\ker f$ corresponds bijectively, under $K \mapsto f(K)$, to the lattice of subgroups of the image $f(G)$, and this correspondence preserves inclusions, normality, and quotients. The same statement holds for rings and modules, where ideals or submodules containing the kernel correspond to those of the image. It reduces the study of a quotient $G/N$ to the subgroups of $G$ that contain $N$.

\term{Coset}
If $H$ is a subgroup of $G$ and $g \in G$, the left coset is $gH = \{gh : h \in H\}$ and the right coset is $Hg = \{hg : h \in H\}$. Left cosets partition $G$ into disjoint sets of equal size $|H|$. Lagrange's theorem states $[G:H] = |G|/|H|$. A subgroup is normal iff left and right cosets coincide.

\term{Cotangent Bundle}
The cotangent bundle $T^*M$ of a manifold $M$ is the disjoint union of cotangent spaces $T^*_p M = (T_p M)^*$ over all points $p \in M$. It is a $2n$-dimensional manifold when $\dim M = n$. The cotangent bundle carries a canonical symplectic form $d\theta$ where $\theta$ is the tautological 1-form. It is the natural phase space in classical mechanics.

\term{Cotangent Complex}
The cotangent complex $L_{B/A}$ is a complex of $B$-modules governing the deformation theory of a morphism of algebras (or schemes) $A \to B$; it is a derived replacement of the module of Kähler differentials $\Omega_{B/A} = H^0(L_{B/A})$. Its cohomology modules $H^i(L_{B/A})$ govern derivations and extensions: obstructions to lifting $B$ over square-zero extensions of $A$ live in $H^2(L_{B/A})$, which also computes André--Quillen homology. For smooth morphisms $L_{B/A}$ is concentrated in degree zero, and under étale localisation it behaves like classical differentials. It was constructed by Quillen and developed by Illusie.

\term{Cotangent Space}
The dual space of the tangent space, whose elements are covectors and which carries the differentials of functions.

\term{Countable Set}
A set that is in bijection with a subset of the natural numbers. Finite sets and $\mathbb{N}$ itself are countable; $\mathbb{R}$ is not.

\term{Counterexample}
An example that disproves a conjecture or proposition.

\term{Covariance}
The measure $\operatorname{Cov}(X,Y) = E[(X - EX)(Y - EY)]$ of the joint variability of two random variables.

\term{Covariant Derivative}
A derivative operator on a manifold that is linear, satisfies the Leibniz rule, and is compatible with the connection.

\term{Covering Space}
A space $\tilde{X}$ with a surjective local homeomorphism $p: \tilde{X} \to X$ over which every point has an evenly covered neighbourhood.

\term{Cramer's Rule}
For a linear system $Ax = b$ with $\det(A) \neq 0$, the unique solution is $x_i = \frac{\det(A_i)}{\det(A)}$, where $A_i$ is $A$ with column $i$ replaced by $b$. Cramer's rule is theoretically important but numerically impractical for large systems due to the cost of computing determinants. It is useful for symbolic computation and proving theoretical results.

\term{Critical Point}
A point where the gradient of a differentiable function vanishes, or where its derivative does not exist.

\term{Crofton Formula}
The Crofton formula relates the length of a curve to the measure of lines intersecting it. In the plane, the length of a rectifiable curve $\gamma$ is $\frac{1}{2} \iint n(\gamma \cap \ell)\, d\mu(\ell)$, where $n(\gamma \cap \ell)$ counts intersection points and $d\mu$ is the invariant measure on lines. This is a special case of the kinematic formula in integral geometry.

\term{Cross Product}
The vector product $a \times b$ of two vectors in $\mathbb{R}^3$, perpendicular to both with magnitude $|a||b|\sin\theta$.

\term{Cross Ratio}
For four distinct points $z_1, z_2, z_3, z_4$ on the Riemann sphere, the cross ratio is $(z_1, z_2; z_3, z_4) = \frac{(z_1-z_3)(z_2-z_4)}{(z_1-z_4)(z_2-z_3)}$. It is invariant under Möbius transformations. Four points are concyclic or collinear iff their cross ratio is real. The cross ratio characterises projective equivalence of quadruples of points.

\term{Cubic Reciprocity}
Cubic reciprocity is the reciprocity law governing solutions of $x^3 \equiv a \pmod{p}$ for rational primes $p \equiv 1 \pmod{3}$. Working in the Eisenstein integers $\mathbb{Z}[\omega]$, where $\omega$ is a primitive cube root of unity, one defines the cubic residue symbol $(\alpha/\beta)_3$ on primary elements; the reciprocity law states $(\alpha/\beta)_3 = (\beta/\alpha)_3$ for primary $\alpha, \beta$. It is an analogue of quadratic reciprocity and was used by Eisenstein in his proof that $x^3 + y^3 = z^3$ has no non-trivial integer solutions. The theory is due to Eisenstein and Jacobi.

\term{Curry's Paradox}
In a system with self-reference and ordinary logical rules, the sentence "If this sentence is true, then X" forces X to be true for any proposition X. It shows that naive self-reference combined with implication (unlike the Liar) can prove arbitrary statements, and plays a role in discussions of paraconsistent logic.

\term{Curvature}
The rate at which a curve or surface deviates from being straight or flat, measured by the radius of the osculating circle.

\term{Curvature Tensor}
The Riemann curvature tensor $R^a_{\ bcd}$ measuring the failure of parallel transport to commute, from which Ricci and scalar curvature are derived.

\term{CW Complex}
A CW complex is a space built inductively by attaching cells $e^n_\alpha \cong D^n$ along their boundary spheres to the $(n-1)$-skeleton $X^{n-1}$, beginning with a discrete 0-skeleton, and equipped with the weak topology and the closure-finiteness condition. CW complexes form the natural combinatorial category for homotopy theory: by the Whitehead theorem a weak equivalence between CW complexes is a homotopy equivalence, and cellular approximation models maps between them. Cellular homology computes $H_*(X)$ from the attaching maps, and every space is weakly equivalent to a CW complex.

\term{Cycle}
A closed path in a graph with no repeated vertices, or an element of the quotient group $Z_n/B_n$ in homology.

\term{Cyclotomic Field}
The $n$-th cyclotomic field is $\mathbb{Q}(\zeta_n)$ where $\zeta_n = e^{2\pi i/n}$ is a primitive $n$-th root of unity. It is a Galois extension of $\mathbb{Q}$ with Galois group $(\mathbb{Z}/n\mathbb{Z})^\times$, abelian of order $\varphi(n)$. Cyclotomic fields are central to algebraic number theory; Kummer used them to study Fermat's Last Theorem via cyclotomic units and ideal class groups.

\term{Cyclotomic Polynomial}
The $n$-th cyclotomic polynomial $\Phi_n(x)$ is the minimal polynomial of a primitive $n$-th root of unity over $\mathbb{Q}$. It satisfies $x^n - 1 = \prod_{d|n} \Phi_d(x)$. The degree is $\Phi_n(x) = \prod_{1 \leq k \leq n, \gcd(k,n)=1} (x - e^{2\pi i k/n})$, with degree $\varphi(n)$. For prime $p$, $\Phi_p(x) = x^{p-1} + x^{p-2} + \cdots + x + 1$.

\term{Cylinder (Topology)}
The topological cylinder is $S^1 \times \mathbb{R}$, homeomorphic to $\mathbb{R}^2 \setminus \{0\}$. More generally, the cylinder over a space $X$ is $X \times [0,1]$. The cylinder is used to define homotopy: a homotopy between $f$ and $g$ is a continuous map $H: X \times [0,1] \to Y$ with $H(x,0) = f(x)$ and $H(x,1) = g(x)$.

\term{Cyclic Group}
A cyclic group is a group generated by a single element: $G = \langle g \rangle = \{g^n : n \in \mathbb{Z}\}$. Every cyclic group is abelian. A finite cyclic group of order $n$ is isomorphic to $\mathbb{Z}/n\mathbb{Z}$. The subgroups of $\mathbb{Z}/n\mathbb{Z}$ are precisely $\langle d \rangle$ for each divisor $d$ of $n$. Cyclic groups are the simplest infinite family of groups.

\term{Cyclic Homology}
Cyclic homology $HC_*(A)$ is a homology theory for associative algebras, introduced by Connes (1985). It is related to Hochschild homology via the Connes exact sequence: $\cdots \to HC_{n-2}(A) \to HH_n(A) \to HC_n(A) \to HC_{n-1}(A) \to \cdots$. Cyclic homology computes invariants of non-commutative spaces and appears in index theory and deformation quantisation.

\term{Cylinder Set}
In measure theory on product spaces, a cylinder set is a set of the form $\{x \in X^I : (x_{i_1}, \ldots, x_{i_n}) \in B\}$ where $B$ is measurable in the finite product. Cylinder sets generate the product $\sigma$-algebra. Kolmogorov's extension theorem uses cylinder sets to construct measures on infinite product spaces from consistent finite-dimensional distributions.
\term{Adjacency Matrix}
A square matrix $A = (a_{ij})$ representing a graph with $n$ vertices, where $a_{ij} = 1$ if vertices $i$ and $j$ are adjacent and $0$ otherwise. For weighted graphs, $a_{ij}$ is the weight of the edge. The eigenvalues of the adjacency matrix form the spectrum of the graph and determine many combinatorial properties.
\term{Categorical Product}
The product $A \times B$ of two objects in a category, equipped with projection morphisms to $A$ and $B$, satisfying the universal property that for any object $X$ with morphisms to $A$ and $B$, there is a unique morphism from $X$ to $A \times B$ making the diagram commute. In the category of sets, it is the Cartesian product; in the category of groups, it is the direct product.
\term{Clique}
A subset of vertices in a graph such that every two distinct vertices are adjacent; equivalently, a complete subgraph. The clique number $\omega(G)$ is the maximum size of a clique in $G$. Finding maximum cliques is NP-hard, and cliques appear in extremal graph theory and in the study of perfect graphs.
\term{Connectivity}
A measure of how connected a graph is. The vertex connectivity $\kappa(G)$ is the minimum number of vertices whose removal disconnects the graph or leaves a trivial graph. The edge connectivity $\kappa'(G)$ is defined analogously for edges. A graph is $k$-connected if its connectivity is at least $k$. By Menger's theorem, $k$-connectivity is equivalent to the existence of $k$ internally vertex-disjoint paths between every pair of vertices.
\term{Covariant Functor}
A functor $F: \mathcal{C} \to \mathcal{D}$ that preserves the direction of morphisms: $F(f \circ g) = F(f) \circ F(g)$ and $F(\operatorname{id}_X) = \operatorname{id}_{F(X)}$. The hom-functor $\operatorname{Hom}(X, -): \mathcal{C} \to \mathbf{Set}$ is a fundamental example. Covariant functors are the default notion of functor.
\term{Cut Vertex}
A vertex whose removal increases the number of connected components of a graph. Also called an articulation point. A graph with no cut vertices is 2-connected. Cut vertices play a key role in the decomposition of graphs into blocks.
\term{Generating Function}
A formal power series $\sum_{n=0}^\infty a_n x^n$ whose coefficients encode a sequence $a_n$. Generating functions transform combinatorial problems into analytic ones. Ordinary generating functions count objects by size; exponential generating functions $\sum a_n x^n/n!$ are natural for labelled structures. The method is due to Euler and de Moivre and is central to enumerative combinatorics.
\term{Chromatic Number}
The smallest number of colours needed to colour the vertices of a graph so that no adjacent vertices share a colour. It is denoted $\chi(G)$. A graph is $k$-chromatic if $\chi(G) = k$. By Brooks' theorem, $\chi(G) \leq \Delta(G)$ for connected graphs other than complete graphs and odd cycles, where $\Delta$ is the maximum degree.

\chapter*{D}
\label{chap:d}

\term{Darboux's Theorem}
If $f$ is differentiable on $[a,b]$, then $f'$ has the intermediate value property: for any $y$ between $f'(a)$ and $f'(b)$, there exists $c \in (a,b)$ with $f'(c) = y$. Unlike general functions, derivatives need not be continuous, but they cannot have jump discontinuities. This theorem is fundamental in real analysis and demonstrates that derivatives, while not necessarily continuous, still satisfy a mean value property.

\term{Darboux Sum}
An approximation of an integral using step functions that bound the function from above or below; their common refinement is used in Riemann integration.

\term{Decidable}
A property or problem for which an algorithm always terminates with the correct yes/no answer.

\term{Decision Problem}
A problem whose answer is yes or no, studied in computability and complexity theory.

\term{Decomposable}
Capable of being split into simpler components, as a vector or module written as a direct sum of non-zero pieces.

\term{Dedekind Cut}
A Dedekind cut partitions the rationals $\mathbb{Q}$ into two non-empty sets $(A, B)$ such that every element of $A$ is less than every element of $B$, and $A$ has no greatest element. Each cut defines a real number. This construction provides a rigorous foundation for the real numbers, establishing completeness without appealing to geometric intuition.

\term{Dedekind Domain}
An integral domain in which every nonzero proper ideal factors uniquely into prime ideals. Equivalently, a Noetherian integrally closed domain of Krull dimension 1. The ring of integers $\mathcal{O}_K$ of a number field $K$ is always a Dedekind domain. Unique factorisation of elements may fail, but unique factorisation of ideals holds.

\term{Dedekind Eta Function}
The Dedekind eta function is $\eta(\tau) = e^{\pi i \tau/12} \prod_{n=1}^{\infty}(1 - e^{2\pi i n \tau})$ for $\tau$ in the upper half-plane. It is a modular form of weight $1/2$ satisfying $\eta(-1/\tau) = \sqrt{\tau/i}\,\eta(\tau)$. The eta function appears in partition theory, string theory, and the theory of vertex operator algebras.

\term{Dedekind's Axiom}
The completeness property of the real numbers: every non-empty bounded-above subset has a least upper bound.

\term{Definite Integral}
The quantity $\int_a^b f(x)\,dx$ measuring the signed area under $f$ between $a$ and $b$.

\term{Deformation Retract}
A subspace $A$ of $X$ is a deformation retract if there exists a homotopy $H: X \times [0,1] \to X$ with $H(x,0) = x$, $H(x,1) \in A$, and $H(a,t) = a$ for all $a \in A$, $t \in [0,1]$. The circle $S^1$ is a deformation retract of the punctured plane $\mathbb{R}^2 \setminus \{0\}$ via radial projection. Homotopy equivalent spaces share the same homotopy groups.

\term{Delaunay Triangulation}
Given a finite set of points in the plane, the Delaunay triangulation is a triangulation such that no point lies inside the circumcircle of any triangle. Equivalently, the Voronoi cells form a dual graph. Delaunay triangulations maximise the minimum angle, avoiding skinny triangles, making them ideal for finite element methods and surface reconstruction.

\term{Delta Method}
In statistics, the delta method approximates the distribution of a function of an asymptotically normal random variable. If $\sqrt{n}(X_n - \theta) \xrightarrow{d} N(0, \sigma^2)$ and $g$ is differentiable at $\theta$, then $\sqrt{n}(g(X_n) - g(\theta)) \xrightarrow{d} N(0, [g'(\theta)]^2\sigma^2)$. It is widely used for deriving standard errors of transformed estimators.

\term{Delta System Lemma}
Also called the Sunflower Lemma. Every uncountable family of finite sets contains an uncountable delta system: a subfamily where pairwise intersections are all equal to a fixed set (the kernel). This combinatorial principle, due to Erdős and Rado, is used in set theory, topology, and combinatorics. A delta system of size $\kappa$ with kernel $K$ consists of sets $\{A_i\}_{i \in \kappa}$ with $A_i \cap A_j = K$ for $i \neq j$.

\term{De Moivre's Formula}
The identity $(\cos \theta + i \sin \theta)^n = \cos(n\theta) + i\sin(n\theta)$.

\term{Dense Set}
A subset $D$ of a topological space $X$ is dense if its closure is $X$: every non-empty open set intersects $D$. The rationals $\mathbb{Q}$ are dense in $\mathbb{R}$; the polynomials are dense in $C[a,b]$ by the Weierstrass approximation theorem. A separable space is one containing a countable dense subset.

\term{Derivation (Algebra)}
A derivation on a commutative algebra $A$ over a ring $R$ is an $R$-linear map $d: A \to A$ satisfying the Leibniz rule: $d(ab) = d(a)b + ad(b)$. The set of all derivations forms a Lie algebra under the bracket $[d_1, d_2] = d_1 \circ d_2 - d_2 \circ d_1$. Vector fields on manifolds are derivations on the algebra of smooth functions.

\term{Derivation (Differential Geometry)}
A derivation on the algebra $C^\infty(M)$ of smooth functions on a manifold $M$ is a linear map $D$ satisfying $D(fg) = D(f)g + fD(g)$. The space of derivations is naturally identified with the tangent bundle $TM$. At each point $p$, derivations correspond to tangent vectors, giving an intrinsic definition of the tangent space.

\term{Derivative}
The derivative $f'(x_0) = \lim_{h \to 0} \frac{f(x_0+h) - f(x_0)}{h}$ measures the instantaneous rate of change of $f$ at $x_0$. For functions of several variables, the derivative is the linear map represented by the Jacobian matrix. The derivative is fundamental to calculus, appearing in the Fundamental Theorem, Taylor's theorem, and the study of differential equations.

\term{Derived Functor}
The functor $R^n F$ obtained by applying a left-exact functor to a resolution, measuring the failure of exactness.

\term{Desargues' Theorem}
Two triangles are in perspective from a point if and only if they are in perspective from a line. Given triangles $ABC$ and $A'B'C'$, if the lines $AA'$, $BB'$, $CC'$ are concurrent, then the intersections of corresponding sides $AB \cap A'B'$, $BC \cap B'C'$, $CA \cap C'A'$ are collinear. This fundamental theorem of projective geometry holds in any projective space of dimension at least 3.

\term{Descartes' Rule of Signs}
For a polynomial $p(x)$ with real coefficients, the number of positive real roots (counting multiplicity) is at most the number of sign changes in the sequence of coefficients, and differs from it by an even number. Applying the rule to $p(-x)$ gives information about negative roots. This provides a bound on the number of real zeros without finding them explicitly.

\term{Descent (Algebraic Geometry)}
Descent theory studies when objects defined locally on a cover glue to form a global object. For a morphism $f: Y \to X$, a descent datum on a sheaf or bundle over $Y$ consists of isomorphisms satisfying a cocycle condition on $Y \times_X Y$. Descent is effective for faithfully flat morphisms (Grothendieck's theorem), forming the basis of étale cohomology.

\term{Determinant}
The determinant $\det(A)$ of a square matrix $A$ is a scalar value characterising the linear transformation represented by $A$. It equals the product of eigenvalues, measures the volume scaling factor, and determines invertibility: $A$ is invertible iff $\det(A) \neq 0$. Key properties: $\det(AB) = \det(A)\det(B)$, $\det(A^T) = \det(A)$, $\det(A^{-1}) = 1/\det(A)$.

\term{Deterministic}
A process whose outcome is fully determined by its initial state, with no randomness involved.

\term{Deterministic Algorithm}
An algorithm is deterministic if its behaviour is completely determined by its input, with no randomness. Every step produces a unique output for a given input. This contrasts with probabilistic algorithms, which use random bits. Deterministic algorithms are analysed for worst-case and average-case complexity; many problems admit efficient deterministic solutions (P) while others may require randomness.

\term{Diagonal Argument}
Cantor's diagonal argument proves that $\mathbb{R}$ is uncountable. Suppose $r_1, r_2, \ldots$ lists all reals in $[0,1]$. Construct $x$ by changing the $n$-th decimal digit of $r_n$. Then $x$ differs from every $r_n$, contradicting the assumption. The same argument shows $|\mathcal{P}(\mathbb{N})| > |\mathbb{N}|$ and appears in Turing's proof of the undecidability of the halting problem.

\term{Diagonalisation (Logic)}
Diagonalisation in computability theory constructs an object that differs from every member of a given list. Turing's diagonal argument shows the halting problem is undecidable: assume a machine $H$ decides halting; construct $D$ that does the opposite of $H$ on its own description. $D$ cannot be decided by $H$, contradiction. This technique underlies Gödel's incompleteness theorems.

\term{Diagonal Matrix}
A square matrix whose only non-zero entries lie on the main diagonal.

\term{Diagonalizability}
The property of a linear operator of having a basis of eigenvectors, so that it is similar to a diagonal matrix.

\term{Difference Equation}
A recurrence relation expressing a sequence term as a function of earlier terms, the discrete analogue of a differential equation.

\term{Differential}
The linear part $\mathrm{d}f$ of the change of a differentiable function, mapping tangent vectors to tangent vectors.

\term{Differential Equation}
An equation relating an unknown function to its derivatives, such as $y'' + y = 0$.

\term{Differential Form}
An antisymmetric tensor field $\omega$ of type $(0,k)$, integrated over oriented submanifolds.

\term{Differential Geometry}
The study of smooth manifolds and the calculus, metrics, and curvatures defined on them.

\term{Differential Topology}
The study of differentiable manifolds and smooth maps between them, focusing on properties invariant under diffeomorphism.

\term{Differentiation}
The process of computing the derivative, measuring the instantaneous rate of change of a function.

\term{Dilworth's Theorem}
In a finite partially ordered set, the minimum number of chains needed to cover the elements equals the size of the largest antichain. For a poset $P$, a chain is a totally ordered subset and an antichain is a set of pairwise incomparable elements. The dual form (Mirsky's theorem) states that the minimum number of antichains needed to cover $P$ equals the size of the largest chain. Dilworth's theorem applies to scheduling problems where tasks are subject to precedence constraints.

\term{Digraph}
A directed graph, whose edges have an orientation from one vertex to another.

\term{Dihedral Group}
The dihedral group $D_n$ is the symmetry group of a regular $n$-gon, of order $2n$. It is generated by a rotation $r$ of order $n$ and a reflection $s$ of order 2, with relation $srs = r^{-1}$. For $n = 3$, $D_3 \cong S_3$. Dihedral groups are the simplest non-abelian groups (for $n \geq 3$) and appear in geometry, chemistry (molecular symmetry), and group theory.

\term{Dimension}
The number of coordinates needed to describe a space; for a vector space, the size of a basis.

\term{Directed Graph}
A directed graph (digraph) $D = (V, A)$ consists of vertices $V$ and ordered pairs of vertices called arcs or directed edges. The out-degree $d^+(v)$ counts arcs leaving $v$; the in-degree $d^-(v)$ counts arcs entering $v$. A directed graph is strongly connected if there is a directed path between any ordered pair of vertices. Digraphs model networks, precedence constraints, and state machines.

\term{Direct Sum}
The coproduct construction for abelian groups or vector spaces, $V \oplus W$, whose elements are pairs $(v, w)$ with componentwise operations.

\term{Dirac Delta Function}
The Dirac delta $\delta(x)$ is a generalised function with $\delta(x) = 0$ for $x \neq 0$ and $\int_{-\infty}^{\infty} \delta(x)\,dx = 1$. Formally, $\delta(x) = \lim_{\epsilon \to 0} \frac{1}{\epsilon}\phi(x/\epsilon)$ for any nascent delta function $\phi$. It satisfies $\int f(x)\delta(x-a)\,dx = f(a)$. The delta function is the identity element for convolution and the derivative of the Heaviside step function.

\term{Dirac Measure}
The probability measure $\delta_a$ concentrated at a single point $a$, with $\delta_a(\{a\}) = 1$.

\term{Dirichlet Boundary Condition}
A Dirichlet boundary condition specifies the value of a solution on the boundary of a domain: $u(x) = g(x)$ for $x \in \partial\Omega$. For the Laplace equation $\Delta u = 0$ on $\Omega$, the Dirichlet problem has a unique solution if $g$ is continuous (Perron's method). Dirichlet conditions model fixed temperature on a boundary in heat flow, or fixed displacement in elasticity.

\term{Dirichlet L-function}
The series $L(s, \chi) = \sum_{n=1}^\infty \chi(n)/n^s$ built from a Dirichlet character $\chi$, central to the theory of primes in arithmetic progressions.

\term{Dirichlet's Approximation Theorem}
For any real number $\alpha$ and any positive integer $N$, there exist integers $p, q$ with $1 \leq q \leq N$ such that $|\alpha - p/q| < 1/(qN) \leq 1/N^2$. This implies infinitely many rational approximations satisfying $|\alpha - p/q| < 1/q^2$. The theorem follows from the pigeonhole principle applied to fractional parts $\{\alpha\}, \{2\alpha\}, \ldots, \{N\alpha\}$.

\term{Dirichlet's Theorem on Arithmetic Progressions}
For any coprime integers $a, d$, the arithmetic progression $a, a+d, a+2d, \ldots$ contains infinitely many primes. Dirichlet proved this (1837) using $L$-functions: $L(s, \chi) = \sum_{n=1}^{\infty} \chi(n)n^{-s}$ for Dirichlet characters $\chi$ modulo $d$, showing $L(1, \chi) \neq 0$. The proof introduced analytic methods into number theory.

\term{Dirichlet's Unit Theorem}
In a number field $K$ with $r_1$ real embeddings and $2r_2$ complex embeddings, the unit group $\mathcal{O}_K^*$ of the ring of integers is isomorphic to $\mu_K \times \mathbb{Z}^{r_1 + r_2 - 1}$, where $\mu_K$ is the finite cyclic group of roots of unity in $K$. The rank $r_1 + r_2 - 1$ is the unit rank. For $\mathbb{Q}(\sqrt{2})$, the fundamental unit is $1 + \sqrt{2}$.

\term{Discontinuity}
A point at which a function is not continuous, classified as removable, jump, or essential.

\term{Discrete Topology}
The discrete topology on a set $X$ declares every subset open. Every function from a discrete space is continuous. Every subset is also closed, so the space is zero-dimensional and totally disconnected. The discrete topology is the finest possible topology; it makes $X$ a metric space with the discrete metric $d(x,y) = 1$ for $x \neq y$.

\term{Discriminant (Number Field)}
The discriminant $d_K$ of a number field $K$ is the determinant of the trace form: $d_K = \det(\operatorname{Tr}_{K/\mathbb{Q}}(\omega_i \omega_j))$ for an integral basis $\{\omega_i\}$. It measures the arithmetic complexity of $K$. The discriminant of $\mathbb{Q}(\sqrt{d})$ for square-free $d$ is $d$ if $d \equiv 1 \pmod 4$ and $4d$ otherwise. Only finitely many fields have discriminant below any given bound.

\term{Discriminant (Polynomial)}
The discriminant of a polynomial $p(x) = a_n \prod(x - \alpha_i)$ is $\Delta = a_n^{2n-2} \prod_{i < j}(\alpha_i - \alpha_j)^2$. It vanishes iff $p$ has a repeated root. For a quadratic $ax^2 + bx + c$, $\Delta = b^2 - 4ac$. For a cubic $x^3 + px + q$, $\Delta = -4p^3 - 27q^2$. The discriminant measures the "separation" of roots and appears in Galois theory.

\term{Disk (Topology)}
The closed disk $D^n = \{x \in \mathbb{R}^n : \|x\| \leq 1\}$ is the standard model for an $n$-dimensional ball. $D^2$ is the unit disk in the plane, homeomorphic to any compact convex set with non-empty interior. The disk is contractible: the identity map is homotopic to a constant map. By the Brouwer fixed-point theorem, every continuous map $D^n \to D^n$ has a fixed point.

\term{Displacement Map}
A displacement map in numerical analysis refers to a matrix representation where a matrix $A$ satisfies $A - S A S^{-1} = uv^T$ for a shift matrix $S$ and vectors $u, v$. Displacement rank measures how far a matrix is from being Toeplitz or circulant. Structured matrices with low displacement rank admit fast algorithms for inversion and eigenvalue computation.

\term{Distribution}
(Probability) The rule assigning probabilities to the values or intervals of a random variable.

\term{Distribution Theory}
The theory of generalised functions, or distributions, constructed by duality with spaces of test functions.

\term{Distributional Solution}
A solution of a PDE in the sense of distributions, tested against smooth compactly supported functions.

\term{Divergence}
The divergence $\operatorname{div} \mathbf{F} = \nabla \cdot \mathbf{F} = \sum \frac{\partial F_i}{\partial x_i}$ measures the net outflow of a vector field per unit volume at a point. The divergence theorem (Gauss's theorem) relates the flux through a closed surface to the integral of divergence over the enclosed volume: $\int_V \nabla \cdot \mathbf{F}\,dV = \oint_{\partial V} \mathbf{F} \cdot d\mathbf{S}$.

\term{Divergence Theorem}
Also called Gauss's theorem or Ostrogradsky's theorem. For a smoothly bounded region $V \subset \mathbb{R}^n$ and a $C^1$ vector field $\mathbf{F}$ on $V$: $\int_V \nabla \cdot \mathbf{F}\,dV = \oint_{\partial V} \mathbf{F} \cdot \mathbf{n}\,dS$, where $\mathbf{n}$ is the outward unit normal. This generalises the fundamental theorem of calculus to higher dimensions and is fundamental to physics (electromagnetism, fluid dynamics).

\term{Divergent Series}
A series whose partial sums do not converge to a finite limit.

\term{Divisible Group}
An abelian group $G$ is divisible if for every $g \in G$ and positive integer $n$, there exists $h \in G$ with $nh = g$. Examples: $\mathbb{Q}/\mathbb{Z}$, $\mathbb{Q}$, and the circle group $\mathbb{R}/\mathbb{Z}$. Divisible groups are injective objects in the category of abelian groups: every abelian group embeds in a divisible group. The classification of divisible groups is simple: they are direct sums of copies of $\mathbb{Q}$ and $\mathbb{Z}(p^\infty)$.

\term{Division Algorithm}
For integers $a$ and $b > 0$, there exist unique integers $q$ (quotient) and $r$ (remainder) such that $a = bq + r$ with $0 \leq r < b$. This fundamental result underlies the Euclidean algorithm for computing GCDs. The division algorithm generalises to polynomials over a field: for $f, g \neq 0$, there exist unique $q, r$ with $\deg(r) < \deg(g)$ and $f = gq + r$.

\term{Dixon's Theorem (Hypergeometric Series)}
Dixon's theorem evaluates the well-poised hypergeometric series $_3F_2(\alpha, \beta, \gamma; 1+\alpha-\beta, 1+\alpha-\gamma; 1)$ when $\operatorname{Re}(\alpha + 1 - \beta - \gamma) > 0$: $_3F_2 = \frac{\Gamma(1+\alpha/2)\Gamma(1+\alpha-\beta)\Gamma(1+\alpha-\gamma)\Gamma(1+\alpha/2-\beta-\gamma)}{\Gamma(1+\alpha)\Gamma(1+\alpha/2-\beta)\Gamma(1+\alpha/2-\gamma)\Gamma(1+\alpha-\beta-\gamma)}$. It is a cornerstone of hypergeometric identity theory.

\term{Doctrine of Natural Deduction}
Natural deduction is a proof calculus developed by Gerhard Gentzen (1934) and Jakob Jakobsson. It uses introduction and elimination rules for each logical connective, organising proofs in tree form. The deduction theorem states that $\Gamma, A \vdash B$ iff $\Gamma \vdash A \to B$. Natural deduction is the basis for proof assistants (Coq, Isabelle) and type theory.

\term{Dodecahedron}
The regular dodecahedron is one of the five Platonic solids, with 12 regular pentagonal faces, 20 vertices, and 30 edges. Its symmetry group is $A_5 \times \mathbb{Z}/2\mathbb{Z}$ of order 120. The dodecahedron is dual to the icosahedron. Its vertices can be placed at coordinates involving the golden ratio $\phi = (1+\sqrt{5})/2$, e.g., $(\pm 1, \pm 1, \pm 1)$ and $(0, \pm \phi, \pm 1/\phi)$ and permutations.

\term{Doob's Optional Stopping Theorem}
For a martingale $X_n$ and a stopping time $T$, under appropriate conditions we have $\mathbb{E}[X_T] = \mathbb{E}[X_0]$. The theorem holds when $T$ is bounded, or when $X$ is uniformly integrable and $T$ is almost surely finite. This formalises the fair-games interpretation of martingales: no strategy of stopping at a random time can improve expected winnings. The theorem is central in martingale theory, harmonic analysis, and mathematical finance.

\term{Domain}
The set of inputs on which a function is defined, or an integral domain when discussing rings.

\term{Dominated Convergence Theorem}
If $f_n \to f$ pointwise almost everywhere and $|f_n| \leq g$ for an integrable $g$, then $f$ is integrable and $\lim_{n \to \infty} \int f_n\,d\mu = \int f\,d\mu$. It is the standard interchange theorem of limit and integral in measure theory, used whenever one must pass a limit under the integral sign. The dominating function $g$ is essential: without it, e.g. $f_n = n\chi_{(0,1/n)}$ on $[0,1]$, the conclusion fails.

\term{Dominator Tree}
In a directed graph with start vertex $s$, vertex $d$ dominates vertex $v$ if every path from $s$ to $v$ passes through $d$. The dominator tree represents the dominance relation: the parent of $v$ is its immediate dominator. Applications include compiler optimisation (program dependence graphs), flow analysis, and finding critical points in control flow graphs.

\term{Donsker Invariance Principle}
For independent identically distributed random variables $X_1, X_2, \ldots$ with mean $0$ and variance $\sigma^2$, the scaled random walk $W_n(t) = \frac{1}{\sigma\sqrt{n}} \sum_{k=1}^{\lfloor nt \rfloor} X_k$, linearly interpolated to a continuous process on $[0,1]$, converges weakly to standard Brownian motion. This is the functional central limit theorem. Convergence holds in distribution on $C[0,1]$, so by the continuous mapping theorem the maximum $\sup_{0 \le t \le 1} W_n(t)$ and other continuous functionals converge to the corresponding functionals of Brownian motion.

\term{Dot Product}
The scalar $a \cdot b = \sum a_i b_i = |a||b|\cos\theta$ in Euclidean space.

\term{Douady--Earle Extension}
The Douady--Earle extension is a canonical extension of a continuous map $f: S^1 \to S^1$ to a diffeomorphism of the disk $D^2$, constructed using the centre of mass in the Poincaré disk model. It is equivariant under Möbius transformations and provides a smooth solution to the extension problem in Teichmüller theory. Used in the construction of the Euler class of the group of diffeomorphisms of the circle.

\term{Double Cover}
A double cover of a topological space $X$ is a 2-to-1 continuous map $p: Y \to X$ that is a covering map. The universal double cover of $\mathbb{R}P^n$ is $S^n$. The Möbius strip double covers the Klein bottle. Double covers correspond to index-2 subgroups of the fundamental group via the Galois correspondence for covering spaces.

\term{Duality}
A correspondence between two theories or objects that reverses structure, such as the duality between a vector space and its dual.

\term{Dual Norm}
The dual norm of a norm $\|\cdot\|$ on a vector space $V$ is $\|f\|_* = \sup\{f(v) : \|v\| \leq 1\}$ for $f \in V^*$. By the Hahn--Banach theorem, the dual norm is attained. The dual of the $\ell^p$ norm is the $\ell^q$ norm where $1/p + 1/q = 1$. For $p = 2$, the dual norm equals the original norm (self-dual). Dual norms appear in convex analysis and optimisation.

\term{Dual Operator}
The operator $T'$ on the dual space induced by $T$, with $(T'\phi)(x) = \phi(Tx)$.

\term{Dual Problem}
The optimization problem associated with the primal, whose optimum bounds the primal optimum under duality.

\term{Dual Space}
The dual space $V^*$ of a vector space $V$ is the space of all linear functionals $f: V \to \mathbb{F}$. For finite-dimensional $V$, $\dim V^* = \dim V$, and $V$ is canonically isomorphic to $V^{**}$ (the double dual). For infinite-dimensional spaces, the algebraic dual is much larger than $V$. The continuous dual $X'$ of a normed space $X$ is fundamental in functional analysis.

\term{Duplication Formula (Gamma Function)}
Legendre's duplication formula states $\Gamma(z)\Gamma(z + 1/2) = 2^{1-2z}\sqrt{\pi}\,\Gamma(2z)$. This generalises to the multiplication theorem: $\prod_{k=0}^{n-1} \Gamma(z + k/n) = (2\pi)^{(n-1)/2} n^{1/2-nz} \Gamma(nz)$. The duplication formula is essential in evaluating Gamma function values and appears in the theory of beta integrals and special functions.

\term{Dynamical System}
A space together with a rule, such as an iteration or differential equation, describing how points evolve over time.


\chapter*{E}
\label{chap:e}

\term{Eccentricity}
For a conic section, eccentricity measures how far the shape is from a circle: for an ellipse with semiaxes $a\ge b$, it is $e=\sqrt{1-b^2/a^2}$. A graph-theory eccentricity is different: it is the greatest shortest-path distance from one vertex to any other. \par\noindent\textbf{Applications:} conic eccentricity classifies ellipses, parabolas, and hyperbolas, while graph eccentricity identifies central and peripheral locations in networks.

\term{Edge}
An edge is a connection between two vertices in a graph; it may have a direction or a weight. \par\noindent\textbf{Applications:} edges represent roads, communication links, relationships, and transitions between states.

\term{Eigenvalue}
An eigenvalue $\lambda$ of a linear operator $T: V \to V$ on a vector space $V$ is a scalar such that $T v = \lambda v$ for some non-zero $v \in V$ (an eigenvector). The eigenvalues are the roots of the characteristic polynomial $\det(T - \lambda I) = 0$. The spectral theorem states that a symmetric matrix is orthogonally diagonalisable with real eigenvalues. Eigenvalues determine stability of dynamical systems and appear throughout physics and engineering.

\term{Eigenvalue Problem}
The eigenvalue problem seeks scalars $\lambda$ and nonzero vectors $v$ satisfying $Av=\lambda v$. \par\noindent\textbf{Applications:} eigenvalues identify natural frequencies, growth rates, stability, and important directions in data.

\term{Eigenvector}
An eigenvector of $T$ is a nonzero vector $v$ satisfying $Tv=\lambda v$; the transformation changes its size or direction only by the scalar $\lambda$. \par\noindent\textbf{Applications:} eigenvectors describe vibration modes, principal components, quantum states, and long-term behavior of dynamical systems.

\term{Eilenberg--Steenrod Axioms}
The Eilenberg--Steenrod axioms characterise ordinary homology and cohomology theories on the category of CW pairs: (1) homotopy invariance, (2) exactness of the long exact sequence of a pair, (3) excision, (4) additivity over disjoint unions, (5) normalisation on a point. Any two theories satisfying these axioms are naturally isomorphic. These axioms (1952) provided the first rigorous framework for homology.

\term{Eilenberg--MacLane Space}
For an abelian group $G$ and integer $n \geq 1$, the Eilenberg--MacLane space $K(G,n)$ is a CW complex whose homotopy groups vanish except $\pi_n(K(G,n)) \cong G$. It represents cohomology: for any CW complex $X$ there is a natural bijection $H^n(X;G) \cong [X, K(G,n)]$ between cohomology classes and homotopy classes of maps into $K(G,n)$. These spaces exist and are unique up to homotopy, and they play a central role in obstruction theory, where the obstructions to extending maps and sections live in cohomology with coefficients in the relevant homotopy groups, and in the theory of cohomology operations.

\term{Eisenstein Series}
The Eisenstein series $E_k(z) = \sum_{(c,d) \neq (0,0)} (cz+d)^{-k}$ on the upper half-plane converges absolutely for $k > 2$ and is a modular form of weight $k$ for $\mathrm{SL}_2(\mathbb{Z})$. Its Fourier expansion $E_k(z) = 2\zeta(k) + \frac{2(2\pi i)^k}{(k-1)!}\sum_{n \geq 1} \sigma_{k-1}(n) e^{2\pi i n z}$ is written in terms of the divisor sums $\sigma_{k-1}$. Together with cusp forms they generate the space of modular forms; the constant terms are algebraic multiples of $\zeta(k)$, and the series connect to the Weierstrass zeta function in the theory of elliptic functions.

\term{Elementary Function}
A function built from polynomials, exponentials, logarithms, trigonometric functions, and their inverses using finitely many operations and compositions. \par\noindent\textbf{Applications:} elementary functions form the standard vocabulary of calculus, modeling, engineering, and scientific computation.

\term{Elementary Symmetric Polynomials}
The elementary symmetric polynomials in $x_1, \ldots, x_n$ are $e_k = \sum_{1 \leq i_1 < \cdots < i_k \leq n} x_{i_1} \cdots x_{i_k}$ for $k=1,\ldots,n$. The Fundamental Theorem of Symmetric Functions states that every symmetric polynomial is a unique polynomial in the $e_k$. The coefficients of a monic polynomial are (up to sign) the elementary symmetric polynomials in its roots.

\term{Ellipse}
An ellipse is the set of points whose sum of distances to two fixed foci is constant. A circle is the special case in which the foci coincide. \par\noindent\textbf{Applications:} ellipses model planetary orbits, optical mirrors, acoustic reflectors, and uncertainty regions in statistics.

\term{Elliptic Curve}
A smooth curve of genus 1 with a chosen point at infinity. Over a field of characteristic different from $2$ and $3$, it can be written as $y^2=x^3+ax+b$ with $4a^3+27b^2\ne0$. Its points form an abelian group under a geometric addition rule. \par\noindent\textbf{Applications:} elliptic curves are used in public-key cryptography, integer factorization, and the study of rational solutions to equations.

\term{Elliptic Curve Cryptography}
Elliptic curve cryptography is a public-key cryptosystem whose security rests on the elliptic curve discrete logarithm problem: given points $P$ and $Q = [k]P$ on an elliptic curve over a finite field, recover $k$. For a general curve no sub-exponential algorithm is known, whereas the index calculus attacks finite fields, so elliptic curve systems achieve security comparable to RSA with much smaller keys: a 256-bit curve offers the security of a 3072-bit RSA modulus. Standardised constructions include key exchange (ECDH) and digital signatures (ECDSA), together with efficient curve models such as Montgomery and twisted Edwards curves that resist side-channel attacks.

\term{Elliptic Function}
A meromorphic function on the complex plane with two independent periods, such as the Weierstrass $\wp$ function. \par\noindent\textbf{Applications:} elliptic functions describe periodic phenomena in two directions and appear in complex analysis, mechanics, and the geometry of elliptic curves.

\term{Elliptic Integral}
An elliptic integral is an integral such as $\int R(x,\sqrt{P(x)})\,dx$, where $P$ is a cubic or quartic and $R$ is rational; it is not generally expressible using elementary functions. \par\noindent\textbf{Applications:} elliptic integrals calculate pendulum periods, arc lengths, planetary motion, and electromagnetic quantities.

\term{Elliptic PDE}
An elliptic PDE has a principal part that behaves like a positive-definite quadratic form; Laplace's equation is the standard example. Under suitable boundary conditions, elliptic equations smooth their solutions. \par\noindent\textbf{Applications:} they model steady temperature, electrostatics, elasticity, fluid potential, and image smoothing.

\term{Ellsberg Paradox}
People strongly prefer gambles with known probabilities over those with ambiguous or unknown probabilities, even when subjective expected utility would treat them the same. It is a classic violation of expected utility theory and motivates ambiguity aversion and the distinction between risk and uncertainty (Knightian uncertainty).

\term{Embedding Theorem (Nash)}
Nash's embedding theorem (1954, 1956) states that every Riemannian manifold can be isometrically embedded into some Euclidean space $\mathbb{R}^N$. The $C^1$ embedding requires $N \geq n+2$; the $C^\infty$ embedding requires $N = n(3n+11)/2$ for compact manifolds. This resolved a long-standing question: abstract Riemannian manifolds arise as submanifolds of Euclidean space.

\term{Endomorphism}
An endomorphism is a structure-preserving map from an object to itself. \par\noindent\textbf{Applications:} endomorphisms describe repeated transformations, linear systems, dynamical updates, and symmetries of algebraic structures.

\term{Enumeration}
Enumeration is the process of listing objects or assigning them natural-number labels. \par\noindent\textbf{Applications:} enumeration generates combinatorial objects, indexes data, tests algorithms, and distinguishes countable from uncountable sets.

\term{Entire Function}
An entire function is a holomorphic function defined on all of $\mathbb{C}$. Liouville's theorem states that a bounded entire function is constant, giving the fundamental theorem of algebra. Every entire function admits a Weierstrass factorization as a product over its zeros times an exponential factor; for functions of finite order, the Hadamard factorization theorem makes this canonical, and the growth is quantified by the order and type. Examples include polynomials, $\exp$, $\sin$, and $\cos$. Entire functions are central to complex analysis and value distribution theory (Nevanlinna theory), and their zeros control the gamma and zeta functions.

\term{Envelope}
An envelope is a curve tangent to each member of a family of curves, at least locally. \par\noindent\textbf{Applications:} envelopes describe wavefronts, caustics in optics, reachable boundaries, and families of solutions to differential equations.

\term{Epimorphism}
In category theory, an epimorphism is a morphism $f$ such that $g\circ f=h\circ f$ implies $g=h$. It need not be surjective in every category, although epimorphisms of sets are surjective. \par\noindent\textbf{Applications:} epimorphisms formalize quotient and “onto-like” constructions in algebra and geometry.

\term{Equality}
Equality means that two expressions denote the same object; it is reflexive, symmetric, and transitive. \par\noindent\textbf{Applications:} equality is the basic relation used in equations, symbolic computation, formal proof, and program verification.

\term{Equivalence Relation}
An equivalence relation on a set $X$ is a binary relation $\sim$ that is reflexive ($x \sim x$), symmetric ($x \sim y \implies y \sim x$), and transitive ($x \sim y \land y \sim z \implies x \sim z$). Equivalence relations partition $X$ into equivalence classes $[x] = \{y \in X : y \sim x\}$. The quotient set $X/{\sim}$ is the set of all equivalence classes. Examples: congruence modulo $n$, isomorphism of structures, homotopy of maps.

\term{Equivalence of Categories}
An equivalence of categories consists of functors $F: \mathcal{C} \to \mathcal{D}$ and $G: \mathcal{D} \to \mathcal{C}$ together with natural isomorphisms $GF \cong \mathrm{id}_{\mathcal{C}}$ and $FG \cong \mathrm{id}_{\mathcal{D}}$. Equivalent categories are interchangeable for all categorical purposes: any construction preserved by composition with $F$ or $G$ holds in both, such as the existence of limits, coproducts, and adjoints. Finite-dimensional vector spaces over a field are equivalent to the category of matrices over that field, and rings with equivalent module categories are Morita equivalent. By Yoneda's lemma, a category is essentially determined up to equivalence by its representable functors.

\term{Equivalent}
Two objects are equivalent when they are related by an equivalence relation; two statements are logically equivalent when each implies the other. \par\noindent\textbf{Applications:} equivalence lets mathematicians replace a difficult object or statement with a simpler one having the same relevant behavior.

\term{Erdős--Ko--Rado Theorem}
If $n \geq 2k$, a family of $k$-element subsets of an $n$-element set with pairwise intersecting members has size at most $\binom{n-1}{k-1}$; equality is achieved by all sets containing a fixed element.

\term{Erdős--Szekeres}
The Erdős--Szekeres theorem asserts that any sequence of $(r-1)(s-1)+1$ distinct real numbers contains an increasing subsequence of length $r$ or a decreasing subsequence of length $s$; equivalently, every sufficiently large point set in the plane contains a large monotone subset. The bound is sharp and the theorem is proved by a pigeonhole argument. It is a foundational result of Ramsey-type combinatorics and a precursor to the Erdős--Szekeres cups-and-caps theorem on convex configurations.

\term{Ergodic Theorem (Birkhoff)}
Birkhoff's ergodic theorem states that for a measure-preserving transformation $T$ on a probability space $(X, \mathcal{F}, \mu)$ and $f \in L^1(\mu)$, the time averages $\frac{1}{n}\sum_{k=0}^{n-1} f(T^k x)$ converge almost everywhere to a $T$-invariant function $\bar{f}$. If $T$ is ergodic, $\bar{f} = \int f\,d\mu$ almost everywhere. This justifies the physical principle that time averages equal space averages for ergodic systems.

\term{Ergodic Theory}
Ergodic theory studies long-term statistical behavior in systems that preserve a measure. \par\noindent\textbf{Applications:} it connects time averages with ensemble averages in statistical mechanics, dynamical systems, information theory, and data modeling.

\term{Ergodicity}
Ergodicity is the property that long-run averages along almost every trajectory equal averages over the whole invariant state space. \par\noindent\textbf{Applications:} it justifies replacing repeated time measurements by ensemble averages in physics, probability, and dynamical simulations.

\term{Essential Singularity}
An essential singularity is an isolated singularity that is neither removable nor a pole. Near it, the function behaves extremely irregularly and approaches almost every complex value. \par\noindent\textbf{Applications:} essential singularities classify complex functions and determine local behavior in contour integration and complex models.

\term{Étale Cohomology}
Étale cohomology $H_{\text{\'et}}^*(X, \mathcal{F})$ is a cohomology theory for schemes developed by Grothendieck and Artin (1960s). It uses the étale site (étale morphisms as covers) instead of the Zariski or analytic topology. It satisfies the Weil conjectures: rationality, functional equation, and Riemann hypothesis for varieties over finite fields. The $\ell$-adic cohomology $H_{\text{\'et}}^*(X_{\bar{k}}, \mathbb{Q}_\ell)$ carries a Galois action.

\term{Étale Morphism}
A morphism of schemes $f: X \to Y$ is étale if it is smooth of relative dimension $0$, equivalently flat and unramified. Étale morphisms are the algebro-geometric analogue of local isomorphisms of complex manifolds: they induce isomorphisms on completions of local rings and lift tangent vectors, so the formal inverse function theorem holds. The étale site, whose covers are surjective families of étale morphisms, underlies étale cohomology, and the étale fundamental group classifies finite étale covers of a connected scheme.

\term{Euclid's Algorithm}
Euclid's algorithm repeatedly divides with remainders to compute the greatest common divisor of two integers. \par\noindent\textbf{Applications:} it supports fraction reduction, modular inverses, polynomial arithmetic, and cryptographic algorithms.

\term{Euclid's Theorem}
There are infinitely many prime numbers. Euclid's proof multiplies any proposed finite list of primes and adds one, producing a number with a prime divisor absent from the list. \par\noindent\textbf{Applications:} the theorem establishes the basic supply of prime building blocks used in arithmetic and cryptography.

\term{Euclidean Domain}
An Euclidean domain is an integral domain with a division algorithm based on a size function. It is therefore a principal ideal domain and a unique factorization domain. \par\noindent\textbf{Applications:} Euclidean domains support efficient gcd calculations and factorization in integers and polynomial rings.

\term{Euler Characteristic}
The Euler characteristic $\chi(X)$ of a topological space is $\chi = \sum_{k=0}^\infty (-1)^k \operatorname{rank} H_k(X; \mathbb{Q})$. For a finite CW complex, $\chi = \sum (-1)^k (\text{number of } k\text{-cells})$. For polyhedra, $\chi = V - E + F$. The Gauss--Bonnet theorem relates $\chi(S)$ of a surface to its Gaussian curvature: $\int_S K\,dA = 2\pi\chi(S)$. The Euler characteristic is multiplicative under fibrations.

\term{Euler's Formula (Polyhedra)}
For a convex polyhedron, $V - E + F = 2$ where $V$ is vertices, $E$ edges, and $F$ faces. More generally, for a connected planar graph, $V - E + F = 2$ including the unbounded face. This is the Euler characteristic of the sphere. For surfaces of genus $g$, $V - E + F = 2 - 2g$. The formula is a cornerstone of topology and graph theory.

\term{Euler's Identity}
Euler's identity $e^{i\pi}+1=0$ follows from $e^{i\theta}=\cos\theta+i\sin\theta$ at $\theta=\pi$. It connects exponential, trigonometric, and complex notation. \par\noindent\textbf{Applications:} this relationship powers Fourier analysis, alternating-current calculations, signal processing, and quantum mechanics.

\term{Euler's Number}
Euler's number is the transcendental constant $e=\lim_{n\to\infty}(1+1/n)^n\approx2.71828$, the base of the natural logarithm. \par\noindent\textbf{Applications:} $e^x$ models continuous growth and decay, compound interest, radioactive decay, and probability distributions.

\term{Euler's Theorem}
(Number theory) If $\gcd(a,n)=1$, then $a^{\varphi(n)}\equiv1\pmod n$. \par\noindent\textbf{Applications:} Euler's theorem supports modular inverses, primality tests, and the RSA cryptosystem.

\term{Euler--Lagrange Equation}
The Euler--Lagrange equation $\partial L/\partial f-d(\partial L/\partial f')/dx=0$ is satisfied by functions that make a variational quantity stationary. \par\noindent\textbf{Applications:} it gives the equations of motion in mechanics and is used in optics, optimal control, and PDEs.

\term{Euler--Maclaurin Formula}
The Euler--Maclaurin formula relates sums to integrals: $\sum_{k=a}^b f(k) = \int_a^b f(x)\,dx + \frac{f(a)+f(b)}{2} + \sum_{k=1}^m \frac{B_{2k}}{(2k)!} \left(f^{(2k-1)}(b) - f^{(2k-1)}(a)\right) + R_m$, where $B_{2k}$ are Bernoulli numbers. It provides asymptotic expansions for sums, computes zeta values, and is used in numerical integration and analytic number theory.

\term{Eulerian Graph}
A connected graph is Eulerian if it has a closed walk using every edge exactly once; equivalently, every vertex has even degree. \par\noindent\textbf{Applications:} Eulerian graphs model inspection routes, street sweeping, circuit testing, and the Chinese postman problem.

\term{Eulerian Path}
An Eulerian path in a graph is a trail that uses every edge exactly once; if it returns to its starting vertex it is an Eulerian circuit. Euler's theorem characterises when these exist: a connected graph admits an Eulerian circuit iff every vertex has even degree, and an Eulerian path iff at most two vertices have odd degree. The problem originated with the bridges of Königsberg. Hierholzer's algorithm finds an Eulerian circuit in linear time. Eulerian circuits are in a sense dual to Hamiltonian cycles.

\term{Event}
An event is a subset of the sample space to which a probability is assigned. \par\noindent\textbf{Applications:} events represent outcomes of experiments, alerts in monitoring systems, medical conditions, and statements tested by data.

\term{Exact Sequence}
An exact sequence is a sequence of homomorphisms in which the image of each map equals the kernel of the next. \par\noindent\textbf{Applications:} exact sequences track information through algebraic constructions and are used in homology, differential equations, and extension problems.

\term{Excision}
Excision is the theorem that for subspaces $Z \subset A \subset X$ with $\overline{Z} \subset \operatorname{int}(A)$, the inclusion $(X-Z, A-Z) \hookrightarrow (X, A)$ induces isomorphisms $H_n(X-Z, A-Z) \cong H_n(X, A)$ on relative homology. It is one of the Eilenberg--Steenrod axioms which, together with homotopy invariance, exactness, additivity, and normalisation, characterise homology theories. Excision enables the computation of homology by attaching cells and yields the Mayer--Vietoris long exact sequence as a consequence.

\term{Existence}
Existence means that at least one object satisfies a condition, written with the quantifier $\exists$. \par\noindent\textbf{Applications:} existence theorems guarantee that a model, solution, equilibrium, or optimum actually exists before it is computed.

\term{Existence and Uniqueness Theorem}
For ordinary differential equations, the Picard--Lindelöf theorem guarantees a unique local solution to $y' = f(x, y)$ with $y(x_0) = y_0$ when $f$ is continuous in $x$ and Lipschitz in $y$.

\term{Expected Value}
The expected value $\mathbb E[X]=\int x\,dP$ is the probability-weighted average of a random variable. \par\noindent\textbf{Applications:} expected values describe long-run averages in risk, finance, statistical estimation, games, and decision-making.

\term{Explicit Formula}
A formula is explicit when it directly gives the quantity of interest rather than defining it indirectly or recursively. \par\noindent\textbf{Applications:} explicit formulas allow direct evaluation, symbolic analysis, and efficient numerical implementation.

\term{Exponential Distribution}
The exponential distribution has density $f(x)=\lambda e^{-\lambda x}$ for $x\ge0$ and models waiting times between events of a Poisson process. \par\noindent\textbf{Applications:} it describes service times, radioactive decay, reliability, queueing, and time-to-failure models.

\term{Expectation (Probability)}
The expectation $\mathbb{E}[X]$ of a random variable $X$ is $\int X\,dP$ on a probability space, or $\sum x P(X=x)$ for discrete, or $\int x f(x)\,dx$ for continuous with density $f$. Linearity: $\mathbb{E}[aX + bY] = a\mathbb{E}[X] + b\mathbb{E}[Y]$. The expectation of a product is $\mathbb{E}[XY] = \operatorname{Cov}(X,Y) + \mathbb{E}[X]\mathbb{E}[Y]$. The law of the unconscious statistician: $\mathbb{E}[g(X)] = \int g(x) f(x)\,dx$.

\term{Exponential Map (Lie Theory)}
The exponential map $\exp: \mathfrak{g} \to G$ from a Lie algebra $\mathfrak{g}$ to its Lie group $G$ sends $X$ to the value at $1$ of the unique one-parameter subgroup with tangent $X$. For matrix groups, $\exp(X) = \sum_{k=0}^\infty X^k/k!$. The exponential map is a local diffeomorphism near $0$; its derivative is the identity. It provides the link between Lie algebras and Lie groups.

\term{Exponential Function}
The exponential function $e^x = \sum_{n=0}^\infty x^n/n!$ is the unique function equal to its own derivative with $e^0 = 1$. For complex $z$, $e^z = e^x(\cos y + i\sin y)$ where $z = x+iy$. The complex exponential is periodic: $e^{z+2\pi i} = e^z$. It satisfies $e^{a+b} = e^a e^b$ and maps lines to rays, strips to annuli. The inverse is the complex logarithm $\log z = \ln|z| + i\arg z$.

\term{Ext Functor}
For an abelian category with enough projectives or injectives, the derived functor $\operatorname{Ext}_R^n(A,B)$ of $\mathrm{Hom}_R$ is computed as the cohomology of $\mathrm{Hom}_R(P_\bullet, B)$ for a projective resolution $P_\bullet$ of $A$, or symmetrically with an injective resolution of $B$. $\operatorname{Ext}^1_R(A,B)$ classifies extensions $0 \to B \to E \to A \to 0$ up to equivalence. Higher groups give sheaf cohomology via $\operatorname{Ext}^n_{\mathcal{O}_X}(\mathcal{O}_X, \mathcal{F}) = H^n(X, \mathcal{F})$, and universal-coefficient duality relates $\operatorname{Ext}$ to $\operatorname{Tor}$.

\term{Extension}
A field $K$ containing a field $F$, written $K/F$; also any structure containing a given structure.

\term{Extension Field}
A field extension $L/K$ is a pair of fields with $K \subseteq L$. The degree $[L:K]$ is $\dim_K L$ as a $K$-vector space. An extension is algebraic if every element of $L$ is a root of a polynomial over $K$. A finite extension is algebraic. The minimal polynomial of $\alpha \in L$ over $K$ is the monic irreducible polynomial with $\alpha$ as a root; its degree divides $[L:K]$. Splitting fields and Galois extensions are key constructions.

\term{Exterior Algebra}
The exterior algebra $\bigwedge V$ of a vector space $V$ is the algebra generated by $V$ with $v \wedge v = 0$ for all $v \in V$ (hence $v \wedge w = -w \wedge v$). It is graded: $\bigwedge V = \bigoplus_{k=0}^n \bigwedge^k V$ where $\bigwedge^k V$ has basis $v_{i_1} \wedge \cdots \wedge v_{i_k}$. The wedge product is associative and bilinear. The exterior algebra is the free graded-commutative algebra on $V$.

\term{Exterior Derivative}
The unique anti-derivation $d$ on differential forms satisfying $d^2 = 0$ and $df$ equal to the ordinary differential.

\term{Exterior Product}
The antisymmetric product $\wedge$ on differential forms, with $dx \wedge dy = -dy \wedge dx$.

\term{Extremal Graph Theory}
Extremal graph theory studies the maximum or minimum number of edges a graph can have while avoiding a given subgraph. Turán's theorem: the maximum edges in an $n$-vertex graph with no $K_{r+1}$ is $\left(1 - \frac{1}{r}\right)\frac{n^2}{2}$, achieved by the Turán graph $T(n,r)$. The extremal number $\operatorname{ex}(n, H)$ is the maximum edges avoiding $H$. The Erdős--Stone theorem gives $\operatorname{ex}(n, H) = \left(1 - \frac{1}{\chi(H)-1} + o(1)\right)\frac{n^2}{2}$.

\term{Extreme Value Theorem}
A continuous real-valued function on a compact space attains its maximum and minimum. On a closed bounded interval $[a,b]$, continuity is enough; compactness captures the essential property that every open cover has a finite subcover.

\term{Extremum}
A maximum or minimum of a function, either local or global.
\term{Cayley Graph}
A graph representing a group: the vertices are group elements and edges correspond to multiplication by generators. Cayley graphs encode the structure of groups geometrically and are central to geometric group theory. For example, the Cayley graph of $\mathbb{Z}^n$ with standard generators is the infinite grid, and the Cayley graph of $S_3$ with generators $(12)$ and $(123)$ is a triangle with diagonals.
\term{Euler Totient Function}
The function $\phi(n)$ counting positive integers $\leq n$ that are coprime to $n$. It is multiplicative: $\phi(mn) = \phi(m)\phi(n)$ when $\gcd(m,n) = 1$, and $\phi(p^k) = p^k - p^{k-1}$ for primes $p$. Euler's theorem states $a^{\phi(n)} \equiv 1 \pmod{n}$ for $\gcd(a,n) = 1$, generalising Fermat's little theorem.
\term{Eulerian Circuit}
A closed walk in a graph that traverses every edge exactly once. A connected graph has an Eulerian circuit if and only if every vertex has even degree, by Euler's theorem (1736). Such a graph is called Eulerian. The Chinese postman problem asks for the shortest closed walk covering every edge of a weighted graph.
\term{Equalizer}
In category theory, the equalizer of two morphisms $f, g: A \to B$ is the object $E$ with a morphism $e: E \to A$ such that $f \circ e = g \circ e$, and which is universal with this property. In sets, the equalizer is $\{a \in A : f(a) = g(a)\}$. Equalizers and coequalizers are instances of limits and colimits, respectively.

\chapter*{F}
\label{chap:f}

\term{Faà di Bruno's Formula}
The general chain rule for higher derivatives: $\frac{d^n}{dx^n}f(g(x)) = \sum \frac{n!}{k_1!(m_1!) \cdots k_n!(m_n!)} f^{(m_1+\cdots+m_n)}(g(x)) \prod_{j=1}^n \left(\frac{g^{(j)}(x)}{j!}\right)^{m_j}$ where the sum is over all non-negative integers $k_1, \ldots, k_n, m_1, \ldots, m_n$ satisfying $\sum j k_j = n$ and $\sum k_j = m$. This formula appears in combinatorics, probability (Bell polynomials), and differential equations.

\term{Face (Polytope)}
A face of a convex polytope is the part cut out by a supporting hyperplane. Faces include vertices, edges, facets, and the whole polytope. \par\noindent\textbf{Applications:} faces describe the geometry of polyhedra and are used in linear programming, computer graphics, and computational geometry.

\term{Factorial Function}
The factorial is $n!=1\cdot2\cdots n$ for a nonnegative integer, with $0!=1$. Stirling's approximation gives $n!\sim\sqrt{2\pi n}(n/e)^n$, and $n!=\Gamma(n+1)$. \par\noindent\textbf{Applications:} factorials count arrangements, enter probability formulas, and support series expansions and numerical approximations.

\term{Fadeev--Popov Determinant}
In gauge theory, the Fadeev--Popov determinant arises when fixing a gauge in the path integral. For a gauge group $G$ acting on a field configuration space, the Fadeev--Popov procedure inserts $1 = \Delta_{FP}[A] \int \mathcal{D}g \, \delta(G(A^g))$ into the path integral. The resulting ghost fields are anticommuting scalars and are essential for maintaining unitarity in non-abelian gauge theories.

\term{Fan (Toric Geometry)}
A fan $\Sigma$ in $\mathbb{R}^n$ is a collection of strongly convex rational polyhedral cones such that every face of a cone in $\Sigma$ is also in $\Sigma$, and the intersection of any two cones is a face of each. Fans correspond to toric varieties: each fan $\Sigma$ defines a complex algebraic variety $X_\Sigma$ with an effective torus action. The fan for $\mathbb{CP}^n$ consists of cones generated by subsets of the standard basis and their negatives.

\term{Fan Graph}
A fan graph joins one central vertex to every vertex of a path. \par\noindent\textbf{Applications:} fan graphs model hub-and-spoke networks and appear in chemical graph theory, connectivity studies, and graph spectra.

\term{Fatou's Lemma}
For a sequence of non-negative measurable functions $f_n \geq 0$: $\int \liminf f_n \, d\mu \leq \liminf \int f_n \, d\mu$. This fundamental result in measure theory is a key ingredient in the proof of the Dominated Convergence Theorem and the Monotone Convergence Theorem. It reflects the lower semicontinuity of the integral under pointwise convergence.

\term{Fault-Tolerant Quantum Computation}
A fault-tolerant quantum computation uses quantum error-correcting codes to protect against errors during computation. The threshold theorem states that if the error rate per gate is below a constant threshold, arbitrary long quantum computations can be performed with polynomial overhead. Surface codes and the $[[7,1,3]]$ Steane code are prominent examples. Fault tolerance requires entanglement, syndrome measurement, and adaptive correction.

\term{Fermat's Last Theorem}
Fermat's Last Theorem states that $x^n+y^n=z^n$ has no positive integer solutions when $n>2$. \par\noindent\textbf{Applications:} its proof connects elliptic curves, modular forms, and number theory and illustrates how elementary equations can require deep mathematics.

\term{Fermat's Little Theorem}
If $p$ is prime and $p\nmid a$, then $a^{p-1}\equiv1\pmod p$. \par\noindent\textbf{Applications:} Fermat's little theorem supports primality tests, modular inverses, and public-key cryptography.

\term{Fiber Bundle}
A fiber bundle is a space $E$ with a projection $\pi:E\to B$ that locally looks like a product $U\times F$, where $B$ is the base and $F$ is the fiber. Globally, the fibers may twist, as in a Möbius strip. \par\noindent\textbf{Applications:} bundles model fields, orientations, configuration spaces, and gauge potentials in geometry and physics.

\term{Fibrant Object}
In a model category, an object is fibrant when its map to the terminal object has the required lifting property. Informally, a fibrant object is one on which homotopy-theoretic constructions behave smoothly. In simplicial sets, the fibrant objects are Kan complexes. Weak equivalences between objects that are both fibrant and cofibrant have homotopy inverses. \par\noindent\textbf{Applications:} fibrant replacements let algebraic models represent the homotopy type of spaces without changing their essential information.

\term{Fibration}
A fibration is a map with a homotopy lifting property: a deformation in the base can be lifted after an initial lift is chosen. Fiber bundles are common examples. \par\noindent\textbf{Applications:} fibrations organize homotopy groups, fibered spaces, and the deformation of geometric or topological models.

\term{Fibonacci Sequence}
The Fibonacci sequence satisfies $F_0=0$, $F_1=1$, and $F_n=F_{n-1}+F_{n-2}$. The ratio of consecutive terms approaches the golden ratio $\phi=(1+\sqrt5)/2$. \par\noindent\textbf{Applications:} Fibonacci numbers appear in combinatorics, algorithms, continued fractions, coding, and models of branching and growth.

\term{Field (Algebra)}
A field is a commutative ring in which every nonzero element has a multiplicative inverse. Examples include $\mathbb Q$, $\mathbb R$, $\mathbb C$, and finite fields. \par\noindent\textbf{Applications:} fields provide the scalars for linear algebra and are central to coding theory, cryptography, and polynomial equations.

\term{Field Extension}
A field extension $L/K$ is a pair of fields with $K \subseteq L$. The degree $[L:K]$ is the dimension of $L$ as a $K$-vector space. A finite extension is algebraic if every element of $L$ is a root of a polynomial over $K$. The splitting field of a polynomial $f$ over $K$ is the smallest extension in which $f$ factors into linear factors. Galois theory studies the correspondence between intermediate fields and subgroups of the Galois group.

\term{Fierz Identity}
In quantum field theory, the Fierz identity is a rearrangement formula for products of spinor bilinears. For Dirac spinors $\psi_i$: $(\bar{\psi}_1 \Gamma^A \psi_2)(\bar{\psi}_3 \Gamma_A \psi_4) = \sum_B c_B (\bar{\psi}_1 \Gamma^B \psi_4)(\bar{\psi}_3 \Gamma_B \psi_2)$ where the sum runs over the 16 Dirac matrices and $c_B$ are numerical coefficients. Fierz identities are used in deriving four-fermion interactions and checking supersymmetry transformations.

\term{Filtration}
A filtration is a nested sequence $\mathcal F_0\subseteq\mathcal F_1\subseteq\cdots$ that records how an object is built in stages. In probability, a filtration represents the information available over time; in algebra, it records layers of a module or complex. \par\noindent\textbf{Applications:} filtrations are used in stochastic processes, spectral sequences, persistent homology, and multiscale data analysis.

\term{Final Object}
In category theory, a final (or terminal) object is an object $T$ such that for every object $X$ there exists a unique morphism $X \to T$. In the category of sets, the final object is any singleton $\{*\}$. In the category of groups, it is the trivial group. Final objects are unique up to unique isomorphism. Dually, an initial object has a unique morphism from it to every object.

\term{Finite Field}
A finite field is a field with finitely many elements; its cardinality is necessarily a prime power $p^n$ and its characteristic is $p$. For every prime power $q = p^n$ there exists a field $\mathbb{F}_q$ with $q$ elements, unique up to isomorphism, namely the splitting field of $x^q - x$ over $\mathbb{F}_p$. The multiplicative group $\mathbb{F}_q^\times$ is cyclic, and a generator is a primitive element. The Frobenius automorphism $x \mapsto x^p$ generates the Galois group of $\mathbb{F}_q$ over $\mathbb{F}_p$. Finite fields underlie coding theory (Reed--Solomon and BCH codes), cryptography (elliptic curves, discrete logarithms), and pseudorandom generation.

\term{Finite Group}
A finite group has finitely many elements, and its order is the number of elements. \par\noindent\textbf{Applications:} finite groups describe symmetries of molecules, crystals, codes, puzzles, and finite computational systems.

\term{Finite Set}
A finite set has a fixed number of elements and is in bijection with $\{1,2,\ldots,n\}$ for some $n$. \par\noindent\textbf{Applications:} finite sets model stored data, finite state machines, discrete experiments, and combinatorial objects.

\term{Finite-difference Method}
A finite-difference method approximates derivatives by differences of function values on a grid. \par\noindent\textbf{Applications:} it solves heat, wave, fluid, and structural equations numerically.

\term{Finite-element Method}
A finite-element method approximates a PDE solution by piecewise polynomial functions on a mesh, usually through a variational formulation. \par\noindent\textbf{Applications:} it is widely used for stress, heat, fluid, electromagnetic, and structural simulations.

\term{Finite-volume Method}
A finite-volume method integrates a conservation law over small control volumes and balances fluxes across their boundaries. \par\noindent\textbf{Applications:} its conservation property makes it effective for computational fluid dynamics, weather, combustion, and traffic flow.

\term{Finitely Presented Group}
A finitely presented group is given by generators and relations: $G = \langle S \mid R \rangle$ where $S$ is a finite set of generators and $R$ is a finite set of relations (words in $S$). Every finite group is finitely presented. However, the word problem is undecidable in general: there is no algorithm to determine whether two words represent the same group element. The Higman embedding theorem characterises recursively presented groups.

\term{First Isomorphism Theorem}
For a group homomorphism $\phi: G \to H$, the quotient group $G/\ker(\phi)$ is isomorphic to the image of $\phi$: $G/\ker(\phi) \cong \operatorname{im}(\phi)$. The analogous statements hold for rings, where $R/\ker(\phi) \cong \operatorname{im}(\phi)$ as rings, and for modules, where $M/\ker(\phi) \cong \operatorname{im}(\phi)$ as modules. This is the first of the three isomorphism theorems, the other two following from it by applying it to suitable maps.

\term{First-order Logic}
A formal system using variables, relations, logical connectives, and quantifiers over elements of a specified domain. \par\noindent\textbf{Applications:} first-order logic underlies mathematical foundations, database queries, automated reasoning, and formal verification.

\term{Fischer--Burmeister Function}
The Fischer--Burmeister function $\phi(a,b) = \sqrt{a^2 + b^2} - a - b$ provides a smooth reformulation of the complementarity problem: $\phi(a,b) = 0$ iff $a \geq 0$, $b \geq 0$, $ab = 0$. It is used in numerical methods for variational inequalities and mathematical programs with equilibrium constraints. The function is continuously differentiable everywhere except at the origin.

\term{Fixed Point}
A fixed point of a function $f: X \to X$ is a point $x$ with $f(x) = x$. The Brouwer fixed-point theorem guarantees fixed points for continuous maps on compact convex sets in $\mathbb{R}^n$. In dynamical systems, fixed points are equilibrium solutions; their stability is determined by the eigenvalues of the Jacobian. Fixed-point iterations $x_{n+1} = f(x_n)$ converge under contraction conditions.

\term{Fixed-Point Theorem}
A fixed-point theorem guarantees a point $x$ with $f(x)=x$ under specified conditions, as in the Banach and Brouwer theorems. \par\noindent\textbf{Applications:} fixed points represent equilibria in economics, games, dynamical systems, and nonlinear equations.

\term{Flat Module}
An $R$-module $M$ is flat if tensoring with $M$ preserves exact sequences. Flatness is weaker than being free: over a local ring, a finitely presented flat module is free, but an arbitrary flat module need not be. Projective modules are flat, and flatness makes the algebraic behavior of fibers change without sudden jumps. \par\noindent\textbf{Applications:} flat modules and flat maps are central in algebraic geometry, where they describe families whose fibers vary consistently.

\term{Flatness}
Flatness means that tensoring a module with $M$ preserves exact sequences; over an integral domain it generalizes being torsion-free. \par\noindent\textbf{Applications:} flatness describes algebraic families whose fibers change without sudden jumps in algebraic geometry.

\term{Floating-Point Arithmetic}
Floating-point arithmetic represents real numbers approximately as $m\times\beta^e$ using a finite mantissa and exponent. Rounding and overflow are unavoidable and are standardized by IEEE 754. \par\noindent\textbf{Applications:} understanding floating point is essential for reliable scientific computing, graphics, finance, and machine learning.

\term{Flow}
A flow is a family of maps describing how states evolve with time; its paths are integral curves of a vector field. \par\noindent\textbf{Applications:} flows model motion, fluid transport, population change, and continuous-time control systems.

\term{Focal Point (Optics)}
A focal point of a lens or mirror is the point where parallel rays converge (or appear to diverge from). For a thin lens with focal length $f$, the lens equation is $\frac{1}{f} = \frac{1}{d_o} + \frac{1}{d_i}$ where $d_o$ and $d_i$ are object and image distances. Focal properties are described by ray transfer matrices in paraxial optics. The focal length determines the magnification and field of view.

\term{Fold Catastrophe}
A fold catastrophe is the simplest singularity in Thom's classification of elementary catastrophes. The normal form is $V(x; \lambda) = x^3 + \lambda x$, with catastrophe set given by the discriminant $27\lambda^2 + 4x^3 = 0$. At the fold, a stable and unstable branch of equilibria coalesce and annihilate. Folds appear in optics (rainbow formation), bifurcation theory, and elasticity.

\term{Foata Transform}
The Foata transform is a bijection on permutations that maps the descent set of a permutation to its cycle structure in a controlled way. Defined by Désiré Foata (1968), it preserves the Eulerian distribution and relates to the Robinson--Schensted correspondence. The transform is used in combinatorics on words, pattern avoidance, and the study of permutation statistics.

\term{Formula}
A formula is a symbolic rule relating quantities, such as the quadratic formula. \par\noindent\textbf{Applications:} formulas make models computable, communicate mathematical relationships, and support prediction and parameter estimation.

\term{Four Color Theorem}
Every planar graph has a vertex coloring using at most four colors such that adjacent vertices receive different colors. The theorem was proved in 1976 by Kenneth Appel and Wolfgang Haken, the first major theorem to be proved with substantial computer assistance. Four colors are sometimes necessary, since the complete graph $K_4$ is planar.

\term{Four-Square Theorem}
Also called Lagrange's four-square theorem: every natural number is a sum of four integer squares, $n = a_1^2 + a_2^2 + a_3^2 + a_4^2$, proved by Lagrange in 1770. Related results about sums of squares include Jacobi's four-square theorem, which counts the number of representations of $n$, and Legendre's three-square theorem, which characterises when three squares suffice.

\term{Fourier Coefficients}
For a $2\pi$-periodic function $f \in L^1([-\pi,\pi])$, the Fourier coefficients are $\hat{f}(n) = \frac{1}{2\pi} \int_{-\pi}^{\pi} f(x) e^{-inx} \, dx$ for $n \in \mathbb{Z}$. The Fourier series is $f(x) \sim \sum_{n=-\infty}^{\infty} \hat{f}(n) e^{inx}$. For $f \in L^2$, Parseval's identity states $\|f\|_2^2 = \sum |\hat{f}(n)|^2$. The coefficients measure the frequency content of $f$.

\term{Fourier Series}
The representation of a periodic function as a sum of sinusoids, $f(x) = a_0 + \sum (a_n \cos nx + b_n \sin nx)$.

\term{Fourier Transform}
The Fourier transform of $f \in L^1(\mathbb{R})$ is $\hat{f}(\xi) = \int_{-\infty}^{\infty} f(x) e^{-2\pi i x \xi} \, dx$. It extends to an isometry on $L^2(\mathbb{R})$ (Plancherel theorem). Key properties: $\widehat{f'}(\xi) = 2\pi i \xi \hat{f}(\xi)$, $\hat{\hat{f}}(x) = f(-x)$, and $\widehat{f*g} = \hat{f} \cdot \hat{g}$. The Fourier transform diagonalises convolution and is fundamental to analysis, signal processing, and quantum mechanics.

\term{Fourier--Motzkin Elimination}
Fourier--Motzkin elimination is a method for eliminating variables from a system of linear inequalities. Given inequalities in variables $x_1, \ldots, x_n$, one can eliminate $x_n$ by combining pairs of inequalities of the form $a x_n \leq b$ and $c x_n \geq d$ to produce an inequality not involving $x_n$. The procedure is the linear programming analogue of Gaussian elimination and is used in constraint projection.

\term{Fourier's Law}
Fourier's law of heat conduction states that the heat flux $\mathbf{q}$ is proportional to the negative temperature gradient: $\mathbf{q} = -k \nabla T$ where $k > 0$ is the thermal conductivity. Combined with conservation of energy, this yields the heat equation $\frac{\partial T}{\partial t} = \alpha \nabla^2 T$ where $\alpha = k/(\rho c_p)$. The law is empirical but derivable from kinetic theory.

\term{Fractal}
A fractal is a set with detail or patterns that repeat across scales, often with a non-integer dimension. Examples include the Cantor and Mandelbrot sets. \par\noindent\textbf{Applications:} fractals model coastlines, branching networks, porous materials, textures, and financial or biological patterns.

\term{Fractal Dimension}
The fractal (Hausdorff) dimension of a set $E \subset \mathbb{R}^n$ is $s = \inf\{t : \mathcal{H}^t(E) = 0\} = \sup\{t : \mathcal{H}^t(E) = \infty\}$ where $\mathcal{H}^t$ is the $t$-dimensional Hausdorff measure. For the Cantor set, $s = \log 2 / \log 3$. The Sierpiński triangle has dimension $\log 3 / \log 2$. Fractal dimensions quantify the "roughness" or "space-filling" property of self-similar sets.

\term{Frame}
A frame is a possibly redundant family of vectors that stably represents every vector in a Hilbert space. Unlike a basis, it may contain more vectors than necessary. \par\noindent\textbf{Applications:} frames enable robust signal reconstruction, image compression, sampling, and noise-resistant measurements.

\term{Frattini Subgroup}
The Frattini subgroup $\Phi(G)$ of a group $G$ is the intersection of all maximal subgroups of $G$. Equivalently, it is the set of nongenerator elements: $g \in \Phi(G)$ iff whenever $G = \langle g, S \rangle$ for some $S$, then $G = \langle S \rangle$. For a finite $p$-group, $\Phi(G) = G^p [G,G]$ (the product of $p$-th powers and commutators). The Frattini subgroup is nilpotent and characteristic.

\term{Fréchet Derivative}
The derivative of a map between normed spaces as a bounded linear operator, generalising the total derivative.

\term{Fredholm Alternative}
For a compact operator $T$ on a normed space, either the equation $(I - T)x = y$ is uniquely solvable for every $y$, or the homogeneous equations $(I - T)x = 0$ and $(I - T^*)y = 0$ have the same positive finite dimension. This dichotomy is the core of Fredholm theory and is the basis for solving integral equations with compact kernels. It reduces solvability to orthogonality conditions against solutions of the adjoint homogeneous equation.

\term{Fredholm Theory}
Fredholm theory studies equations of the form $(I-K)x=y$, where $K$ is often a compact operator. It gives a solvability alternative: an equation is solvable precisely when $y$ satisfies compatibility conditions coming from the corresponding adjoint equation. \par\noindent\textbf{Applications:} the theory helps analyze integral equations, boundary-value problems, and linearized partial differential equations.

\term{Free Group}
A free group is generated by symbols with no relations except cancellation of a symbol with its inverse. Every group is a quotient of a free group. \par\noindent\textbf{Applications:} free groups describe paths, symmetries, automata, and the presentations used in computational group theory.

\term{Free Module}
A free module has a basis, so every element has a unique finite linear combination of basis elements. Equivalently, it is a direct sum of copies of its ring. \par\noindent\textbf{Applications:} free modules provide coordinates for algebraic systems and are used in linear algebra, coding, and algebraic geometry.

\term{Frenet--Serret Formulas}
For a smooth curve $\alpha(s)$ parametrised by arc length in $\mathbb{R}^3$, the Frenet--Serret frame consists of the tangent $\mathbf{T}$, normal $\mathbf{N}$, and binormal $\mathbf{B} = \mathbf{T} \times \mathbf{N}$. The formulas are: $\mathbf{T}' = \kappa \mathbf{N}$, $\mathbf{N}' = -\kappa \mathbf{T} + \tau \mathbf{B}$, $\mathbf{B}' = -\tau \mathbf{N}$ where $\kappa$ is curvature and $\tau$ is torsion. These equations characterise curves up to rigid motion.

\term{Fremlin Measure}
The Fremlin measure is a construction in measure theory that produces a localisable measure from a premeasure on an algebra. Defined by David Fremlin, it is used in the theory of Riesz spaces and vector lattices. The Fremlin construction is essential for extending measures to larger $\sigma$-algebras while preserving locality and completeness properties.

\term{Frequency}
Frequency is the number of cycles per unit time of a periodic phenomenon; it is the reciprocal of the period. \par\noindent\textbf{Applications:} frequency is fundamental in sound, radio, vibration, image analysis, and Fourier-based signal processing.

\term{Freudenthal's Theorem (Homotopy)}
Freudenthal's theorem on the stability of homotopy groups states that for a connected CW-complex $X$, the suspension homomorphism $\pi_k(X) \to \pi_{k+1}(\Sigma X)$ is an isomorphism for $k < 2n-1$ and surjective for $k = 2n-1$ when $X$ is $(n-1)$-connected. This is the first step in the Freudenthal suspension theorem, foundational for stable homotopy theory.

\term{Frobenius Coin Problem}
The Frobenius coin problem asks for the largest integer that cannot be expressed as a non-negative integer combination of given coprime positive integers $a_1, \ldots, a_n$. For two numbers $a, b$, the Frobenius number is $g(a,b) = ab - a - b$ (Sylvester, 1884). For three or more numbers, no closed formula is known; the problem is NP-hard. The number of non-representable integers is $\frac{1}{(n-1)!} \prod a_i$ for $n=2$.

\term{Frobenius Endomorphism}
In characteristic $p > 0$, the Frobenius endomorphism is the map $x \mapsto x^p$ on a commutative ring or field. It is a ring homomorphism: $(x+y)^p = x^p + y^p$ in characteristic $p$. For a perfect field $\mathbb{F}_q$, the Frobenius generates the Galois group $\operatorname{Gal}(\mathbb{F}_q/\mathbb{F}_p) \cong \mathbb{Z}/n\mathbb{Z}$. In algebraic geometry, the Frobenius acts on cohomology and is central to the Weil conjectures.

\term{Frobenius Reciprocity}
Frobenius reciprocity relates induction and restriction of representations. If $H \leq G$, $\chi$ is a character of $H$, and $\psi$ is a character of $G$, then $\langle \operatorname{Ind}_H^G \chi, \psi \rangle_G = \langle \chi, \operatorname{Res}_H^G \psi \rangle_H$. This adjunction between induction and restriction is fundamental in representation theory and appears in many categories (modules, sheaves, Hilbert spaces).

\term{Frobenius Theorem (Distributions)}
The Frobenius theorem characterises integrable distributions on manifolds. A smooth distribution $D \subset TM$ is integrable (tangent to a foliation) iff it is involutive: $[X, Y] \in D$ for all $X, Y \in D$. In Lie group theory, a Lie subalgebra integrates to a Lie subgroup. The theorem provides the bridge between infinitesimal (Lie algebra) and global (manifold) structures.

\term{Fubini Property}
(Measure theory) The interchange of order of integration in iterated integrals for positive measurable functions, i.e. $\int\int f \, dx\, dy = \int\int f \, dy\, dx$.

\term{Fubini's Theorem}
If $f$ is integrable on a product space $X \times Y$, then $\int_{X \times Y} f \, d(\mu \times \nu) = \int_X \left(\int_Y f(x,y) \, d\nu\right) d\mu = \int_Y \left(\int_X f(x,y) \, d\mu\right) d\nu$. This justifies iterated integration for Lebesgue integrals and is foundational in analysis. Tonelli's theorem applies to non-negative functions without integrability assumptions. Fubini extends to countable products and product measures.

\term{Fuchsian Group}
A Fuchsian group is a discrete subgroup of $\operatorname{PSL}(2, \mathbb{R})$, acting on the upper half-plane $\mathbb{H}$ by Möbius transformations. The quotient $\mathbb{H}/\Gamma$ is a Riemann surface (possibly with orbifold points). Fuchsian groups generalise modular groups; the modular group $\operatorname{PSL}(2, \mathbb{Z})$ is the prototypical example. Uniformization theorem states every hyperbolic Riemann surface is a quotient by a Fuchsian group.

\term{Fuchsian Differential Equation}
A Fuchsian differential equation is a linear ODE whose singular points are all regular (solutions have at most polynomial growth near each singularity). The hypergeometric equation $z(1-z)y'' + [c - (a+b+1)z]y' - aby = 0$ has three regular singular points at $0, 1, \infty$. Riemann's $P$-function formalism classifies Fuchsian equations by their singularities and exponents.

\term{Fullerene}
A fullerene is a molecular graph representing a carbon molecule with 20 or more vertices, where each vertex has degree 3 and every face is a pentagon or hexagon. Euler's formula forces exactly 12 pentagonal faces; the number of hexagons can vary. The buckyball $C_{60}$ (truncated icosahedron) is the most symmetric fullerene. Fullerenes model nanotubes and zero-dimensional carbon allotropes.

\term{Function Space}
A space whose points are functions, such as $C(X)$ or $L^p$, with a topology or metric measuring the closeness of functions.

\term{Functional}
A map from a vector space to its scalar field, especially a continuous linear one such as integration.

\term{Functional Analysis}
Functional analysis studies infinite-dimensional vector spaces equipped with additional structure (norms, inner products, topologies). Key objects: Banach spaces (complete normed spaces), Hilbert spaces (complete inner product spaces), and operator algebras. Fundamental theorems include Hahn--Banach (extension of functionals), Banach--Steinhaus (uniform boundedness), Open Mapping, and Closed Graph theorems. Applications span PDEs, quantum mechanics, and harmonic analysis.

\term{Functional Calculus}
The construction of functions of an operator, such as $f(A)$, via spectral measures or power series.

\term{Functional Determinant}
The functional determinant generalises the matrix determinant to operators. For a differential operator $L$ on a domain with boundary conditions, $\det L$ is defined via zeta-function regularisation: $\det L = \exp(-\zeta_L'(0))$ where $\zeta_L(s) = \sum \lambda_n^{-s}$ sums over eigenvalues. The determinant appears in path integrals, spectral geometry, and the study of quantum anomalies.

\term{Functor}
A structure-preserving map between categories, sending objects to objects and morphisms to morphisms while preserving composition and identities.

\term{Fundamental Group}
The fundamental group $\pi_1(X, x_0)$ consists of homotopy classes of loops based at $x_0$, with composition given by concatenation. For the circle, $\pi_1(S^1) \cong \mathbb{Z}$ (winding number). The fundamental group detects topological features: it is trivial for simply connected spaces and non-trivial for spaces with "holes." Covering space theory establishes a correspondence between subgroups of $\pi_1$ and covering spaces.

\term{Fundamental Solution}
A distribution solving a linear PDE with a delta-function source, from which general solutions are built by convolution.

\term{Fundamental Theorem of Algebra}
Every non-constant polynomial with complex coefficients has at least one complex root. Equivalently, $\mathbb{C}$ is algebraically closed: every polynomial of degree $n$ has exactly $n$ roots counting multiplicity. Proofs exist via complex analysis (Liouville's theorem), algebraic topology (degree theory), and real analysis (infimum argument). The theorem was first stated by Gauss (1799) and is foundational to algebra and analysis.

\term{Fundamental Theorem of Algebraic Topology}
Homology (and cohomology) is a functor from the homotopy category of topological spaces to the category of groups: a continuous map $f: X \to Y$ induces homomorphisms $f_*: H_n(X) \to H_n(Y)$, and a homotopy equivalence induces an isomorphism for every $n$. Spaces of the same homotopy type therefore have isomorphic homology groups, making homology a topological invariant of homotopy type.

\term{Fundamental Theorem of Arithmetic}
Every integer greater than 1 factors uniquely, up to the order of the factors, into a product of primes. This uniqueness is the reason primes are described as the atoms of integer multiplication. The theorem also underlies the Euclidean algorithm and the computation of greatest common divisors.

\term{Fundamental Theorem of Calculus}
If $f$ is continuous on $[a,b]$, then $F(x) = \int_a^x f(t) \, dt$ is differentiable with $F'(x) = f(x)$ (Part 1). If $F$ is any antiderivative of $f$, then $\int_a^b f(x) \, dx = F(b) - F(a)$ (Part 2). This theorem links differentiation and integration, the two central operations of calculus. It generalises to manifolds via Stokes' theorem: $\int_M d\omega = \int_{\partial M} \omega$.

\term{Fundamental Theorem of Finite Abelian Groups}
Every finite abelian group is isomorphic to a direct product of cyclic groups of prime-power order, $G \cong C_{p_1^{a_1}} \times \cdots \times C_{p_k^{a_k}}$, and the decomposition is unique up to the order of the factors. Equivalently, a finite abelian group is the direct product of its Sylow $p$-subgroups, each of which splits into cyclic factors; this is the primary decomposition. The theorem classifies all finite abelian groups.

\term{Fundamental Theorem of Galois Theory}
For a finite Galois extension $L/K$ with group $G = \operatorname{Gal}(L/K)$, there is an inclusion-reversing bijection between intermediate fields $K \subseteq M \subseteq L$ and subgroups $H \leq G$: $M \mapsto \operatorname{Gal}(L/M)$ and $H \mapsto L^H$ (fixed field). Normal subgroups correspond to Galois extensions. This theorem is the cornerstone of Galois theory, linking field theory to group theory.

\term{Fundamental Theorem of Lie Algebras}
A finite-dimensional Lie algebra over $\mathbb{C}$ decomposes as $\mathfrak{g} = \mathfrak{s} \ltimes \mathfrak{r}$ where $\mathfrak{s}$ is semisimple (direct sum of simple algebras) and $\mathfrak{r}$ is the solvable radical (maximal solvable ideal). Cartan's criterion characterises semisimplicity by the non-degeneracy of the Killing form $B(x,y) = \operatorname{tr}(\operatorname{ad}_x \operatorname{ad}_y)$. This structural theorem parallels the Jordan--Hölder theorem for groups.

\term{Fundamental Theorem of Linear Algebra}
For an $m \times n$ matrix $A$, the four fundamental subspaces satisfy: (1) $\dim \operatorname{Nul}(A) + \dim \operatorname{Col}(A) = n$ (rank-nullity); (2) $\operatorname{Nul}(A) = \operatorname{Col}(A^T)^\perp$; (3) $\operatorname{Nul}(A^T) = \operatorname{Col}(A)^\perp$; (4) $\operatorname{rank}(A) = \operatorname{rank}(A^T)$. This theorem, due to Gilbert Strang, provides the geometric foundation for matrix analysis and is essential for understanding systems of linear equations.

\term{Fundamental Theorem of Symmetric Functions}
Every symmetric polynomial in $x_1, \ldots, x_n$ can be uniquely expressed as a polynomial in the elementary symmetric polynomials $e_k = \sum_{i_1 < \cdots < i_k} x_{i_1} \cdots x_{i_k}$. The ring of symmetric polynomials is $\mathbb{Z}[e_1, \ldots, e_n]$. This theorem underlies the theory of resultants, discriminants, and invariant theory. It is also the algebraic basis for expressing coefficients of polynomials in terms of their roots.

\term{Future Cone}
The future light cone at a point $p$ in Minkowski spacetime is the set $\{q : q - p \text{ is future-directed timelike or null}\}$. In coordinates with signature $(-,+,+,+)$, the future cone is $\{t > 0, t^2 > x^2 + y^2 + z^2\}$. The future cone defines the causal structure of spacetime: events in the future cone can be influenced by $p$. Light cones generalise to curved spacetimes in general relativity.
\term{Coequalizer}
The dual notion to the equalizer. The coequalizer of $f, g: A \to B$ is an object $C$ with $q: B \to C$ such that $q \circ f = q \circ g$, universal among such morphisms. In sets, the coequalizer is the quotient of $B$ by the equivalence relation generated by $f(a) \sim g(a)$ for all $a \in A$.
\term{Commutative Diagram}
A diagram in a category in which any two paths with the same endpoints give the same composed morphism. Commutative diagrams are the primary language for expressing universal properties, naturality, and exactness in algebra and topology.
\term{Feasible Region}
In optimization, the set of all points satisfying the constraints of a problem. For a linear program, the feasible region is a convex polyhedron, and the optimum occurs at a vertex. In nonlinear programming, the feasible region may be a more general subset of $\mathbb{R}^n$.
\term{Finite Difference}
An approximation to a derivative using values of a function at discrete points: the forward difference $\Delta f(x) = f(x+h) - f(x)$, the backward difference, and the central difference. Finite difference methods replace differential equations by recurrence relations on grids. They are foundational to numerical analysis and the discretization of PDEs.

\chapter{G}
\label{chap:g}

\term{Gabriel's Horn}
The surface of revolution of $y = 1/x$ for $x \geq 1$ about the $x$-axis. Its volume is finite (equal to $\pi$) but its surface area is infinite, so the horn can be filled with a finite amount of paint while no finite amount suffices to cover it. It illustrates that volume and area behave independently and connects to improper integrals and convergence.

\term{Galois Extension}
A normal, separable field extension; the fixed field of its full automorphism group.

\term{Galois Group}
The Galois group $\operatorname{Gal}(L/K)$ of a field extension $L/K$ is the group of field automorphisms of $L$ that fix $K$ pointwise. For a finite Galois extension, $|\operatorname{Gal}(L/K)| = [L:K]$. The Fundamental Theorem of Galois Theory gives an inclusion-reversing bijection between intermediate fields and subgroups. The Galois group of a polynomial is the Galois group of its splitting field; solvability by radicals corresponds to solvability of the Galois group.

\term{Galois Theory}
Galois theory, developed by \'Evariste Galois (1811--1832), links field theory and group theory. For a finite Galois extension $L/K$ with group $G = \operatorname{Gal}(L/K)$, intermediate fields $K \subseteq M \subseteq L$ correspond bijectively to subgroups $H \leq G$ via $M \mapsto \operatorname{Gal}(L/M)$ and $H \mapsto L^H$ (fixed field). Normal subgroups correspond to Galois subextensions. A polynomial is solvable by radicals iff its Galois group is solvable. This resolved the classical problem of solving polynomials of degree $\geq 5$.

\term{Game Theory}
Game theory studies strategic interactions between rational agents. A game consists of players, strategies, and payoffs. Nash equilibrium is a strategy profile where no player can improve by unilaterally deviating. Every finite game has a mixed-strategy Nash equilibrium. The minimax theorem for zero-sum games states $\max_x \min_y u(x,y) = \min_y \max_x u(x,y)$. Cooperative game theory studies coalitions and the core, Shapley value. Applications span economics, biology, computer science, and political science.

\term{Gamma Function}
The function $\Gamma(s) = \int_0^\infty t^{s-1} e^{-t}\, dt$ generalising the factorial: $\Gamma(n+1) = n!$.

\term{Gâteaux Derivative}
The directional derivative of a map between normed spaces, defined by a limit along a line.

\term{Gaussian Curvature}
The Gaussian curvature $K$ of a surface at a point is the product of the principal curvatures $\kappa_1 \kappa_2$. Equivalently, $K = \det(dN)$ where $N$ is the Gauss map. For a surface in $\mathbb{R}^3$, $K = \frac{LN - M^2}{EG - F^2}$ using the first and second fundamental forms. By Theorema Egregium, $K$ is intrinsic (depends only on the first fundamental form). The Gauss--Bonnet theorem: $\int_S K\,dA = 2\pi\chi(S)$. Positive curvature: sphere; zero: plane, cylinder; negative: saddle, hyperbolic plane.

\term{Gaussian Distribution (Normal Distribution)}
The normal distribution $\mathcal{N}(\mu, \sigma^2)$ has density $\frac{1}{\sqrt{2\pi\sigma^2}} e^{-(x-\mu)^2/(2\sigma^2)}$. It is the unique distribution with maximum entropy given mean and variance. The Central Limit Theorem: sums of i.i.d. random variables converge to normal. Standard normal has $\mu=0, \sigma=1$. Properties: symmetric, $68\%/95\%/99.7\%$ within $1/2/3$ standard deviations. Multivariate normal: $\mathcal{N}(\boldsymbol{\mu}, \Sigma)$ with density involving $\det(\Sigma)$ and Mahalanobis distance.

\term{Gaussian Elimination}
Gaussian elimination solves linear systems $Ax = b$ by reducing $A$ to row echelon form using elementary row operations: (1) swap rows, (2) multiply row by non-zero scalar, (3) add multiple of one row to another. The LU decomposition $A = LU$ (lower and upper triangular) captures the elimination steps. Pivoting (partial, complete) improves numerical stability. Complexity: $O(n^3)$ for $n \times n$ matrices. Back substitution then solves $Ux = L^{-1}b$.

\term{Gaussian Integers}
The Gaussian integers $\mathbb{Z}[i] = \{a+bi : a,b \in \mathbb{Z}\}$ form a Euclidean domain with norm $N(a+bi) = a^2+b^2$. They are a UFD; primes are either rational primes $p \equiv 3 \pmod{4}$ (remain prime), $p \equiv 1 \pmod{4}$ (split as $p = \pi \bar{\pi}$), or $2 = -i(1+i)^2$ (ramifies). The units are $\{\pm 1, \pm i\}$. Gaussian integers are used to prove Fermat's two-square theorem: $p = x^2+y^2$ iff $p \equiv 1 \pmod{4}$.

\term{Gauss--Bonnet Theorem}
The Gauss--Bonnet theorem for a compact 2-manifold $S$ without boundary: $\int_S K\,dA = 2\pi\chi(S)$, where $K$ is Gaussian curvature and $\chi$ is Euler characteristic. For a surface with boundary, $\int_S K\,dA + \int_{\partial S} k_g\,ds = 2\pi\chi(S)$ where $k_g$ is geodesic curvature. This links local geometry (curvature) to global topology (Euler characteristic). For polyhedral surfaces, curvature concentrates at vertices: angle defect $2\pi - \sum \theta_i$.

\term{Gauss's Law (Electrostatics)}
Gauss's law states that the electric flux through a closed surface is proportional to the enclosed charge: $\oint_S \mathbf{E} \cdot d\mathbf{A} = Q_{\text{enc}}/\varepsilon_0$. In differential form: $\nabla \cdot \mathbf{E} = \rho/\varepsilon_0$. This is one of Maxwell's equations. It implies the field of a point charge $Q$ is $\mathbf{E} = \frac{Q}{4\pi\varepsilon_0 r^2} \hat{\mathbf{r}}$ (Coulomb's law). Gauss's law is a powerful tool for calculating fields of symmetric charge distributions.

\term{Gauss's Theorem (Divergence Theorem)}
The divergence theorem (Gauss's theorem): $\int_V (\nabla \cdot \mathbf{F})\,dV = \oint_{\partial V} \mathbf{F} \cdot d\mathbf{S}$ for a vector field $\mathbf{F}$ on a volume $V$ with boundary $\partial V$. This relates volume integral of divergence to surface integral of flux. In coordinates: $\int_V \frac{\partial F^i}{\partial x^i}\,dV = \int_{\partial V} F^i n_i\,dS$. Special cases: Green's identities in PDE theory, Gauss's law in electromagnetism. Generalises to manifolds via Stokes' theorem.

\term{Generalized Eigenspace}
The subspace $\ker (A - \lambda I)^m$ for sufficiently large $m$, spanned by the generalized eigenvectors of $A$ for the eigenvalue $\lambda$.

\term{Generator}
An element whose orbit under a group action produces the whole set; also the single element generating a cyclic group.

\term{Genus}
The number of handles of a surface, or the dimension of the space of holomorphic differentials on a Riemann surface.

\term{Geodesic}
A geodesic is a curve $\gamma(t)$ that parallel-transports its own tangent vector: $\nabla_{\dot{\gamma}} \dot{\gamma} = 0$. In coordinates: $\ddot{\gamma}^k + \Gamma^k_{ij} \dot{\gamma}^i \dot{\gamma}^j = 0$. Geodesics are locally length-minimising curves (but not globally). On a sphere, great circles are geodesics; on a plane, straight lines. The exponential map $\exp_p(v) = \gamma_v(1)$ maps tangent vectors to geodesic endpoints. The geodesic flow is Hamiltonian on the cotangent bundle.

\term{Geometric Distribution}
The distribution of the number of trials needed for the first success in repeated Bernoulli trials, $P(X = k) = (1-p)^{k-1}p$.

\term{Geometric Mean}
The $n$-th root of the product of $n$ numbers, $(\prod a_i)^{1/n}$, appropriate for multiplicative data.

\term{Geometric Measure Theory}
Geometric measure theory studies geometric properties of sets via measure theory. Key objects: rectifiable sets, currents (generalised surfaces), varifolds (unoriented surfaces). Federer--Fleming deformation theorem: any integral current can be deformed into a polyhedral current. The coarea formula relates integrals over level sets to integrals over the domain. Applications: minimal surfaces, Plateau's problem, image processing, calculus of variations. The isoperimetric inequality: among sets of given volume, the ball minimises surface area.

\term{Geometric Series}
A series of the form $\sum_{n=0}^\infty ar^n$, converging to $a/(1-r)$ when $|r| < 1$.

\term{Geometry of Numbers}
Geometry of numbers, initiated by Minkowski, studies lattice points in convex bodies. Minkowski's theorem: a centrally symmetric convex body in $\mathbb{R}^n$ with volume $> 2^n \det(\Lambda)$ contains a non-zero lattice point of $\Lambda$. Applications: Minkowski's bound for class numbers, Dirichlet's unit theorem, lattice basis reduction (LLL algorithm). The successive minima $\lambda_i$ are the radii of the smallest balls containing $i$ linearly independent lattice points.

\term{Gergonne Point}
The point of concurrency of the lines from each vertex of a triangle to the touch point of the incircle on the opposite side.

\term{Gibbs Measure}
The Gibbs measure (or Gibbs distribution) on a configuration space is $\mu(\sigma) = \frac{1}{Z} e^{-\beta H(\sigma)}$ where $H$ is a Hamiltonian (energy function), $\beta = 1/kT$ is inverse temperature, and $Z = \sum e^{-\beta H}$ is the partition function. It maximises entropy for given expected energy. The Gibbs measure is the equilibrium state of a statistical mechanical system. Phase transitions correspond to non-analyticities in the free energy $-\frac{1}{\beta} \log Z$. Examples: Ising model, Potts model, lattice gas.

\term{Girard's Theorem}
Girard's theorem (spherical trigonometry): the area of a spherical triangle on a unit sphere is its spherical excess $A = \alpha + \beta + \gamma - \pi$, where $\alpha, \beta, \gamma$ are interior angles. On a sphere of radius $R$, area $= R^2(\alpha + \beta + \gamma - \pi)$. The sum of angles exceeds $\pi$; the excess is proportional to area. This contrasts with Euclidean geometry where $\alpha + \beta + \gamma = \pi$ always.

\term{Global Minimum}
The smallest value attained by a function over its entire domain.

\term{Gödel's Completeness Theorem}
Gödel's completeness theorem (1929): first-order logic is complete. A formula is provable from a set of axioms iff it is true in every model of those axioms. Equivalently, a theory is consistent iff it has a model. This is distinct from the incompleteness theorem (1931). The proof constructs a model from a maximal consistent set of formulas using the Lindenbaum--Henkin method. Compactness theorem: a set of sentences has a model iff every finite subset does.

\term{Gödel's Incompleteness Theorems}
First Incompleteness Theorem (1931): Any consistent formal system containing Peano arithmetic is incomplete: there are true statements about natural numbers unprovable in the system. Second Incompleteness Theorem: No such system can prove its own consistency. Gödel constructed a statement ``I am not provable'' using self-reference via G\"odel numbering. These theorems showed the limitations of Hilbert's program to formalise all mathematics.

\term{Goldbach's Conjecture}
Goldbach's conjecture (1742): every even integer $n \geq 4$ is the sum of two primes. The weak Goldbach conjecture (every odd $n \geq 7$ is sum of three primes) was proved by Helfgott (2013). The strong conjecture remains open. Vinogradov (1937) proved all sufficiently large odd integers are sum of three primes. Chen (1973) proved every sufficiently large even integer is sum of a prime and a number with at most two prime factors. Computational verification up to $4 \times 10^{18}$.

\term{Golden Ratio}
The number $\varphi = (1+\sqrt{5})/2 \approx 1.618$ satisfying $\varphi = 1 + 1/\varphi$, appearing throughout geometry and art.

\term{Gradient}
The gradient $\nabla f$ of a scalar function $f: \mathbb{R}^n \to \mathbb{R}$ is the vector of partial derivatives $(\partial f/\partial x_1, \ldots, \partial f/\partial x_n)$. It points in the direction of steepest ascent and its magnitude is the maximum directional derivative: $D_u f = \nabla f \cdot u$. For a differentiable function, $df = \nabla f \cdot dx$. The gradient is orthogonal to level sets. On a Riemannian manifold, $\nabla f$ is the vector field dual to $df$ via the metric: $\langle \nabla f, X \rangle = df(X)$.

\term{Gradient Descent}
The iterative optimization method $x_{k+1} = x_k - \alpha \nabla f(x_k)$ that moves opposite the gradient.

\term{Gradient Vector}
The vector of first partial derivatives $\nabla f = (\partial f/\partial x_1, \ldots, \partial f/\partial x_n)$, pointing in the direction of steepest ascent.

\term{Graph}
A mathematical structure of vertices and edges, either undirected or directed, modelling pairwise relations.

\term{Graph Laplacian}
The matrix $L = D - A$ for a graph, whose spectrum encodes connectivity and is used in spectral graph theory and clustering.

\term{Graph Theory}
The branch of mathematics studying graphs and their properties, applications, and algorithms.

\term{Grassmannian}
The Grassmannian $\operatorname{Gr}(k,n)$ (or $\operatorname{Gr}(k,V)$) is the space of all $k$-dimensional subspaces of an $n$-dimensional vector space $V$. It is a smooth projective variety of dimension $k(n-k)$. The Pl\"ucker embedding $\operatorname{Gr}(k,n) \hookrightarrow \mathbb{P}(\bigwedge^k V)$ maps a subspace to its Pl\"ucker coordinates (minors of a $k \times n$ matrix). The cohomology ring $H^*(\operatorname{Gr}(k,n); \mathbb{Z})$ is generated by Schubert classes. Grassmannians classify vector bundles and appear in quantum cohomology and mirror symmetry.

\term{Gray Code}
A Gray code is an ordering of $2^n$ binary strings of length $n$ such that adjacent strings differ in exactly one bit. The binary-reflected Gray code is constructed recursively: $G_1 = (0,1)$; $G_{n+1} = (0G_n, 1 \text{reverse}(G_n))$. Gray codes minimise transitions in analog-to-digital conversion, solve the Chinese rings puzzle, and generate Hamiltonian cycles on hypercubes. The $k$-th Gray code is $k \oplus \lfloor k/2 \rfloor$.

\term{Green's Function}
The Green's function $G(x,y)$ for a linear differential operator $L$ satisfies $L G(x,y) = \delta(x-y)$ with appropriate boundary conditions. The solution to $L u = f$ is $u(x) = \int G(x,y) f(y)\,dy$. For the Laplacian on $\mathbb{R}^n$, $G(x,y) = \frac{1}{n(n-2)\omega_n} |x-y|^{2-n}$ ($n \geq 3$) or $\frac{1}{2\pi} \log |x-y|$ ($n=2$). Green's functions are fundamental in PDE theory, potential theory, and quantum field theory (propagators).

\term{Green's Identity}
Green's identity is obtained by combining the divergence theorem with the product rule: $\iint_D (u \Delta v)\,dA = \oint_{\partial D} u (\nabla v) \cdot \mathbf{n}\,ds - \iint_D \nabla u \cdot \nabla v\,dA$. It relates a double integral of a Laplacian to boundary integrals, converting domain data into boundary data. This makes it central to boundary value problems for the Laplacian, where it produces weak formulations and uniqueness arguments. It also underlies the construction and use of Green's functions.

\term{Green's Theorem}
Green's theorem in the plane: $\oint_C (P\,dx + Q\,dy) = \iint_D \left(\frac{\partial Q}{\partial x} - \frac{\partial P}{\partial y}\right) dx\,dy$ for a simple closed curve $C$ bounding region $D$. This is the 2D case of Stokes' theorem. The circulation form relates line integral of a vector field to area integral of its curl: $\oint_C \mathbf{F} \cdot d\mathbf{r} = \iint_D (\nabla \times \mathbf{F}) \cdot \mathbf{k}\,dA$. The flux form: $\oint_C \mathbf{F} \cdot \mathbf{n}\,ds = \iint_D \nabla \cdot \mathbf{F}\,dA$. Fundamental in vector calculus and fluid dynamics.

\term{Gromov–Hausdorff Distance}
A distance between compact metric spaces, measuring how far two spaces are from being isometric.

\term{Group Action}
A group action of $G$ on a set $X$ is a map $G \times X \to X$, $(g,x) \mapsto g \cdot x$, satisfying $e \cdot x = x$ and $(gh) \cdot x = g \cdot (h \cdot x)$. Equivalently, a homomorphism $G \to \operatorname{Sym}(X)$. Orbit of $x$: $G \cdot x = \{g \cdot x : g \in G\}$; stabiliser: $G_x = \{g : g \cdot x = x\}$. Orbit-stabiliser theorem: $|G \cdot x| = [G : G_x]$. The class equation: $|X| = \sum [G : G_{x_i}]$. Burnside's lemma: number of orbits $= \frac{1}{|G|} \sum_{g \in G} |X^g|$.

\term{Group (Algebra)}
A group $(G, \cdot)$ is a set with an associative binary operation, an identity element $e$, and inverses: $\forall g \in G,\; \exists g^{-1} : g \cdot g^{-1} = e$. Groups capture symmetry. Examples: symmetric group $S_n$, cyclic group $\mathbb{Z}/n\mathbb{Z}$, general linear group $\operatorname{GL}(n,\mathbb{R})$, fundamental group $\pi_1(X)$. Finite simple groups classified into 18 infinite families and 26 sporadic groups. The classification theorem (2004) is a monumental achievement.

\term{Group Cohomology}
Group cohomology $H^*(G, M)$ for a group $G$ and $G$-module $M$ is derived functors of the fixed-point functor $M \mapsto M^G$. It can be computed from the bar resolution or the standard resolution. $H^1(G, M) = Z^1/B^1$ classifies crossed homomorphisms modulo principal ones. $H^2(G, M)$ classifies group extensions $1 \to M \to E \to G \to 1$ with $M$ abelian. Tate cohomology $\hat{H}^*(G, M)$ combines homology and cohomology for finite groups.

\term{Group Homomorphism}
A map between groups preserving the operation, $f(ab) = f(a)f(b)$.

\term{Group Representation}
A representation of a group $G$ on a vector space $V$ is a homomorphism $\rho: G \to \operatorname{GL}(V)$. Characters $\chi(g) = \operatorname{tr}(\rho(g))$ are central. Irreducible representations of finite groups are finite-dimensional and completely reducible (Maschke's theorem). The character table encodes all irreducible characters. Orthogonality relations: $\frac{1}{|G|} \sum_{g \in G} \chi_i(g) \overline{\chi_j(g)} = \delta_{ij}$. Representations of Lie groups and algebras are central to physics (standard model, particle classification).

\term{Group Theory}
The study of groups and their actions, symmetries, and representations.

\term{Grothendieck Duality}
The generalisation of Serre duality to coherent sheaves on arbitrary schemes, via a dualising complex.

\term{Grothendieck Group}
The Grothendieck group $K_0(\mathcal{C})$ of an additive category $\mathcal{C}$ is the abelian group generated by isomorphism classes $[X]$ of objects with relations $[X] = [Y] + [Z]$ for every short exact sequence $0 \to Y \to X \to Z \to 0$. For a ring $R$, $K_0(R)$ is built from finitely generated projective modules. $K_0(\text{Var}_k)$ is generated by varieties with $[X] = [Y] + [X \setminus Y]$ for closed $Y \subset X$. Grothendieck groups are universal for additive invariants (Euler characteristic, rank, Chern classes).

\term{Grothendieck Topology}
A Grothendieck topology on a category $\mathcal{C}$ specifies covering families $\{U_i \to U\}$ satisfying: (1) isomorphisms are covers, (2) pullbacks of covers are covers, (3) composites of covers are covers. This generalises the notion of open covers from topological spaces. A site is a category with a Grothendieck topology. Sheaves on sites define cohomology (étale, fppf, fpqc, crystalline). Grothendieck topologies are the foundation of modern algebraic geometry and the Weil conjectures.

\term{Grothendieck's Six Operations}
Grothendieck's six operations in derived categories of sheaves: $f_*$ (pushforward), $f^*$ (pullback), $f_!$ (proper pushforward), $f^!$ (exceptional pullback), $\otimes$ (tensor product), $\mathcal{H}om$ (internal hom). These functors satisfy adjunctions: $f^* \dashv f_*$, $f_! \dashv f^!$. The formalism unifies Poincar\'e duality, Verdier duality, and the Grothendieck trace formula. It is the language of modern cohomology theories (étale, $\ell$-adic, perverse sheaves).

\term{Grunsky Coefficients}
The Grunsky coefficients $c_{mn}$ of a univalent function $f(z) = z + a_2 z^2 + \cdots$ on the unit disk are defined by $\log \frac{f(z) - f(w)}{z-w} = -\sum_{m,n \geq 1} c_{mn} z^m w^n$. The Grunsky inequalities state that the matrix $(c_{mn})$ is positive semi-definite. They provide necessary and sufficient conditions for univalence and are used in the proof of the Bieberbach conjecture ($|a_n| \leq n$ for univalent functions).

\term{G\"urtler's Theorem}
G\"urtler's theorem in differential geometry: a complete Riemannian manifold with non-positive sectional curvature is diffeomorphic to $\mathbb{R}^n$ (Hadamard--Cartan theorem). More generally, if the manifold has no conjugate points, the universal cover is diffeomorphic to $\mathbb{R}^n$. This relates global topology to curvature assumptions. The theorem is fundamental in comparison geometry and the study of CAT(0) spaces.

\term{Gyrovector Space}
A gyrovector space (Ungar, 2001) is an algebraic structure for hyperbolic geometry analogous to vector spaces for Euclidean geometry. It has ``gyrovector addition'' $\oplus$ (not associative but gyroassociative) and scalar multiplication. The gyrogroup axioms include the gyrocommutative law $a \oplus b = \operatorname{gyr}[a,b](b \oplus a)$ where $\operatorname{gyr}$ is the Thomas precession. Gyrovector spaces provide a vector-like formalism for Einstein's velocity addition in special relativity.

\chapter{H}
\label{chap:h}

\term{Haar Measure}
A translation-invariant measure on a locally compact topological group. Such a measure exists and is unique up to a positive scalar factor.

\term{Hadamard Factorization Theorem}
The theorem stating that any entire function of finite order can be expressed as a product involving its zeros times an exponential polynomial. It generalises the Weierstrass product theorem.

\term{Hadamard Matrix}
A square matrix whose entries are either $+1$ or $-1$ and whose rows are mutually orthogonal. Such matrices are used in error-correcting codes and signal processing.

\term{Hadamard Product}
The element-wise product of two matrices of the same dimensions, where $(A \circ B)_{ij} = A_{ij} B_{ij}$. It is also known as the Schur product.

\term{Hahn-Banach Theorem}
A fundamental result in functional analysis stating that a bounded linear functional on a subspace of a normed vector space can be extended to the whole space without increasing its norm.

\term{Hairy Ball Theorem}
Every continuous tangent vector field on an even-dimensional sphere must vanish somewhere; equivalently, there is no continuous non-vanishing field of tangent vectors on $S^{2n}$. It is often summarized as ``you cannot comb a hairy ball flat.'' A consequence is that at any moment the wind must be calm at some point on Earth. Odd-dimensional spheres do admit non-vanishing tangent vector fields.

\term{Half-plane}
A region in $\mathbb{R}^2$ bounded by a line, consisting of all points on one side of the line.

\term{Half-space}
A region in $\mathbb{R}^n$ bounded by a hyperplane.

\term{Hall's Marriage Theorem}
A combinatorial theorem giving a necessary and sufficient condition for a finite family of sets to have a system of distinct representatives: the union of any $k$ sets must contain at least $k$ elements.

\term{Halting Problem}
The problem of deciding, given a program and its input, whether the program eventually halts. Turing proved that it is undecidable: no algorithm solves it in all cases. The proof uses a diagonalization argument, constructing a program that halts exactly when the purported decider claims it does not. It is the classic undecidable problem and the starting point for the theory of computability.

\term{Hamiltonian (Mechanics)}
A function $H(p, q, t)$ representing the total energy of a system. The evolution of the system is governed by Hamilton's equations: $\dot{q} = \frac{\partial H}{\partial p}$ and $\dot{p} = -\frac{\partial H}{\partial q}$.

\term{Hamiltonian Operator}
In quantum mechanics, the operator $\hat{H}$ corresponding to the total energy of the system. The time-independent Schrödinger equation is $\hat{H} \psi = E \psi$.

\term{Hamiltonian System}
A dynamical system governed by Hamilton's equations $\dot{q} = \frac{\partial H}{\partial p}$ and $\dot{p} = -\frac{\partial H}{\partial q}$.

\term{Hamming Code}
A family of linear error-correcting codes that can detect up to two-bit errors and correct one-bit errors. The (7,4) Hamming code is a prototypical example.

\term{Hamming Distance}
The number of positions at which the corresponding symbols of two strings of equal length are different. It is used in information theory to measure error rates.

\term{Hankel Transform}
An integral transform of the form $F(r) = \int_0^\infty f(t) J_n(rt)\, t\, dt$, where $J_n$ is a Bessel function of the first kind. It is the Fourier transform adapted to radial symmetry.

\term{Harmonic Analysis}
The study of how signals or functions can be represented as sums of basic waves (sinusoids), primarily via Fourier series and Fourier transforms.

\term{Harmonic Function}
A function $u$ that satisfies Laplace's equation $\Delta u = 0$. Such functions are smooth and satisfy the mean value property.

\term{Harmonic Mean}
The reciprocal of the arithmetic mean of the reciprocals of a set of values: $H = n / (\sum 1/x_i)$. It is typically used for average rates.

\term{Harmonic Number}
The sum of the reciprocals of the first $n$ positive integers: $H_n = \sum_{k=1}^n 1/k$. The harmonic numbers grow like $\ln n + \gamma$, where $\gamma$ is Euler's constant.

\term{Harmonic Series}
The divergent infinite series $\sum_{n=1}^\infty \frac{1}{n}$. Although the terms tend to zero, the sum grows logarithmically.

\term{Harnack's Inequality}
An inequality bounding the ratio of values of a positive harmonic function at two points in terms of a constant depending only on the domain. It implies that positive harmonic functions on connected domains satisfy a strong form of rigidity.

\term{Hausdorff Dimension}
A way of assigning a dimension to a set (including fractals) that can be non-integer. It is defined via the $s$-dimensional Hausdorff measure.

\term{Hausdorff Distance}
A measure of how far two subsets of a metric space are from each other. It is the maximum distance from a point in one set to the nearest point in the other set.

\term{Hausdorff Paradox}
A unit sphere $S^2$ can be decomposed into finitely many disjoint pieces which, by rotations alone, can be reassembled into two copies of the sphere (one possibly together with a countable set). It is a precursor of the Banach--Tarski paradox and similarly depends on the axiom of choice and non-measurable sets.

\term{Hausdorff Space}
A topological space in which any two distinct points can be separated by disjoint open neighborhoods. Most "natural" spaces in analysis and geometry are Hausdorff.

\term{Heat Equation}
The partial differential equation $\frac{\partial u}{\partial t} = \alpha \nabla^2 u$ describing the diffusion of heat through a medium.

\term{Heaviside Step Function}
A function $\text{H}(x)$ that is $0$ for $x < 0$ and $1$ for $x \geq 0$. It is used to model switches in circuit theory and is the integral of the Dirac delta function.

\term{Hedgehog Space}
A topological space formed by attaching several copies of the unit interval to a single common endpoint. It is used as a counterexample in the study of metric spaces.

\term{Heegner Number}
A number $d$ such that the imaginary quadratic field $\mathbb{Q}(\sqrt{-d})$ has class number 1. There are exactly nine such numbers: 1, 2, 3, 7, 11, 19, 43, 67, 163.

\term{Heisenberg Group}
A Lie group of $3 \times 3$ upper triangular matrices with ones on the diagonal. It is the simplest non-abelian nilpotent Lie group and is central to quantum mechanics.

\term{Helly's Theorem}
A result in convex geometry stating that if every $d+1$ elements of a collection of convex sets in $\mathbb{R}^d$ have a common intersection, then the entire collection has a common intersection.

\term{Hensel's Lemma}
A result in number theory stating that if a polynomial has a simple root modulo a prime $p$, then it has a unique root modulo $p^k$ for every $k$, obtained by successive lifting.

\term{Hereditary Property}
A property of a class of graphs or structures such that if a structure has the property, then all its substructures (e.g., induced subgraphs) also have it.

\term{Hereditary Ring}
A ring where every ideal is a projective module. In such rings, the global dimension is at most 1.

\term{Hermite Normal Form}
A canonical integer matrix form obtained by column operations, in which the matrix is upper triangular with non-negative off-diagonal entries. It is used in lattice theory and integer linear programming.

\term{Hermite Polynomials}
A sequence of orthogonal polynomials $H_n(x)$ that arise in the solution of the quantum harmonic oscillator and in probability theory.

\term{Hermitian Form}
A sesquilinear form $h$ on a complex vector space satisfying $h(u,v) = \overline{h(v,u)}$. A positive definite Hermitian form is an inner product.

\term{Hermitian Matrix}
A complex square matrix $A$ that is equal to its own conjugate transpose: $A = A^*$. All eigenvalues of a Hermitian matrix are real.

\term{Hessian Matrix}
The square matrix of second-order partial derivatives of a scalar-valued function. It is used to determine if a critical point is a local maximum, minimum, or saddle point.

\term{Heuristics}
Practical methods or "rules of thumb" used in problem solving that may not be guaranteed to find an optimal solution but often find a "good enough" one in a reasonable time.

\term{Hilbert Basis Theorem}
The theorem that the polynomial ring $R[x_1, \ldots, x_n]$ over a Noetherian ring $R$ is itself Noetherian.

\term{Hilbert Curve}
A space-filling curve that maps the unit interval $[0,1]$ onto the unit square $[0,1]^2$ continuously. It is a fractal curve.

\term{Hilbert Space}
A complete inner product space. They generalize Euclidean space to infinite dimensions and are the primary setting for quantum mechanics.

\term{Hilbert Transform}
An integral transform that shifts the phase of every frequency component of a signal by $90^\circ$. It is used to create the analytic signal in communications.

\term{Hilbert's Nullstellensatz}
A fundamental theorem in algebraic geometry relating ideals in polynomial rings to the zero-sets (varieties) they define. It states that every ideal in $k[x_1, \ldots, x_n]$ is the radical of the ideal of its zero-set.

\term{Hilbert's Tenth Problem}
The problem of determining whether a given Diophantine equation has a solution in integers. This was proven undecidable by Yuri Matiyasevich in 1970.

\term{Hilbert-Schmidt Norm}
The norm of a linear operator $T$ on a Hilbert space defined by $\|T\|_{HS}^2 = \sum \|Te_i\|^2$ for any orthonormal basis $\{e_i\}$.

\term{Hilbert-Schmidt Operator}
A bounded linear operator whose Hilbert-Schmidt norm is finite. Such operators are always compact.

\term{Hille-Yosida Theorem}
A theorem characterising the generators of strongly continuous semigroups of contractions on a Banach space in terms of resolvent bounds.

\term{H-index}
A metric that measures both the productivity and citation impact of a publication, where a scientist has index $h$ if $h$ of their papers have at least $h$ citations each.

\term{Hodge Star Operator}
A linear isomorphism between $k$-forms and $(n-k)$-forms on an $n$-dimensional oriented Riemannian manifold. It is used to define the adjoint of the exterior derivative.

\term{Hodge Theory}
The study of the cohomology of compact Kähler manifolds. The Hodge decomposition $H^k(X, \mathbb{C}) = \bigoplus_{p+q=k} H^{p,q}(X)$ relates the topology of the manifold to its complex structure.

\term{Hodge-Laplacian}
The operator $\Delta = d\delta + \delta d$ acting on differential forms. The kernel of this operator consists of the harmonic forms.

\term{Hofstadter's Butterfly}
A fractal pattern representing the energy spectrum of an electron in a 2D lattice subjected to a magnetic field.

\term{Hölder Space}
A space of functions whose differences are controlled by a power of the distance, i.e.\ functions satisfying $|f(x) - f(y)| \leq C |x - y|^\alpha$ for some $0 < \alpha \leq 1$.

\term{Hölder's Inequality}
The inequality $\int |fg| \leq \|f\|_p \|g\|_q$ for $1/p + 1/q = 1$, $p, q \geq 1$. It generalises the Cauchy--Schwarz inequality, which is the case $p = q = 2$.

\term{Holomorphic Function}
A complex-valued function of a complex variable that is complex-differentiable in a neighborhood of every point in its domain.

\term{Holonomy}
The group of linear transformations obtained by parallel transport of tangent vectors around closed loops on a manifold. It encodes curvature.

\term{Homeomorphism}
A continuous bijection $f: X \to Y$ whose inverse $f^{-1}$ is also continuous. Two spaces are homeomorphic if there exists a homeomorphism between them, meaning they are topologically identical.

\term{Homogeneous Coordinate}
A system of coordinates used in projective geometry to represent a point in $\mathbb{P}^n$ as a vector in $\mathbb{R}^{n+1}$, where two vectors represent the same point if they are scalar multiples.

\term{Homogeneous Polynomial}
A polynomial where every term has the same total degree. For example, $x^2 + xy + y^2$ is homogeneous of degree 2.

\term{Homogeneous Space}
A space with a transitive action of a Lie group, so that any point can be moved to any other by a group element. Examples include spheres and Grassmannians.

\term{Homological Algebra}
The branch of algebra studying chain complexes, exact sequences, and derived functors, unifying algebraic topology and homological invariants.

\term{Homology}
A tool in algebraic topology used to count the "holes" of different dimensions in a space. For a CW complex, the $n$-th homology group $H_n(X)$ describes the $n$-dimensional cycles that are not boundaries.

\term{Homomorphism}
A map between two algebraic structures (like groups, rings, or vector spaces) that preserves the operations. For groups $G$ and $H$, $f(ab) = f(a)f(b)$.

\term{Homotopy}
A continuous deformation between two continuous maps $f, g: X \to Y$. Two maps are homotopic if one can be continuously transformed into the other.

\term{Homotopy Category}
The category whose objects are topological spaces and whose morphisms are homotopy classes of continuous maps.

\term{Homotopy Group}
The group $\pi_n(X)$ of based homotopy classes of maps $S^n \to X$. The fundamental group $\pi_1$ is the case $n = 1$.

\term{Hopf Algebra}
A bialgebra equipped with an antipode, generalizing the group algebra of a group. They are central to the study of quantum groups and tensor categories.

\term{Hopf Fibration}
A fiber bundle $S^3 \to S^2$ where the fibers are circles $S^1$. It provides a way to view the 3-sphere as a collection of circles over a 2-sphere.

\term{Hopf Link}
The simplest non-trivial link, consisting of two circles interlocked once.

\term{Hopf-Rinow Theorem}
A theorem stating that for a Riemannian manifold, geodesic completeness is equivalent to metric completeness.

\term{Hopfian Group}
A group $G$ that is not isomorphic to any of its proper quotients. Every finite group is Hopfian.

\term{Horned Sphere}
A topological space that is homeomorphic to a 3-sphere but whose boundary is not a smooth 2-sphere, serving as a counterexample to several intuitive conjectures in topology.

\term{H-space}
A topological space $X$ equipped with a continuous binary operation that is homotopic to a group operation. Examples include loop spaces $\Omega Y$.

\term{Hurewicz Theorem}
A theorem relating the first non-vanishing homotopy group of a space to its homology group: if $\pi_i(X) = 0$ for $i < n$ with $n \geq 2$, then $\pi_n(X) \cong H_n(X)$.

\term{Hurwitz Zeta Function}
The function $\zeta(s, a) = \sum_{n=0}^\infty (n + a)^{-s}$, which generalises the Riemann zeta function $\zeta(s) = \zeta(s, 1)$.

\term{Hyperbolic Cosine (cosh)}
The hyperbolic function defined as $\cosh x = \frac{e^x + e^{-x}}{2}$. It satisfies $\cosh^2 x - \sinh^2 x = 1$.

\term{Hyperbolic Equation}
A second-order PDE whose characteristic equation has two distinct real characteristic directions, such as the wave equation. Hyperbolic problems admit well-posed initial value problems.

\term{Hyperbolic Function}
The functions $\sinh$, $\cosh$, and $\tanh$, defined in terms of exponentials and analogous to the trigonometric functions. They are the functions of a unit hyperbola rather than a circle.

\term{Hyperbolic Geometry}
A non-Euclidean geometry where the parallel postulate is replaced: through a point not on a line, there are at least two distinct lines parallel to the given line.

\term{Hyperbolic Group}
A finitely generated group whose Cayley graph is a hyperbolic metric space in the sense of Gromov.

\term{Hyperbolic Paraboloid}
A doubly ruled surface that looks like a saddle. It is defined by the equation $z = x^2/a^2 - y^2/b^2$.

\term{Hyperbolic Sine (sinh)}
The hyperbolic function defined as $\sinh x = \frac{e^x - e^{-x}}{2}$. It is the analogue of the sine function for hyperbolic geometry.

\term{Hyperbolic Space}
The standard model for hyperbolic geometry, such as the Poincaré disk or the upper half-plane. It has constant sectional curvature $-1$.

\term{Hyperbolic Tangent (tanh)}
The hyperbolic function $\tanh x = \frac{\sinh x}{\cosh x}$, used extensively as an activation function in neural networks.

\term{Hypercube}
The $n$-dimensional generalisation of a square and a cube, defined as the product of $n$ copies of the unit interval.

\term{Hyperelliptic Curve}
An algebraic curve of genus at least 2 that admits a degree-two map to the projective line. They generalise elliptic curves.

\term{Hypergeometric Function}
The function ${}_2F_1(a, b; c; z) = \sum_{n=0}^\infty \frac{(a)_n (b)_n}{(c)_n n!} z^n$, where $(x)_n$ is the Pochhammer symbol. It specialises to many elementary functions and orthogonal polynomials.

\term{Hypergraph}
A generalisation of a graph in which edges, called hyperedges, may connect any number of vertices.

\term{Hyperplane}
An $(n-1)$-dimensional subspace of an $n$-dimensional vector space. In $\mathbb{R}^3$, a hyperplane is a 2D plane.

\term{Hyperplane Arrangement}
A finite collection of hyperplanes in a vector space. The study of the topology of the complement of such an arrangement is central to the theory of braid groups.

\term{Hyperreal Number}
An extension of the real numbers that includes infinitesimals and infinite numbers, used in non-standard analysis.

\term{Hypersphere}
The set of all points in $\mathbb{R}^n$ at a constant distance from a fixed center.

\term{Hypersurface}
A manifold of dimension $n-1$ embedded in a manifold of dimension $n$.

\term{Hypervolume}
The $n$-dimensional volume of a region in $\mathbb{R}^n$ for $n > 3$.

\chapter*{I}
\label{chap:i}

\term{Ideal (Ring Theory)}
An ideal $I$ of a ring $R$ is an additive subgroup closed under multiplication by elements of $R$. A left ideal requires $ra\in I$, a right ideal requires $ar\in I$, and a two-sided ideal satisfies both. The quotient $R/I$ has a ring structure exactly when $I$ is two-sided. \par\noindent\textbf{Applications:} ideals describe divisibility, kernels of homomorphisms, and solution sets of polynomial equations.

\term{Ideal Class Group}
The ideal class group is the group of fractional ideals of a number field modulo principal fractional ideals. It measures the failure of unique factorization of elements: a trivial class group implies unique factorization of ideals and, in important cases, of elements. Its finite order is the class number. \par\noindent\textbf{Applications:} class groups organize arithmetic in number fields and appear in class field theory, Diophantine equations, and cryptography.

\term{Ideal Sheaf}
The sheaf $\mathcal{I}_Y$ on a space $X$ associated with a closed subscheme $Y$, whose sections over an open set $U$ are the functions in $\mathcal{O}_X(U)$ that vanish on $Y \cap U$. It is the kernel of the restriction map $\mathcal{O}_X \to i_* \mathcal{O}_Y$.

\term{Idempotent Element}
An element $e$ is idempotent if $e^2=e$. Besides $0$ and $1$, a ring may have idempotents that split its structure: for a central idempotent, $R\cong eR\times(1-e)R$. An idempotent matrix is a projection. \par\noindent\textbf{Applications:} idempotents decompose rings and modules and represent repeated decisions or projections in algebra and computation.

\term{Identity Element}
An element $e$ in a set $S$ with a binary operation $\cdot$ such that $e \cdot a = a \cdot e = a$ for all $a \in S$. In group theory, the identity is unique. In category theory, an identity morphism $\text{id}_X: X \to X$ exists for every object.

\term{Identity Theorem}
If two holomorphic functions on a connected domain agree on a set having a limit point inside the domain, they agree everywhere. Thus a holomorphic function is determined by surprisingly sparse data, and analytic continuation is unique. \par\noindent\textbf{Applications:} the theorem proves uniqueness in complex analysis and helps extend solutions of analytic differential equations.

\term{Idempotent Matrix}
A square matrix $P$ is idempotent when $P^2=P$. It projects vectors onto its image along its kernel; its only possible eigenvalues are $0$ and $1$, so it is diagonalizable. \par\noindent\textbf{Applications:} projection matrices are used in least squares, statistics, numerical linear algebra, and decomposing vector spaces.

\term{Iff}
Abbreviation for "if and only if", denoting logical equivalence between two statements.

\term{Image}
The image of $A$ under $f$ is $f(A)=\{f(x):x\in A\}$, the set of values actually attained. It is different from the whole codomain when $f$ is not onto. \par\noindent\textbf{Applications:} images describe reachable states, solution sets, and the range of linear transformations.

\term{Imaginary Number}
A purely imaginary number has the form $bi$ with $b\in\mathbb R$ and $i^2=-1$. It lies on the imaginary axis of the complex plane. Together with real numbers, imaginary numbers form the complex numbers $a+bi$. \par\noindent\textbf{Applications:} complex numbers solve polynomial equations and represent oscillations, AC circuits, waves, and Fourier modes.

\term{Imaginary Unit}
The number $i$ defined by $i^2 = -1$. It allows for the construction of the field of complex numbers $\mathbb{C}$. In physics, $i$ often appears in the Schrödinger equation and Fourier transforms.

\term{Immersion}
A smooth map $f:M\to N$ is an immersion if its differential $df_p$ is injective at every point. Locally it preserves tangent directions, but its image may cross itself; an embedding additionally requires suitable global one-to-one behavior. \par\noindent\textbf{Applications:} immersions describe curves and surfaces in differential geometry, computer graphics, and the kinematics of constrained systems.

\term{Implication}
A logical connective $P \Rightarrow Q$ true except when $P$ is true and $Q$ is false.

\term{Improper Fraction}
A fraction whose numerator is greater than or equal to its denominator.

\term{Improper Integral}
An improper integral has an unbounded interval or an unbounded integrand. For example, $\int_a^\infty f(x)\,dx=\lim_{b\to\infty}\int_a^b f(x)\,dx$ when the limit exists. Convergence must be checked rather than assumed. \par\noindent\textbf{Applications:} improper integrals measure total probability, energy, mass, and long-time effects in applied models.

\term{Implicit Function Theorem}
If $F:\mathbb R^{n+m}\to\mathbb R^m$ is smooth, $F(x_0,y_0)=0$, and the Jacobian $\partial F/\partial y$ is invertible there, then locally the equation $F(x,y)=0$ defines a unique smooth function $y=g(x)$. It explains when dependent variables can be solved from constraints. \par\noindent\textbf{Applications:} it supports coordinate descriptions of surfaces, equilibrium analysis, constrained optimization, and nonlinear equations.

\term{In-degree}
In a directed graph, the in-degree of a vertex $v$ is the number of edges pointing into $v$. In a regular tournament, all vertices have the same in-degree. The sum of in-degrees equals the sum of out-degrees, both equaling the total number of edges.

\term{Incenter}
The center of the incircle of a triangle, formed by the intersection of the internal angle bisectors. It is always equidistant from the three sides. The incenter is always located in the interior of the triangle.

\term{Incircle}
The largest circle that can be inscribed within a triangle, tangent to all three sides. Its radius $r$ is given by $r = \text{Area}/s$ where $s$ is the semi-perimeter.

\term{Independence}
Events $A$ and $B$ are independent when $P(A\cap B)=P(A)P(B)$, meaning learning one does not change the probability of the other. Random variables are independent when their joint distribution factors into the product of their marginal distributions. \par\noindent\textbf{Applications:} independence simplifies probability models, sampling, reliability calculations, and statistical inference.

\term{Inscribed Angle}
An angle formed by two chords that share an endpoint on a circle's circumference. An inscribed angle is always half the measure of the central angle subtending the same arc.

\term{Inscribed Circle}
Also known as the incircle. A circle that is tangent to all sides of a polygon. Only some polygons (like all triangles) possess an inscribed circle.

\term{Index (Group Theory)}
The index $[G:H]$ is the number of left or right cosets of a subgroup $H$ in $G$. If $G$ is finite, Lagrange's theorem gives $[G:H]=|G|/|H|$. \par\noindent\textbf{Applications:} index counts symmetry classes, orbit sizes, and the number of states obtained after identifying transformations in a subgroup.

\term{Index (Complex Analysis)}
The index, or winding number, of a closed curve $\gamma$ around $a\notin\gamma$ is
\[
\operatorname{ind}_\gamma(a)=\frac{1}{2\pi i}\oint_\gamma\frac{dz}{z-a}.
\]
It counts signed turns around the point. \par\noindent\textbf{Applications:} winding numbers support the argument principle, contour integration, and the study of zeros and poles of complex functions.

\term{Index (Fixed Point)}
A topological invariant assigned to a fixed point of a map $f: M \to M$. The sum of indices of all fixed points equals the Euler characteristic $\chi(M)$ by the Poincaré--Hopf Theorem.

\term{Indicator Function}
The indicator function $\mathbb 1_A(x)$ is $1$ when $x\in A$ and $0$ otherwise. Its integral records the measure of the set: $\int\mathbb1_A\,d\mu=\mu(A)$. \par\noindent\textbf{Applications:} indicators encode events, define piecewise functions, and simplify probability, measure theory, statistics, and optimization models.

\term{Inductive Proof}
A method of proof where a statement $P(n)$ is proved for a base case $n_0$, and then the implication $P(k) \implies P(k+1)$ is proved. This establishes $P(n)$ for all $n \geq n_0$.

\term{Induction (Mathematical)}
Mathematical induction proves a statement for every natural number by establishing a base case and showing $P(n)\Rightarrow P(n+1)$. Strong induction permits using all earlier cases $P(m)$ for $m<n$. \par\noindent\textbf{Applications:} induction proves formulas, recursive algorithms, divisibility results, and correctness of constructions defined step by step.

\term{Inductive Limit}
In category theory, the colimit of a directed system of objects. For a sequence of nested spaces $X_1 \subseteq X_2 \subseteq \cdots$, the inductive limit is their union $\bigcup X_i$. It is the categorical dual of the projective limit.

\term{Inertia Tensor}
The inertia tensor
\[
I_{ij}=\int\rho(r^2\delta_{ij}-r_ir_j)\,dV
\]
describes how mass is distributed relative to rotation axes. It relates angular momentum and angular velocity by $\mathbf L=I\boldsymbol\omega$; its eigenvectors are principal axes. \par\noindent\textbf{Applications:} it predicts rotation of rigid bodies, spacecraft attitude, robotics motion, and molecular dynamics.

\term{Inequation}
A statement asserting that two expressions are not equal, typically using $\leq$ or $\geq$. Solving an inequation involves finding the set of values for which the inequality holds.

\term{Inference}
The derivation of a conclusion from premises according to the rules of logic.

\term{Infimum}
The infimum $\inf S$ is the greatest number that is no larger than every element of $S$. It need not belong to $S$: the infimum of $(0,1)$ is $0$. \par\noindent\textbf{Applications:} infima define limits and optimization objectives when a minimum is not attained and are central to completeness and variational methods.

\term{Infinite Set}
A set that is not finite; it cannot be placed in bijection with any $\{1, \ldots, n\}$. A set is infinite iff it has a proper subset with the same cardinality (Dedekind-infinite).

\term{Infinite Series}
An infinite series $\sum_{n=1}^\infty a_n$ is defined through its partial sums $S_N=\sum_{n=1}^N a_n$. It converges when $S_N$ has a finite limit. The ratio, root, comparison, and integral tests help determine convergence. \par\noindent\textbf{Applications:} series approximate functions, solve differential equations, and quantify accumulated probabilities or errors.

\term{Infinite Product}
The limit of a sequence of partial products $\prod_{n=1}^\infty (1+a_n)$. It converges to a non-zero value iff $\sum a_n$ converges absolutely. Used in the Weierstrass product theorem.

\term{Infinite Dimensional Space}
A vector space is infinite-dimensional if no finite list of vectors spans it; examples include $L^p$ spaces and all polynomials. In such spaces, convergence and continuity require a chosen norm or topology. \par\noindent\textbf{Applications:} infinite-dimensional spaces model functions, signals, quantum states, and solutions of differential equations.

\term{Infinitesimal}
A quantity considered to be closer to zero than any standard real number, but not zero. In non-standard analysis, infinitesimals are formal elements of the hyperreal numbers ${}^*\mathbb{R}$.

\term{Infinitesimal Generator}
For a strongly continuous semigroup $T_t$ on a Banach space, the operator $A = \lim_{t \to 0} \frac{T_t - I}{t}$ defined on the domain where the limit exists. It represents the rate of change at $t=0$.

\term{Infinity}
A quantity larger than every finite number, denoted $\infty$, handled rigorously through limits and cardinalities.

\term{Information Entropy}
For a discrete random variable with probabilities $p_i$, Shannon entropy is $H(X)=-\sum_i p_i\log p_i$. It measures average uncertainty: rare outcomes carry more information, and entropy is largest for a uniform distribution with a fixed number of outcomes. \par\noindent\textbf{Applications:} entropy guides data compression, coding, decision trees, privacy, and statistical modeling.

\term{Information Theory}
The mathematical study of the quantification, storage, and communication of information, founded by Claude Shannon (1948). Key concepts include entropy, mutual information, and channel capacity.

\term{Information Gain}
Information gain is the decrease in entropy produced by observing or splitting on a variable. For a decision tree, the preferred split is often the one with the largest expected reduction in class uncertainty. \par\noindent\textbf{Applications:} it selects features in decision trees and helps quantify the usefulness of observations in data analysis.

\term{Information Geometry}
Information geometry treats a family of probability distributions as points on a manifold. The Fisher information matrix supplies a metric that measures how distinguishable nearby distributions are. \par\noindent\textbf{Applications:} it explains natural-gradient optimization, parameter estimation, exponential families, and statistical learning.

\term{Initial Condition}
Data specifying the state of a dynamical system at an initial time, used to select a unique solution of an ODE.

\term{Initial-value Problem}
An initial-value problem consists of a differential equation and data specifying the solution, and possibly its derivatives, at an initial point or time. Existence and uniqueness conditions determine whether those data select one solution. \par\noindent\textbf{Applications:} initial-value problems model motion, circuits, population change, heat flow, and control systems.

\term{Injective}
A function mapping distinct inputs to distinct outputs, $f(a) = f(b) \Rightarrow a = b$. Also called one-to-one.

\term{Injective Function}
A function $f:X\to Y$ is injective, or one-to-one, if $f(x_1)=f(x_2)$ implies $x_1=x_2$. Distinct inputs have distinct outputs, so $f$ has an inverse on its image even if it is not onto $Y$. \par\noindent\textbf{Applications:} injectivity guarantees recoverable encodings, unique parameter identification, and cancellation in algebra.

\term{Injective Module}
A module $M$ such that any embedding $M \subseteq N$ splits. Equivalently, the functor $\operatorname{Hom}(-, M)$ is exact. Every module can be embedded into an injective module.

\term{Injective Hull}
The smallest injective module containing a given module $M$, denoted $E(M)$. It is the unique (up to isomorphism) essential extension of $M$ that is injective.

\term{Injective Limit}
Another name for the inductive limit (colimit) of a directed system. It describes the smallest object that contains all objects in the system consistently.

\term{Inner Automorphism}
An inner automorphism of a group is conjugation by a fixed element: $g\mapsto aga^{-1}$. It preserves the group operation and describes a symmetry internal to the group. \par\noindent\textbf{Applications:} inner automorphisms organize conjugacy classes, normal subgroups, and the structure of symmetry groups.

\term{Inner Product Space}
An inner-product space is a vector space with a positive-definite form $\langle u,v\rangle$ that defines lengths $\|u\|=\sqrt{\langle u,u\rangle}$, angles, and orthogonality. In complex spaces the form is conjugate-linear in one argument. \par\noindent\textbf{Applications:} inner products support projections, Fourier expansions, least squares, and geometric models of data.

\term{Inner Product}
A mapping $\langle \cdot, \cdot \rangle: V \times V \to \mathbb{F}$ that is linear in the first argument, conjugate symmetric, and positive definite. It generalises the dot product in $\mathbb{R}^n$.

\term{Inseparable Extension}
An algebraic extension $L/K$ where some element has a minimal polynomial with multiple roots in its splitting field. This occurs only in characteristic $p > 0$.

\term{Inseparable Polynomial}
A polynomial that has multiple roots in its splitting field. For a field of characteristic $p$, a polynomial is inseparable iff it is a function of $x^p$.

\term{Inspection Paradox}
The inspection paradox is the tendency for a randomly encountered interval to be longer than a typical interval. For example, arriving at a random time is more likely to occur during a long bus gap than a short one. It is caused by length-biased sampling. \par\noindent\textbf{Applications:} it matters in queueing, reliability, renewal processes, traffic, and service-wait analysis.

\term{Instanton}
A solution of the Yang--Mills equations on a principal $G$-bundle over a Riemannian four-manifold with finite action, equivalently a connection whose curvature $F$ satisfies the self-duality equation $F = *F$ (or anti-self-duality $F = -*F$), where $*$ is the Hodge star. Instantons minimise the Yang--Mills functional in their topological class, and their classification by Chern classes and Pontryagin number uses index theory; over $\mathbb{R}^4$ the ADHM construction parametrises all instantons. Donaldson invariants of smooth four-manifolds are built from moduli spaces of instantons, making them central to gauge theory and the differential topology of four-manifolds.

\term{Integers}
The integers are $\mathbb Z=\{\ldots,-2,-1,0,1,2,\ldots\}$. They form a commutative ring and Euclidean domain, so division with remainder and unique prime factorization are available. The only multiplicative units are $\pm1$. \par\noindent\textbf{Applications:} integers model discrete quantities and form the basis of number theory, modular arithmetic, coding, and exact computation.

\term{Integer Partition}
An integer partition writes a positive integer $n$ as a sum of positive integers, with order ignored; for example, $4=3+1=2+2=2+1+1=1+1+1+1$ gives four nontrivial forms plus $4$ itself. The partition number $p(n)$ grows rapidly. \par\noindent\textbf{Applications:} partitions count combinations in number theory, statistical mechanics, representation theory, and algorithmic enumeration.

\term{Integer Programming}
Integer programming optimizes a linear or nonlinear objective subject to constraints while requiring some or all variables to be integers. Integrality makes many problems combinatorial and NP-hard. \par\noindent\textbf{Applications:} it models scheduling, facility location, routing, production planning, portfolio selection, and network design.

\term{Integral}
An integral assigns a limit to accumulated quantities. The Riemann integral uses partitions of the domain, while the Lebesgue integral groups points by function values and handles more general functions. \par\noindent\textbf{Applications:} integrals compute area, mass, probability, energy, work, and expected values.

\term{Integral Calculus}
The branch of calculus dealing with the computation of integrals and the area under curves, linked to differential calculus via the Fundamental Theorem of Calculus.

\term{Integral Equation}
An integral equation contains an unknown function under an integral sign. Fredholm equations have fixed limits, while Volterra equations have a variable limit such as $\int_a^x$. Kernels describe how values at different points interact. \par\noindent\textbf{Applications:} integral equations model boundary-value problems, inverse problems, scattering, population dynamics, and transport.

\term{Integral Curve}
An integral curve of a vector field $V$ satisfies $\dot\gamma(t)=V(\gamma(t))$, so its tangent agrees with the field at every point. It is the trajectory of the corresponding first-order ODE. \par\noindent\textbf{Applications:} integral curves describe particle motion, fluid streamlines, dynamical systems, and control trajectories.

\term{Integral Manifold}
A submanifold whose tangent space at every point is the subspace specified by a distribution. By Frobenius' Theorem, a distribution is integrable iff it is involutive.

\term{Integrability}
In analysis, a function is integrable when its absolute integral is finite, for example $\int|f|<\infty$ in the Lebesgue sense. In dynamics, a system is integrable when enough independent conserved quantities reduce its motion to a solvable form. \par\noindent\textbf{Applications:} integrability controls whether totals are finite and whether differential equations admit explicit or structured solutions.

\term{Integrable Function}
A function $f$ for which the integral $\int f$ exists and is finite. In the Lebesgue sense, $f$ is integrable iff $\int |f| \, d\mu < \infty$.

\term{Integrable System}
A dynamical system with as many independent constants of motion as degrees of freedom. Their trajectories are confined to invariant tori (Liouville-Arnold theorem).

\term{Integrating Factor}
An integrating factor is a nonzero function $\mu$ that turns a differential equation into an exact one. For $y'+P(x)y=Q(x)$, the factor $\mu(x)=e^{\int P(x)dx}$ gives $(\mu y)'=\mu Q$. \par\noindent\textbf{Applications:} integrating factors solve first-order ODEs and appear in conservation laws and differential-form calculations.

\term{Integration by Parts}
Integration by parts is the integral form of the product rule:
\[
\int u\,dv=uv-\int v\,du.
\]
It transfers a derivative from one factor to another. \par\noindent\textbf{Applications:} it evaluates products of functions, derives weak PDE formulations, proves orthogonality relations, and establishes estimates.

\term{Integration by Substitution}
A method for simplifying integrals by changing variables, based on the chain rule. It transforms the integral $\int f(g(x))g'(x)\,dx$ into $\int f(u)\,du$.

\term{Integral Operator}
An operator $T$ defined by $(Tf)(x) = \int_a^b K(x,y) f(y)\,dy$. The function $K(x,y)$ is the kernel. Many such operators are compact on $L^2$.

\term{Integral Transform}
An integral transform maps a function to another function through an integral, as in the Fourier, Laplace, or Mellin transform. A suitable transform can turn differentiation or convolution into multiplication. \par\noindent\textbf{Applications:} transforms solve differential equations, analyze signals, model systems, and study asymptotic behavior.

\term{Integral Domain}
An integral domain is a commutative ring with $1\ne0$ and no zero divisors: $ab=0$ implies $a=0$ or $b=0$. It supports cancellation and embeds in a field of fractions. \par\noindent\textbf{Applications:} integral domains provide the algebraic setting for polynomial factorization, divisibility, and number theory.

\term{Integral Basis}
A basis for the ring of integers $\mathcal{O}_K$ of a number field $K$ as a $\mathbb{Z}$-module. If $\{\omega_i\}$ is an integral basis, every $\alpha \in \mathcal{O}_K$ is $\sum a_i \omega_i$ with $a_i \in \mathbb{Z}$.

\term{Integral Closure}
For a subring $R \subseteq S$, the set of elements in $S$ that are roots of monic polynomials in $R[x]$. $S$ is integrally closed if $S$ equals its closure in its fraction field.

\term{Integrally Closed Domain}
An integral domain that is equal to its integral closure in its field of fractions. Examples include UFDs and Dedekind domains.

\term{Integer-valued Polynomial}
A polynomial $p(x) \in \mathbb{Q}[x]$ such that $p(n) \in \mathbb{Z}$ for all $n \in \mathbb{Z}$. The binomial coefficients $\binom{x}{k}$ form a basis for such polynomials.

\term{Interior (Topology)}
The interior of a set $S$, denoted $\operatorname{int}(S)$, is the union of all open sets contained in $S$. It consists of all interior points of $S$.

\term{Interior Point}
A point $x \in S$ such that there exists an open neighborhood $U$ with $x \in U \subseteq S$.

\term{Interior Product}
An operation that contracts a vector field $X$ with a differential $k$-form $\omega$, resulting in a $(k-1)$-form: $\iota_X \omega$. It is a derivation on the exterior algebra.

\term{Interior Point Method}
Interior-point methods solve optimization problems by staying inside the feasible region and approaching its boundary through barrier functions. For linear programming they can achieve polynomial-time complexity under standard models. \par\noindent\textbf{Applications:} they solve large-scale resource allocation, network, engineering, and machine-learning optimization problems.

\term{Intermediate Value Theorem}
If $f$ is continuous on $[a,b]$ and $y$ lies between $f(a)$ and $f(b)$, then some $c\in[a,b]$ satisfies $f(c)=y$. A sign change therefore guarantees a zero. \par\noindent\textbf{Applications:} the theorem justifies bisection methods, existence of equilibria, and root-finding in scientific computation.

\term{Intercept (Curve)}
The points where a curve intersects the coordinate axes. For a function $y=f(x)$, the $y$-intercept is $f(0)$ and the $x$-intercepts are the roots of $f(x)=0$.

\term{Interpolation}
Interpolation constructs a function that matches given data points, using lines, polynomials, splines, or other bases. It estimates values between measured samples, unlike extrapolation outside the data range. \par\noindent\textbf{Applications:} interpolation supports numerical tables, computer graphics, signal reconstruction, and scientific data processing.

\term{Interpolating Polynomial}
The unique polynomial of degree at most $n-1$ that passes through $n$ given points. It can be constructed using Lagrange or Newton forms.

\term{Interpolation Error}
The difference between the actual function and its interpolant. For polynomials, it is bounded by a term involving the $(n+1)$-th derivative of the function.

\term{Interpolation Formula}
A general expression for an interpolant, such as the Lagrange formula $\sum y_i L_i(x)$ or the Newton divided-difference formula.

\term{Interpolation Space}
Spaces arising in interpolation theory, the study of how a linear operator bounded on two Banach spaces behaves on intermediate spaces. The real method, built on the $K$-functional, and the complex method construct the intermediate spaces $(A_0, A_1)_{\theta,q}$ and $[A_0, A_1]_\theta$ between a pair of Banach spaces. Interpolation inequalities such as the Riesz--Thorin theorem bound the norm of an operator on the interpolant in terms of its endpoint norms. Sobolev spaces and $L^p$ spaces arise as real interpolation spaces, and the theory is essential in harmonic analysis and the regularity theory of PDEs.

\term{Interquartile Range}
The interquartile range is $IQR=Q_3-Q_1$, the difference between the 75th and 25th percentiles. It measures the spread of the middle half of a dataset and is resistant to extreme outliers. \par\noindent\textbf{Applications:} it summarizes skewed data and is used to flag potential outliers with box plots and robust statistics.

\term{Intersection Theory}
The study of the intersection of subvarieties and cycles on a scheme, whose equivalence classes form the Chow ring, graded by codimension, with the intersection product corresponding to proper intersection of cycles. Rational equivalence, Chow groups, and the pushforward and pullback of cycles are the basic structures. Chern classes of vector bundles are computed via intersections with zero loci, and the theory governs enumerative geometry and birational geometry, connecting to algebraic K-theory through the relation between Chow groups and K-cohomology.

\term{Interval (Analysis)}
A connected subset of $\mathbb{R}$. Closed intervals $[a,b]$ are compact; open intervals $(a,b)$ are open. Intervals are the building blocks of the Borel $\sigma$-algebra.

\term{Interval Arithmetic}
A method for computing with intervals $[a,b]$ rather than point values, used to bound rounding errors and prove inequalities rigorously.

\term{Inverse Function}
A function $f^{-1}$ such that $f^{-1}(f(x)) = x$ and $f(f^{-1}(y)) = y$. $f$ must be bijective for $f^{-1}$ to exist on the whole codomain.

\term{Inverse Map}
A synonym for the inverse function in category theory or topology. A map $f: X \to Y$ is an isomorphism iff it has an inverse map $f^{-1}: Y \to X$.

\term{Inverse Element}
In a group, the element $g^{-1}$ such that $g \cdot g^{-1} = e$. In a ring, the multiplicative inverse $u^{-1}$ such that $u u^{-1} = 1$.

\term{Inverse Operator}
An operator $T^{-1}$ such that $T T^{-1} = T^{-1} T = I$. For matrices, this is the inverse matrix. In functional analysis, the inverse is often a bounded operator.

\term{Inverse Matrix}
A square matrix $A^{-1}$ such that $A A^{-1} = I$. It exists iff $\det(A) \neq 0$. The inverse is computed via the adjugate matrix or Gaussian elimination.

\term{Inverse Transform}
An operation that recovers the original function from its transform, e.g., the Inverse Fourier Transform $f(x) = \int \hat{f}(\xi) e^{2\pi i x \xi} \, d\xi$.

\term{Inverse Problem}
A problem where the goal is to determine the cause (e.g., a source) from the observed effects (e.g., the field). Often ill-posed and requires regularisation.

\term{Inverse Limit}
The categorical limit of a projective system of objects. It is the dual of the inductive limit. Examples include the $p$-adic integers $\mathbb{Z}_p = \varprojlim \mathbb{Z}/p^n\mathbb{Z}$.

\term{Invariant (Algebra)}
An expression or property that remains unchanged under a group action. For example, the determinant of a matrix is invariant under basis change.

\term{Invariant (Geometry)}
A property of a geometric object that is preserved under a group of transformations, such as distance under isometry or angle under conformal mapping.

\term{Invariant (Topology)}
A property of a space preserved under homeomorphism, such as the Euler characteristic or fundamental group. If two spaces have different invariants, they are not homeomorphic.

\term{Invariant Subspace}
A subspace $W \subseteq V$ such that $T(W) \subseteq W$ for a linear operator $T$. The goal of many decomposition theorems (e.g., Jordan form) is to find a basis of invariant subspaces.

\term{Invariant Measure}
A measure $\mu$ on a space $X$ that is preserved by a transformation $T$: $\mu(T^{-1} A) = \mu(A)$ for all measurable $A$. Ergodic theory studies these measures.

\term{Invariant Set}
A set $S$ in a dynamical system such that if $x \in S$, then $f(x) \in S$ for all time. Examples include attractors and unstable manifolds.

\term{Invariant Ring}
The ring of all polynomials $p \in K[x_1, \ldots, x_n]$ that are invariant under the action of a group $G$. The study of such rings is the core of invariant theory.

\term{Invariant Theory}
The branch of algebra studying polynomial invariants. Classical invariant theory focuses on SL(n) and GL(n) acting on tensors.

\term{Inverse Square Law}
A physical law stating that a specified physical quantity is inversely proportional to the square of the distance from the source, e.g., Newton's law of gravity.

\term{Involutive Distribution}
A distribution $D$ such that for any two vector fields $X, Y \in D$, their Lie bracket $[X, Y]$ is also in $D$. By Frobenius' Theorem, such distributions are integrable.

\term{Involution (Group Theory)}
An element $g$ of a group such that $g^2 = e$ (and $g \neq e$). The set of involutions often plays a key role in the classification of finite simple groups.

\term{Involution (Algebra)}
An antiautomorphism $\ast: A \to A$ of an algebra such that $(a^\ast)^\ast = a$ and $(ab)^\ast = b^\ast a^\ast$. C*-algebras are defined via such involutions.

\term{Involutory Matrix}
A square matrix $A$ such that $A^2 = I$. Such matrices represent reflections in geometry and have eigenvalues $\pm 1$.

\term{Icosahedron}
A regular polyhedron with 20 equilateral triangular faces, 12 vertices, and 30 edges. It is one of the five Platonic solids.

\term{Icosahedral Group}
The symmetry group of the regular icosahedron, which is isomorphic to the alternating group $A_5$ (for rotational symmetries) or $A_5 \times \mathbb{Z}/2\mathbb{Z}$.

\term{Irrational Number}
A real number that is not rational, i.e. not expressible as a ratio of integers, such as $\sqrt{2}$ or $\pi$.

\term{Irreducible}
An element of a ring that cannot be written as a product of two non-units; an irreducible polynomial cannot be factored over the ground field.

\term{Isolated Point}
A point of a set having a neighbourhood containing no other points of the set.

\term{Isomorphism}
A bijective map $f: X \to Y$ that preserves the structure of the objects. In groups, $f(ab) = f(a)f(b)$. In topology, it is a homeomorphism.

\term{Isometry}
A distance-preserving map between metric spaces: $d(f(x), f(y)) = d(x,y)$. Isometries are always injective. Examples include translations, rotations, and reflections.

\term{Isoperimetric Inequality}
The inequality $4\pi A \leq L^2$ relating the area $A$ of a plane region to its perimeter $L$, with equality for circles.

\term{Isothermal Coordinates}
Local coordinates in which a surface's metric takes the conformal form $\lambda^2(dx^2 + dy^2)$.

\term{Iteration}
The repeated application of a function or process; in numerical analysis, the basis of fixed-point methods.

\term{Iterative Method}
A numerical method that generates successive approximations converging to the solution.

\term{Itô Calculus}
The stochastic calculus for semimartingales, centered on the Itô integral of predictable processes with respect to martingales such as Brownian motion. Itô's formula is the change-of-variables rule, whose second-order term is the quadratic variation $[X,X]_t$, and it is the source of stochastic differential equations, whose solutions are diffusions governed by their generators. Unlike the Stratonovich integral, the Itô integral is an isometry and a martingale. Itô calculus is the foundation of the theory of diffusions and of mathematical finance through the Itô isometry and Itô's lemma.
\term{Finite Element Method}
A numerical technique for solving PDEs by partitioning the domain into simple elements and approximating the solution on each element by basis functions, typically polynomials. The method converts a PDE into a large sparse linear or nonlinear system. It is the dominant method in engineering simulation and has a rigorous foundation in Sobolev space theory and Céa's lemma.
\term{Initial Object}
An object $0$ in a category such that for every object $X$ there is a unique morphism $0 \to X$. Dually, a terminal object has a unique morphism from every object. In **Set**, the initial object is the empty set and the terminal objects are singleton sets. In **Grp**, the trivial group is both initial and terminal.
\term{Independent Set}
A set of vertices in a graph no two of which are adjacent. The independence number $\alpha(G)$ is the maximum size of an independent set. Independent sets are dual to vertex covers: a set is a vertex cover if and only if its complement is an independent set. The problem of finding maximum independent sets is NP-hard.

\chapter*{J}
\label{chap:j}

\term{Jacobi Identity}
For a Lie algebra, the bracket satisfies
\[
[X,[Y,Z]]+[Y,[Z,X]]+[Z,[X,Y]]=0.
\]
This identity expresses the compatibility between the bracket and the action of one infinitesimal symmetry on another. \par\noindent\textbf{Applications:} it guarantees consistent Lie-algebra calculations and appears in rotations, differential equations, gauge theory, and mechanics.

\term{Jacobi Symbol}
The Jacobi symbol $(\frac an)$ extends the Legendre symbol to an odd positive composite denominator. If $n=\prod p_i^{e_i}$, it is defined as $\prod(\frac a{p_i})^{e_i}$. It is useful for testing some quadratic-residue conditions, although a value $1$ does not by itself guarantee that $a$ is a square modulo $n$. \par\noindent\textbf{Applications:} it supports modular arithmetic, primality tests, and computational number theory.

\term{Jacobian Matrix}
For $f:\mathbb R^n\to\mathbb R^m$, the Jacobian matrix contains the first partial derivatives, $J_f=(\partial f_i/\partial x_j)$. Near a point, it gives the best first-order linear approximation $f(x+h)\approx f(x)+J_f(x)h$. \par\noindent\textbf{Applications:} Jacobians are used in nonlinear equation solving, sensitivity analysis, coordinate changes, robotics, and machine learning.

\term{Jacobian Determinant}
The Jacobian determinant is $\det J_f$ for a map $f:\mathbb R^n\to\mathbb R^n$. Its absolute value gives the local volume-scaling factor. In a change of variables, $\int f(x)\,dx=\int f(g(u))|\det J_g(u)|\,du$. \par\noindent\textbf{Applications:} it is essential in multiple integration, probability-density transformations, mechanics, and coordinate geometry.

\term{Jacobian Variety}
The Jacobian of a compact Riemann surface or algebraic curve is an abelian variety whose points represent degree-zero line bundles, or divisor classes. For a genus-$g$ curve it has dimension $g$. \par\noindent\textbf{Applications:} Jacobians encode curve arithmetic, support the study of Abelian integrals, and are important in algebraic geometry and cryptography.

\term{Jacobson Radical}
The Jacobson radical $J(R)$ is the intersection of all maximal left ideals of a ring; equivalently, it consists of elements acting as zero on every simple module. It measures the part of a ring invisible to simple representations. A ring is semisimple when its radical is zero. \par\noindent\textbf{Applications:} it separates solvable or nilpotent behavior from the semisimple part of an algebra.

\term{Jaccard Index}
A measure of similarity between two sets $A$ and $B$, defined as the size of the intersection divided by the size of the union: $J(A,B) = |A \cap B| / |A \cup B|$. Used extensively in ecology, data science, and clustering.

\term{Jensen's Formula}
Jensen's formula relates the average of $\log|f|$ on a circle to the value at the center and the zeros inside it. For a suitable holomorphic $f$ with zeros $a_k$ in the unit disk,
\[
\int_0^1\log|f(e^{2\pi it})|\,dt=\log|f(0)|+\sum_k\log\frac1{|a_k|}.
\]
\par\noindent\textbf{Applications:} it counts zeros and connects growth of holomorphic functions with their zero sets.

\term{Jensen's Inequality}
If $f$ is convex and the expectations exist, Jensen's inequality gives $f(\mathbb E[X])\leq\mathbb E[f(X)]$. Convexity means that averaging inputs cannot produce a larger value than averaging their function values. \par\noindent\textbf{Applications:} it proves bounds in probability and statistics, supports risk analysis, and explains why convex optimization is tractable.

\term{Jet}
Two functions have the same $k$-jet at a point if all their derivatives through order $k$ agree there. A $k$-jet records the corresponding Taylor data without recording the functions away from the point. \par\noindent\textbf{Applications:} jets describe local behavior, formulate differential equations geometrically, and organize higher-order approximations.

\term{Jet Bundle}
A jet bundle has as its fiber the $k$-jets of local sections of another bundle. A differential equation can then be viewed as a subset or submanifold in a jet bundle, with solutions represented by prolonged sections. \par\noindent\textbf{Applications:} jet bundles provide a coordinate-independent language for PDEs, symmetries, and variational calculus.

\term{Jet Space}
The space of all possible jets of a given order. It generalizes the notion of tangent space by including information about higher-order derivatives.

\term{J-invariant}
The $j$-invariant is a modular invariant of an elliptic curve. In one convention,
\[
j=1728\frac{g_2^3}{g_2^3-27g_3^2}.
\]
Over $\mathbb C$, two elliptic curves are isomorphic exactly when their $j$-invariants agree. \par\noindent\textbf{Applications:} it classifies complex elliptic curves and appears in modular forms, Diophantine equations, and elliptic-curve cryptography.

\term{Joint Distribution}
A joint distribution describes the simultaneous behavior of two or more random variables. It may be given by a joint mass function in the discrete case or a joint density in the continuous case. Marginal distributions are obtained by summing or integrating over the other variables. \par\noindent\textbf{Applications:} joint models describe dependence in experiments, finance, measurement, and statistical prediction.

\term{Joint Entropy}
The joint entropy of discrete variables is $H(X,Y)=-\sum_{x,y}p(x,y)\log p(x,y)$. It measures uncertainty about the pair, while conditional entropy satisfies $H(X,Y)=H(X)+H(Y\mid X)$. \par\noindent\textbf{Applications:} joint entropy supports data compression, mutual information, communication theory, and feature-dependence analysis.

\term{Joint Probability Distribution}
The probability distribution of a set of random variables. For variables $X$ and $Y$, it is $P(X \leq x, Y \leq y)$. Marginal distributions are obtained by integrating (or summing) over the other variables.

\term{Joint Spectral Radius}
For a finite set of matrices $\mathcal M$, the joint spectral radius is the maximal asymptotic growth rate of products:
\[
\rho(\mathcal M)=\lim_{k\to\infty}\max_{A_i\in\mathcal M}\|A_k\cdots A_1\|^{1/k}.
\]
It is independent of the chosen compatible matrix norm. \par\noindent\textbf{Applications:} it tests stability of switched systems, iterated function systems, wavelets, and constrained control processes.

\term{Jordan Algebra}
A Jordan algebra is a commutative, generally non-associative algebra satisfying $(xy)x^2=x(yx^2)$. The product records compatible observables rather than the full order of physical operations. \par\noindent\textbf{Applications:} Jordan algebras model observables in quantum mechanics and appear in symmetric cones, optimization, and Lie theory.

\term{Jordan Block}
A Jordan block $J_k(\lambda)$ has $\lambda$ on the diagonal, $1$ on the superdiagonal, and zeros elsewhere. It represents a chain of generalized eigenvectors when an operator cannot be diagonalized. \par\noindent\textbf{Applications:} Jordan blocks determine powers and exponentials of matrices and therefore solutions of linear differential systems.

\term{Jordan Canonical Form}
Jordan canonical form writes a square matrix over an algebraically closed field as a block-diagonal matrix of Jordan blocks. It exposes eigenvalues and the sizes of generalized-eigenvector chains. \par\noindent\textbf{Applications:} it analyzes matrix powers, exponentials, recurrence relations, and linear dynamical systems.

\term{Jordan Normal Form}
A block-diagonal matrix where each block is a Jordan block. Every square matrix over an algebraically closed field is similar to a matrix in Jordan normal form, unique up to the ordering of the blocks.

\term{Jordan Curve Theorem}
The Jordan curve theorem states that every simple closed curve in the plane separates the plane into exactly two connected components, one bounded inside and one unbounded outside, with the curve as their common boundary. \par\noindent\textbf{Applications:} it underlies planar geometry, contour algorithms, region filling, and topological reasoning about boundaries.

\term{Jordan Decomposition}
For a suitable linear operator, Jordan decomposition writes $T=S+N$, where $S$ is diagonalizable, $N$ is nilpotent, and $SN=NS$. The two parts separately describe eigenvalue behavior and the finite-length generalized-eigenvector effect. \par\noindent\textbf{Applications:} it simplifies matrix functions, differential equations, and representation-theoretic calculations.

\term{Jordan Domain}
A Jordan domain is a bounded planar region whose boundary is a Jordan curve. Its boundary is a simple closed curve, so the region has a well-defined inside. If the domain is simply connected and proper, the Riemann mapping theorem maps it conformally to the disk. \par\noindent\textbf{Applications:} Jordan domains are used in contour integration, conformal mapping, and boundary-value problems.

\term{Jordan Measure}
Jordan measure defines volume using finite unions of rectangles and their inner and outer approximations. A bounded set is Jordan measurable exactly when its boundary has Lebesgue measure zero. \par\noindent\textbf{Applications:} it formalizes elementary area and volume calculations and clarifies when Riemann integration applies.

\term{Jordan--Brouwer Separation Theorem}
The Jordan--Brouwer separation theorem states that an embedded sphere $S^{n-1}$ in $\mathbb R^n$ separates the space into exactly two connected components, an inside and an outside. This generalizes the Jordan curve theorem. \par\noindent\textbf{Applications:} it supports topology of boundaries, solid geometry, mesh algorithms, and the classification of embedded surfaces.

\term{Jordan--Hölder Theorem}
The Jordan--Hölder theorem says that any two composition series of a finite group have the same length and the same simple factor groups, up to order and isomorphism. The factors need not determine the group completely, but they provide a canonical list of building blocks. \par\noindent\textbf{Applications:} it supports classification and comparison of finite groups and modules.

\term{Jordan--von Neumann Theorem}
A norm comes from an inner product if and only if it satisfies the parallelogram law
\[
2(\|x\|^2+\|y\|^2)=\|x+y\|^2+\|x-y\|^2.
\]
Thus the theorem identifies when a normed space has genuine Euclidean geometry. \par\noindent\textbf{Applications:} it distinguishes Hilbert-space methods from general normed-space methods in approximation, optimization, and functional analysis.

\term{J-homomorphism}
The $J$-homomorphism maps homotopy information from a Lie group, especially an orthogonal group, into the stable homotopy groups of spheres. It compares vector-bundle data with sphere-shaped topological data. \par\noindent\textbf{Applications:} it is a central tool in stable homotopy theory and the study of vector bundles and characteristic classes.

\term{Join (Lattice Theory)}
In a lattice, the join $a\vee b$ is the least upper bound of $a$ and $b$: it is the smallest element $x$ with $a\leq x$ and $b\leq x$. In a power-set lattice it is ordinary union. \par\noindent\textbf{Applications:} joins combine subspaces, propositions, sets, and states in ordered systems.

\term{Join (Geometry)}
The geometric join of two subspaces is the smallest subspace containing both. For distinct points $P,Q$, it is the line through them; for two compatible lines, it may be their containing plane. \par\noindent\textbf{Applications:} joins are basic in projective geometry, incidence computations, interpolation, and computational geometry.

\term{Join Semilattice}
A join semilattice is a partially ordered set in which every pair $a,b$ has a least upper bound, written $a\vee b$. It captures a notion of combining two pieces of information while remaining minimal. \par\noindent\textbf{Applications:} join semilattices model logical combination, dataflow analysis, type systems, and hierarchical classification.

\term{Joukowski Transform}
The Joukowski transform is the complex map $w=z+1/z$. Away from its critical points it is conformal, and it maps suitable circles in the $z$-plane to airfoil-like curves in the $w$-plane. \par\noindent\textbf{Applications:} it gives an idealized model for airflow around wings and illustrates conformal mapping in complex analysis.

\term{Julia Set}
For iteration of a complex map such as $f(z)=z^2+c$, the Julia set is the boundary separating stable and unstable behavior. Equivalently, it is where the iterates fail to form a normal family. \par\noindent\textbf{Applications:} Julia sets illustrate fractal geometry, sensitive dependence, dynamical systems, and computer visualization.

\term{Jump Discontinuity}
A function has a jump discontinuity at $c$ when the one-sided limits exist but differ: $\lim_{x\to c^-}f(x)\ne\lim_{x\to c^+}f(x)$. The difference is the jump size. \par\noindent\textbf{Applications:} jumps model switching, shocks, digital signals, and piecewise physical laws.

\term{Jump Function}
A step function whose value changes discontinuously at a set of points.

\term{Jump Process}
A jump process changes state at random times rather than continuously. A Poisson process counts random arrivals, while a Markov jump process changes among states with transition rates. \par\noindent\textbf{Applications:} jump processes model queues, failures, epidemics, financial events, chemical reactions, and communication traffic.

\term{Junction}
A junction is a vertex or point where three or more edges, paths, or channels meet. In a graph it is a high-connectivity location, while in a physical network it may be a branching node. \par\noindent\textbf{Applications:} junctions represent intersections in road networks, switches in circuits, and branching in transport or communication systems.

\term{J-spectrum}
For an operator self-adjoint with respect to an indefinite inner product defined by a fundamental symmetry $J$, the $J$-spectrum records spectral behavior compatible with that symmetry. The precise definition depends on the operator class and the chosen indefinite metric. \par\noindent\textbf{Applications:} $J$-spectral methods study stability and wave or quantum systems with non-positive energy forms.

\chapter*{K}
\label{chap:k}

\term{K-theory (Topological)}
A cohomology-like theory that assigns a ring $K(X)$ to a topological space $X$, constructed from the Grothendieck group of vector bundles over $X$. It is a powerful tool for studying the fact that some vector bundles are not trivial.

\term{K-theory (Algebraic)}
A sequence of functors $K_n$ from rings to abelian groups. $K_0(R)$ is the Grothendieck group of finitely generated projective modules over $R$. Algebraic K-theory bridges ring theory, topology, and number theory.

\term{Kähler Manifold}
A complex manifold equipped with a Hermitian metric whose associated 2-form $\omega$ is closed, so that $[\omega]$ defines a Kähler class in $H^2(X, \mathbb{R})$. The condition $d\omega = 0$ makes the Levi--Civita connection and the Chern connection coincide, and it implies the $\partial\bar{\partial}$-lemma, the Hodge decomposition into Dolbeault cohomology groups, and vanishing theorems such as Kodaira's. Complex projective space and smooth projective varieties are Kähler. Kähler geometry is central to complex and algebraic geometry, mirror symmetry, and the study of Calabi--Yau manifolds.

\term{Kantorovich Inequality}
The inequality $(\sum a_i x_i)(\sum a_i / x_i) \leq \frac{(M+m)^2}{4Mm}$ for positive $a_i$ and $m \leq x_i \leq M$, used in optimization and numerical analysis.

\term{Karush--Kuhn--Tucker (KKT) Conditions}
First-order necessary conditions for a solution in nonlinear programming to be optimal, provided that some regularity conditions are met. They generalize the method of Lagrange multipliers to inequality constraints.

\term{Kazhdan--Lusztig Polynomials}
Polynomials associated with Coxeter groups used to describe the singularities of Schubert varieties. They play a central role in the representation theory of Lie algebras and the Hecke algebra.

\term{Keller's Conjecture}
The conjecture that any tiling of $\mathbb{R}^n$ by equal hypercubes has a pair of cubes meeting full-face-to-full-face; true for $n \leq 6$ and disproved in dimension 10 or more.

\term{Kepler Conjecture}
The theorem that the densest packing of equal spheres in three dimensions is the face-centred cubic lattice, with density $\pi/\sqrt{18}$.

\term{Kernel (Linear Algebra)}
The set of all vectors $v$ such that $Tv = 0$ for a linear operator $T$. Also known as the null space. The dimension of the kernel is the nullity of the operator.

\term{Kernel (Group Theory)}
The set of elements in the domain of a group homomorphism that are mapped to the identity element of the codomain. The kernel is always a normal subgroup of the domain.

\term{Kernel (Integral Equation)}
The function $K(x,y)$ in an integral operator $(Tf)(x) = \int K(x,y)f(y)\,dy$. The properties of the kernel (e.g., symmetry, continuity) determine the properties of the operator.

\term{Kernel (Convolution)}
A function used in convolution to extract specific features from a signal or image. In the context of the Fourier transform, convolution in the time domain corresponds to multiplication in the frequency domain.

\term{Kernel Method}
The technique of implicitly mapping data into a high-dimensional feature space via a positive definite kernel, used in support vector machines.

\term{Killing Form}
A symmetric bilinear form on a Lie algebra $\mathfrak{g}$ defined by $B(X, Y) = \operatorname{tr}(\operatorname{ad}_X \operatorname{ad}_Y)$. A Lie algebra is semisimple iff its Killing form is non-degenerate.

\term{Killing Vector Field}
A vector field $X$ on a Riemannian manifold such that the Lie derivative of the metric $g$ along $X$ vanishes: $\mathcal{L}_X g = 0$. These fields generate isometries of the manifold.

\term{Kirchhoff's Laws}
Conservation laws for electric circuits: the current law (sum of currents into a node is zero) and the voltage law (sum of potential differences around a loop is zero).

\term{Kirchhoff's Theorem}
The formula for the number of spanning trees of a graph as any cofactor of its Laplacian matrix.

\term{Klee's Measure Problem}
The problem of computing the volume of the union of $n$ axis-parallel hyperrectangles in $d$-dimensional space. The most efficient algorithms have complexity $O(n \log n)$ for $d=2$ and $O(n^{d/2})$ for larger $d$.

\term{Klein Bottle}
A non-orientable surface of genus 1, formed by gluing the opposite edges of a rectangle with one pair reversed. It cannot be embedded in $\mathbb{R}^3$ without self-intersection, though it can be embedded in $\mathbb{R}^4$.

\term{Klein Four-Group}
The group $\mathbb{Z}/2 \times \mathbb{Z}/2$, the smallest non-cyclic group, with four elements each of order 2.

\term{Klein Quadric}
The Grassmannian $\operatorname{Gr}(2, 4)$ of lines in $\mathbb{P}^3$, which can be embedded as a smooth quadric hypersurface in $\mathbb{P}^5$. It is central to the study of line geometry.

\term{Kneser Graph}
The graph $KG(n,k)$ whose vertices are the $k$-subsets of an $n$-set, adjacent when they are disjoint; used in Ramsey theory.

\term{Knot Theory}
The study of mathematical knots, which are embeddings of the circle $S^1$ into $\mathbb{R}^3$ or $S^3$. Two knots are equivalent if one can be deformed into the other without cutting.

\term{Knot Invariants}
Quantities (numbers, polynomials, or groups) assigned to a knot that remain unchanged under ambient isotopy. Examples include the Jones polynomial, the Alexander polynomial, and the knot group $\pi_1(S^3 \setminus K)$.

\term{Knot Sum}
The connected sum of two knots $K_1$ and $K_2$, performed by removing a small arc from each and joining the endpoints. This operation makes the set of knots into a commutative monoid.

\term{Knuth's Algorithm}
Algorithms of Donald Knuth for combinatorial generation and sorting, notably his efficient permutation generation algorithms.

\term{Konigsberg Bridges}
The historical problem, solved by Euler in 1736, of whether one can walk across all seven bridges of Königsberg exactly once, founding graph theory.

\term{König's Theorem}
In a bipartite graph, the size of a maximum matching equals the size of a minimum vertex cover. It is equivalent to Hall's marriage theorem and follows from the max-flow min-cut theorem.

\term{K-tuple}
An ordered list of $k$ elements. The set of all $k$-tuples of a set $S$ is the Cartesian product $S^k$.

\term{Koopman Operator}
An infinite-dimensional linear operator that describes the evolution of functions (observables) on a state space under a dynamical system, mapping $f \to f \circ T$.

\term{Kolmogorov Complexity}
The length of the shortest computer program that produces a given string as output. It provides a formal measure of the algorithmic randomness of a string.

\term{Kolmogorov--Savage Divergence}
Also known as Kullback--Leibler (KL) divergence. A measure of how one probability distribution $P$ differs from a second, reference probability distribution $Q$.

\term{Kolmogorov Continuity Theorem}
A criterion for a stochastic process to have a continuous modification. It requires the increments to be bounded by a power of the time difference: $\mathbb{E}[|X_t - X_s|^\alpha] \leq C|t-s|^{1+\beta}$.

\term{Kolmogorov--Smirnov Test}
A non-parametric test used to determine if a sample comes from a specific probability distribution by comparing the empirical distribution function to the theoretical one.

\term{Kolmogorov Three-Series Theorem}
A series of independent random variables $\sum X_n$ converges almost surely if and only if three series converge: $\sum P(|X_n| > 1)$, the series of the means of the truncated variables $X_n 1_{\{|X_n| \leq 1\}}$, and the series of their variances.

\term{K-nearest neighbors (KNN)}
A non-parametric method used for classification and regression, where the label of a point is determined by the majority vote of its $k$ closest neighbors in the feature space.

\term{K-means Clustering}
An iterative algorithm that partitions a dataset into $k$ clusters by minimizing the variance within each cluster (the sum of squared distances to the centroid).

\term{Koszul Complex}
A chain complex used in commutative algebra to study the regularity of a sequence of elements in a ring. It is the primary tool for defining and computing the depth of a module.

\term{Koszul Duality}
A duality between certain categories of modules over quadratic algebras. It relates the cohomology of a space to the algebra of its loop space.

\term{Krull Dimension}
The supremum of the lengths of all chains of prime ideals in a commutative ring. For a polynomial ring $k[x_1, \ldots, x_n]$, the Krull dimension is $n$.

\term{Krull's Intersection Theorem}
In a Noetherian ring, if $I$ is an ideal and $M$ is a finitely generated module, the intersection $\bigcap_{n=1}^\infty I^n M$ consists of elements $m \in M$ such that $(1-a)m = 0$ for some $a \in I$.

\term{Krull--Schmidt Theorem}
The theorem that modules of finite length decompose uniquely into indecomposable direct summands.

\term{Kronecker Delta}
A function of two variables $\delta_{ij}$ that is $1$ if $i=j$ and $0$ otherwise. It is the identity element for the Hadamard product and is used to represent the identity matrix.

\term{Kronecker Product}
An operation on two matrices $A$ and $B$ resulting in a block matrix $A \otimes B$ where each element $a_{ij}$ of $A$ is multiplied by the entire matrix $B$.

\term{Kronecker Lemma}
A result in real analysis stating that if $\sum a_n$ is a convergent series of real numbers and $0 < b_n \uparrow \infty$, then $\frac{1}{b_n} \sum_{k=1}^n a_k b_k \to 0$.

\term{Kronecker--Weber Theorem}
The theorem that every abelian extension of the rationals is contained in a cyclotomic field $\mathbb{Q}(\zeta_n)$; equivalently, the maximal abelian extension of $\mathbb{Q}$ is the union $\bigcup_n \mathbb{Q}(\zeta_n)$ of cyclotomic fields. By the irreducibility of the cyclotomic polynomials, the Galois group $\operatorname{Gal}(\mathbb{Q}(\zeta_n)/\mathbb{Q}) \cong (\mathbb{Z}/n\mathbb{Z})^\times$ is abelian, so the cyclotomic tower realises the whole abelian part of the absolute Galois group. The theorem is the prototype of explicit class field theory: it gives an explicit description of abelian extensions of $\mathbb{Q}$, generalised by Kronecker's Jugendtraum and the Langlands program to other number fields.

\term{Krylov Subspace}
The subspace $\mathrm{span}\{b, Ab, A^2b, \ldots\}$ used in iterative methods for linear systems such as GMRES and the conjugate gradient method.

\term{K-matrix}
In scattering theory, the K-matrix is a representation of the S-matrix that ensures unitarity. It is related to the phase shifts of the scattering process.

\term{K-theoretic Euler Characteristic}
The alternating sum of the ranks of the K-theory groups of a space, providing a generalized notion of the classical Euler characteristic.

\term{Kummer's Theorem}
A theorem giving the exponent of a prime $p$ in a binomial coefficient in terms of carries in base-$p$ addition.

\term{Kurtosis}
The fourth standardized moment of a distribution, measuring the heaviness of its tails.
\term{Opposite Category}
For a category $\mathcal{C}$, the opposite category $\mathcal{C}^{\operatorname{op}}$ has the same objects but all morphisms reversed: $\operatorname{Hom}_{\mathcal{C}^{\operatorname{op}}}(A, B) = \operatorname{Hom}_{\mathcal{C}}(B, A)$. Duality is the principle of obtaining a new theorem by reversing all arrows in a proof. The opposite of an abelian category is abelian, and $(\mathcal{C}^{\operatorname{op}})^{\operatorname{op}} \cong \mathcal{C}$.


\chapter*{L}
\label{chap:l}

\term{Lagrange's Theorem}
For a finite group $G$ and subgroup $H$, the group is partitioned into cosets of $H$, so $|G|=[G:H]|H|$. Thus the order of $H$ divides the order of $G$. \par\noindent\textbf{Applications:} the theorem restricts possible subgroup orders and supports the study of symmetries, permutation groups, and modular arithmetic.

\term{Lagrange Multipliers}
For an equality constraint $g(x)=c$, a constrained extremum of $f$ often satisfies $\nabla f=\lambda\nabla g$, together with $g(x)=c$. The gradients are parallel because the objective cannot change to first order along the constraint surface. \par\noindent\textbf{Applications:} the method supports constrained design, economics, mechanics, and statistical estimation.

\term{Lagrange Interpolation}
Given distinct nodes $x_i$ and values $y_i$, Lagrange interpolation constructs the unique polynomial of degree at most $n-1$ through the data using $p(x)=\sum_i y_iL_i(x)$, where $L_i(x_j)=\delta_{ij}$. \par\noindent\textbf{Applications:} it estimates values from tables and measurements and is used in numerical integration and scientific computing.

\term{Lagrange Inversion Theorem}
The Lagrange inversion theorem gives coefficients of the local inverse of an analytic function. If $w=z\phi(w)$ with $\phi(0)\ne0$, then coefficients of $w(z)$ can be computed from powers of $\phi$. \par\noindent\textbf{Applications:} it derives generating functions, counts combinatorial structures, and analyzes implicit equations.

\term{Lagrangian (Mechanics)}
The Lagrangian is often $L=T-V$, kinetic minus potential energy. The motion follows the Euler--Lagrange equations $\frac d{dt}(\partial L/\partial\dot q)-\partial L/\partial q=0$. The formulation handles generalized coordinates and constraints naturally. \par\noindent\textbf{Applications:} it models mechanics, optics, field theory, robotics, and control.

\term{Lambda Calculus}
Lambda calculus is a formal language for computation built from variables, function abstraction, and function application, with substitution as its main operation. It separates the idea of a function from a particular programming syntax. \par\noindent\textbf{Applications:} it underlies functional programming, type theory, interpreters, and the mathematical study of computability.

\term{Lambda-term}
A lambda term is a variable, an abstraction $\lambda x.M$ representing a function, or an application $(MN)$ applying one term to another. Bound variables can be renamed, and substitution models computation. \par\noindent\textbf{Applications:} lambda terms describe programs, symbolic reduction, and formal proofs in logic and type systems.

\term{Langlands Program}
A vast network of conjectures, initiated by Robert Langlands in 1967, connecting automorphic forms and representations of reductive groups over global fields to Galois representations and motives. The central reciprocity conjecture, the Langlands correspondence, predicts that $L$-functions of automorphic representations match those of Galois representations, and functoriality asserts functorial transfer between reductive groups. Its local version concerns representations of $p$-adic groups, and it generalizes Artin's conjecture on holomorphy, providing the framework for modularity results such as the proof of Fermat's Last Theorem.

\term{Laplace Equation}
The equation $\Delta u = 0$ for the Laplacian $\Delta$, whose solutions are the harmonic functions. Harmonic functions satisfy the maximum principle and the mean value property and are real analytic. Boundary value problems for the equation prescribe the values of $u$ or of its normal derivative on the boundary, the Dirichlet and Neumann problems; the inhomogeneous Poisson equation $\Delta u = f$ is solved using fundamental solutions and Green's functions. The Laplace equation is the model elliptic equation, governing electrostatics, steady heat flow, and potential flow.

\term{Laplace Transform}
The Laplace transform is $\mathcal L\{f\}(s)=\int_0^\infty e^{-st}f(t)\,dt$. Differentiation becomes multiplication by $s$ together with initial-value terms, so differential equations become algebraic equations in transform space. \par\noindent\textbf{Applications:} it solves ODEs for circuits, control systems, mechanical vibrations, and transient responses.

\term{Laplacian}
The Laplacian is the divergence of the gradient: $\Delta f=\nabla\cdot\nabla f=\sum_i\partial^2f/\partial x_i^2$. It measures local curvature or imbalance of flux. \par\noindent\textbf{Applications:} it governs heat and diffusion, electrostatic potential, wave models, image smoothing, and graph-based learning.

\term{Large Cardinals}
Very large infinite cardinals whose existence cannot be proved in ZFC and whose consistency implies the consistency of ZFC, arranged in a hierarchy of consistency strength. They include inaccessible cardinals, measurable cardinals, Woodin cardinals, and supercompact cardinals. Their theory measures the strength of set-theoretic statements, interacts with the axiom of determinacy and with forcing, and features in inner model theory, where the existence of such cardinals reflects deep structure of the universe of sets beyond the usual cumulative hierarchy.

\term{Large Deviations Principle}
The asymptotic description of rare events: a sequence of probability measures $P_n$ satisfies a large deviations principle with rate function $I$ if $P_n(A) \approx \exp(-n \inf_{x \in A} I(x))$ for suitable sets $A$, where $I \geq 0$ is lower semicontinuous with compact sublevel sets. Cramér's theorem gives the rate function for empirical means of i.i.d. variables, and Sanov's theorem for empirical measures. Varadhan's lemma computes limits of normalized logarithmic expectations. The theory underpins statistics and statistical mechanics.

\term{Lattice (Algebra)}
A lattice is a partially ordered set in which every pair has a least upper bound $a\vee b$ and greatest lower bound $a\wedge b$. Distributive and modular lattices are important special cases. \par\noindent\textbf{Applications:} lattices organize subspaces, ideals, logical propositions, and data or information hierarchies.

\term{Lattice (Geometry)}
A geometric lattice is a discrete subgroup $\Lambda=\{\sum_i m_i\mathbf b_i:m_i\in\mathbb Z\}$ of $\mathbb R^n$. A basis generates a fundamental parallelepiped whose volume is $|\det(\mathbf b_1,\ldots,\mathbf b_n)|$ in full rank. \par\noindent\textbf{Applications:} lattices model crystals, error-correcting codes, cryptography, optimization, and integer approximation.

\term{Laurent Series}
A Laurent series represents a complex function near $z_0$ as $f(z)=\sum_{n=-\infty}^{\infty}a_n(z-z_0)^n$. Negative powers describe the singular part: finitely many indicate a pole, while infinitely many may indicate an essential singularity. \par\noindent\textbf{Applications:} Laurent series calculate residues and evaluate contour integrals.

\term{L-function}
An $L$-function is a Dirichlet series, often with an Euler product, that generalizes the Riemann zeta function; for a character $\chi$, $L(s,\chi)=\sum\chi(n)n^{-s}$. Its analytic behavior encodes arithmetic information. \par\noindent\textbf{Applications:} $L$-functions study primes, elliptic curves, modular forms, and number-field extensions.

\term{L-space}
In Heegaard Floer homology, an L-space is a 3-manifold $Y$ whose Heegaard Floer homology $\widehat{HF}(Y)$ is as small as possible (rank equal to the order of $H_1(Y; \mathbb{Z})$).

\term{Law of Large Numbers}
The law of large numbers says that averages of many independent identically distributed variables approach their common expectation. The weak law gives convergence in probability, while the strong law gives almost-sure convergence under suitable conditions. \par\noindent\textbf{Applications:} it justifies sample averages, Monte Carlo computation, polling, insurance, and statistical estimation.

\term{Law of Total Probability}
If $\{A_i\}$ is a partition, then $P(B)=\sum_iP(B\mid A_i)P(A_i)$. It decomposes an event according to mutually exclusive cases. \par\noindent\textbf{Applications:} it supports Bayesian inference, reliability calculations, risk analysis, and prediction with hidden categories.

\term{Least Common Multiple}
The least common multiple $\operatorname{lcm}(a,b)$ is the smallest positive integer divisible by both $a$ and $b$. For positive integers, $\gcd(a,b)\operatorname{lcm}(a,b)=ab$. \par\noindent\textbf{Applications:} LCM handles synchronized cycles, common denominators, modular schedules, and periodic processes.

\term{Lax Equivalence Theorem}
For a consistent finite-difference scheme approximating a well-posed linear initial-value problem, stability is equivalent to convergence. Consistency controls local truncation error, while stability prevents errors from growing without bound. \par\noindent\textbf{Applications:} the theorem guides reliable numerical schemes for wave, transport, and other evolution equations.

\term{Lax--Milgram Theorem}
If $a(u,v)$ is bounded and coercive on a Hilbert space $H$, then for every bounded linear functional $f$ there is a unique $u\in H$ with $a(u,v)=f(v)$ for all $v$. The solution depends continuously on $f$. \par\noindent\textbf{Applications:} it proves existence and stability of weak solutions to elliptic PDEs and supports finite-element methods.

\term{Least Squares}
Least squares chooses parameters to minimize the sum of squared residuals between observations and model predictions. For a linear model this leads to the normal equations, or a more stable QR or singular-value computation. \par\noindent\textbf{Applications:} least squares is used in regression, calibration, signal reconstruction, inverse problems, and data fitting.

\term{Lebesgue Integral}
The Lebesgue integral groups points according to function values and measures the corresponding sets, rather than relying only on intervals in the domain. It extends the Riemann integral and handles limits of functions and many discontinuities effectively. \par\noindent\textbf{Applications:} it is the foundation of modern probability, Fourier analysis, PDEs, and $L^p$ spaces.

\term{Lebesgue Measure}
Lebesgue measure assigns length, area, or volume to measurable subsets of $\mathbb R^n$. It is translation invariant and agrees with ordinary volume on boxes, with the unit cube having measure $1$. \par\noindent\textbf{Applications:} it defines probability densities, integration, volume in analysis, and the size of exceptional or fractal sets.

\term{Lebesgue's Dominated Convergence Theorem}
If measurable $f_n\to f$ almost everywhere and $|f_n|\leq g$ for an integrable $g$, then $\int f_n\to\int f$. The common integrable bound permits the limit to pass through the integral. \par\noindent\textbf{Applications:} it justifies limiting calculations in probability, Fourier analysis, numerical approximation, and PDEs.

\term{Lusin's Theorem}
On a finite-measure space, a measurable function is nearly continuous: for every $\varepsilon>0$, there is a closed set whose complement has measure below $\varepsilon$ and on which the function is continuous. \par\noindent\textbf{Applications:} Lusin's theorem connects measurable functions with continuous approximation and supports density results in $L^p$ analysis.

\term{Lemma}
A "helping theorem"; a minor result that is proved as a stepping stone toward a larger, more significant theorem.

\term{Leibniz Rule (Generalised)}
The generalized Leibniz rule is
\[
(fg)^{(n)}=\sum_{k=0}^n\binom nk f^{(k)}g^{(n-k)}.
\]
It extends the product rule to higher derivatives. \par\noindent\textbf{Applications:} it is used in Taylor expansions, differential equations, symbolic computation, and estimates for products of smooth functions.

\term{Leibniz Integral Rule}
The Leibniz rule differentiates an integral with variable limits:
\[
\frac d{dx}\int_{a(x)}^{b(x)}f(x,t)dt=f(x,b(x))b'(x)-f(x,a(x))a'(x)+\int_{a(x)}^{b(x)}f_x(x,t)dt.
\]
\par\noindent\textbf{Applications:} it derives conservation laws, sensitivity equations, moving-boundary models, and gradients of integral objectives.

\term{Levi--Civita Connection}
The Levi--Civita connection is the unique connection that is torsion-free and preserves the Riemannian metric. It defines covariant differentiation, parallel transport, curvature, and geodesics. \par\noindent\textbf{Applications:} it provides the basic language for curved surfaces, general relativity, mechanics, and geometric data analysis.

\term{L'Hôpital's Rule}
Under appropriate differentiability and nonzero-denominator hypotheses, an indeterminate limit of $f/g$ of type $0/0$ or $\infty/\infty$ may be found from $f'/g'$ when that latter limit exists. \par\noindent\textbf{Applications:} it evaluates difficult limits in calculus and asymptotic error analysis; its hypotheses must be checked before use.

\term{Liar Paradox}
The self-referential sentence "this statement is false" cannot consistently be assigned a truth value, since it is true iff it is false. It is a foundational example of self-reference in logic, motivating Tarski's hierarchy of truth and discussions of paraconsistent and gap-based approaches.

\term{Lie Algebra}
A Lie algebra is a vector space with a bilinear bracket satisfying $[X,X]=0$ and the Jacobi identity. The bracket measures the failure of infinitesimal operations to commute. \par\noindent\textbf{Applications:} Lie algebras linearize Lie groups and describe rotations, differential equations, symmetries, gauge fields, and conservation laws.

\term{Lie Bracket}
The Lie bracket measures noncommutativity. For matrices or operators it is $[X,Y]=XY-YX$; for vector fields it records their commutator as differential operators. It is bilinear, antisymmetric, and satisfies Jacobi. \par\noindent\textbf{Applications:} brackets describe infinitesimal symmetries, controllability, curvature, and quantum commutators.

\term{Lie Group}
A Lie group is both a group and a smooth manifold, with smooth multiplication and inversion. Its tangent space at the identity is a Lie algebra describing infinitesimal motions. \par\noindent\textbf{Applications:} Lie groups model continuous symmetries in robotics, mechanics, computer vision, relativity, and physics.

\term{Lie Derivative}
The Lie derivative measures how a tensor field changes when carried along the flow of a vector field. For a scalar it is directional differentiation; for vector fields it is related to the Lie bracket. \par\noindent\textbf{Applications:} it tests symmetries, expresses conservation laws, and is used in fluid mechanics, differential geometry, and relativity.

\term{Lie Superalgebra}
A generalization of a Lie algebra that is $\mathbb{Z}_2$-graded (even and odd parts) and satisfies a graded version of the Jacobi identity. Used in supersymmetry in physics.

\term{Limit}
The limit of a function or sequence is the value it approaches as its input approaches a point or infinity. The $\varepsilon$--$\delta$ definition makes this idea precise without requiring the function to attain the limiting value. \par\noindent\textbf{Applications:} limits define derivatives, integrals, continuity, convergence, and asymptotic approximations.

\term{Limit (Category Theory)}
A limit of a diagram $F: J \to C$ is a cone $(L, \pi_j)$ over $F$ that is universal: every other cone factors uniquely through $L$. Limits encompass products, pullbacks, equalizers, and inverse limits. A category with all small limits is complete, and limits are computed objectwise in functor categories. Right adjoints preserve limits, and their existence is governed by the adjoint functor theorems. Limits are the dual of colimits and central to universal constructions; this is distinct from the analytic notion of a limit in calculus.

\term{Limit Point}
A point $x$ is a limit or accumulation point of a set $S$ if every neighborhood of $x$ contains a point of $S$ different from $x$. In metric spaces this is equivalent to being the limit of a sequence of distinct points of $S$. \par\noindent\textbf{Applications:} limit points describe closure, continuity, convergence, and the boundary structure of sets.

\term{Lindenbaum--Tarski Algebra}
The algebra formed by the equivalence classes of formulas in a logical theory, where the equivalence relation is provable equivalence.

\term{Lindeberg--Lévy Central Limit Theorem}
The Lindeberg--Lévy central limit theorem states that normalized sums of independent identically distributed variables with finite mean and nonzero finite variance converge in distribution to a standard normal variable. \par\noindent\textbf{Applications:} it explains normal approximations for sample means, measurement errors, polling, and statistical confidence intervals.

\term{Lindelöf Space}
A Lindelöf space is one in which every open cover has a countable subcover. This property provides countable control over potentially large collections of open sets; every second-countable space is Lindelöf. \par\noindent\textbf{Applications:} it supports separability and countable constructions in analysis and topology.

\term{Line}
A one-dimensional object with zero thickness extending without end in both directions; in a projective plane, the dual of a point.

\term{Line Bundle}
A line bundle is a vector bundle with one-dimensional fibers. Local sections look like functions whose descriptions change by nonzero scaling on overlaps. Divisors and line bundles are closely related, and the first Chern class records topological twisting. \par\noindent\textbf{Applications:} line bundles define projective embeddings, describe phases and magnetic potentials, and organize divisors in algebraic and complex geometry.

\term{Line Integral}
A line integral accumulates a scalar field along arc length, $\int_Cf\,ds$, or measures work and circulation of a vector field, $\int_C\mathbf F\cdot d\mathbf r$. Its value can depend on the path unless the field is conservative. \par\noindent\textbf{Applications:} line integrals compute work, circulation, mass of wires, and quantities along trajectories.

\term{Linear Algebra}
Linear algebra studies vector spaces, linear maps, matrices, and linear equations. Concepts such as bases, rank, eigenvalues, and orthogonality turn complicated systems into structured calculations. \par\noindent\textbf{Applications:} it is fundamental in engineering, data science, graphics, quantum mechanics, differential equations, and optimization.

\term{Linear Combination}
A linear combination of vectors is $\sum_i a_iv_i$, where the $a_i$ are scalars. The span of a set is all its linear combinations. \par\noindent\textbf{Applications:} linear combinations describe coordinates, approximations, mixtures, signals, and solutions of linear systems.

\term{Linear Dependence}
A set of vectors $\{v_i\}$ is linearly dependent if there exist scalars $a_i$, not all zero, such that $\sum a_i v_i = 0$. Otherwise, they are linearly independent.

\term{Linear Functional}
A linear functional is a linear map $\ell:V\to\mathbb F$, satisfying $\ell(au+bv)=a\ell(u)+b\ell(v)$. All such maps form the dual space $V^*$. \par\noindent\textbf{Applications:} functionals represent measurements, constraints, gradients, weak derivatives, and observations in optimization and PDEs.

\term{Linear Independence}
A set of vectors is linearly independent if the only linear combination equal to zero has all coefficients zero.

\term{Linear Map}
A linear map satisfies $T(au+bv)=aT(u)+bT(v)$. Once bases are chosen, it is represented by a matrix, so kernels, images, rank, and eigenvalues can be studied algebraically. \par\noindent\textbf{Applications:} linear maps model transformations, networks, discretized physical systems, and changes of coordinates.

\term{Linear Programming}
Linear programming optimizes a linear objective subject to linear equalities and inequalities. The feasible set is a convex polyhedron, and when an optimum exists it can be found at an extreme point. \par\noindent\textbf{Applications:} it allocates resources, schedules production, routes goods, designs networks, and supports machine learning.

\term{Linear Recurrence}
A linear recurrence expresses each term as a fixed linear combination of earlier terms, such as $a_n=c_1a_{n-1}+\cdots+c_ka_{n-k}$. Its characteristic polynomial often gives a closed form. \par\noindent\textbf{Applications:} recurrences model population growth, algorithms, digital filters, and discrete dynamical systems.

\term{Linear Regression}
Linear regression models a response as $y\approx\beta_0+\beta_1x_1+\cdots+\beta_px_p+\varepsilon$. Parameters are often estimated by least squares, with assumptions used for uncertainty and inference. \par\noindent\textbf{Applications:} regression supports prediction, calibration, trend estimation, and experimental analysis.

\term{Linear Operator}
A linear map from a vector space to itself. The study of these operators includes eigenvalues, eigenvectors, and spectral theory.

\term{Lipschitz Continuity}
A function is Lipschitz continuous if $|f(x)-f(y)|\leq L|x-y|$ for all $x,y$. It cannot change faster than a fixed linear rate. Lipschitz continuity implies uniform continuity. \par\noindent\textbf{Applications:} it guarantees uniqueness of ODE solutions under standard hypotheses and controls approximation error, stability, and learning models.

\term{Lipschitz Constant}
The Lipschitz constant is the smallest admissible $L$ in $|f(x)-f(y)|\leq L|x-y|$. It bounds the maximum global rate of change. \par\noindent\textbf{Applications:} it quantifies sensitivity, controls propagation of errors, and sets step-size or robustness bounds in numerical methods and optimization.

\term{Liouville's Theorem}
Liouville's theorem states that every bounded entire holomorphic function is constant. Cauchy's estimates force all positive-degree Taylor coefficients to vanish. \par\noindent\textbf{Applications:} it proves the fundamental theorem of algebra and gives powerful uniqueness and growth arguments in complex analysis.

\term{Local Compactness}
A space is locally compact if every point has a neighborhood whose closure is compact, or equivalently a compact neighborhood in common Hausdorff settings. Local compactness allows useful integration and compactness arguments without the whole space being compact. \par\noindent\textbf{Applications:} it underlies Haar measure, Fourier analysis, potential theory, and one-point compactification.

\term{Local Field}
A locally compact, non-discrete topological field: the real or complex numbers, a finite extension of the field $\mathbb{Q}_p$ of $p$-adic numbers, or a finite extension of the field $\mathbb{F}_q((t))$ of formal Laurent series over a finite field. Every non-archimedean local field carries a discrete valuation and a finite residue field. Local class field theory describes abelian extensions of local fields, Tate duality governs local Galois cohomology, and Haar measure on the additive group provides the local structure underlying the adèles.

\term{Local Homeomorphism}
A map is a local homeomorphism if every point has a neighborhood on which the map is a homeomorphism onto an open subset of the target. It need not be globally one-to-one. \par\noindent\textbf{Applications:} local homeomorphisms describe covering maps, coordinate charts, and locally invertible transformations.

\term{Local Minimum}
A point $x_0$ is a local minimum if $f(x_0)\leq f(x)$ for all $x$ in some neighborhood of $x_0$. It is a local comparison and may not be the smallest value globally. \par\noindent\textbf{Applications:} local minima represent stable designs and equilibria and are targets of many optimization algorithms.

\term{Local Maximum}
A point $x_0$ is a local maximum if $f(x)\leq f(x_0)$ for all $x$ in some neighborhood of $x_0$. It need not be the largest value on the whole domain. \par\noindent\textbf{Applications:} local maxima describe nearby best responses, peaks in data, and stable or unstable equilibria.

\term{Local Ring}
A local ring has exactly one maximal ideal, representing the functions that vanish at the point under study. In algebraic geometry, localizing at a prime ideal focuses on behavior near that point. \par\noindent\textbf{Applications:} local rings study singularities, tangent spaces, divisors, and local solvability of equations.

\term{Localization}
The process of formally adjoining inverses of a multiplicatively closed subset $S$ of a commutative ring $R$, producing the ring of fractions $S^{-1}R$; localization at the complement of a prime ideal $\mathfrak{p}$ yields the local ring $R_{\mathfrak{p}}$. For modules, localization is exact and compatible with tensor products. Geometrically, localization corresponds to restricting sheaves to an open set, expressing the local behaviour of a variety at a point. The notion generalizes to the localization of a category, formally inverting a class of morphisms, as in the construction of derived categories.

\term{Logarithm}
The logarithm is the inverse of exponentiation: $\log_bx=y$ when $b^y=x$. The natural logarithm satisfies $\ln x=\int_1^xdt/t$ and converts products into sums. \par\noindent\textbf{Applications:} logarithms measure growth, scale data, solve exponential equations, and appear in entropy and decibel units.

\term{Logarithmic Function}
The inverse of the exponential function, $\log_b x = y$ when $b^y = x$, satisfying $\log(xy) = \log x + \log y$.

\term{Loop Space}
The loop space $\Omega X$ consists of paths in a pointed space that start and end at the basepoint. Concatenation gives it a group-like structure up to homotopy, and $\pi_n(X)\cong\pi_{n-1}(\Omega X)$. \par\noindent\textbf{Applications:} loop spaces convert questions about maps and homotopy groups into structured algebraic-topology problems and are central to stable homotopy theory.

\term{Lorentz Transformation}
A Lorentz transformation is a linear change between inertial frames preserving the spacetime interval $\Delta s^2=-c^2\Delta t^2+\Delta x^2+\Delta y^2+\Delta z^2$. It combines rotations with boosts and replaces ordinary Euclidean distance by spacetime geometry. \par\noindent\textbf{Applications:} Lorentz transformations describe relativistic motion, time dilation, length contraction, and electromagnetic-field changes.

\term{Löwenheim--Skolem Theorem}
A countable first-order theory that has an infinite model has a model of any infinite cardinality (downward and upward forms). It yields the Skolem paradox: set theory has a countable model, yet proves the existence of uncountable sets.

\term{Lower Bound}
A lower bound of a set $S$ is a number $m$ with $m\leq s$ for every $s\in S$. The greatest lower bound is $\inf S$, whether or not it belongs to $S$. \par\noindent\textbf{Applications:} lower bounds certify performance limits, prove convergence, and establish minimum possible costs or errors.

\term{L\textasciicircum p Space}
The space $L^p$ contains measurable functions with $\int|f|^p<\infty$ for $1\leq p<\infty$, with an essential-supremum norm when $p=\infty$. For $p=2$ it is a Hilbert space. \par\noindent\textbf{Applications:} $L^p$ spaces measure signal size, define probability moments, and provide the standard setting for Fourier analysis and PDEs.

\term{L-polynomial}
The polynomial associated with an L-function, often appearing in the context of the Birch and Swinnerton-Dyer conjecture for elliptic curves.

\term{Lefschetz Fixed-Point Theorem}
The Lefschetz fixed-point theorem relates fixed points of a continuous map to its action on homology. The Lefschetz number is $L(f)=\sum_k(-1)^k\operatorname{tr}(f_*|H_k(X))$; if it is nonzero, $f$ has a fixed point. \par\noindent\textbf{Applications:} it proves existence of equilibria and fixed points in topology, dynamics, and geometry.

\term{Lefschetz Hyperplane Theorem}
A result in complex geometry stating that the cohomology of a smooth projective variety $X$ is determined by the cohomology of a generic hyperplane section $H$ up to a certain dimension.

\term{L-system (Lindenmayer system)}
An L-system is a parallel rewriting system: every symbol in a string is replaced simultaneously according to production rules. Interpreting symbols as drawing commands produces branching patterns. \par\noindent\textbf{Applications:} L-systems model plant growth, generate fractals, and create procedural graphics.

\term{Lattice Gauge Theory}
Lattice gauge theory approximates spacetime by a discrete lattice while preserving gauge symmetry through variables on edges. It provides a non-perturbative formulation of quantum field theories. \par\noindent\textbf{Applications:} numerical lattice calculations study strong interactions, confinement, phase transitions, and quantum statistical systems.

\term{Lumpability}
Lumpability is a property of a Markov chain under which grouped states form an exact lower-dimensional Markov chain. Transition probabilities from states in the same group must agree after aggregation. \par\noindent\textbf{Applications:} lumpability enables model reduction, state aggregation, faster simulation, and simplified queueing or reliability models.

\chapter*{M}
\label{chap:m}

\term{Maclaurin Series}
The Maclaurin series is the Taylor expansion at zero:
\[
f(x)=\sum_{n=0}^{\infty}\frac{f^{(n)}(0)}{n!}x^n,
\]
when the series converges to $f$. \par\noindent\textbf{Applications:} it approximates functions, solves differential equations, and supports error estimates in scientific computing.

\term{Manifold}
A manifold is a space that locally looks like $\mathbb R^n$: every point has a neighborhood with coordinates in an open subset of $\mathbb R^n$. Extra structure can make it smooth, Riemannian, complex, or symplectic. \par\noindent\textbf{Applications:} manifolds model curved surfaces, configuration spaces, spacetime, and data with nonlinear geometry.

\term{Metric Tensor}
A metric tensor is a smoothly varying inner product on tangent spaces. It measures lengths, angles, distances, areas, and volumes and determines geodesics and curvature. \par\noindent\textbf{Applications:} metric tensors describe physical spacetime, curved surfaces, mechanics, and geometric statistics.

\term{Map}
A map or function assigns each input in a domain exactly one output in a codomain. Depending on context it may preserve structure, such as addition, distance, or multiplication. \par\noindent\textbf{Applications:} maps represent transformations, models, coordinate changes, and input-output relationships.

\term{Markov Chain}
A Markov chain is a sequence of random states whose next-state distribution depends only on the current state. Its transition matrix $P$ satisfies $P_{ij}=P(X_{n+1}=j\mid X_n=i)$. \par\noindent\textbf{Applications:} Markov chains model queues, web ranking, weather, genetics, inventory, and discrete decision systems.

\term{Markov Chain Monte Carlo}
Markov chain Monte Carlo (MCMC) samples from a target distribution $\pi$ by simulating a Markov chain whose stationary distribution is $\pi$. Metropolis--Hastings and Gibbs sampling are standard methods; burn-in, mixing, and effective sample size must be assessed. \par\noindent\textbf{Applications:} MCMC is central to Bayesian inference, statistical physics, uncertainty quantification, and models with intractable integrals.

\term{Markov Process}
A Markov process has the property that, conditional on its present state, its future is independent of its past. Markov chains use discrete time and states, while diffusion and jump processes give continuous-time examples. \par\noindent\textbf{Applications:} Markov processes model evolving systems in finance, biology, queues, physics, and control.

\term{Markov Property}
The property that the future state of a process depends only on its present state, not on its history. This "memoryless" property is the core of Markov processes.

\term{Martingale}
A process $(X_n)$ is a martingale with respect to information $(\mathcal F_n)$ if $\mathbb E[|X_n|]<\infty$ and $\mathbb E[X_{n+1}\mid\mathcal F_n]=X_n$. Its current value is the best conditional prediction of the next value. \par\noindent\textbf{Applications:} martingales model fair games, financial prices, stochastic integrals, and convergence arguments.

\term{Martingale Convergence Theorem}
A submartingale that is bounded in $L^1$ converges almost surely to an integrable limit. For martingales, $L^1$-boundedness guarantees almost sure convergence, with convergence in $L^1$ under uniform integrability. It underlies the almost sure convergence of branching processes, reversed martingales, and many limit theorems in probability.

\term{Mathematics}
Mathematics studies quantity, structure, space, change, and uncertainty through precise definitions, logical deduction, and proof. It provides both abstract theories and practical models. \par\noindent\textbf{Applications:} mathematics supports every quantitative science, including engineering, computing, economics, medicine, and physics.

\term{Matroid}
A matroid abstracts independence. It consists of a finite ground set and independent subsets satisfying heredity and an exchange rule. Vector-space linear independence and edge sets of forests in graphs are examples. \par\noindent\textbf{Applications:} matroids explain why greedy algorithms solve certain optimization problems and connect combinatorics, network design, and algebra.

\term{Matrix}
A matrix is a rectangular array representing a linear map once bases are chosen. Matrix operations encode composition, systems of equations, and changes of coordinates. \par\noindent\textbf{Applications:} matrices are used in simulations, graphics, statistics, networks, control, and machine learning.

\term{Matrix Multiplication}
For compatible matrices, $(AB)_{ij}=\sum_kA_{ik}B_{kj}$. Multiplication is associative but usually not commutative because it represents composition of linear maps. \par\noindent\textbf{Applications:} it combines transformations, propagates network relationships, and drives numerical algorithms and neural-network layers.

\term{Matrix Norm}
A matrix norm measures matrix size and satisfies positivity, homogeneity, the triangle inequality, and usually $\|AB\|\leq\|A\|\|B\|$. The Frobenius norm sums squared entries; the operator norm measures maximum stretching. \par\noindent\textbf{Applications:} norms quantify errors, conditioning, stability, and convergence of matrix algorithms.

\term{Maximal Ideal}
A maximal ideal is a proper ideal not contained in any larger proper ideal. The quotient $R/I$ is a field exactly when $I$ is maximal. \par\noindent\textbf{Applications:} maximal ideals identify points of algebraic varieties, construct fields, and analyze ring homomorphisms.

\term{Maximum Flow Problem}
The maximum-flow problem sends as much material as possible from a source to a sink while respecting edge capacities and flow conservation at intermediate vertices. The max-flow min-cut theorem equates the optimum with the capacity of a smallest separating cut. \par\noindent\textbf{Applications:} routing, communication bandwidth, transportation, matching, and supply chains.

\term{Maximum Likelihood Estimator (MLE)}
The maximum likelihood estimator chooses the parameter $\hat\theta$ maximizing $L(\theta)=p(data\mid\theta)$. Under regularity conditions it is often consistent and approximately normal for large samples, but it can be biased or unstable in small samples. \par\noindent\textbf{Applications:} MLE fits statistical, biological, reliability, and machine-learning models.

\term{Maximum Modulus Principle}
If a holomorphic function is nonconstant, its modulus cannot have a local maximum at an interior point of its domain. If a maximum exists on the closure of a bounded domain, it occurs on the boundary. \par\noindent\textbf{Applications:} the principle gives bounds for analytic functions and helps prove uniqueness results such as the Schwarz lemma.

\term{Maximum Principle}
The maximum principle says that a nonconstant harmonic or holomorphic function cannot attain a maximum at an interior point under the usual hypotheses. Boundary data therefore control the interior. \par\noindent\textbf{Applications:} it proves uniqueness and stability for boundary-value problems and bounds physical potentials.

\term{Mayer--Vietoris}
The long exact sequence relating the homology of a union $X = A \cup B$ to the homology of the pieces and their intersection: $\cdots \to H_n(A \cap B) \to H_n(A) \oplus H_n(B) \to H_n(X) \to H_{n-1}(A \cap B) \to \cdots$. Together with excision, the Mayer--Vietoris sequence is the fundamental computational tool for computing homology of spaces built from simpler pieces; it has analogous formulations for cohomology and for sheaves.

\term{Mean Value Theorem}
If $f$ is continuous on $[a,b]$ and differentiable inside, some $c\in(a,b)$ satisfies $f'(c)=\frac{f(b)-f(a)}{b-a}$. At one point, instantaneous rate equals average rate. \par\noindent\textbf{Applications:} it proves error bounds, monotonicity, uniqueness, and estimates for numerical approximations.

\term{Measure}
A measure assigns nonnegative size to measurable sets, with $\mu(\emptyset)=0$ and countable additivity for disjoint unions. It generalizes length, area, probability, and volume. \par\noindent\textbf{Applications:} measures provide the foundation for Lebesgue integration, probability theory, ergodic theory, and geometric analysis.

\term{Measure Theory}
Measure theory studies measurable sets, measures, and integrals. It extends length, area, and volume to abstract spaces and handles limits and discontinuous functions systematically. \par\noindent\textbf{Applications:} it underlies probability, modern analysis, stochastic processes, PDEs, and mathematical statistics.

\term{Mellin Transform}
The Mellin transform is $F(s)=\int_0^\infty x^{s-1}f(x)\,dx$. It converts multiplicative scaling into translation in transform space and is closely related to the Laplace transform. \par\noindent\textbf{Applications:} it analyzes asymptotic behavior, scale-invariant systems, special functions, and zeta functions.

\term{Member}
An element of a set; the objects contained in a set.

\term{Menger's Theorem}
Menger's theorem states that the maximum number of internally vertex-disjoint paths between two vertices equals the minimum number of vertices whose removal separates them. An edge version equates edge-disjoint paths with minimum edge cuts. \par\noindent\textbf{Applications:} it measures network connectivity and identifies the number of independent routes or failures a network can tolerate.

\term{Meromorphic Function}
A meromorphic function is holomorphic except at isolated poles, where it may tend to infinity like a finite negative power. Locally it is a quotient of holomorphic functions. \par\noindent\textbf{Applications:} meromorphic functions support residue calculus, differential equations, complex dynamical systems, and analytic number theory.

\term{Method of Characteristics}
The method of characteristics solves certain first-order PDEs by following curves along which the PDE becomes an ODE. Information is transported along these characteristic curves; crossing characteristics can create shocks or loss of smoothness. \par\noindent\textbf{Applications:} it models wave and transport equations, fluid flow, traffic, and conservation laws.

\term{Metric Space}
A metric space is a set with a distance $d$ satisfying nonnegativity, identity of indiscernibles, symmetry, and the triangle inequality. These axioms define convergence, continuity, compactness, and completeness beyond Euclidean space. \par\noindent\textbf{Applications:} metric spaces model physical distance, function spaces, data sets, and optimization landscapes.

\term{Microlocal Analysis}
Microlocal analysis studies functions and distributions simultaneously by location and frequency in phase space. It tracks singularities and their propagation rather than only measuring global smoothness. \par\noindent\textbf{Applications:} it analyzes wave fronts, PDE regularity, scattering, inverse problems, and imaging.

\term{Minimal Polynomial}
The minimal polynomial of an operator $A$ is the unique monic polynomial of least degree satisfying $p(A)=0$. It divides every annihilating polynomial and has the same distinct roots as the characteristic polynomial. \par\noindent\textbf{Applications:} it determines diagonalizability, reduces matrix powers, and simplifies linear differential systems.

\term{Minimal Spanning Tree}
In a connected weighted graph, a minimum spanning tree connects every vertex without cycles and has the smallest total edge weight. Kruskal's and Prim's algorithms find one efficiently. \par\noindent\textbf{Applications:} it designs inexpensive communication, electrical, road, and pipeline networks.

\term{Minkowski Space}
Minkowski space is four-dimensional spacetime with a flat Lorentzian metric, commonly $(-,+,+,+)$. Lorentz transformations preserve the interval between events, distinguishing timelike, lightlike, and spacelike separations. \par\noindent\textbf{Applications:} it is the geometric framework of special relativity and electromagnetism.

\term{Minkowski Sum}
The Minkowski sum of sets is $A\oplus B=\{a+b:a\in A,b\in B\}$. Adding a disk to a shape creates an offset or thickened shape, and convexity is preserved. \par\noindent\textbf{Applications:} it supports collision detection, robot motion planning, mathematical morphology, and convex geometry.

\term{Möbius Inversion Formula}
If $g(n)=\sum_{d\mid n}f(d)$, Möbius inversion recovers $f(n)=\sum_{d\mid n}\mu(d)g(n/d)$, where $\mu$ is the Möbius function. \par\noindent\textbf{Applications:} it separates divisor-sum data, counts primitive objects, and appears in number theory and combinatorial enumeration.

\term{Möbius Strip}
A Möbius strip is formed by joining the ends of a strip after a half-twist. It has one side and one boundary component and is non-orientable. \par\noindent\textbf{Applications:} it is a standard topological example and illustrates orientation, boundary, embeddings, and quotient spaces.

\term{Model Category}
A category with three distinguished classes of morphisms, the weak equivalences, cofibrations, and fibrations, satisfying the model axioms of Quillen. It provides an abstract framework for homotopy theory in which the localization of the category at the weak equivalences is the homotopy category. Model structures exist for topological spaces, simplicial sets, and chain complexes, and model categories formalize derived functors.

\term{Model (Logic)}
A model of a formal theory is a mathematical structure in which all the theory's axioms hold. A sentence is true in the model when its interpretation is true there. \par\noindent\textbf{Applications:} models connect formal logic with algebra, geometry, databases, and the study of what axioms can or cannot determine.

\term{Model Theory}
The branch of mathematical logic that studies structures and the first-order theories describing them. Central tools include the compactness theorem, the Löwenheim--Skolem theorems, elementary equivalence, saturated models, and type spaces. Morley's categoricity theorem and stability theory classify theories through the spectrum of their models. Model theory has deep applications in algebra, real and $p$-adic geometry, and analytic number theory.

\term{Modular Arithmetic}
Modular arithmetic identifies integers that have the same remainder modulo $n$; calculations effectively wrap around after $n$. For example, $17\equiv5\pmod{12}$. \par\noindent\textbf{Applications:} it supports clocks, checksums, cryptography such as RSA, pseudorandom generators, and computer algorithms.

\term{Modular Form}
A modular form of weight $k$ is a holomorphic function on the upper half-plane satisfying $f((az+b)/(cz+d))=(cz+d)^kf(z)$ for matrices in the modular group, together with controlled behavior at cusps. Its Fourier coefficients often carry arithmetic information. \par\noindent\textbf{Applications:} modular forms connect elliptic curves, $L$-functions, partition formulas, and number theory.

\term{Modular Group}
The modular group is $\operatorname{PSL}(2,\mathbb{Z})=\operatorname{SL}(2,\mathbb{Z})/\{I,-I\}$. Here $\operatorname{SL}(2,\mathbb{Z})$ consists of the $2\times2$ integer matrices with determinant $1$, and the quotient identifies a matrix with its negative. It acts on the upper half-plane by $z\mapsto(az+b)/(cz+d)$. \par\noindent\textbf{Applications:} this symmetry is fundamental in modular forms, elliptic curves, and arithmetic geometry.

\term{Module}
A module is like a vector space whose scalars come from a ring rather than a field. It is an abelian group with scalar multiplication satisfying distributive and associative laws. Modules may lack bases or unique coordinates. \par\noindent\textbf{Applications:} modules unify linear algebra, integer equations, representations, coding, and algebraic geometry.

\term{Moduli Space}
A space, scheme, stack, or analytic space whose points parametrize isomorphism classes of geometric objects, such as algebraic curves, vector bundles, or varieties. A coarse moduli space may lack a universal family, whereas a fine moduli space carries one; stacks are needed when objects have nontrivial automorphisms. Deformation theory describes the local structure of moduli spaces. Examples include the moduli $\mathcal{M}_g$ of curves and the moduli of elliptic curves.

\term{Modulus of Continuity}
The modulus of continuity is $\omega(\delta)=\sup_{|x-y|<\delta}|f(x)-f(y)|$. It quantifies the largest change in $f$ over inputs separated by less than $\delta$; $\omega(\delta)\to0$ is uniform continuity. \par\noindent\textbf{Applications:} it gives approximation errors, regularity estimates, and stability bounds in numerical analysis and PDEs.

\term{Moment (Statistics)}
A moment is a numerical summary of a probability distribution. The raw $k$th moment is $\mathbb{E}[X^k]$, while the $k$th central moment is $\mathbb{E}[(X-\mu)^k]$. The second central moment is the variance; standardized third and fourth central moments describe skewness and kurtosis. \par\noindent\textbf{Applications:} moments summarize location, spread, asymmetry, and tail behavior in data analysis.

\term{Moment Generating Function}
The moment-generating function is $M_X(t)=\mathbb E[e^{tX}]$ when finite near $0$. Its derivatives satisfy $M_X^{(n)}(0)=\mathbb E[X^n]$. \par\noindent\textbf{Applications:} it identifies distributions, computes sums of independent variables, and derives means, variances, and concentration bounds.

\term{Monad}
In category theory, a monad on a category $\mathcal{C}$ is an endofunctor $T$ with natural transformations $\eta: \mathrm{id} \to T$ (the unit) and $\mu: T^2 \to T$ (the multiplication) satisfying associativity and unit laws. Monads formalize algebraic theories and, in computer science, computational effects such as exceptions and state. Every adjunction induces a monad; the Eilenberg--Moore and Kleisli categories recover algebras and free algebras, respectively.

\term{Monoid}
A monoid is a set with an associative operation and an identity element. Unlike a group, its elements need not have inverses. Examples include strings under concatenation and functions under composition. \par\noindent\textbf{Applications:} monoids model sequences, computation, formal languages, and repeated processes.

\term{Monoidal Category}
A monoidal category has a tensor-like product $\otimes$, a unit object, and coherent associativity and unit isomorphisms. It abstracts tensor products of vector spaces and composition of systems. \par\noindent\textbf{Applications:} monoidal categories organize quantum processes, programming semantics, representation theory, and algebraic topology.

\term{Monomial}
A monomial is a coefficient times variables raised to nonnegative integer powers, such as $3x^2y$. A polynomial is a finite sum of monomials. \par\noindent\textbf{Applications:} monomials form the basis for polynomial algebra, Gröbner-basis computations, numerical approximation, and multivariable models.

\term{Monomorphism}
A monomorphism is a morphism $f:X\to Y$ that can be cancelled on the left: $f\circ g=f\circ h$ implies $g=h$. In sets it is exactly an injective function, though in other categories it need not look like literal inclusion. \par\noindent\textbf{Applications:} monomorphisms formalize embeddings and subobjects in algebra and geometry.

\term{Monotone Convergence Theorem}
For measurable functions $0\leq f_n\uparrow f$, the monotone convergence theorem states $\int f_n\,d\mu\to\int f\,d\mu$. This is stronger than the elementary theorem about bounded increasing real sequences. \par\noindent\textbf{Applications:} it justifies exchanging limits and integrals in probability and measure theory.

\term{Monotone Function}
A monotone function is nondecreasing or nonincreasing: $x\leq y$ implies $f(x)\leq f(y)$ or $f(x)\geq f(y)$. Monotone functions have one-sided limits and only countably many jumps on an interval. \par\noindent\textbf{Applications:} they model cumulative quantities, threshold rules, optimization constraints, and distribution functions.

\term{Monty Hall Problem}
In the Monty Hall problem, after a contestant chooses one of three doors, a host who knows the prize opens a different losing door. Switching then wins with probability $2/3$, while staying wins with $1/3$, because the host's action is informative rather than random. \par\noindent\textbf{Applications:} it illustrates conditional probability, Bayesian updating, and the importance of modeling information.

\term{Moore--Penrose Pseudoinverse}
The Moore--Penrose pseudoinverse $A^+$ extends inversion to rectangular or singular matrices. It gives the least-squares solution of $Ax\approx b$ with minimum norm and is characterized by four Penrose equations, including $AA^+A=A$. \par\noindent\textbf{Applications:} it supports regression, inverse problems, signal reconstruction, and rank-deficient systems.

\term{Morphism}
A morphism is a structure-preserving arrow between objects in a category. It may be a function, linear map, continuous map, or group homomorphism depending on the category. \par\noindent\textbf{Applications:} morphisms provide a common language for transformations, equivalence, composition, and universal constructions.

\term{Morse Theory}
Morse theory studies a smooth manifold $M$ through a Morse function $f: M \to \mathbb{R}$ whose critical points are non-degenerate. The Morse lemma gives a normal form near each critical point, and $M$ admits a handle decomposition producing a CW structure; the Morse inequalities bound its homology in terms of the critical points of each index. Infinite-dimensional variants, including Morse--Witten and Floer theory, underlie modern symplectic and gauge theory.

\term{Motives}
Grothendieck's conjectural universal cohomology theory underlying all Weil cohomology theories. The category of pure motives is constructed from smooth projective varieties by adjoining correspondences as morphisms, and the Chow and Künneth standard conjectures concern its structure. Mixed motives extend the construction to arbitrary varieties, motivic cohomology provides the associated cohomology theory, and motives are central to the Langlands program.

\term{Multigrid Method}
A multigrid method solves discretized PDEs on several grid resolutions. Fine grids remove local errors, while coarse grids efficiently remove slowly varying errors. \par\noindent\textbf{Applications:} multigrid accelerates elliptic PDE solvers in fluid flow, heat transfer, elasticity, imaging, and engineering simulation.

\term{Multilinear Map}
A multilinear map is linear in each argument when the others are held fixed. The determinant is multilinear in the columns of a matrix. \par\noindent\textbf{Applications:} multilinear maps describe tensors, volume, interactions among variables, differential forms, and higher-order data models.

\term{Multinomial Distribution}
The multinomial distribution counts outcomes in $n$ independent trials with $k$ categories and probabilities $p_1,\ldots,p_k$. Its probability is $\frac{n!}{x_1!\cdots x_k!}\prod_i p_i^{x_i}$ when $\sum x_i=n$. \par\noindent\textbf{Applications:} it models categorical data, survey counts, genetics, text frequencies, and contingency tables.

\term{Multiple Integral}
A multiple integral accumulates a function over a multidimensional region, such as $\iint_Df\,dA$ or $\iiint_Vf\,dV$. Fubini's theorem and changes of variables give practical ways to compute it under suitable hypotheses. \par\noindent\textbf{Applications:} multiple integrals calculate mass, probability, volume, work, and fields.

\term{Multiplicative Function}
An arithmetic function $f$ such that $f(mn) = f(m)f(n)$ for any two coprime integers $m$ and $n$. Examples include the Euler totient function $\varphi(n)$ and the Möbius function $\mu(n)$.

\term{Multiplicative Inverse}
An element $a^{-1}$ is a multiplicative inverse when $aa^{-1}=1$. Every nonzero field element has one; in $\mathbb Z/n\mathbb Z$, $a$ is invertible exactly when $\gcd(a,n)=1$. \par\noindent\textbf{Applications:} inverses solve equations, support modular cryptography, and define division in algebraic systems.

\term{Multiplicity}
The multiplicity of a polynomial root is the number of repeated factors $(x-r)$ it contributes. For an eigenvalue, algebraic multiplicity is its multiplicity in the characteristic polynomial. \par\noindent\textbf{Applications:} multiplicities distinguish repeated solutions, determine diagonalizability, and describe resonance and degeneracy.

\term{Multiset}
A multiset allows an element to occur several times; the occurrence count is its multiplicity. Unlike an ordinary set, multiplicity is retained. \par\noindent\textbf{Applications:} multisets model repeated measurements, inventory, words, chemical compositions, and counting problems.

\term{Myhill--Nerode Theorem}
The Myhill--Nerode theorem states that a language is regular exactly when the relation identifying strings with the same possible continuations has finitely many equivalence classes. Those classes are the states of the minimal deterministic finite automaton. \par\noindent\textbf{Applications:} it proves non-regularity and constructs minimal automata for compilers, pattern matching, and verification.

\chapter*{N}
\label{chap:n}

\term{Nash Equilibrium}
A Nash equilibrium is a strategy profile in which no player can improve their payoff by changing strategy alone while others keep theirs fixed. Equilibria may be pure or mixed and need not be socially optimal. \par\noindent\textbf{Applications:} Nash equilibrium models competition, bargaining, auctions, evolution, and network behavior.

\term{Natural Transformation}
A natural transformation assigns a morphism between the outputs of two functors for every object, in a way that commutes with every map in the category. The commuting condition makes the construction independent of arbitrary choices. \par\noindent\textbf{Applications:} natural transformations formalize equivalence of constructions in algebra, topology, geometry, and theoretical computer science.

\term{Noether's Theorem}
Noether's theorem states that every suitable continuous symmetry of an action produces a conservation law. Time-translation symmetry gives energy conservation, spatial translation gives momentum, and rotational symmetry gives angular momentum. \par\noindent\textbf{Applications:} it organizes classical mechanics, field theory, fluid dynamics, and particle physics.

\term{Navier--Stokes Equations}
The Navier--Stokes equations describe conservation of mass and momentum in a moving viscous fluid. They balance acceleration and pressure against viscous and external forces; incompressibility adds $\nabla\cdot u=0$. \par\noindent\textbf{Applications:} they model air flow, water, weather, blood, combustion, and industrial fluid systems. Global regularity for 3D incompressible flow remains a Millennium Prize problem.

\term{Necessary and Sufficient Condition}
A condition $A$ is necessary and sufficient for $B$ when $A\iff B$: $A$ implies $B$, and $B$ implies $A$. This gives an exact characterization rather than only a test or consequence. \par\noindent\textbf{Applications:} such conditions simplify proofs, algorithms, classification, and diagnostic criteria.

\term{Necessary Condition}
A is necessary for $B$ when $B\Rightarrow A$. If $A$ fails, $B$ is impossible, but $A$ alone may not guarantee $B$. \par\noindent\textbf{Applications:} necessary conditions eliminate infeasible designs, rule out solutions, and provide preliminary tests in optimization and differential equations.

\term{Neighborhood (Topology)}
A neighborhood of $x$ is a set containing an open set that contains $x$. Neighborhoods define local concepts such as convergence, continuity, interior, and differentiability without requiring a distance. \par\noindent\textbf{Applications:} they provide the basic language for topology, manifolds, analysis, and local approximation.

\term{Néron Model}
The smooth group scheme over the ring of integers, or more generally over a Dedekind base, that extends an abelian variety and satisfies the Néron mapping property: morphisms from the generic fiber extend uniquely to morphisms from the whole base scheme. Its special fiber carries the group of connected components. Néron models are central to arithmetic geometry and the theory of abelian varieties over number fields.

\term{Nested Set}
A nested family has successive inclusion, such as $S_1\subseteq S_2\subseteq\cdots$ or decreasing inclusion. Nested sets support limiting constructions and completeness arguments. \par\noindent\textbf{Applications:} they describe feasible-region refinement, interval methods, fractals, probability events, and approximation schemes.

\term{Net (Topology)}
A net generalizes a sequence by allowing its indices to range over any directed set. Nets characterize convergence and closure in arbitrary topological spaces, including spaces where sequences are insufficient. \par\noindent\textbf{Applications:} they support general compactness arguments, functional analysis, and convergence of approximations.

\term{Network Flow}
Network-flow theory studies quantities moving through directed edges with capacities. The max-flow min-cut theorem says the greatest source-to-sink flow equals the capacity of a smallest separating cut. \par\noindent\textbf{Applications:} flow algorithms solve transportation, bandwidth allocation, matching, scheduling, and supply-chain problems.

\term{Neumann Boundary Condition}
A Neumann boundary condition prescribes the normal derivative on a boundary: $\partial u/\partial n=g$. It often represents flux through the boundary, such as heat flow or electric current. \par\noindent\textbf{Applications:} Neumann conditions model insulated or flux-controlled boundaries in PDEs, heat transfer, and electrostatics.

\term{Newcomb's Paradox}
A puzzle about an infallible predictor who has placed money in boxes: one box always contains \$1000 and a second contains \$1,000,000 exactly when the predictor expects you to take only the second box. Two-boxing dominates in payoff, but the predictor's accuracy makes one-boxing optimal, exposing a tension between causal and evidential decision theory.

\term{Newton Polytopes}
The Newton polytope of a polynomial is the convex hull of the exponent vectors of its nonzero monomials. Its geometry records which terms dominate in different directions. \par\noindent\textbf{Applications:} Newton polytopes study polynomial roots, sparse systems, tropical geometry, mixed volumes, and computational algebra.

\term{Newton--Raphson Method}
Newton's method finds a root using $x_{n+1}=x_n-f(x_n)/f'(x_n)$. It replaces the function locally by its tangent line and is usually quadratically convergent near a simple root, but poor initial guesses or small derivatives can cause failure. \par\noindent\textbf{Applications:} it solves nonlinear equations in engineering, calibration, optimization, and scientific computing.

\term{Nilpotent Element}
An element $x$ of a ring is nilpotent if $x^n=0$ for some positive integer $n$. The least such $n$ is its nilpotency index. Nilpotents behave like algebraic infinitesimals. \par\noindent\textbf{Applications:} they reveal singular or nonreduced structure in rings and algebraic varieties.

\term{Nilpotent Group}
A group is nilpotent if its central series reaches the whole group in finitely many steps. Nilpotent groups generalize abelian groups and are always solvable, though not every solvable group is nilpotent. \par\noindent\textbf{Applications:} they provide manageable models of noncommutativity and appear in Lie theory, matrix groups, and local group structure.

\term{Nilpotent Ideal}
An ideal $I$ is nilpotent if $I^n=0$ for some $n$. Every element of $I$ is then nilpotent, and the ideal records a part of the ring that disappears under repeated multiplication. \par\noindent\textbf{Applications:} nilpotent ideals describe infinitesimal thickenings and help separate reduced structure from algebraic multiplicity.

\term{Nilpotent Ring}
A ring where every element is nilpotent, or equivalently, the ring itself is a nilpotent ideal of itself.

\term{Nodal Domain}
A nodal domain is a connected component of the region where an eigenfunction is positive or negative, separated by its zero set. For a vibrating membrane, nodal domains are regions oscillating with the same sign. \par\noindent\textbf{Applications:} nodal domains study eigenvalue problems, vibrations, quantum states, and spectral geometry.

\term{Node (Graph Theory)}
A node or vertex is a basic object in a graph, connected to other nodes by edges. Nodes may represent entities, locations, states, or variables. \par\noindent\textbf{Applications:} graph nodes model routers, people, cities, tasks, molecules, and states in a finite automaton.

\term{Noetherian Module}
A module is Noetherian if every ascending chain of submodules eventually stabilizes, equivalently if every submodule is finitely generated. Noetherian modules prevent endlessly increasing algebraic complexity. \par\noindent\textbf{Applications:} they guarantee finite descriptions in algebra, algebraic geometry, and module computations.

\term{Noetherian Ring}
A ring is Noetherian if every ascending chain of ideals stabilizes. Equivalently, every ideal is finitely generated. \par\noindent\textbf{Applications:} Noetherianity guarantees finite algebraic calculations and is fundamental to polynomial rings, algebraic varieties, and Gröbner-basis methods.

\term{Non-archimedean Field}
A non-archimedean field has an absolute value satisfying $|x+y|\leq\max(|x|,|y|)$. The $p$-adic field $\mathbb Q_p$ is the main example; its geometry is ultrametric, so triangles are often isosceles with very strong nesting. \par\noindent\textbf{Applications:} non-archimedean fields support number theory, coding, arithmetic geometry, and $p$-adic computation.

\term{Non-commutative Geometry}
Noncommutative geometry studies spaces through algebras of functions that may not commute. The algebra is treated as the primary object, even when no ordinary point-set space exists. \par\noindent\textbf{Applications:} it connects operator algebras, quantum physics, index theory, foliations, and mathematical models of quantum spaces.

\term{Nonconvex Optimization}
Nonconvex optimization has a nonconvex objective or feasible set, so it may contain many local optima, saddle points, or disconnected feasible regions. Global optimization is generally difficult. \par\noindent\textbf{Applications:} it appears in neural networks, engineering design, matrix factorization, robotics, and physical inverse problems.

\term{Non-degenerate}
A bilinear or sesquilinear form is nondegenerate if $B(x,y)=0$ for every $y$ implies $x=0$. For a square matrix this corresponds to full rank and a nonzero determinant. \par\noindent\textbf{Applications:} nondegeneracy guarantees solvable coordinate changes, well-defined metrics, unique pairings, and stable geometric structures.

\term{Non-Euclidean Geometry}
Non-Euclidean geometries alter Euclid's parallel postulate. Hyperbolic geometry has infinitely many parallels through an external point, while elliptic geometry has none. \par\noindent\textbf{Applications:} these geometries model curved surfaces, relativity, navigation, computer graphics, and geometric group theory.

\term{Nonlinear}
A system is nonlinear when its response does not satisfy superposition: scaling or adding solutions need not produce another solution. Nonlinear systems can have bifurcations, chaos, shocks, and multiple equilibria. \par\noindent\textbf{Applications:} they model fluids, climate, biological populations, circuits, and complex engineering systems.

\term{Non-orientable Surface}
A non-orientable surface has no globally consistent choice of clockwise orientation or side; a loop can reverse orientation. The Möbius strip and Klein bottle are examples. \par\noindent\textbf{Applications:} non-orientability is a central topological property and appears in surface classification, geometry, and physical models with orientation reversal.

\term{Non-standard Analysis}
Nonstandard analysis extends the real numbers with infinitesimal and infinite elements and uses a transfer principle to preserve suitable first-order statements. It gives rigorous versions of intuitive infinitesimal calculus. \par\noindent\textbf{Applications:} it provides alternative proofs and models for limits, derivatives, integrals, probability, and mathematical economics.

\term{Normal Bundle}
For a submanifold $M$ inside a Riemannian manifold, the normal bundle has fiber $N_pM$, the orthogonal complement of the tangent space at $p$. It records directions leaving the submanifold. \par\noindent\textbf{Applications:} normal bundles describe tubular neighborhoods, curvature, embeddings, perturbations, and intersection theory.

\term{Normal Distribution}
The normal distribution $N(\mu,\sigma^2)$ has density $\frac1{\sqrt{2\pi\sigma^2}}e^{-(x-\mu)^2/(2\sigma^2)}$. It is centered at $\mu$ with variance $\sigma^2$ and arises as a limit of normalized sums in the central limit theorem. \par\noindent\textbf{Applications:} it models measurement errors, noise, regression residuals, and confidence intervals.

\term{Normal Extension}
An algebraic field extension $K/F$ is normal if every irreducible polynomial over $F$ having a root in $K$ splits completely in $K$, equivalently if $K$ is stable under all $F$-automorphisms in an algebraic closure. Splitting fields are normal, and a finite extension is normal and separable exactly when it is Galois. Composita of normal extensions are again normal, making normality a key ingredient of the fundamental theorem of Galois theory.

\term{Normal Family}
A family of holomorphic functions on a domain is normal if it is precompact in the topology of locally uniform convergence, so that every sequence contains a locally uniformly convergent subsequence. Montel's theorem characterizes normality through local boundedness. Normality is the engine behind proofs of the Picard theorems and is fundamental in value distribution theory and the iteration of holomorphic maps.

\term{Normal Form}
A normal form is a standard representative of an object under an allowed equivalence, such as Jordan form or Smith form. It exposes structural information and makes comparison or computation easier. \par\noindent\textbf{Applications:} normal forms solve classification, matrix, polynomial, rewriting, and symbolic-computation problems.

\term{Normality}
In topology, normality means that any two disjoint closed sets have disjoint open neighborhoods. This separation property supports continuous functions that distinguish closed sets. \par\noindent\textbf{Applications:} normality underlies extension and metrization theorems and is automatic for metric spaces.

\term{Normalization}
Normalization rescales an object to satisfy a chosen standard, such as $\|v\|=1$ for a vector or total mass $1$ for a probability distribution. The exact scaling depends on context. \par\noindent\textbf{Applications:} it compares data on common scales, creates unit directions, and converts weights into probabilities.

\term{Normalizer (Group Theory)}
The normalizer of $H\leq G$ is $N_G(H)=\{g\in G:gHg^{-1}=H\}$. It is the largest subgroup of $G$ in which $H$ is normal and measures the symmetries preserving $H$ under conjugation. \par\noindent\textbf{Applications:} normalizers analyze subgroup structure, group actions, and symmetry reductions.

\term{Normal Map}
In homotopy theory, a map $f: X \to Y$ that is a "normal" map in the sense of surgery theory, typically involving a bundle map from the normal bundle of $X$ to a bundle over $Y$.

\term{Normal Operator}
A bounded operator is normal when $TT^*=T^*T$. On finite-dimensional complex inner-product spaces, the spectral theorem says it is unitarily diagonalizable. Hermitian, unitary, and diagonal operators are normal. \par\noindent\textbf{Applications:} normal operators simplify spectral analysis, quantum observables, Fourier methods, and numerical computations.

\term{Normal Space}
A topological space in which every two disjoint closed sets have disjoint open neighbourhoods. Such spaces are also called $T_4$ spaces. Every metric space is normal. Normality is a key property in the proof of the Urysohn Metrization Theorem and the Tietze Extension Theorem.

\term{Normal Space (Differential Geometry)}
The orthogonal complement of the tangent space at a point $p$ on a manifold $M \subset \mathbb{R}^n$.

\term{Normal Subgroup}
A subgroup $N\trianglelefteq G$ is normal if $gNg^{-1}=N$ for every $g\in G$. This makes the quotient $G/N$ a group, and every homomorphism's kernel is normal. \par\noindent\textbf{Applications:} normal subgroups produce quotient symmetries and support classification and homomorphism theorems.

\term{Normed Space}
A normed space is a vector space with a norm, and therefore a metric $d(x,y)=\|x-y\|$. If every Cauchy sequence converges in the space, it is a Banach space. \par\noindent\textbf{Applications:} normed spaces measure function and signal size and provide the setting for functional analysis, approximation, and PDEs.

\term{Norm}
See \term{Norm (Vector Space)}.

\term{Norm (Vector Space)}
A norm is a function satisfying $\|v\|=0\iff v=0$, $\|\alpha v\|=|\alpha|\|v\|$, and $\|u+v\|\leq\|u\|+\|v\|$. It measures vector size and induces distance and convergence. \par\noindent\textbf{Applications:} norms quantify error, regularization, signal magnitude, and stability in numerical and functional analysis.

\term{Nuclear Space}
A locally convex space whose topology is generated by Hilbert seminorms satisfying the nuclear (trace-class) property. Fréchet nuclear spaces, such as the Schwartz space, are the domain of the Schwartz kernel theorem, and the class of nuclear spaces is stable under tensor products. Nuclear spaces provide the natural setting for distribution theory.

\term{Null Hypothesis}
The null hypothesis $H_0$ is a baseline claim, often that there is no effect or difference. A statistical test measures how surprising the data would be under $H_0$; failing to reject it is not the same as proving it true. \par\noindent\textbf{Applications:} null hypotheses structure experiments, clinical studies, quality control, and scientific inference.

\term{Nullity}
The nullity of a linear map is the dimension of its kernel. In finite dimensions, rank--nullity gives $\operatorname{rank}T+\operatorname{nullity}T=\dim V$. \par\noindent\textbf{Applications:} nullity measures non-uniqueness, redundant variables, unobservable modes, and degrees of freedom in linear systems.

\term{Null Set}
A null set has measure zero, meaning it contributes no mass, length, area, or volume under the chosen measure. A property holding almost everywhere may fail on a null set. \par\noindent\textbf{Applications:} null sets allow negligible exceptions in probability, integration, differential equations, and geometric measure theory.

\term{Null Space}
The null space of a matrix $A$ is $\ker A=\{v:Av=0\}$. It consists of directions invisible to the transformation and determines whether solutions of $Ax=b$ are unique. \par\noindent\textbf{Applications:} null spaces identify constraints, redundancy, inverse-problem ambiguity, and degrees of freedom.

\term{Nullstellensatz}
See Hilbert's Nullstellensatz: the correspondence between ideals of a polynomial ring and algebraic varieties.

\term{Null Vector}
A null vector has zero value under a quadratic or indefinite form, $g(v,v)=0$. In a positive-definite inner product this forces $v=0$, but in spacetime a nonzero null vector represents lightlike motion. \par\noindent\textbf{Applications:} null vectors describe light cones, relativistic signals, and degeneracy in bilinear-form geometry.

\term{Number Theory}
Number theory studies integers and related arithmetic structures. Elementary, analytic, algebraic, and computational branches investigate divisibility, primes, equations, fields, and algorithms. \par\noindent\textbf{Applications:} number theory supports cryptography, coding, pseudorandom generation, secure protocols, and error detection.

\term{Numerical Analysis}
Numerical analysis develops algorithms for approximate solutions of mathematical problems. It studies convergence, stability, conditioning, discretization, and error bounds. \par\noindent\textbf{Applications:} it enables simulation and computation in engineering, physics, finance, weather prediction, imaging, and scientific research.

\term{Numerical Optimization}
Numerical optimization designs algorithms for approximate minima or maxima of functions, possibly with constraints. Methods include gradient, Newton, trust-region, and derivative-free approaches. \par\noindent\textbf{Applications:} it supports parameter estimation, machine learning, engineering design, control, and resource allocation.

\term{Numerical Quadrature}
Numerical quadrature approximates a definite integral by a weighted sum of function values, as in the trapezoidal, Simpson, or Gaussian rules. Accuracy depends on smoothness, spacing, and the chosen rule. \par\noindent\textbf{Applications:} quadrature evaluates integrals in simulation, probability, finite elements, and physics.

\term{Numerical Stability}
An algorithm is numerically stable when rounding and input errors produce controlled output errors, ideally comparable to the problem's conditioning. Stability is distinct from whether the underlying problem is well-conditioned. \par\noindent\textbf{Applications:} stability analysis guides reliable linear algebra, differential-equation solvers, simulation, and scientific software.

\term{Numeric Integration}
Numerical integration approximates an integral using discrete samples and weights, such as the trapezoidal or Simpson rule. Refinement and error estimates determine reliability. \par\noindent\textbf{Applications:} it computes areas, expectations, energies, and model outputs when antiderivatives are unavailable.

\term{Nyquist Frequency}
For a sampling rate $f_s$, the Nyquist frequency is $f_s/2$, the highest frequency that can be represented without aliasing under ideal band-limited assumptions. A signal must be sampled at more than twice its maximum frequency. \par\noindent\textbf{Applications:} this rule guides digital audio, imaging, communications, sensors, and anti-aliasing filter design.
\term{Nonlinear Programming}
Nonlinear programming optimizes a nonlinear objective or constraints, including smooth, nonsmooth, convex, nonconvex, and mixed-integer cases. Local methods may depend strongly on initialization; KKT conditions extend Lagrange multipliers to inequalities. \par\noindent\textbf{Applications:} it models engineering design, energy systems, machine learning, finance, and nonlinear control.

\chapter*{O}
\label{chap:o}

\term{Object}
An object is an entity in a category to which morphisms are connected. In the category of sets, the objects are sets; in other categories they may be groups, spaces, or vector spaces. \par\noindent\textbf{Applications:} objects provide the units that are compared and transformed in abstract mathematics.

\term{Occurrence}
An occurrence is one appearance of a symbol or term at a particular position in a formula or expression. The same symbol may have several occurrences with different meanings or scopes. \par\noindent\textbf{Applications:} occurrences are important in parsing, formal logic, compilers, and symbolic algebra.

\term{Octonion}
An octonion is an element of an eight-dimensional real division algebra that extends the complex numbers and quaternions. Multiplication is not associative, but every nonzero octonion has an inverse. \par\noindent\textbf{Applications:} octonions appear in exceptional geometry, rotations, theoretical physics, and some models of string theory.

\term{Octahedron}
An octahedron is a Platonic solid with eight equilateral triangular faces, six vertices, and twelve edges. It is dual to the cube, meaning their faces and vertices correspond. \par\noindent\textbf{Applications:} octahedral geometry appears in crystallography, molecular structure, graphics, and spatial design.

\term{Octave}
In some mathematical contexts, an octave is an informal name for an octonion. More broadly, the term can refer to a doubling step in the construction of normed algebras. \par\noindent\textbf{Applications:} this terminology occurs in discussions of composition algebras and exceptional algebraic structures.

\term{Odd Permutation}
An odd permutation is a product of an odd number of transpositions. Its sign is $-1$, whereas even permutations have sign $+1$ and form the alternating group $A_n$. \par\noindent\textbf{Applications:} permutation parity is used in determinants, sorting, combinatorics, and group theory.

\term{One-Point Compactification}
The one-point compactification adds one point at infinity to a locally compact Hausdorff space $X$. A neighborhood of infinity is the complement of a compact subset of $X$. \par\noindent\textbf{Applications:} it turns unbounded spaces into compact ones and is useful in topology, complex analysis, and dynamical systems.

\term{One-Sided Limit}
A one-sided limit describes the value approached as $x$ tends to $a$ from the right or left. A function is continuous at $a$ when both one-sided limits exist, agree, and equal $f(a)$. \par\noindent\textbf{Applications:} one-sided limits analyze jumps, boundary behavior, piecewise functions, and numerical approximations.

\term{Operad}
An operad records operations with several inputs and one output, together with rules for composing those operations. It separates the pattern of composition from the particular objects on which the operations act. \par\noindent\textbf{Applications:} operads organize homotopy theory, deformation theory, algebraic topology, and structures in mathematical physics.

\term{Open Cover}
An open cover of a space is a collection of open sets whose union contains the whole space. Compactness means that every open cover has a finite subcover. \par\noindent\textbf{Applications:} open covers define compactness, manifolds, sheaves, partitions of unity, and local-to-global arguments.

\term{Open Interval}
An open interval $(a,b)$ contains all real numbers strictly between $a$ and $b$ but excludes the endpoints. It is an open set in the usual topology on the real line. \par\noindent\textbf{Applications:} open intervals describe domains, neighborhoods, convergence ranges, and feasible strict inequalities.

\term{Open Mapping Theorem}
The open mapping theorem states that a surjective bounded linear operator between Banach spaces maps open sets to open sets. It also implies that a bijective bounded linear operator has a bounded inverse. \par\noindent\textbf{Applications:} it establishes stability and continuous solvability in functional analysis and operator equations.

\term{Open Problem}
An open problem is a precisely stated mathematical question for which no accepted proof or disproof is known. Its status can change when new research resolves it. \par\noindent\textbf{Applications:} open problems guide research and often motivate new methods, conjectures, algorithms, and theories.

\term{Open Set}
An open set contains a neighborhood of each of its points. In a metric space, every point of an open set has a small open ball contained in the set. \par\noindent\textbf{Applications:} open sets define continuity, convergence, differentiability, manifolds, and the topology of function spaces.

\term{Operation}
An operation is a rule that combines one or more inputs to produce an output, such as addition, multiplication, or composition. Its properties, such as associativity or commutativity, determine the algebraic structure. \par\noindent\textbf{Applications:} operations model combination, transformation, computation, and composition of processes.

\term{Operator (Linear)}
A linear operator is a linear map $T:V\to W$ satisfying $T(au+bv)=aT(u)+bT(v)$. On Hilbert spaces, operators may be bounded or unbounded, and their domains matter. \par\noindent\textbf{Applications:} operators represent derivatives, integrations, matrices, quantum observables, and evolution equations.

\term{Operator Algebra}
Operator algebra studies algebras of bounded operators on Hilbert spaces, including C*-algebras and von Neumann algebras. Algebraic operations and analytic norms interact in these structures. \par\noindent\textbf{Applications:} operator algebras support quantum mechanics, functional analysis, noncommutative geometry, and statistical mechanics.

\term{Operator Norm}
The operator norm is $\|T\|=\sup_{\|x\|=1}\|Tx\|$. It is the largest factor by which $T$ can stretch a unit vector and satisfies $\|Tx\|\leq\|T\|\|x\|$. \par\noindent\textbf{Applications:} it measures sensitivity, bounds errors, and analyzes stability and convergence of operators and numerical algorithms.

\term{Optimal Control}
Optimal control chooses an input $u(t)$ to minimize a cost while satisfying system dynamics and constraints. Pontryagin's principle gives necessary conditions, while dynamic programming gives Bellman or Hamilton--Jacobi--Bellman equations. \par\noindent\textbf{Applications:} optimal control is used in robotics, aerospace, energy systems, economics, and finance.

\term{Orbit (Group Action)}
The orbit of $x$ under a group action is $\operatorname{Orb}(x)=\{g\cdot x:g\in G\}$. It contains every point reachable from $x$ by the allowed symmetries. For finite groups, orbit--stabilizer gives $|\operatorname{Orb}(x)|=[G:G_x]$. \par\noindent\textbf{Applications:} orbits classify symmetry classes and count configurations in geometry, chemistry, and combinatorics.

\term{Ordinary Differential Equation (ODE)}
An ordinary differential equation relates an unknown function of one independent variable to its derivatives. Initial or boundary data select particular solutions. \par\noindent\textbf{Applications:} ODEs model motion, circuits, population dynamics, chemical reactions, and control systems.

\term{Orthogonal Complement}
For a subspace $V$ of an inner-product space $W$, the orthogonal complement is $V^\perp=\{w:\langle w,v\rangle=0\text{ for every }v\in V\}$. In finite-dimensional Euclidean spaces, $W=V\oplus V^\perp$. \par\noindent\textbf{Applications:} orthogonal complements support projections, least squares, Fourier expansions, and constraint decomposition.

\term{Orthogonal Group}
The orthogonal group $O(n)$ consists of matrices $Q$ satisfying $Q^{\mathsf T}Q=I$. Its transformations preserve inner products, lengths, angles, and distances. \par\noindent\textbf{Applications:} orthogonal groups describe rotations and reflections in geometry, mechanics, computer graphics, and physics.

\term{Orthogonal Matrix}
An orthogonal matrix satisfies $Q^{\mathsf T}Q=I$. Its columns form an orthonormal basis, its inverse is $Q^{\mathsf T}$, and its determinant is $1$ or $-1$. \par\noindent\textbf{Applications:} orthogonal matrices provide stable coordinate changes, rotations, reflections, and QR-based numerical algorithms.

\term{Orthogonal Projection}
An orthogonal projection maps a vector to its closest point in a subspace. Its matrix satisfies $P^2=P$ and $P^*=P$. \par\noindent\textbf{Applications:} projections are used in least squares, signal filtering, approximation, and separating components of data.

\term{Orthogonality}
Two vectors are orthogonal when their inner product is zero. In Euclidean space this means they are perpendicular, although the definition also works in function spaces. \par\noindent\textbf{Applications:} orthogonality simplifies coordinates, Fourier series, regression, signal separation, and numerical computation.

\term{Orthonormal Basis}
An orthonormal basis consists of mutually orthogonal unit vectors that span the space. Every finite-dimensional inner-product space has one, for example by Gram--Schmidt. \par\noindent\textbf{Applications:} orthonormal bases simplify projections, Fourier expansions, quantum states, and stable numerical representations.

\term{Osculating Circle}
The osculating circle is the circle that matches a curve's position, tangent, and curvature at a point. Its radius is $1/|\kappa|$ when the curvature $\kappa$ is nonzero. \par\noindent\textbf{Applications:} it measures local bending and is used in computer-aided design, robotics, road geometry, and motion planning.

\term{Outer Measure}
An outer measure assigns a nonnegative size to every subset, is monotone, and is countably subadditive. Carathéodory's criterion selects the sets on which the outer measure is additive, producing a measurable space. \par\noindent\textbf{Applications:} outer measures construct Lebesgue measure and handle size questions for irregular sets.

\term{Overdetermined System}
An overdetermined system has more equations than unknowns. It may have no exact solution, so least squares finds the vector that minimizes the residual $\|Ax-b\|$. \par\noindent\textbf{Applications:} overdetermined systems arise in measurement, regression, calibration, tomography, and parameter estimation.

\term{Oval}
An oval is a smooth, closed, convex plane curve. Every ellipse is an oval, but many ovals are not ellipses. \par\noindent\textbf{Applications:} ovals model rounded shapes in geometry, optics, mechanical design, and computer graphics.

\term{O-minimal Structure}
An o-minimal structure is an ordered mathematical structure in which every definable subset of the line is a finite union of points and intervals. This prevents uncontrolled one-dimensional oscillation. \par\noindent\textbf{Applications:} o-minimality supports tame geometry, real algebraic geometry, optimization, and differential equations.

\term{Order (Group Theory)}
The order of a finite group is its number of elements, written $|G|$. The order of an element $g$ is the least positive integer $n$ for which $g^n=e$, if such an integer exists. \par\noindent\textbf{Applications:} order restricts subgroup structure, periodic behavior, and symmetries in finite groups.

\term{Order (Field)}
A field order is a total order compatible with addition and multiplication: $a\leq b$ implies $a+c\leq b+c$, and nonnegative elements have nonnegative products. The real numbers admit such an order, while the complex numbers do not. \par\noindent\textbf{Applications:} ordered fields support inequalities, limits, optimization, and real analysis.

\term{Order of Growth}
Order of growth describes how a quantity behaves as its input becomes large. Big-O notation gives an upper asymptotic bound, while $\Theta$ describes matching growth. \par\noindent\textbf{Applications:} growth rates compare algorithm efficiency, resource use, convergence speed, and model scaling.

\term{Order-Preserving Map}
An order-preserving map satisfies $x\leq y\Rightarrow f(x)\leq f(y)$. It carries the order information from one partially ordered set to another. \par\noindent\textbf{Applications:} such maps represent monotone processes, lattice homomorphisms, ranking systems, and fixed-point algorithms.

\term{Order-Statistic}
The $k$th order statistic is the $k$th smallest observation after sorting a sample. The median is a central order statistic, and the minimum and maximum are the first and last. \par\noindent\textbf{Applications:} order statistics support robust summaries, quantiles, outlier detection, and nonparametric inference.

\term{Ordered Pair}
An ordered pair $(a,b)$ records both entries and their positions. Equality means $(a,b)=(c,d)$ exactly when $a=c$ and $b=d$. \par\noindent\textbf{Applications:} ordered pairs represent coordinates, relations, input-output pairs, and points in a Cartesian product.

\term{Ordering}
An ordering is a relation that organizes elements. A partial order is reflexive, antisymmetric, and transitive; a total order also compares every pair. \par\noindent\textbf{Applications:} orderings support scheduling, dependency analysis, sorting, lattices, and optimization.

\term{Ordinary Least Squares (OLS)}
Ordinary least squares estimates regression coefficients by minimizing the sum of squared residuals between observed and fitted responses. It provides coefficient estimates, fitted values, and uncertainty measures under statistical assumptions. \par\noindent\textbf{Applications:} OLS is used for prediction, calibration, trend analysis, and experimental modeling.

\term{Orthocenter}
The orthocenter is the common intersection of the three altitudes of a triangle. It lies outside an obtuse triangle and inside an acute triangle. \par\noindent\textbf{Applications:} the orthocenter is used in triangle geometry and illustrates concurrency and perpendicularity.

\term{Orthogonal Polynomials}
A sequence of polynomials is orthogonal with respect to a weight $w$ when $\int w(x)p_n(x)p_m(x)dx=0$ for $n\ne m$. Legendre, Chebyshev, and Hermite polynomials are standard examples. \par\noindent\textbf{Applications:} orthogonal polynomials support approximation, Gaussian quadrature, spectral PDE methods, and probability.

\term{Orlicz Space}
An Orlicz space generalizes $L^p$ by measuring size through a Young function $\Phi$ rather than only a power. It can describe growth between or beyond standard power laws. \par\noindent\textbf{Applications:} Orlicz spaces are used in nonlinear PDEs, harmonic analysis, elasticity, and models with nonstandard growth.

\term{Orthonormal Set}
An orthonormal set consists of unit vectors that are pairwise orthogonal. Such a set is automatically linearly independent, although it may not span the whole space. \par\noindent\textbf{Applications:} orthonormal sets provide efficient coordinates, signal decompositions, and basis construction.

\term{Orthonormality}
Orthonormality means that vectors are pairwise orthogonal and each has norm one. An orthonormal set becomes an orthonormal basis when it also spans the space. \par\noindent\textbf{Applications:} orthonormality simplifies projections, numerical computation, Fourier analysis, and quantum-state representations.

\term{Outlier}
An outlier is an observation unusually far from the main pattern of a dataset. It may result from error, a rare event, or a genuine subpopulation, so it should not be removed automatically. \par\noindent\textbf{Applications:} outlier analysis supports quality control, fraud detection, robust statistics, and anomaly monitoring.

\term{Outer Product}
The outer product of vectors $u$ and $v$ is the matrix $uv^{\mathsf T}$ with entries $(uv^{\mathsf T})_{ij}=u_iv_j$. Unlike an inner product, it produces a matrix. \par\noindent\textbf{Applications:} outer products create rank-one updates, covariance estimates, tensor decompositions, and attention or feature interactions.

\term{Outer Automorphism}
An outer automorphism is a group automorphism not obtained by conjugation within the group. The quotient $\operatorname{Out}(G)=\operatorname{Aut}(G)/\operatorname{Inn}(G)$ measures automorphisms external to the group's internal symmetries. \par\noindent\textbf{Applications:} outer automorphisms support classification and extension problems in group theory.

\term{Overlap}
An overlap is the part common to two or more sets, intervals, or regions. For measurable sets $A$ and $B$, it is represented by $A\cap B$ and its size by $\mu(A\cap B)$. \par\noindent\textbf{Applications:} overlap measures intersections in geometry, image matching, scheduling, databases, and probability.

\term{Ordinal Number}
An ordinal number describes the order type of a well-ordered set. The first infinite ordinal is $\omega$, the order type of the natural numbers, followed by larger transfinite ordinals. \par\noindent\textbf{Applications:} ordinals support transfinite induction, recursion, and the analysis of well-founded hierarchies in set theory.

\term{Overline}
An overline can denote the closure $\overline A$ of a set, the complex conjugate $\overline z$, or an average depending on context. The closure contains the set and all of its limit points. \par\noindent\textbf{Applications:} overlines describe limits and boundaries in topology, conjugation in complex algebra, and averages in statistics.

\chapter{P}
\label{chap:p}

\term{P-adic Number}
Elements of the field $\mathbb{Q}_p$, constructed as the completion of the rational numbers $\mathbb{Q}$ with respect to the $p$-adic valuation $|x|_p = p^{-v_p(x)}$. These numbers are fundamental in modern number theory and arithmetic geometry.

\term{P-adic Valuation}
A function $v_p: \mathbb{Q} \to \mathbb{Z} \cup \{\infty\}$ where $v_p(n)$ is the exponent of the highest power of $p$ that divides $n$. It satisfies the ultrametric inequality: $v_p(x+y) \geq \min(v_p(x), v_p(y))$.

\term{Painter's Paradox}
The observation connected with Gabriel's horn: the horn can be filled with a finite volume of paint, yet covering its surface seems to require an infinite amount. Under an idealized model with paint of no thickness the two statements conflict; in practice a finite film of paint on an infinite area still uses infinite volume, and the paradox highlights the difference between surface area and volume.

\term{Palindromic Sequence}
A sequence of symbols or numbers that reads the same forwards and backwards. In combinatorics on words, the study of palindromes relates to the structure of free monoids and Sturmian words.

\term{Parabolic PDE}
A partial differential equation of heat type, with one distinguished direction of time and smoothing evolution.

\term{Parallelepiped}
A three-dimensional figure whose six faces are all parallelograms. Its volume is given by the absolute value of the scalar triple product: $V = | \mathbf{a} \cdot (\mathbf{b} \times \mathbf{c}) |$.

\term{Parallel Transport}
A way of moving a vector along a curve on a manifold such that the vector "remains parallel" to itself, meaning its covariant derivative along the curve is zero. The result depends on the path taken if the manifold has curvature.

\term{Parameter}
A quantity that can be varied and appears as a constant in an equation, or the coordinate of a parametric curve.

\term{Parametric Equation}
An equation describing a curve through functions of one or more parameters, such as $x = \cos t$, $y = \sin t$.

\term{Parity}
The evenness or oddness of an integer; also the property of a function that is even or odd.

\term{Parrondo's Paradox}
Two games that are each losing can be combined, played alternately or by a random rule, into a winning strategy. It arises from capital-dependent or history-dependent strategies and connects to the theory of Brownian ratchets in physics; the winning arises from the nonlinearity of the combined process.

\term{Parseval's Identity}
The equality $\sum |a_n|^2 = \int |f|^2$ relating the Fourier coefficients of a function to its energy.

\term{Partial Derivative}
The derivative of a multivariable function with respect to one variable, while holding all other variables constant. It measures the rate of change along a specific coordinate axis.

\term{Partial Fraction}
The decomposition of a rational function into a sum of simpler rational terms, useful for integration.

\term{Partial Sum}
The sum of the first $n$ terms of a series, $S_n = a_1 + \cdots + a_n$.

\term{Partial Order}
A binary relation $\leq$ on a set $S$ that is reflexive ($a \leq a$), antisymmetric ($a \leq b \land b \leq a \implies a = b$), and transitive ($a \leq b \land b \leq c \implies a \leq c$).

\term{Partition (Integer)}
A way of writing a positive integer $n$ as a sum of positive integers, where the order of the summands does not matter. The partition function $p(n)$ grows rapidly and is studied using generating functions and Hardy-Ramanujan asymptotics.

\term{Path-Connected Space}
A topological space $X$ where any two points $x, y \in X$ can be joined by a continuous path $\gamma: [0,1] \to X$ such that $\gamma(0)=x$ and $\gamma(1)=y$. Path-connectedness implies connectedness.

\term{Path Integral}
A formulation of quantum mechanics developed by Richard Feynman, where the probability amplitude for a particle to travel from $a$ to $b$ is the sum over all possible paths, weighted by $e^{iS/\hbar}$ where $S$ is the action.

\term{Pattern}
A regular, repeating or predictable structure, studied formally in combinatorics and formal language theory.

\term{Pearson Correlation Coefficient}
A measure of the linear correlation between two variables, $\rho_{XY} = \operatorname{Cov}(X,Y) / (\sigma_X \sigma_Y)$, taking values in $[-1, 1]$. It is the sample correlation used in linear regression.

\term{Perfect Number}
A positive integer that is equal to the sum of its proper positive divisors. Examples include 6, 28, and 496. All known perfect numbers are even and given by the formula $2^{p-1}(2^p - 1)$ where $2^p-1$ is a Mersenne prime.

\term{Perfect Field}
A field $K$ such that every irreducible polynomial over $K$ has distinct roots in its splitting field. All fields of characteristic 0 and all finite fields are perfect.

\term{Perfect Matching}
In graph theory, a set of edges such that every vertex of the graph is incident to exactly one edge in the set. Also called a 1-factor.

\term{Permutation}
A bijection from a set to itself. A permutation of a finite set of $n$ elements is one of the $n!$ possible linear orderings.

\term{Permutation Group}
A group whose elements are permutations of a given set and whose operation is the composition of permutations. The symmetric group $S_n$ is the group of all permutations of $n$ elements.

\term{Permutation Matrix}
A square binary matrix that has exactly one entry of 1 in each row and each column and 0s elsewhere. Multiplying a matrix by a permutation matrix reorders its rows or columns.

\term{Perron--Frobenius Theorem}
A non-negative irreducible square matrix has a positive real eigenvalue equal to its spectral radius, with a strictly positive eigenvector; this eigenvalue is simple and dominates all others in modulus. It is the foundation for the convergence of Markov chains and for stable states in population dynamics, such as those modeled by Leslie matrices.

\term{Phase Space}
In classical mechanics, a space in which all possible states of a system are represented, with plot coordinates asCanadian position and momentum. The volume of a region in phase space is preserved by Hamiltonian flow (Liouville's theorem).

\term{Phi ($\phi$) (Golden Ratio)}
The irrational number $\phi = \frac{1+\sqrt{5}}{2} \approx 1.618$, which is the positive solution to $x^2 - x - 1 = 0$. It appears in the geometry of pentagons, the Fibonacci sequence, and phyllotaxis.

\term{Picard--Lindelöf Theorem}
For the initial value problem $y' = f(x,y)$, $y(x_0) = y_0$, if $f$ is continuous in $x$ and Lipschitz continuous in $y$, then a unique solution exists on some interval about $x_0$. The proof typically proceeds via the contraction mapping principle applied to a space of continuous functions.

\term{Pigeonhole Principle}
A combinatorial principle stating that if $n$ items are put into $m$ containers, with $n > m$, then at least one container must contain more than one item.

\term{Pointwise Convergence}
A sequence of functions $f_n$ converges pointwise to $f$ on a set $S$ if for every $x \in S$, the sequence of values $f_n(x)$ converges to $f(x)$. This is a weaker condition than uniform convergence.

\term{Poisson Distribution}
A discrete probability distribution that expresses the probability of a given number of events occurring in a fixed interval of time or space, provided these events occur with a known constant mean rate and independently of the time since the last event.

\term{Poisson Equation}
A partial differential equation of the form $\Delta u = f$, where $\Delta$ is the Laplacian. It describes the potential $u$ generated by a source density $f$, appearing in electrostatics and gravity.

\term{Poisson Process}
A stochastic process $N(t)$ that counts the number of events occurring up to time $t$, where the increments are independent and the number of events in any interval follows a Poisson distribution.

\term{Poisson Summation Formula}
A relationship between the sum of values of a function at integers and the sum of values of its Fourier transform at integers: $\sum_{n=-\infty}^\infty f(n) = \sum_{k=-\infty}^\infty \hat{f}(k)$.

\term{Polytope}
The generalization of a polygon (2D) and polyhedron (3D) to $n$ dimensions. A convex polytope is the convex hull of a finite set of points in $\mathbb{R}^n$.

\term{Polynomial}
An expression consisting of variables and coefficients, using only addition, subtraction, multiplication, and non-negative integer exponents of variables.

\term{Polynomial Ring}
The ring $R[x]$ consisting of all polynomials with coefficients in a ring $R$. If $R$ is a field, $R[x]$ is a Euclidean domain.

\term{Polar Coordinates}
A two-dimensional coordinate system in which each point is determined by a distance $r$ from a reference point (the origin) and an angle $\theta$ from a reference direction.

\term{Polar Decomposition}
The factorization of a complex square matrix $A$ as $A = UP$, where $U$ is a unitary matrix and $P$ is a positive semi-definite Hermitian matrix.

\term{Pontryagin Duality}
A duality between a locally compact abelian (LCA) group $G$ and its dual group $\widehat{G}$, which consists of all continuous characters of $G$. The double dual is canonically isomorphic to $G$.

\term{Poincaré Disk Model}
A model of hyperbolic geometry where the hyperbolic plane is represented by the interior of the unit disk in $\mathbb{R}^2$. Geodesics are arcs of circles orthogonal to the boundary.

\term{Poincaré Recurrence Theorem}
A theorem in dynamical systems stating that for a measure-preserving transformation on a finite-measure space, almost every point in any set $A$ will eventually return to $A$ infinitely often.

\term{Poincaré Conjecture}
A famous theorem stating that every simply connected, closed 3-manifold is homeomorphic to the 3-sphere. It was proven by Grigori Perelman (2003) using Ricci flow.

\term{Poincaré Duality}
A fundamental result in manifold topology: for a compact oriented $n$-manifold $M$, there is an isomorphism $H^k(M; \mathbb{Z}) \cong H_{n-k}(M; \mathbb{Z})$.

\term{Population}
The entire set of objects of interest in a statistical study, from which samples are drawn.

\term{Power Set}
The set of all subsets of a set $S$, denoted $\mathcal{P}(S)$ or $2^S$. By Cantor's theorem, the cardinality of the power set is strictly greater than the cardinality of $S$.

\term{Power Series}
An infinite series of the form $\sum_{n=0}^\infty a_n (x-c)^n$. Each power series has a radius of convergence $R$, within which the series converges absolutely.

\term{Preconditioner}
A matrix or operator that transforms a linear system into one with better conditioning for iterative solvers.

\term{Preimage}
For a function $f: X \to Y$ and a subset $B \subseteq Y$, the preimage $f^{-1}(B)$ is the set of all $x \in X$ such that $f(x) \in B$.

\term{Pre-order}
A binary relation that is reflexive and transitive, but not necessarily antisymmetric. A partial order is a pre-order that is also antisymmetric.

\term{Presheaf}
An assignment of data to open sets together with restriction maps, compatible with inclusions of open sets.

\term{Primal Problem}
The original optimization problem, to which a dual problem is associated.

\term{Prime Ideal}
An ideal $P$ such that $ab \in P$ implies $a \in P$ or $b \in P$; the quotient ring is an integral domain.

\term{Prime Number Theorem}
The number $\pi(x)$ of primes at most $x$ is asymptotically $\pi(x) \sim x / \ln x$. An equivalent statement is that $\pi(x)$ is asymptotic to the logarithmic integral $\operatorname{li}(x)$, which approximates $\pi(x)$ even more closely. The size of the error term depends on the zeros of the Riemann zeta function, and the Riemann hypothesis would imply a sharp square-root error bound.

\term{Principal Ideal}
An ideal $I$ in a ring $R$ that is generated by a single element $a$: $I = (a) = \{ra : r \in R\}$.

\term{Principal Component Analysis (PCA)}
A statistical procedure that uses an orthogonal transformation to convert a set of observations of possibly correlated variables into a set of values of linearly uncorrelated variables called principal components.

\term{Principal Fiber Bundle}
A fiber bundle $\pi: P \to B$ together with a Lie group $G$ that acts freely and transitively on each fiber. These are used to define gauge theories in physics.

\term{Principal Ideal Domain (PID)}
An integral domain in which every ideal is a principal ideal. Examples include the ring of integers $\mathbb{Z}$ and the ring of polynomials $F[x]$ over a field $F$.

\term{Principal Value (Cauchy)}al
The Cauchy principal value is a method for assigning a value to certain divergent integrals, such as $\int_{-a}^a \frac{1}{x}\,dx$, by taking a symmetric limit around the singularity.

\term{Probability Density Function}
The function $f$ such that $P(a \leq X \leq b) = \int_a^b f(x)\,dx$ for a continuous random variable $X$.

\term{Probability Space}
A mathematical triple $(\Omega, \mathcal{F}, P)$ consisting of a sample space $\Omega$, a $\sigma$-algebra $\mathcal{F}$ of events, and a probability measure $P$.

\term{Product Space}
The Cartesian product of two or more sets, spaces, or groups, equipped with the product structure (e.g., the product topology or product measure).

\term{Product Topology}
The coarsest topology on a product $\prod X_i$ such that all projection maps $\pi_i$ are continuous. A basis for this topology consists of products of open sets, where all but finitely many are the whole space $X_i$.

\term{Projection (Linear Algebra)}
A linear operator $P: V \to V$ such that $P^2 = P$. It decomposes the space $V$ into the direct sum of the image of $P$ and the kernel of $P$.

\term{Projective Geometry}
The study of geometric properties that are invariant under projective transformations. It extends Euclidean geometry by adding "points at infinity."

\term{Projective Space}
The set of all lines passing through the origin in a vector space $\mathbb{R}^{n+1}$. It is denoted $\mathbb{P}^n(\mathbb{R})$ and can be viewed as the quotient of $\mathbb{R}^{n+1} \setminus \{0\}$ by the equivalence $v \sim \lambda v$.

\term{Projective Variety}
The set of common zeros of a collection of homogeneous polynomials in a projective space. Projective varieties are the central objects of study in classical algebraic geometry.

\term{Proper Map}
A continuous map $f: X \to Y$ between topological spaces is proper if the preimage of every compact set in $Y$ is compact in $X$.

\term{Primitive Root}
An integer $g$ is a primitive root modulo $n$ if every integer $a$ coprime to $n$ is congruent to a power of $g$ modulo $n$. This means $g$ generates the multiplicative group $(\mathbb{Z}/n\mathbb{Z})^\times$.

\term{Primitive Element Theorem}
A theorem in field theory stating that every finite separable extension $L/K$ is a simple extension, meaning $L = K(\alpha)$ for some $\alpha \in L$.

\term{Proportionality}
A relationship between two quantities where their ratio remains constant. If $y$ is proportional to $x$, then $y = kx$ for some constant $k$.

\term{Proximal Operator}
The map $\operatorname{prox}_f(x) = \arg\min_y (f(y) + \frac{1}{2}\|x-y\|^2)$, central to proximal methods in optimization.

\term{Pullback}
The map $f^*$ induced by a function $f$, sending objects on the codomain to objects on the domain, such as differential forms or bundles.

\term{Purely Imaginary}
A complex number of the form $bi$ with real part zero.

\term{Pushforward}
The map $f_*$ induced by a function $f$, pushing vectors from the tangent space of the domain to that of the codomain.

\chapter*{Q}
\label{chap:q}

\term{Q-analog}
A $q$-analog introduces a parameter $q$ into a formula so that the classical expression is recovered as $q\to1$. The $q$-binomial coefficient is a standard example. \par\noindent\textbf{Applications:} $q$-analogs occur in combinatorics, number theory, statistical mechanics, and quantum algebra.

\term{Q-matrix}
For a continuous-time Markov chain, a $Q$-matrix has nonnegative off-diagonal transition rates and row sums zero, so $q_{ii}=-\sum_{j\ne i}q_{ij}$. \par\noindent\textbf{Applications:} $Q$-matrices describe queues, reliability systems, population models, and chemical reaction networks.

\term{Q-series}
A $q$-series is a power series in $q$, often with coefficients encoding arithmetic or combinatorial data. The generating function for integer partitions is a famous example. \par\noindent\textbf{Applications:} $q$-series support partition theory, modular forms, combinatorics, and statistical mechanics.

\term{Quadratic Equation}
A quadratic equation has the form $ax^2+bx+c=0$ with $a\ne0$. Its roots are $x=(-b\pm\sqrt{b^2-4ac})/(2a)$, and the discriminant determines whether the roots are real or complex. \par\noindent\textbf{Applications:} quadratic equations model trajectories, optimization, geometry, and basic physical laws.

\term{Quadratic Form}
A quadratic form is a homogeneous degree-two expression $Q(v)=v^{\mathsf T}Av$, with $A$ taken symmetric. Its definiteness describes curvature and its level sets include ellipses, hyperbolas, and higher-dimensional analogues. \par\noindent\textbf{Applications:} quadratic forms appear in optimization, energy, statistics, geometry, and number theory.

\term{Quadratic Reciprocity}
For distinct odd primes $p$ and $q$, quadratic reciprocity states that $\left(\frac{p}{q}\right)\left(\frac{q}{p}\right)=(-1)^{\frac{p-1}{2}\frac{q-1}{2}}$. It relates whether $p$ is a square modulo $q$ to whether $q$ is a square modulo $p$. \par\noindent\textbf{Applications:} it makes quadratic-residue calculations efficient and supports number theory and cryptography.

\term{Quadratic Residue}
An integer $a$ is a quadratic residue modulo a prime $p$ if $x^2\equiv a\pmod p$ has a solution. Otherwise it is a nonresidue. \par\noindent\textbf{Applications:} quadratic residues support modular arithmetic, primality tests, cryptography, and number-theoretic reciprocity.

\term{Quadrature}
Quadrature means evaluating a definite integral, often by a weighted numerical sum. Gaussian quadrature chooses nodes and weights to integrate polynomials of high degree exactly. \par\noindent\textbf{Applications:} quadrature is used in simulation, probability, finite elements, and engineering computation.

\term{Quadrature Domain}
A quadrature domain is a bounded complex domain for which integrals of suitable holomorphic functions equal a finite combination of function values and derivatives at specified points. \par\noindent\textbf{Applications:} quadrature domains connect complex analysis with exact integration and potential-flow models.

\term{Quadrature Rule}
A quadrature rule approximates an integral by a weighted sum of function values, such as the trapezoidal, Simpson, or Gaussian rule. Accuracy depends on smoothness and node placement. \par\noindent\textbf{Applications:} quadrature rules evaluate integrals when elementary antiderivatives are unavailable.

\term{Qualitative}
Qualitative analysis describes properties such as stability, boundedness, or long-term behavior without finding an exact formula. \par\noindent\textbf{Applications:} it provides useful conclusions for dynamical systems, differential equations, and models that cannot be solved explicitly.

\term{Qualitative Analysis}
Qualitative analysis studies features such as stability, equilibria, limit cycles, and long-term behavior without finding an explicit formula. \par\noindent\textbf{Applications:} it analyzes ODEs, dynamical systems, population models, and physical systems that are not exactly solvable.

\term{Quantifier}
A quantifier specifies how widely a statement applies. The universal quantifier $\forall$ means “for every,” and the existential quantifier $\exists$ means “for at least one.” \par\noindent\textbf{Applications:} quantifiers express mathematical definitions, theorems, specifications, and computer-verification statements.

\term{Quantile}
A quantile is a cutoff dividing a distribution so that a specified fraction of observations lies below it. The $0.5$ quantile is the median, and the $0.25$ and $0.75$ quantiles define the quartiles. \par\noindent\textbf{Applications:} quantiles summarize skewed data, define thresholds, and support robust statistics.

\term{Quantile Function}
The quantile function is $Q(p)=\inf\{x:F(x)\geq p\}$, where $F$ is a cumulative distribution function. It returns a value below which a fraction $p$ of the distribution lies. \par\noindent\textbf{Applications:} quantile functions support simulation, risk thresholds, confidence intervals, and robust summaries.

\term{Quantitative}
Quantitative information is expressed through numbers, measurements, or numerical relationships. \par\noindent\textbf{Applications:} quantitative descriptions support modeling, comparison, prediction, experiments, and data-driven decision-making.

\term{Quantitative Analysis}
Quantitative analysis uses numerical data, mathematical models, and statistical methods to describe patterns and make predictions. \par\noindent\textbf{Applications:} it supports finance, engineering, epidemiology, experimentation, and data-driven policy.

\term{Quantity}
A quantity is a measurable or countable amount, such as length, mass, probability, or number. Quantities may be represented by numbers with units or by abstract algebraic elements. \par\noindent\textbf{Applications:} quantities provide the variables and parameters used in scientific models and measurements.

\term{Quantum State}
A quantum state is a unit vector, or physically a ray, in a Hilbert space. Measurement probabilities are determined by squared inner products, such as $|\langle\psi\mid\phi\rangle|^2$. \par\noindent\textbf{Applications:} quantum states describe qubits, atoms, molecules, quantum computation, and quantum communication.

\term{Quarter}
A quarter is one fourth of a whole. In the plane, a quadrant is one of the four regions determined by the coordinate axes. \par\noindent\textbf{Applications:} quarters and quadrants organize fractions, angles, coordinate geometry, and seasonal or financial data.

\term{Quarter-turn Symmetry}
Quarter-turn symmetry means that a $90$-degree rotation leaves an object unchanged. Repeating the rotation gives the cyclic group of order four. \par\noindent\textbf{Applications:} this symmetry appears in tilings, logos, crystal patterns, image processing, and geometric design.

\term{Quasicoherent Sheaf}
A quasicoherent sheaf is a sheaf of modules that locally corresponds to a module over the coordinate ring of an affine piece. It is the natural sheaf-theoretic form of algebraic data on a scheme. \par\noindent\textbf{Applications:} quasicoherent sheaves describe functions, ideals, vector bundles, and cohomology in algebraic geometry.

\term{Quasiconcave Function}
A function is quasiconcave when every superlevel set $\{x:f(x)\geq\alpha\}$ is convex. It is useful for functions whose high-value regions have no gaps. \par\noindent\textbf{Applications:} quasiconcavity appears in economics, utility theory, location problems, and optimization.

\term{Quasiconformal Map}
A quasiconformal map is a homeomorphism that may distort angles but has uniformly bounded eccentricity. Infinitesimal circles become ellipses whose axis ratio is controlled. \par\noindent\textbf{Applications:} quasiconformal maps are used in complex analysis, geometric topology, image registration, and shape deformation.

\term{Quasi-convex Function}
A function is quasiconvex when every sublevel set $\{x:f(x)\leq\alpha\}$ is convex. Every convex function is quasiconvex, but not conversely. \par\noindent\textbf{Applications:} quasiconvexity supports optimization when objectives are not strictly convex.

\term{Quasigroup}
A quasigroup is a set with a binary operation for which both $a*x=b$ and $y*a=b$ have unique solutions. It need not be associative or have an identity. \par\noindent\textbf{Applications:} quasigroups appear in algebraic combinatorics, Latin squares, coding, and cryptographic constructions.

\term{Quasi-isomorphism}
A quasi-isomorphism is a map of chain complexes that induces isomorphisms on all homology groups. Derived categories invert these maps, treating complexes with the same homological information as equivalent. \par\noindent\textbf{Applications:} quasi-isomorphisms simplify homological algebra, derived geometry, and algebraic topology.

\term{Quasilinear}
A quasilinear differential equation is linear in its highest-order derivatives but may be nonlinear in lower-order terms or coefficients. \par\noindent\textbf{Applications:} quasilinear equations model fluid flow, waves, transport, and nonlinear continuum mechanics.

\term{Quasi-metric Space}
A quasi-metric space has a distance-like function satisfying nonnegativity, identity, and the triangle inequality but not necessarily symmetry. Thus traveling from $x$ to $y$ may have a different cost from traveling back. \par\noindent\textbf{Applications:} quasi-metrics model directed networks, transportation costs, computer science, and asymmetric geometry.

\term{Quasi-Newton Method}
A quasi-Newton method updates an approximation to the Hessian or its inverse using changes in gradients, rather than computing second derivatives directly. BFGS is a standard example. \par\noindent\textbf{Applications:} quasi-Newton methods solve smooth optimization problems in statistics, machine learning, and engineering.

\term{Quasi-periodic Function}
A quasiperiodic function combines oscillations with incommensurate frequencies. It does not repeat exactly, but its values can return arbitrarily close to previous states. \par\noindent\textbf{Applications:} quasiperiodicity models celestial motion, oscillators, solid-state systems, and dynamical systems.

\term{Quasi-random Sequence}
A quasirandom sequence is a low-discrepancy sequence designed to cover a region uniformly rather than randomly. Sobol and Halton sequences are common examples. \par\noindent\textbf{Applications:} quasirandom sequences improve numerical integration, uncertainty quantification, finance, and computer graphics.

\term{Quaternion Algebra}
A quaternion algebra is a four-dimensional central simple algebra over its center that generalizes the usual quaternions. Its multiplication may be described by generators and relations, and it need not be a division algebra over every field. \par\noindent\textbf{Applications:} quaternion algebras are used in number theory, quadratic forms, arithmetic geometry, and the study of rotations.

\term{Quaternion Norm}
A quaternion norm is defined for $q=a+bi+cj+dk$ by $|q|=\sqrt{a^2+b^2+c^2+d^2}$. It is multiplicative: $|q_1q_2|=|q_1||q_2|$. \par\noindent\textbf{Applications:} the norm supports inversion, numerical stability, and rotation computations.

\term{Quaternions}
A quaternion is an element of the noncommutative division algebra with basis $1,i,j,k$ and relations $i^2=j^2=k^2=ijk=-1$. Unit quaternions represent spatial rotations without the singularities of some angle-coordinate systems. \par\noindent\textbf{Applications:} quaternions are used in computer graphics, robotics, aerospace, and physics.

\term{Quine}
A quine is a program that prints its own source code without reading it as external input. It is constructed using quotation and self-reference. \par\noindent\textbf{Applications:} quines illustrate recursion, fixed-point theorems, programming-language semantics, and computability.

\term{Quintic}
A quintic is a polynomial of degree five. Some quintics are solvable by radicals, but the general quintic is not. \par\noindent\textbf{Applications:} quintics illustrate polynomial classification and the role of Galois groups.

\term{Quintic Equation}
A quintic equation is a polynomial equation of degree five. The Abel--Ruffini theorem states that no formula using radicals solves every quintic, although particular quintics can be solvable. \par\noindent\textbf{Applications:} quintics illustrate the limits of symbolic algebra and the role of Galois theory in polynomial equations.

\term{Quotient Group}
If $N$ is normal in $G$, the quotient group $G/N$ consists of cosets with $(aN)(bN)=(ab)N$. It treats elements differing by $N$ as equivalent. \par\noindent\textbf{Applications:} quotient groups describe reduced symmetries, congruences, and the structure of homomorphisms.

\term{Quotient Manifold}
A quotient manifold is formed by identifying points of a manifold under a free, properly discontinuous group action. The quotient inherits a smooth structure, as in a torus obtained from a plane by lattice translations. \par\noindent\textbf{Applications:} quotient manifolds model periodic geometry, covering spaces, and symmetry reductions.

\term{Quotient Map}
A quotient map is a surjective map $f:X\to Y$ such that $U\subseteq Y$ is open exactly when $f^{-1}(U)$ is open in $X$. It creates the quotient topology and is continuous, but need not be open. \par\noindent\textbf{Applications:} quotient maps construct circles, projective spaces, and glued topological spaces.

\term{Quotient Ring}
For a two-sided ideal $I$ of a ring $R$, the quotient ring $R/I$ consists of cosets $r+I$ with induced addition and multiplication. For example, $\mathbb C\cong\mathbb R[x]/(x^2+1)$. \par\noindent\textbf{Applications:} quotient rings encode congruences, construct fields, and define algebraic varieties.

\term{Quotient Space}
A quotient space $X/{\sim}$ is formed by identifying points related by an equivalence relation. The quotient map sends each point to its equivalence class. \par\noindent\textbf{Applications:} quotient spaces construct tori, projective spaces, and spaces obtained by gluing boundaries.

\term{Quotient Topology}
The quotient topology on $Y=X/{\sim}$ is the finest topology for which the quotient map $q:X\to Y$ is continuous; equivalently, $U\subseteq Y$ is open when $q^{-1}(U)$ is open. \par\noindent\textbf{Applications:} it formalizes gluing and identification in topology and geometry.

\chapter*{R}
\label{chap:r}

\term{Radius of Convergence}
The radius of convergence is the distance $R$ from the center of a power series within which it converges absolutely. For $|z-c|<R$ it converges, and for $|z-c|>R$ it diverges. \par\noindent\textbf{Applications:} it determines where Taylor expansions can be used in analysis, physics, and numerical computation.

\term{Radon--Nikodym Theorem}
The Radon--Nikodym theorem states that if a $\sigma$-finite measure $\nu$ is absolutely continuous with respect to $\mu$, then $\nu(A)=\int_A f\,d\mu$ for a unique function $f$ up to $\mu$-almost everywhere equality. \par\noindent\textbf{Applications:} the derivative $f$ represents densities in probability, statistics, and mathematical finance.

\term{Riemann–Stieltjes Integral}
The Riemann--Stieltjes integral integrates $f$ with respect to a function of bounded variation $\alpha$, written $\int_a^b f\,d\alpha$. Taking $\alpha(x)=x$ gives the ordinary Riemann integral. \par\noindent\textbf{Applications:} it represents expectations, cumulative distributions, and systems with point masses.

\term{Radon Transform}
The Radon transform maps a function on $\mathbb R^n$ to its integrals over hyperplanes. It records projections from many directions and can be inverted under suitable conditions. \par\noindent\textbf{Applications:} it is the mathematical basis of CT reconstruction and other imaging methods.

\term{Radical (Ring Theory)}
In ring theory, a radical is an ideal that records elements with a specified form of algebraic degeneracy. The Jacobson radical is the intersection of maximal ideals, while the nilradical is the set of nilpotent elements. \par\noindent\textbf{Applications:} radicals help classify rings and analyze modules and algebraic varieties.

\term{Radical (Number Theory)}
For a positive integer $n$, the number-theoretic radical $\operatorname{rad}(n)$ is the product of the distinct primes dividing $n$. \par\noindent\textbf{Applications:} it measures the square-free part of an integer and appears in Diophantine equations and the $abc$ conjecture.

\term{Radon Measure}
A Radon measure is a Borel measure that is finite on compact sets and can be approximated regularly by compact and open sets. \par\noindent\textbf{Applications:} Radon measures provide a reliable framework for integration on locally compact spaces and for probability models.

\term{Ramsey's Theorem}
Ramsey's theorem states that every sufficiently large finitely colored complete graph contains a monochromatic complete subgraph of any prescribed size. The numbers $R(k,l)$ give the smallest size needed for red or blue structure. \par\noindent\textbf{Applications:} it explains unavoidable patterns in networks, scheduling, and combinatorial design.

\term{Random}
Random describes outcomes governed by chance and modeled by probability distributions. Randomness does not mean that the behavior is necessarily patternless; it means that individual outcomes cannot be predicted with certainty. \par\noindent\textbf{Applications:} random models support sampling, simulation, risk analysis, and statistical inference.

\term{Random Walk}
A random walk is a stochastic process formed by successive random steps, such as steps on the lattice $\mathbb Z^d$. Its long-term behavior depends on the step distribution and the underlying space. \par\noindent\textbf{Applications:} random walks model diffusion, search algorithms, ranking, and financial movement.

\term{Random Variable}
A random variable is a measurable function from a probability space to numerical values. It assigns a value to each outcome while allowing probabilities to be computed for ranges of values. \par\noindent\textbf{Applications:} random variables are the basic objects of statistics, probability models, and uncertainty quantification.

\term{Rank}
The rank of a linear map is the dimension of its image. For a matrix, it is also the maximum number of linearly independent rows or columns. \par\noindent\textbf{Applications:} rank measures information, identifies solvability, and supports regression and data reduction.

\term{Rank (Linear Algebra)}
The rank of a matrix is the dimension of its column space, equivalently its row space. It counts the independent directions represented by the matrix. \par\noindent\textbf{Applications:} matrix rank is used in linear systems, least squares, compression, and machine learning.

\term{Rank (Group Theory)}
In group theory, rank measures the minimum size of a generating set in settings where that notion is well defined. For a finitely generated abelian group, it counts the independent infinite-order generators. \par\noindent\textbf{Applications:} group rank describes algebraic complexity in topology and geometric group theory.

\term{Rank--Nullity Theorem}
The rank--nullity theorem states that for a linear map $T:V\to W$, $\dim V=\operatorname{rank}(T)+\operatorname{nullity}(T)$. The rank measures the output directions and the nullity measures the lost input directions. \par\noindent\textbf{Applications:} it tests invertibility and explains degrees of freedom in linear models.

\term{Rational Function}
A rational function is a ratio $f(x)=P(x)/Q(x)$ of two polynomials, defined where $Q(x)\ne0$. Its graph may have poles or removable holes. \par\noindent\textbf{Applications:} rational functions model transfer functions, interpolation, asymptotic behavior, and algebraic curves.

\term{Rational Root Theorem}
The rational root theorem states that if an integer-coefficient polynomial has a root $p/q$ in lowest terms, then $p$ divides the constant term and $q$ divides the leading coefficient. \par\noindent\textbf{Applications:} it gives a finite list of candidates for exact polynomial factorization.

\term{Rayleigh Quotient}
For a Hermitian matrix $A$ and nonzero vector $x$, the Rayleigh quotient is $R(A,x)=\frac{x^*Ax}{x^*x}$. Its maximum and minimum are the largest and smallest eigenvalues of $A$. \par\noindent\textbf{Applications:} it estimates eigenvalues and appears in vibration analysis, stability, and principal-component methods.

\term{Ricci Curvature}
Ricci curvature is a contraction of the Riemann curvature tensor, written locally as $\operatorname{Ric}_{ab}=R^c_{\ acb}$. It summarizes how volumes and geodesics deviate from their flat-space behavior. \par\noindent\textbf{Applications:} it is central to general relativity and geometric analysis.

\term{Rice's Theorem}
Rice's theorem states that every nontrivial semantic property of partial computable functions is undecidable. A property is semantic when it concerns what a program computes rather than how it is written. \par\noindent\textbf{Applications:} the theorem sets limits on automatic program analysis and verification.

\term{Richard's Paradox}
Richard's paradox constructs a number by changing the $n$th decimal digit of the $n$th number described by a finite phrase. The apparent contradiction shows that informal language cannot safely enumerate all definitions. \par\noindent\textbf{Applications:} it motivates formal languages, diagonal arguments, and the study of definability.

\term{Real Analysis}
Real analysis studies limits, sequences, functions, integration, and measure on the real numbers. It supplies precise definitions for continuity, convergence, and approximation. \par\noindent\textbf{Applications:} real analysis underlies differential equations, probability, optimization, and numerical analysis.

\term{Real Line}
The real line is the ordered set $\mathbb R$ of real numbers viewed as a one-dimensional continuum. It has no gaps and provides the usual setting for limits and intervals. \par\noindent\textbf{Applications:} it represents measured quantities, time, position, and scalar state variables.

\term{Reciprocal}
The reciprocal of a nonzero number $x$ is its multiplicative inverse $1/x$, satisfying $x(1/x)=1$. In $\mathbb Z/n\mathbb Z$, an element has a reciprocal exactly when it is coprime to $n$. \par\noindent\textbf{Applications:} reciprocals support division, rates, unit conversion, and modular cryptography.

\term{Recurrence Relation}
A recurrence relation defines sequence terms from earlier terms, such as $F_n=F_{n-1}+F_{n-2}$. Initial values are needed to determine a specific sequence. \par\noindent\textbf{Applications:} recurrences model population growth, algorithms, financial schedules, and discrete dynamical systems.

\term{Recursive}
Recursive means that a definition or algorithm refers to itself on smaller inputs and eventually reaches a base case. \par\noindent\textbf{Applications:} recursion expresses divide-and-conquer algorithms, inductive definitions, trees, and generated sequences.

\term{Recursive Function}
A recursive function is a function computable by a finite algorithm, equivalently by standard formal models such as Turing machines. A recursive set has a computable characteristic function, while a recursively enumerable set may only have an algorithm that lists its members. \par\noindent\textbf{Applications:} recursive functions define the limits of effective computation and support logic and programming-language theory.

\term{Reducible Polynomial}
A polynomial is reducible over a field if it factors into two nonconstant polynomials over that field. Otherwise it is irreducible. \par\noindent\textbf{Applications:} factorization supports solving polynomial equations, constructing field extensions, and coding algebraic data.

\term{Reduction (Algebra)}
Reduction is the process of replacing an object by a simpler but equivalent or controlled form, such as reducing a polynomial modulo another polynomial. \par\noindent\textbf{Applications:} reductions simplify algebraic calculations, prove computational hardness, and construct quotient structures.

\term{Reflection (Geometry)}
A reflection is an isometry that sends each point to its mirror image across a fixed line, plane, or hyperplane. It preserves distances and reverses orientation. \par\noindent\textbf{Applications:} reflections describe symmetry, image transformations, crystallography, and geometric algorithms.

\term{Reflexive Space}
A Banach space is reflexive when its natural map into its double dual is onto, so the space is canonically isomorphic to its double dual. Equivalently, its closed unit ball is weakly compact. \par\noindent\textbf{Applications:} reflexivity provides compactness tools for variational problems, optimization, and functional analysis.

\term{Regular Expression}
A regular expression is a symbolic pattern describing a regular language, using operations such as concatenation, choice, and repetition. It can be recognized by a finite automaton. \par\noindent\textbf{Applications:} regular expressions support text search, lexical analysis, validation, and data extraction.

\term{Regular Graph}
A regular graph is a graph in which every vertex has the same degree. A $k$-regular graph has degree $k$ at every vertex. \par\noindent\textbf{Applications:} regular graphs model balanced networks and support designs in coding, communication, and combinatorics.

\term{Regular Language}
A regular language is a set of finite strings recognized by a finite automaton or described by a regular expression. It has limited memory requirements. \par\noindent\textbf{Applications:} regular languages define tokens in compilers, search patterns, and simple communication protocols.

\term{Regular Sequence}
In a ring, a regular sequence $a_1,\ldots,a_n$ consists of elements where each $a_i$ is not a zero-divisor after the preceding elements have been factored out. \par\noindent\textbf{Applications:} regular sequences measure algebraic independence and define complete intersections in algebraic geometry.

\term{Regular Space}
A regular space is a topological space in which a point and a disjoint closed set can be separated by disjoint open sets. \par\noindent\textbf{Applications:} regularity supplies enough separation to construct continuous functions and study compactness.

\term{Relation (Set Theory)}
A relation from $X$ to $Y$ is a subset of $X\times Y$ that records which elements are connected. It is a function when each element of $X$ is related to exactly one element of $Y$. \par\noindent\textbf{Applications:} relations represent order, reachability, database links, and logical constraints.

\term{Relatively Prime}
Two integers are relatively prime, or coprime, when their greatest common divisor is $1$. \par\noindent\textbf{Applications:} coprimality controls modular inverses, reduced fractions, primitive lattice points, and public-key cryptography.

\term{Remainder Theorem}
The remainder theorem states that dividing a polynomial $f(x)$ by $x-a$ leaves remainder $f(a)$. \par\noindent\textbf{Applications:} it tests roots and quickly evaluates remainders during polynomial factorization.

\term{Renormalisation}
Renormalisation is a method for replacing divergent intermediate quantities by scale-dependent parameters that produce finite predictions. \par\noindent\textbf{Applications:} it is central to quantum field theory and also informs critical phenomena and statistical mechanics.

\term{Representation (Group Theory)}
A group representation is a homomorphism $\rho:G\to\operatorname{GL}(V)$ that realizes group elements as invertible linear maps. \par\noindent\textbf{Applications:} representations turn symmetry questions into matrix and linear-algebra problems in physics, chemistry, and harmonic analysis.

\term{Representation Theory}
Representation theory studies how groups, algebras, and related structures act by linear transformations. \par\noindent\textbf{Applications:} it classifies symmetries in quantum mechanics, molecular spectra, geometry, and signal analysis.

\term{Residue}
The residue of a meromorphic function at an isolated singularity is the coefficient of $(z-z_0)^{-1}$ in its Laurent expansion. \par\noindent\textbf{Applications:} residues evaluate contour integrals and many real integrals in complex analysis.

\term{Residue (Complex Analysis)}
The residue at $z_0$ is the coefficient of $(z-z_0)^{-1}$ in the Laurent expansion of $f$. It captures the part of the singularity that contributes to a contour integral. \par\noindent\textbf{Applications:} residues are used in complex integration, asymptotics, and inverse transforms.

\term{Residue Theorem}
The residue theorem states that for a positively oriented closed curve avoiding the poles, $\int_\gamma f(z)\,dz=2\pi i\sum\operatorname{Res}(f,z_k)$ over the enclosed poles. Reversing orientation changes the sign. \par\noindent\textbf{Applications:} it evaluates contour integrals and difficult real integrals.

\term{Resolution of Singularities}
Resolution of singularities replaces a singular algebraic variety by a smooth variety related to it through a birational map. In characteristic zero this can be achieved by controlled blow-ups. \par\noindent\textbf{Applications:} resolutions make intersection theory, cohomology, and geometric classification more manageable.

\term{Resultant}
The resultant is a polynomial in the coefficients of two polynomials that vanishes exactly when they share a root over an algebraic closure. \par\noindent\textbf{Applications:} resultants eliminate variables and test common roots without explicitly solving the polynomials.

\term{Retraction (Topology)}
A retraction is a continuous map $r:X\to A$ that fixes every point of the subspace $A$. Then $A$ is called a retract of $X$. \par\noindent\textbf{Applications:} retractions simplify topology, prove non-embedding results, and define deformation arguments.

\term{Reverse Polish Notation}
Reverse Polish notation writes operators after their operands, so expressions can be evaluated without parentheses. \par\noindent\textbf{Applications:} it simplifies stack-based calculators, compilers, and expression parsers.

\term{Riemann Mapping Theorem}
The Riemann mapping theorem states that every simply connected proper open subset of the complex plane is conformally equivalent to the unit disk. \par\noindent\textbf{Applications:} it transfers complex-analytic problems to a standard domain and supports potential-flow and boundary-value calculations.

\term{Riemann Sum}
A Riemann sum approximates an integral by adding the areas of thin rectangles whose heights are function values. Under suitable hypotheses, the sum converges as the mesh width tends to zero. \par\noindent\textbf{Applications:} Riemann sums provide the foundation for integration and numerical quadrature.

\term{Riemann Surface}
A Riemann surface is a one-dimensional complex manifold. It turns multivalued expressions such as $\sqrt z$ and $\log z$ into single-valued functions on a suitable surface. \par\noindent\textbf{Applications:} Riemann surfaces support complex analysis, algebraic curves, and mathematical physics.

\term{Riemann Zeta Function}
The Riemann zeta function is $\zeta(s)=\sum_{n=1}^\infty n^{-s}$ for $\operatorname{Re}(s)>1$, extended meromorphically to the complex plane. Its zeros encode information about prime numbers. \par\noindent\textbf{Applications:} it connects prime distribution, analytic number theory, and spectral models.

\term{Riemann--Hilbert Problem}
The Riemann--Hilbert problem seeks a holomorphic function with a prescribed jump across a contour and specified behavior at infinity. \par\noindent\textbf{Applications:} it solves problems in integrable systems, orthogonal polynomials, random matrices, and asymptotic analysis.

\term{Riemannian Manifold}
A Riemannian manifold is a smooth manifold with a smoothly varying positive-definite inner product on each tangent space. The metric defines lengths, angles, and geodesics. \par\noindent\textbf{Applications:} Riemannian manifolds model curved physical spaces and are fundamental in geometry and relativity.

\term{Riemannian Metric}
A Riemannian metric assigns a positive-definite inner product smoothly to every tangent space of a manifold. It determines lengths, angles, areas, and distances. \par\noindent\textbf{Applications:} metrics support geometric modeling, optimization on manifolds, and general relativity.

\term{Ring (Algebra)}
A ring is a set with addition and multiplication such that addition forms an abelian group, multiplication is associative, and multiplication distributes over addition. Some conventions also require a multiplicative identity. \par\noindent\textbf{Applications:} rings organize integers, polynomials, matrices, and algebraic structures used in coding and cryptography.

\term{Ring of Integers}
The ring of integers $\mathcal O_K$ of a number field $K$ is the set of elements of $K$ that satisfy monic polynomials with integer coefficients. Its ideals factor uniquely into prime ideals, even when elements do not factor uniquely. \par\noindent\textbf{Applications:} rings of integers support algebraic number theory, Diophantine equations, and algebraic cryptography.

\term{Riesz Representation Theorem}
The Riesz representation theorem states that every continuous linear functional on a Hilbert space has the form $f(x)=\langle x,y\rangle$ for a unique vector $y$. \par\noindent\textbf{Applications:} it identifies dual spaces and supports weak formulations of differential equations.

\term{Riesz Space}
A Riesz space is a real vector space equipped with an order in which vectors have lattice operations such as maximum and minimum. \par\noindent\textbf{Applications:} Riesz spaces model ordered functions and support integration and operator theory.

\term{Root (Polynomial)}
A root of a polynomial $f$ is a value $r$ for which $f(r)=0$. Over the complex numbers, every nonconstant polynomial has a root. \par\noindent\textbf{Applications:} roots determine factors, equilibria, resonances, and solutions of algebraic models.

\term{Root Mean Square}
The root mean square of $x_1,\ldots,x_n$ is $\sqrt{\frac1n\sum_i x_i^2}$. It measures magnitude while giving greater weight to large values. \par\noindent\textbf{Applications:} RMS is used for signal amplitude, vibration, electrical voltage, and error measurement.

\term{Ross--Littlewood Paradox}
The Ross--Littlewood paradox considers an urn in which balls are added and removed during infinitely many steps. The urn contains many balls after every finite step, yet each particular ball is eventually removed. \par\noindent\textbf{Applications:} it illustrates subtleties of infinite processes, limits, and interchange of operations.

\term{Rotational Symmetry}
Rotational symmetry means that a figure is unchanged after rotation through a specified angle about a fixed point. \par\noindent\textbf{Applications:} it guides geometric design, crystallography, robotics, and pattern recognition.

\term{Rotation Matrix}
A rotation matrix is an orthogonal matrix with determinant $1$ that preserves lengths and orientation. In two dimensions it rotates vectors through an angle $\theta$. \par\noindent\textbf{Applications:} rotation matrices describe rigid-body motion, computer graphics, robotics, and coordinate changes.

\term{Rolle's Theorem}
Rolle's theorem states that if $f$ is continuous on $[a,b]$, differentiable on $(a,b)$, and $f(a)=f(b)$, then some $c\in(a,b)$ satisfies $f'(c)=0$. \par\noindent\textbf{Applications:} it proves the mean value theorem and helps locate stationary points.

\term{Rouché's Theorem}
Rouché's theorem states that if $|f|<|g|$ on a closed contour, then $f$ and $f+g$ have the same number of zeros inside it, counted with multiplicity. \par\noindent\textbf{Applications:} it counts complex roots without solving the polynomial explicitly.

\term{Rough Path Theory}
Rough path theory defines integration and differential equations for paths too irregular for ordinary calculus, including many stochastic paths. \par\noindent\textbf{Applications:} it provides a pathwise framework for stochastic differential equations, signal analysis, and control.

\term{Round-off Error}
Round-off error is the difference caused by representing exact real numbers with finite-precision arithmetic. It can accumulate during long computations or cancellation. \par\noindent\textbf{Applications:} analyzing round-off error is essential for reliable numerical algorithms and scientific software.

\term{Row Echelon Form}
Row echelon form is a matrix form in which zero rows are at the bottom and each leading entry lies to the right of the leading entry above it. \par\noindent\textbf{Applications:} Gaussian elimination uses it to solve systems and determine rank.

\term{Rudin--Shapiro Sequence}
The Rudin--Shapiro sequence is a recursively defined binary sequence with controlled autocorrelation. \par\noindent\textbf{Applications:} it is studied in harmonic analysis, signal design, and the behavior of polynomials with signs $\{-1,1\}$.

\term{Ruler}
A ruler is a straightedge used to draw lines and compare or measure lengths. In classical geometry it is paired with a compass for constructions. \par\noindent\textbf{Applications:} rulers support surveying, technical drawing, geometry, and calibration.

\term{Runge--Kutta Method}
Runge--Kutta methods approximate solutions of ordinary differential equations by combining several slope evaluations at each step. Higher-order versions improve accuracy without requiring higher derivatives. \par\noindent\textbf{Applications:} they simulate motion, chemical reactions, circuits, and dynamical systems.

\term{Runge's Theorem}
Runge's theorem says that a function holomorphic near a compact set can be uniformly approximated there by rational functions whose poles lie in the allowed complementary components. If the complement of the compact set is connected, polynomials are enough. \par\noindent\textbf{Applications:} the theorem replaces analytic functions by simpler expressions for approximation and numerical computation.

\term{Russell's Paradox}
Russell's paradox considers the collection of all sets that do not contain themselves. Asking whether this collection contains itself produces a contradiction. \par\noindent\textbf{Applications:} the paradox motivated axiomatic set theory and careful foundations for mathematics.

\term{R-module}
An $R$-module is an abelian group on which a ring $R$ acts compatibly with addition. When $R$ is a field, it is a vector space. \par\noindent\textbf{Applications:} modules generalize vector spaces and organize linear equations, ideals, and algebraic invariants.

\term{R-linear map}
A map between $R$-modules is $R$-linear when $f(ax+by)=af(x)+bf(y)$ for $a,b\in R$. \par\noindent\textbf{Applications:} linear maps preserve module structure and represent transformations in algebra, coding, and systems theory.
\term{Prime Distribution}
Prime distribution studies how prime numbers occur among the integers. The prime number theorem gives the leading estimate $\pi(x)\sim x/\log x$. \par\noindent\textbf{Applications:} prime distribution supports analytic number theory and the security analysis of public-key cryptography.
\term{Riemann Hypothesis}
The Riemann hypothesis conjectures that every nontrivial zero of $\zeta(s)$ lies on $\operatorname{Re}(s)=1/2$. It remains unproved and would imply strong bounds on the error in prime-counting estimates. \par\noindent\textbf{Applications:} it guides research in number theory, spectral theory, and mathematical physics.
\term{Ramsey Theory}
A branch of combinatorics showing that sufficiently large systems contain ordered substructures under any finite coloring or partition. Ramsey's theorem guarantees monochromatic complete subgraphs in large enough complete graphs. \par\noindent\textbf{Applications:} Ramsey theory studies unavoidable patterns in networks, schedules, coding, and discrete mathematics.

\chapter{S}
\label{chap:s}

\term{S-matrix (Scattering Matrix)}
In quantum mechanics, a matrix that relates the initial and final states of a physical system that undergoes a scattering process. The elements represent the probability amplitudes for transitions between different states.

\term{S-adic Space}
A space constructed as the limit of a sequence of spaces, often used in the study of tilings and symbolic dynamics. It generalizes the notion of $p$-adic spaces to a sequence of different bases.

\term{Sample}
A subset of a population selected for statistical study; random samples aim to be representative.

\term{Sample Space}
The set of all possible outcomes of a random experiment, denoted $\Omega$.

\term{Scalar}
A quantity that is described by a single number, as opposed to a vector which has both magnitude and direction. In linear algebra, scalars are the elements of the field over which a vector space is defined.

\term{Scalar Curvature}
The trace of the Ricci curvature, $R = g^{ab}\operatorname{Ric}_{ab}$, a key invariant of Riemannian manifolds.

\term{Scalar Product}
Another name for the inner product. For two vectors $\mathbf{u}$ and $\mathbf{v}$ in $\mathbb{R}^n$, it is the sum $\sum u_i v_i$. It is used to define lengths and angles between vectors.

\term{Scale Factor}
The ratio by which the dimensions of a figure are multiplied under a similarity transformation.

\term{Scale Space}
A multi-scale representation of an image or signal, obtained by repeatedly smoothing the original data with a low-pass filter (typically Gaussian). It is used in computer vision to detect features at different resolutions.

\term{Schauder Basis}
A sequence of vectors $\{e_n\}$ in a Banach space $X$ such that every $x \in X$ has a unique representation $x = \sum_{n=1}^\infty a_n e_n$. Unlike Hamel bases, Schauder bases allow for infinite sums.

\term{Scheme}
The central object of study in modern algebraic geometry. A scheme is a locally ringed space that is locally isomorphic to the spectrum of a commutative ring. It generalizes the notion of an algebraic variety.

\term{Schröder--Bernstein Theorem}
A classical result in set theory, also called the Cantor--Bernstein theorem: if there are injections from $A$ into $B$ and from $B$ into $A$, then there exists a bijection between $A$ and $B$. It follows that cardinal comparison is antisymmetric, so two sets are equinumerous precisely when each is no larger than the other.

\term{Schur Complement}
For a block matrix $M = \begin{pmatrix} A & B \\ C & D \end{pmatrix}$, the Schur complement of the block $D$ is the matrix $A - B D^{-1} C$. It is used in matrix inversion and the study of positive-definite matrices.

\term{Schur's Lemma}
A fundamental result in representation theory: any irreducible representation of a group $G$ over an algebraically closed field has only scalar multiples of the identity as endomorphisms.

\term{Schur Decomposition}
The factorization of a square matrix $A$ as $A = QUQ^*$ where $Q$ is unitary and $U$ is upper triangular. The diagonal elements of $U$ are the eigenvalues of $A$.

\term{Schwarz Lemma}
A result in complex analysis: if $f$ is a holomorphic function from the unit disk to itself such that $f(0)=0$, then $|f(z)| \leq |z|$ and $|f'(0)| \leq 1$ for all $z$ in the disk.

\term{Schwarz Space}
The space of smooth functions whose derivatives all decay faster than any power of $1/|x|$ as $|x| \to \infty$. It is the natural domain for the Fourier transform, as the transform of a Schwarz function is also a Schwarz function.

\term{Schwartz Space}
The space $\mathcal{S}(\mathbb{R}^n)$ of rapidly decreasing smooth functions, on which the Fourier transform is an isomorphism.

\term{S-curve}
A sigmoid-shaped curve used in various fields to describe growth processes that start slowly, accelerate, and then level off as they reach a carrying capacity.

\term{Second Isomorphism Theorem}
If $G$ is a group with a subgroup $H$ and a normal subgroup $N$, then $HN/N \cong H/(H \cap N)$. This is one of the isomorphism theorems; the ring and module analogues state the same isomorphism when $N$ is an ideal or a submodule.

\term{Second-order PDE}
A partial differential equation where the highest derivative is of the second order. These are classified as elliptic, parabolic, or hyperbolic based on the coefficients of the second-order terms.

\term{Section (Fiber Bundle)}
A continuous map $s: B \to E$ from the base space $B$ to the total space $E$ of a fiber bundle such that $\pi \circ s = \text{id}_B$. A section picks one point from each fiber.

\term{Sector (Geometry)}
The portion of a disk enclosed by two radii and an arc. The area of a sector with radius $r$ and angle $\theta$ is $\frac{1}{2} r^2 \theta$.

\term{Seifert--van Kampen Theorem}
If $X = U \cup V$ with $U$, $V$, and $U \cap V$ path-connected, the fundamental group of $X$ is the amalgamated free product of $\pi_1(U)$ and $\pi_1(V)$ over $\pi_1(U \cap V)$. It allows the computation of fundamental groups by decomposing spaces into simpler parts: for instance, the wedge of $k$ circles has free group on $k$ generators, while $\pi_1(S^n)$ is trivial for $n \geq 2$.

\term{Self-adjoint Operator}
An operator equal to its adjoint, $T = T^*$; on Hilbert spaces self-adjoint operators have real spectrum and a spectral decomposition.

\term{Semi-algebraic Set}
A subset of $\mathbb{R}^n$ defined by a finite number of polynomial inequalities and equalities. They are the basic objects of study in real algebraic geometry.

\term{Semi-continuous Function}
A function is lower semi-continuous if its lower limit at any point is at least the function's value at that point; upper semi-continuous if the upper limit is at most the value.

\term{Semi-group}
A set equipped with an associative binary operation. Unlike a group, a semi-group does not require an identity element or inverses.

\term{Semi-norm}
A function that satisfies all the properties of a norm except for the identity of indiscernibles (i.e., $\|v\|=0$ does not imply $v=0$). It is used to define the topology of certain function spaces.

\term{Separable Space}
A topological space that contains a countable dense subset. For example, $\mathbb{R}$ is separable because $\mathbb{Q}$ is a countable dense subset.

\term{Separable Extension}
An algebraic extension $L/K$ where every element of $L$ is the root of a separable polynomial over $K$ (a polynomial with distinct roots in its splitting field).

\term{Separating}
(Topology) A family of functions or sets that distinguishes distinct points, as in the axioms of a Tychonoff space.

\term{Sequence}
An ordered list of elements, typically indexed by the natural numbers. A sequence $(a_n)$ converges to $L$ if for every $\epsilon > 0$ there exists $N$ such that $|a_n - L| < \epsilon$ for all $n \geq N$.

\term{Series}
The sum of the terms of a sequence. A series converges if the sequence of its partial sums has a finite limit.

\term{Set Theory}
The branch of mathematics studying sets and their operations, relations, and cardinalities, forming the foundation of mathematics.

\term{Sheaf}
A presheaf satisfying the locality and gluing axioms, allowing consistent local data to be assembled into global sections.

\term{Sieve Method}
A technique in number theory for counting or estimating the number of elements in a set that are not divisible by any of a given set of primes. The Sieve of Eratosthenes is the most basic example.

\term{Sigma-Algebra}
A collection of subsets of a set closed under complements and countable unions, providing the measurable sets of a measure space.

\term{Sigmoid Function}
A mathematical function having a characteristic "S"-shaped curve. The most common is the logistic function $f(x) = 1/(1+e^{-x})$, used in neural networks and population modeling.

\term{Simplex}
The generalization of a triangle to any dimension. An $n$-simplex is the convex hull of $n+1$ affinely independent points. A 0-simplex is a point, a 1-simplex is a line segment, and a 2-simplex is a triangle.

\term{Simple Group}
A group whose only normal subgroups are the identity and the group itself. Simple groups are the "building blocks" of all finite groups via composition series.

\term{Simple Loop}
A continuous path in a topological space that starts and ends at the same point and does not intersect itself except at the endpoints.

\term{Simple Root}
In the theory of Lie algebras, a root that cannot be written as a sum of two other positive roots. The set of simple roots forms a basis for the root system.

\term{Simpson's Paradox}
A trend that holds in each separate subgroup of data can be reversed when the subgroups are pooled together. It arises from lurking or confounding variables and unequal group sizes; the classic example is a treatment that looks better for every patient group yet worse overall. (Relates to the existing "Simpson's Rule" only by name.)

\term{Singular Matrix}
A square matrix that is not invertible, which occurs if and only if its determinant is zero. A singular matrix has a non-trivial kernel.

\term{Singular Value Decomposition (SVD)}
A factorization of any $m \times n$ matrix $A$ as $A = U \Sigma V^T$ where $U$ and $V$ are orthogonal matrices and $\Sigma$ is a diagonal matrix of singular values. It is used in data compression and noise reduction.

\term{Singular Point}
A point on a curve or surface where the derivative vanishes or is undefined. At such a point, the curve may have a cusp, a node, or a self-intersection.

\term{Singular Value}
The square roots of the eigenvalues of $A^* A$. They measure the "stretch" of the linear transformation $A$ along specific directions.

\term{Sine (Trigonometric)}
The ratio of the opposite side to the hypotenuse in a right triangle. For a unit circle, it is the $y$-coordinate of a point at angle $\theta$.

\term{Skew Lines}
Lines in space that are neither parallel nor intersecting.

\term{Skolem's Paradox}
By the Löwenheim--Skolem theorem, Zermelo--Fraenkel set theory (if consistent) has a countable model, yet it proves that uncountable sets exist. The resolution is that "uncountable" inside the model only means no internal bijection to the naturals; the apparent contradiction highlights the relativity of cardinality.

\term{Sleeping Beauty Problem}
A probability puzzle: Beauty is woken on Monday, and possibly again on Tuesday depending on a coin toss, without knowing which day it is. Whether the probability of heads is $1/2$ or $1/3$ is debated; the halfer and thirder positions rest on different interpretations of conditional probability and self-location.

\term{Smooth Manifold}
A topological manifold equipped with a smooth atlas of charts whose transition maps are infinitely differentiable.

\term{Sobolev Embedding Theorem}
Under suitable conditions on the indices, Sobolev spaces embed continuously into other Sobolev spaces, into $L^q$ spaces, or into spaces of continuous functions. For instance, $W^{k,p}$ embeds into $C^{m,\alpha}$ whenever $k - n/p > m$. Such embeddings play a central role in the regularity theory of PDEs, showing that weak solutions are smoother than initially guaranteed.

\term{Sobolev Solution}
A solution of a PDE that lies in a Sobolev space, defined through weak derivatives.

\term{Sobolev Space}
The space $W^{k,p}$ of functions whose weak derivatives up to order $k$ lie in $L^p$, the natural function space of PDE theory.

\term{Sorgenfrey Line}
The real line equipped with the lower-limit topology, where basis elements are half-open intervals $[a, b)$. It is a classic example of a space that is normal but not second-countable.

\term{Span (Linear Algebra)}
The set of all linear combinations of a given set of vectors. A set of vectors is a basis if it is linearly independent and its span is the entire vector space.

\term{Spectral Method}
A numerical method that expands the solution in global basis functions such as trigonometric or Chebyshev polynomials.

\term{Spectral Radius}
The maximum absolute value of the eigenvalues of a square matrix $A$. It determines the convergence of the power series $(I-A)^{-1} = \sum A^k$.

\term{Spectral Theorem}
A result stating that every symmetric matrix (or more generally, every normal operator on a Hilbert space) can be diagonalized by an orthonormal basis of eigenvectors.

\term{Spectral Theory}
The study of the spectrum (eigenvalues) of linear operators. It is fundamental to quantum mechanics, where observables are represented by self-adjoint operators.

\term{Sphere}
The set of all points in $\mathbb{R}^n$ at a fixed distance $r$ from a center point. An $n$-sphere $S^n$ is the boundary of an $(n+1)$-ball.

\term{Spherical Harmonic}
A set of special functions defined on the surface of a sphere. They are the angular portion of the solution to Laplace's equation in spherical coordinates.

\term{Spherical Geometry}
A non-Euclidean geometry where the "lines" are great circles on a sphere. In this geometry, the sum of the angles of a triangle is always greater than $180^\circ$.

\term{Spanning Tree}
A subgraph of a connected graph that is a tree and includes all the vertices of the original graph. Every connected graph has at least one spanning tree.

\term{Special Linear Group (SLn)}
The group of $n \times n$ matrices with determinant equal to 1. These transformations preserve volume and orientation.

\term{Special Orthogonal Group (SOn)}
The group of $n \times n$ orthogonal matrices with determinant 1. It consists of all rotations of $\mathbb{R}^n$ around the origin.

\term{Spec (Spectrum of a Ring)}
The set of all prime ideals of a commutative ring, equipped with the Zariski topology. It is the basic building block of schemes.

\term{Sperner's Theorem}
The largest antichain in the Boolean lattice of all subsets of an $n$-element set has size $\binom{n}{\lfloor n/2 \rfloor}$, realized by the middle level of subsets of size $\lfloor n/2 \rfloor$ or $\lceil n/2 \rceil$. It is a central result of extremal set theory and can be seen as a sharp version of Dilworth's theorem for the Boolean lattice.

\term{Sphere Packing Problem}
The problem of arranging non-overlapping spheres to occupy the maximum fraction of space. In 3D, the Kepler conjecture states the maximum density is $\pi/\sqrt{18}$.

\term{St. Petersburg Paradox}
A lottery in which the payoff doubles with each successive head and the expected value diverges to infinity, yet people will pay only a small sum to play. It motivated Bernoulli's introduction of diminishing marginal utility of money (utility of the logarithm), and connects to expected value versus expected utility.

\term{Stability}
The property of a solution or numerical scheme that small perturbations remain small; in ODE theory, solutions starting near an equilibrium remain near it for all time.

\term{Stable Marriage Problem}
An algorithmic problem of finding a matching between two sets of elements such that no pair of elements would both prefer each other over their current partners.

\term{Stable Range}
In algebraic K-theory, the range of dimensions where certain stability properties of the general linear group $\operatorname{GL}_n(R)$ hold.

\term{Standard Form}
A conventional way of writing a mathematical expression to make it easy to read or compare, such as the standard form of a linear equation $Ax + By = C$.

\term{State Space}
In dynamical systems and control theory, a space whose coordinates represent the state of the system at a given time.

\term{Stationary Process}
A stochastic process whose statistical properties (mean, variance, autocorrelation) do not change over time.

\term{Stationary Point}
A point where the derivative of a function is zero. It is a candidate for a local maximum, minimum, or saddle point.

\term{Statistics}
The science of collecting, analyzing, interpreting, and presenting empirical data to make inferences about a population.

\term{Steady State}
A solution of a dynamical system that does not change with time, such as a fixed point of an iteration.

\term{Stochastic Integral}
The integral of a stochastic process with respect to a semimartingale, defined in the sense of Itô or Stratonovich.

\term{Stochastic Process}
A collection of random variables indexed by time or space. Examples include Brownian motion and Poisson processes.

\term{Stokes' Theorem}
For a compact oriented manifold $M$ with boundary $\partial M$ and a differential form $\omega$, it states $\int_M d\omega = \int_{\partial M} \omega$, relating the integral of the exterior derivative over the manifold to the integral of the form over its boundary. Green's theorem, Gauss's divergence theorem, and the classical curl theorem are all special cases.

\term{Stone-Cech Compactification}
The largest compact Hausdorff space $\beta X$ that contains a given completely regular space $X$ as a dense subspace.

\term{Stone-Weierstrass Theorem}
A generalization of the Weierstrass approximation theorem: any continuous function on a compact Hausdorff space can be uniformly approximated by elements of a dense subalgebra.

\term{Stratification}
The decomposition of a space into a disjoint union of smooth manifolds (strata) of varying dimensions, used in singularity theory.

\term{Strong Induction}
A form of mathematical induction where the truth of a statement for $n$ is proved by assuming the statement is true for all $k < n$.

\term{Strong Topology}
In the context of dual spaces, the topology on $V^*$ induced by the norm of the original space $V$.

\term{Sturm-Liouville Problem}
The eigenvalue problem $-(p y')' + q y = \lambda w y$ with boundary conditions, yielding systems of orthogonal eigenfunctions.

\term{Subalgebra}
A subset of an algebra that is itself an algebra under the same operations.

\term{Subbasis}
A collection of sets whose finite intersections form a basis of a topology.

\term{Subgradient}
For a convex function, any vector $g$ with $f(y) \geq f(x) + g \cdot (y - x)$ for all $y$; a generalisation of the gradient to non-smooth functions.

\term{Subgroup}
A subset of a group that is itself a group under the same operation.

\term{Submanifold}
A subset of a manifold that is itself a manifold, with the subspace topology and a compatible smooth structure.

\term{Subspace}
A subset of a vector space that is closed under addition and scalar multiplication.

\term{Sub-additive Function}
A function $f$ such that $f(x+y) \leq f(x) + f(y)$ for all $x, y$ in its domain.

\term{Super-additive Function}
A function $f$ such that $f(x+y) \geq f(x) + f(y)$ for all $x, y$ in its domain.

\term{Super-symmetry}
A hypothesized symmetry in physics relating bosons (integer spin) and fermions (half-integer spin).

\term{Support (of a function)}
The closure of the set of points where a function is non-zero: $\operatorname{supp}(f) = \overline{\{x : f(x) \neq 0\}}$.

\term{Supremum}
The least upper bound of a set. If a set has a maximum, the supremum is that maximum.

\term{Surjection}
A function that hits every element of its codomain, also called onto.

\term{Symmetric}
A relation $R$ with $xRy \Rightarrow yRx$; a symmetric matrix satisfies $A^T = A$.

\term{Symmetric Group (Sn)}
The group of all permutations of $n$ elements. It has order $n!$ and is a central object in finite group theory.

\term{Symmetric Matrix}
A square matrix $A$ such that $A = A^T$. All eigenvalues of a real symmetric matrix are real, and it is always orthogonally diagonalizable.

\term{Symmetric Space}
A Riemannian manifold where for every point $p$, there exists an isometry that acts as a geodesic reflection through $p$.

\term{Symmetric Polynomial}
A polynomial $p(x_1, \ldots, x_n)$ that is invariant under any permutation of its variables.

\term{Symmetrization}
The process of making a function or operator symmetric by averaging it over a group action.

\term{Symplectic Geometry}
The study of manifolds equipped with a closed, non-degenerate 2-form $\omega$. It is the natural geometry of phase spaces in classical mechanics.

\term{Symplectic Form}
A closed, non-degenerate differential 2-form $\omega$ on a manifold. It allows for a canonical isomorphism between the tangent and cotangent bundles.

\term{Symplectic Manifold}
A smooth manifold equipped with a symplectic form. All symplectic manifolds are even-dimensional.

\term{Synge's Equation}
A differential equation used in the study of the stability of geodesics in general relativity.

\term{Syntactic Category}
A category whose objects are types and whose morphisms are provable functional dependencies in a logical theory.

\term{S-system}
A system of differential equations used in biochemical systems theory to model the interaction of components.

\term{Sylow Theorems}
A set of theorems describing the $p$-subgroups of a finite group $G$: Sylow $p$-subgroups exist, any two are conjugate, and the number $n_p$ of them divides $|G|$ and satisfies $n_p \equiv 1 \pmod p$. They are indispensable for classifying groups of small order and for analyzing the structure of finite groups generally.

\chapter*{T}
\label{chap:t}

\term{T-distribution}
The t-distribution is a probability distribution used when a sample is small and the population variance is unknown. It has heavier tails than the normal distribution, but approaches it as degrees of freedom increase. \par\noindent\textbf{Applications:} it supports small-sample confidence intervals and hypothesis tests.

\term{T-test}
A t-test compares means using an estimate of variability, especially when samples are small. Its assumptions and independent or paired design determine the appropriate version. \par\noindent\textbf{Applications:} t-tests compare treatments, measurements, and experimental groups.

\term{Tableau (Combinatorics)}
A tableau is a filling of a diagram subject to specified ordering rules. Young tableaux use rows and columns to encode partitions and permutations. \par\noindent\textbf{Applications:} tableaux support representation theory, symmetric functions, and combinatorics.

\term{T-norm (Triangular Norm)}
A t-norm is a commutative, associative, monotone operation on $[0,1]$ with $T(a,1)=a$. It generalizes intersection when truth values are graded. \par\noindent\textbf{Applications:} t-norms model conjunction in fuzzy logic and approximate reasoning.

\term{T-conorm}
A t-conorm is a fuzzy-logic operation that generalizes union and is dual to a t-norm, often by $S(a,b)=1-T(1-a,1-b)$. \par\noindent\textbf{Applications:} t-conorms combine graded alternatives in decision systems and control.

\term{Torsion (Group Theory)}
An element of a group is torsion if it has finite order. A torsion group consists entirely of torsion elements. \par\noindent\textbf{Applications:} torsion identifies periodic behavior in algebra, topology, and symmetry groups.

\term{Torsion}
See \term{Torsion Tensor}. \par\noindent\textbf{Applications:} torsion tensors measure the failure of a connection to be symmetric.

\term{Torsion-free Group}
A torsion-free group has no nonidentity element of finite order. The additive group $\mathbb Z$ is an example. \par\noindent\textbf{Applications:} torsion-free groups simplify algebraic and geometric classification.

\term{Torsion Point}
A torsion point on an elliptic curve is a point $P$ for which $nP=\mathcal O$ for some positive integer $n$. \par\noindent\textbf{Applications:} torsion points are important in elliptic-curve arithmetic and cryptography.

\term{Torsion Tensor}
The torsion tensor of a connection is $T(X,Y)=\nabla_XY-\nabla_YX-[X,Y]$. It measures the failure of parallel transport to be symmetric in two directions. \par\noindent\textbf{Applications:} torsion appears in differential geometry, relativity, and theories of defects in materials.

\term{Torsion Subgroup}
The torsion subgroup of an abelian group contains all elements of finite order. It decomposes into primary parts associated with primes. \par\noindent\textbf{Applications:} it separates periodic components from free components in algebraic classification.

\term{Tacnode}
A tacnode is a singular point where two curve branches touch with the same tangent line. \par\noindent\textbf{Applications:} tacnodes arise in algebraic-curve classification, bifurcation diagrams, and geometric modeling.

\term{Tangent}
A tangent describes either the limiting line along a curve at a point or the trigonometric function $\tan x=\sin x/\cos x$. \par\noindent\textbf{Applications:} tangents approximate curves, define slopes, and model angles and periodic systems.

\term{Tangent Space}
The tangent space $T_pM$ is the vector space of directions of curves passing through $p$ on a smooth manifold. \par\noindent\textbf{Applications:} tangent spaces define derivatives, velocities, gradients, and local linear models.

\term{Tangent Bundle}
The tangent bundle $TM$ is the disjoint union of all tangent spaces of a manifold. If $M$ has dimension $n$, then $TM$ has dimension $2n$. \par\noindent\textbf{Applications:} tangent bundles describe vector fields, velocities, and differential equations on manifolds.

\term{Tangent Line}
A tangent line is the best first-order linear approximation to a curve at a point. Its slope is the derivative when the curve is represented as a function. \par\noindent\textbf{Applications:} tangent lines support local approximation, optimization, and motion analysis.

\term{Tangent Vector}
A tangent vector represents the direction and rate of change of a curve at a point. On a smooth manifold it can be defined as a derivation on smooth functions. \par\noindent\textbf{Applications:} tangent vectors represent velocities and infinitesimal changes in geometry and mechanics.

\term{Tautological Bundle}
The tautological bundle over a Grassmannian has fiber $W$ at the point representing the subspace $W$. \par\noindent\textbf{Applications:} it records the subspaces being parametrized and is used in geometry, topology, and characteristic classes.

\term{Tautology (Logic)}
A tautology is a logical formula true under every assignment of its variables. \par\noindent\textbf{Applications:} tautologies support valid inference rules, circuit simplification, and formal verification.

\term{Taylor Series}
A Taylor series expresses a function as an infinite sum of derivative terms at one point, $\sum_{n\ge0}f^{(n)}(a)(x-a)^n/n!$. \par\noindent\textbf{Applications:} Taylor series support approximation, differential equations, numerical methods, and error analysis.

\term{Taylor's Theorem}
Taylor's theorem approximates a function near a point by a polynomial and gives a remainder that bounds the error. \par\noindent\textbf{Applications:} it justifies numerical approximation, local models, and estimates in calculus and differential equations.

\term{Teichmüller Space}
Teichmüller space parametrizes complex structures on a fixed topological surface up to an appropriate equivalence. \par\noindent\textbf{Applications:} it studies Riemann surfaces, moduli spaces, and deformation of geometric structures.

\term{Teichmüller Theory}
Teichmüller theory studies deformations of complex and hyperbolic structures on surfaces. \par\noindent\textbf{Applications:} it connects mapping class groups, Riemann surfaces, hyperbolic geometry, and moduli theory.

\term{T-matrix (Scattering Matrix)}
The T-matrix describes the transition from an incoming state to a scattered state in a scattering problem. \par\noindent\textbf{Applications:} it is used in quantum mechanics, nuclear physics, and wave scattering.

\term{Temperate Distribution}
A tempered distribution is a continuous linear functional on the Schwartz space of rapidly decreasing test functions. It may grow at most polynomially and has a well-defined Fourier transform. \par\noindent\textbf{Applications:} tempered distributions model signals, point sources, and generalized PDE solutions.

\term{Tempered Distribution}
A tempered distribution is a continuous linear functional on the Schwartz space. Examples include the Dirac delta, principal values, and polynomial derivatives. \par\noindent\textbf{Applications:} tempered distributions provide the natural framework for Fourier analysis and weak solutions of PDEs.

\term{Tensor}
A tensor is a multilinear object that generalizes scalars, vectors, and matrices. Its components transform predictably when coordinates change. \par\noindent\textbf{Applications:} tensors describe stress, curvature, material properties, and multilinear data.

\term{Tensor Product}
The tensor product $V\otimes W$ converts bilinear maps from $V\times W$ into linear maps from one vector space. \par\noindent\textbf{Applications:} tensor products combine state spaces, represent multilinear maps, and support quantum mechanics and algebra.

\term{Tensor Algebra}
The tensor algebra of $V$ is the graded algebra $T(V)=\bigoplus_{k\ge0}V^{\otimes k}$, with multiplication given by concatenating tensors. \par\noindent\textbf{Applications:} it constructs free algebras and leads to symmetric and exterior algebras.

\term{Tensor Field}
A tensor field assigns a tensor smoothly to every point of a manifold. \par\noindent\textbf{Applications:} tensor fields represent metrics, electromagnetic fields, stress, and curvature.

\term{Tensor Analysis}
Tensor analysis studies tensors, their coordinate transformations, and their derivatives. \par\noindent\textbf{Applications:} it provides the language of differential geometry, continuum mechanics, and general relativity.

\term{Term (Algebra)}
A term in a polynomial is one coefficient times variables raised to nonnegative integer powers. \par\noindent\textbf{Applications:} terms organize algebraic expressions and support symbolic manipulation and model construction.

\term{Tessellation}
A tessellation partitions a plane or space into cells with no gaps or overlaps. \par\noindent\textbf{Applications:} tessellations model crystals, meshes, spatial data, image processing, and nearest-neighbor geometry.

\term{Ternary Operation}
A ternary operation takes three inputs and returns one output. \par\noindent\textbf{Applications:} ternary operations appear in algebra, logic, programming languages, and multiway composition.

\term{Ternary Ring}
A ternary ring replaces ordinary multiplication by a three-input operation satisfying ring-like axioms. \par\noindent\textbf{Applications:} ternary rings arise in nonassociative algebra and the coordinatization of geometries.

\term{Theorem}
A theorem is a mathematical statement established by a proof from accepted axioms and earlier results. \par\noindent\textbf{Applications:} theorems provide reliable conclusions used in further mathematics and scientific models.

\term{Theory (Logic)}
A theory is a collection of formal statements closed under logical consequence. \par\noindent\textbf{Applications:} theories formalize mathematical structures, axioms, computation, and scientific reasoning.

\term{Theta Function}
A theta function is a holomorphic function with structured periodic or quasi-periodic behavior, commonly defined by a rapidly convergent series. \par\noindent\textbf{Applications:} theta functions occur in elliptic functions, modular forms, abelian varieties, and mathematical physics.

\term{Theta Series}
A theta series is a generating series whose coefficients count representations by a quadratic form, such as sums of squares. \par\noindent\textbf{Applications:} theta series connect number theory, lattices, modular forms, and coding.

\term{Theorema Egregium}
Theorema Egregium states that Gaussian curvature is intrinsic: it can be determined from measurements within the surface and does not depend on how the surface sits in space. \par\noindent\textbf{Applications:} it underlies surface geometry, cartography, and curvature-based shape analysis.

\term{Third Isomorphism Theorem}
The third isomorphism theorem states that if $N\subseteq M$ are normal in $G$, then $(G/N)/(M/N)\cong G/M$. \par\noindent\textbf{Applications:} it simplifies nested quotient constructions in groups, rings, and modules.

\term{Third Root}
The third root, or cube root, of a real number $x$ is the unique real $y$ satisfying $y^3=x$. \par\noindent\textbf{Applications:} cube roots occur in scaling laws, volume calculations, and solving cubic equations.

\term{Thomson's Lamp}
Thomson's lamp is a thought experiment in which a lamp is switched infinitely many times before a finite time. The example shows that an infinite sequence of actions may not determine a final state. \par\noindent\textbf{Applications:} it illustrates limits, supertasks, and questions about infinite processes.

\term{Three-dimensional}
Three-dimensional means that three independent coordinates are needed to specify a position or state. \par\noindent\textbf{Applications:} three-dimensional models describe physical space, graphics, geometry, and engineering structures.

\term{Tilde (Operator)}
A tilde is a notation mark used to indicate a modified, approximate, or related object, such as $\tilde f$. Its meaning depends on context. \par\noindent\textbf{Applications:} tildes denote transforms, estimates, equivalence relations, and asymptotic forms.

\term{Time Complexity}
Time complexity describes how an algorithm's running time grows with input size, often using Big O notation. \par\noindent\textbf{Applications:} it compares algorithms and guides the design of scalable software.

\term{Tikhonov Regularization}
Tikhonov regularization stabilizes an ill-posed inverse problem by minimizing a data-fit term plus a penalty, often $\lambda\|x\|_2^2$. \par\noindent\textbf{Applications:} it reduces noise and overfitting in inverse problems, imaging, and machine learning.

\term{Tiling}
A tiling covers a plane or space with tiles without gaps or overlaps. \par\noindent\textbf{Applications:} tilings appear in architecture, materials, discrete geometry, and computational meshing.

\term{Tiling Theory}
Tiling theory studies how shapes cover spaces and which local rules force or prevent global tilings. \par\noindent\textbf{Applications:} it informs crystallography, aperiodic materials, combinatorics, and pattern design.

\term{Time-reversal Symmetry}
Time-reversal symmetry means that the governing laws remain unchanged when time is replaced by $-t$, with appropriate changes to velocities. \par\noindent\textbf{Applications:} it tests physical models and distinguishes reversible from dissipative dynamics.

\term{Timelike Curve}
A timelike curve in spacetime has a timelike tangent at every point, so it can represent the path of a massive particle. \par\noindent\textbf{Applications:} timelike curves describe causality and possible observer trajectories in relativity.

\term{Timelike Vector}
A timelike vector has negative Lorentzian squared norm under the $(-,+,+,+)$ convention. \par\noindent\textbf{Applications:} timelike vectors represent physical four-velocities and causal directions in relativity.

\term{Tits Building}
A Tits building is a simplicial complex encoding the incidence structure of subgroups associated with a group of Lie type. \par\noindent\textbf{Applications:} buildings connect group theory, geometry, and representation theory.

\term{Tits Alternative}
The Tits alternative states that every finitely generated subgroup of a linear group is either virtually solvable or contains a free group of rank two. \par\noindent\textbf{Applications:} it reveals a major structural divide in linear groups and geometric group theory.

\term{Tits Group}
The Tits group is a finite simple group arising from constructions related to groups of Lie type. \par\noindent\textbf{Applications:} it provides an example in finite-group classification and the theory of sporadic groups.

\term{Toeplitz Matrix}
A Toeplitz matrix is constant along each diagonal from upper left to lower right. \par\noindent\textbf{Applications:} Toeplitz matrices model convolution and support fast algorithms in signal processing and numerical analysis.

\term{Toeplitz Operator}
A Toeplitz operator multiplies a Hardy-space function by a symbol and then projects back to the Hardy space. \par\noindent\textbf{Applications:} Toeplitz operators connect operator theory, signal processing, and boundary-value problems.

\term{Topological Invariant}
A topological invariant is a property unchanged by homeomorphisms, such as connectedness, compactness, or a fundamental group. \par\noindent\textbf{Applications:} invariants distinguish spaces and support classification in topology and geometry.

\term{Topological Space}
A topological space is a set with a collection of subsets called open sets, closed under arbitrary unions and finite intersections and containing the empty set and whole set. \par\noindent\textbf{Applications:} topological spaces formalize continuity, convergence, and geometric structure.

\term{Topology}
Topology studies properties preserved by continuous deformation. It uses open sets to define continuity, convergence, and neighborhoods. \par\noindent\textbf{Applications:} topology supports geometry, data analysis, dynamical systems, and physics.

\term{Topos}
A topos is a category that shares key structural features with the category of sets, including products, function objects, and a classifier for subobjects. \par\noindent\textbf{Applications:} topoi unify geometry and logic and provide settings for sheaves, cohomology, and constructive mathematics.

\term{Total Derivative}
The total derivative accounts for every way a composite function changes when its variables depend on one parameter. \par\noindent\textbf{Applications:} total derivatives support sensitivity analysis, mechanics, optimization, and differential equations.

\term{Total Differential}
The total differential is the linear approximation $df=\sum_i(\partial f/\partial x_i)dx_i$ to a multivariable change. \par\noindent\textbf{Applications:} it estimates measurement error and supports tangent-plane approximations and optimization.

\term{Total Order}
A total order is a partial order in which every two elements are comparable. The usual order on the real numbers is total. \par\noindent\textbf{Applications:} total orders support sorting, optimization, ranking, and induction arguments.

\term{Total Variation}
The total variation of a function is the supremum of sums of absolute changes over all partitions. Finite total variation means the function has controlled cumulative oscillation. \par\noindent\textbf{Applications:} it supports integration, signal analysis, and functions with jumps.

\term{Total Variation Diminishing (TVD)}
Total variation diminishing means that a numerical method does not increase the total variation of the computed solution. \par\noindent\textbf{Applications:} TVD schemes reduce spurious oscillations near shocks in hyperbolic PDEs.

\term{Touchard Polynomials}
Touchard polynomials encode the numbers of partitions of a set into blocks, using Stirling numbers as coefficients. \par\noindent\textbf{Applications:} they support combinatorics, moments of Poisson variables, and enumerative probability.

\term{Tournament (Graph Theory)}
A tournament is a directed graph obtained by assigning one direction to every edge of a complete graph. \par\noindent\textbf{Applications:} tournaments model competitions, preference comparisons, and ranking systems.

\term{Torsion-free Abelian Group}
A torsion-free abelian group has no nonidentity element of finite order. \par\noindent\textbf{Applications:} torsion-free groups describe lattice-like structures and the free part of abelian-group classifications.

\term{Trace}
The trace of a square matrix is the sum of its diagonal entries. It is invariant under similarity and equals the sum of eigenvalues with multiplicity. \par\noindent\textbf{Applications:} trace is used in linear algebra, covariance analysis, quantum mechanics, and optimization.

\term{Transcendence Degree}
The transcendence degree of $L/K$ is the size of a largest algebraically independent subset of $L$ over $K$. It measures the number of independent transcendental parameters. \par\noindent\textbf{Applications:} it gives the dimension of algebraic varieties and supports field theory and transcendence results.

\term{Transcendental Extension}
A transcendental extension contains an element that does not satisfy any nonzero polynomial with coefficients in the base field. \par\noindent\textbf{Applications:} transcendental extensions describe function fields and dimensions in algebraic geometry.

\term{Transfinite Induction}
Transfinite induction proves statements on a well-ordered set by proving each case from all earlier cases. \par\noindent\textbf{Applications:} it supports proofs about ordinals, cardinals, recursive constructions, and set theory.

\term{Transfinite Number}
A transfinite number describes the size or order type of an infinite set, using cardinals or ordinals such as $\aleph_0$. \par\noindent\textbf{Applications:} transfinite numbers organize infinite hierarchies in set theory and logic.

\term{Transfinite Recursion}
Transfinite recursion defines an object at each ordinal using values already defined at smaller ordinals. \par\noindent\textbf{Applications:} it constructs functions and hierarchies in set theory, logic, and infinitary combinatorics.

\term{Transformation}
A transformation is a function that maps one mathematical object or space to another. \par\noindent\textbf{Applications:} transformations describe coordinate changes, geometric motions, data processing, and symmetries.

\term{Transformation (Linear)}
A linear transformation preserves vector addition and scalar multiplication. \par\noindent\textbf{Applications:} linear transformations represent matrices, physical systems, projections, and changes of coordinates.

\term{Transformation (Möbius)}
A Möbius transformation has the form $f(z)=(az+b)/(cz+d)$ with $ad-bc\ne0$. It acts conformally on the extended complex plane. \par\noindent\textbf{Applications:} Möbius maps support complex analysis, hyperbolic geometry, and fractional-linear models.

\term{Transitive Action}
An action of $G$ on $X$ is transitive when one group element can carry any point to any other point. \par\noindent\textbf{Applications:} transitive actions describe homogeneous spaces, symmetries, and geometric configurations.

\term{Translation}
A translation moves every point by the same vector, preserving distances and directions. \par\noindent\textbf{Applications:} translations model displacement, coordinate changes, image registration, and periodic structures.

\term{Transpose}
The transpose $A^T$ is obtained by exchanging rows and columns, so $(A^T)_{ij}=A_{ji}$. \par\noindent\textbf{Applications:} transposes define symmetry, inner products, least squares, and matrix factorizations.

\term{Transversality}
Two submanifolds are transverse when their tangent spaces together span the tangent space of the ambient manifold at each intersection point. \par\noindent\textbf{Applications:} transversality makes intersections stable and supports differential topology and geometry.

\term{Transverse Intersection}
A transverse intersection satisfies the transversality condition, so the tangent spaces meet in the expected dimension. \par\noindent\textbf{Applications:} transverse intersections are stable under perturbation and are used to define intersection numbers.

\term{Transposition (Permutation)}
A transposition is a permutation that swaps two elements and fixes all others. Every finite permutation is a product of transpositions. \par\noindent\textbf{Applications:} transpositions generate symmetric groups and support permutation algorithms.

\term{Transport (Parallel)}
Parallel transport moves a vector along a curve while keeping its covariant derivative zero. \par\noindent\textbf{Applications:} it compares tangent vectors at different points and defines holonomy, curvature, and geometric phase.

\term{Trapezium}
A trapezium is a quadrilateral whose definition varies by convention: it may have no parallel sides or at least one pair. \par\noindent\textbf{Applications:} trapezia are used in geometry, area calculation, and design.

\term{Trapezoidal Rule}
The trapezoidal rule approximates an integral by replacing each interval with a trapezoid and summing the areas. \par\noindent\textbf{Applications:} it evaluates experimental data and numerical integrals when antiderivatives are unavailable.

\term{Tree}
A tree is a connected graph with no cycles. A tree on $n$ vertices has $n-1$ edges. \par\noindent\textbf{Applications:} trees model hierarchies, file systems, search algorithms, and communication networks.

\term{Triangle Inequality}
The triangle inequality states that $d(x,z)\leq d(x,y)+d(y,z)$ for a metric. A direct route is no longer than a route through an intermediate point. \par\noindent\textbf{Applications:} it controls convergence, error bounds, clustering, and shortest-path algorithms.

\term{Triangulated Category}
A triangulated category is an additive category with a shift operation and distinguished triangles satisfying axioms modeled on mapping cones. \par\noindent\textbf{Applications:} triangulated categories organize derived categories, homological algebra, algebraic geometry, and stable homotopy.

\term{Trivial Group}
The trivial group contains only the identity element and has order one. \par\noindent\textbf{Applications:} it is the base case for group extensions, kernels, and induction arguments.

\term{Trivial Ring}
The trivial ring has one element, so $0=1$. \par\noindent\textbf{Applications:} it serves as a zero object or degenerate case in ring theory and algebraic geometry.

\term{Trivial Topology}
The trivial topology contains only the empty set and the whole space as open sets. \par\noindent\textbf{Applications:} it provides simple examples and counterexamples in topology and continuity.

\term{Truncation Error}
Truncation error is the error introduced when an exact equation or infinite process is replaced by a finite approximation. \par\noindent\textbf{Applications:} it guides step-size selection and accuracy analysis in numerical methods.

\term{Turán's Theorem}
Turán's theorem gives the maximum number of edges in an $n$-vertex graph that contains no complete graph on $r+1$ vertices. \par\noindent\textbf{Applications:} it bounds dense networks and is foundational in extremal graph theory.

\term{Turing Machine}
A Turing machine is an abstract model with a tape, read-write head, finite control, and transition rules. It captures the idea of algorithmic computation and has undecidable problems such as halting. \par\noindent\textbf{Applications:} Turing machines define computability and complexity classes.

\term{Twin Primes}
Twin primes are pairs of primes differing by two, such as $(3,5)$. The twin-prime conjecture says that infinitely many such pairs exist. \par\noindent\textbf{Applications:} they are a central example in prime-distribution research.

\term{Two Envelopes Paradox}
The two-envelopes paradox appears to show that switching envelopes always increases expected money, even though one envelope is already selected. The error comes from using an unjustified probability model for unbounded amounts. \par\noindent\textbf{Applications:} it clarifies conditional expectation and decision-making under uncertainty.

\term{Two's Complement}
Two's complement represents signed binary integers so that negation is performed by inverting bits and adding one. \par\noindent\textbf{Applications:} it allows computer hardware to use the same addition circuit for positive and negative integers.

\term{Tychonoff's Theorem}
Tychonoff's theorem states that any product of compact spaces is compact in the product topology. \par\noindent\textbf{Applications:} it supports infinite-dimensional analysis, functional analysis, and existence proofs.

\term{Type Theory}
Type theory assigns types to expressions to organize computation and proofs. Under the Curry--Howard correspondence, types represent propositions and programs represent proofs. \par\noindent\textbf{Applications:} type theory underlies programming languages, proof assistants, formal verification, and foundations of mathematics.
\term{Stirling Number}
Stirling numbers of the first kind count permutations by cycles, while those of the second kind count partitions into nonempty blocks. \par\noindent\textbf{Applications:} they expand factorial polynomials and support enumerative combinatorics and probability.
\term{Strong Duality}
Strong duality means that a primal optimization problem and its dual have equal optimal values. \par\noindent\textbf{Applications:} it provides certificates of optimality and supports algorithms in linear, convex, and machine-learning optimization.

\chapter*{U}
\label{chap:u}

\term{Unit}
A unit is either the multiplicative identity in a ring or an element that has a multiplicative inverse, depending on context. \par\noindent\textbf{Applications:} units describe reversible scaling and invertibility in algebra.

\term{Unit (Ring Theory)}
A unit of a ring $R$ is an element $u$ with an inverse $v$ satisfying $uv=vu=1$. The units form the group $R^\times$; in $\mathbb Z$ they are $1$ and $-1$. \par\noindent\textbf{Applications:} units identify invertible coefficients and symmetries in algebra and number theory.

\term{Unit Circle}
The unit circle is the circle $x^2+y^2=1$ centered at the origin. Its points are $(\cos\theta,\sin\theta)$. \par\noindent\textbf{Applications:} it defines trigonometric functions and models rotations, periodic motion, and complex numbers.

\term{Unit Group}
The unit group $R^\times$ of a ring consists of all elements with multiplicative inverses. In a field it is the group of all nonzero elements. \par\noindent\textbf{Applications:} unit groups encode invertibility in algebraic number theory and modular arithmetic.

\term{Unit Vector}
A unit vector has norm one. A nonzero vector $v$ can be normalized as $\hat v=v/\|v\|$. \par\noindent\textbf{Applications:} unit vectors represent directions, orientations, and normalized features.

\term{Unitary Matrix}
A unitary matrix satisfies $U^*U=UU^*=I$. It preserves inner products, lengths, and angles in complex space. \par\noindent\textbf{Applications:} unitary matrices describe quantum gates, rotations, Fourier transforms, and stable changes of basis.

\term{Unitary Operator}
A unitary operator is a surjective linear operator on a Hilbert space that preserves inner products. It is the infinite-dimensional analogue of a unitary matrix. \par\noindent\textbf{Applications:} unitary operators model quantum evolution and energy-preserving transformations.

\term{Universal}
Universal means applying to every case in a specified domain. The symbol $\forall$ expresses “for all.” \par\noindent\textbf{Applications:} universal statements formulate definitions, theorems, specifications, and logical rules.

\term{Universal Set}
A universal set is the chosen set containing every object under discussion; all sets in that context are its subsets. \par\noindent\textbf{Applications:} it fixes the complement operation and organizes elementary set calculations.

\term{Universal Mapping Property}
A universal mapping property characterizes an object by a unique map or factorization property rather than by a formula. Tensor products and free groups are examples. \par\noindent\textbf{Applications:} it gives canonical constructions throughout algebra, topology, and category theory.

\term{Uncountable Set}
An uncountable set is an infinite set that cannot be put in bijection with $\mathbb N$. The real numbers are uncountable. \par\noindent\textbf{Applications:} uncountability distinguishes sizes of infinity and limits enumeration and computation.

\term{Undecidable}
An undecidable problem has no algorithm that always halts and returns the correct answer for every input. \par\noindent\textbf{Applications:} undecidability identifies fundamental limits of software verification and computation.

\term{Underdetermined}
An underdetermined system has fewer independent equations than unknowns and usually has many solutions. \par\noindent\textbf{Applications:} underdetermined models arise in inverse problems, data fitting, and parameter estimation.

\term{Uniform}
Uniform means that one bound or choice works throughout a whole domain, rather than changing from point to point. \par\noindent\textbf{Applications:} uniform estimates control approximation, convergence, and numerical error.

\term{Uniform Continuity}
Uniform continuity means that for every $\epsilon>0$ one $\delta>0$ works for all points in the domain. Unlike ordinary continuity, $\delta$ does not depend on the point. \par\noindent\textbf{Applications:} it preserves controlled approximation and ensures uniform limits behave well.

\term{Uniform Convergence}
Uniform convergence means that the largest error between functions and their limit tends to zero. \par\noindent\textbf{Applications:} it allows limits to preserve continuity and supports reliable approximation by function sequences.

\term{Uniform Boundedness Principle}
See \term{Banach--Steinhaus Theorem}. \par\noindent\textbf{Applications:} it controls families of linear operators using pointwise bounds.

\term{Uniform Distribution}
A uniform distribution assigns equal probability to equal-sized regions of its support. On $[a,b]$ its density is $1/(b-a)$. \par\noindent\textbf{Applications:} it models unbiased random choices and supplies a basic simulation distribution.

\term{Uniform Space}
A uniform space is a set with a structure describing when pairs of points are uniformly close, even when no distance function is given. \par\noindent\textbf{Applications:} uniform spaces define completeness and uniform continuity in general topology.

\term{Union}
The union of sets contains every element belonging to at least one set: $A\cup B=\{x:x\in A\text{ or }x\in B\}$. \par\noindent\textbf{Applications:} unions combine alternatives and describe events, regions, and solution sets.

\term{Unimodular Matrix}
A unimodular matrix has integer entries and determinant $\pm1$, so its inverse also has integer entries. \par\noindent\textbf{Applications:} unimodular matrices preserve integer lattices and support integer programming and lattice algorithms.

\term{Uniqueness Theorem}
A uniqueness theorem gives conditions under which a problem has at most one solution. \par\noindent\textbf{Applications:} uniqueness establishes well-defined models and is central to differential equations, optimization, and inverse problems.

\term{Unipotent Element}
A unipotent element satisfies that $u-1$ is nilpotent. A unipotent matrix has only eigenvalue $1$ and is similar to an upper triangular matrix with ones on the diagonal. \par\noindent\textbf{Applications:} unipotent elements appear in algebraic groups, Lie theory, and matrix decompositions.

\term{Upper Half-Plane}
The upper half-plane is $\{z\in\mathbb C:\operatorname{Im}(z)>0\}$. It is a standard domain for complex and hyperbolic geometry. \par\noindent\textbf{Applications:} it supports modular forms, conformal maps, and the Poincaré model.

\term{Upper Bound}
An upper bound $M$ of a set $S$ satisfies $s\leq M$ for every $s\in S$. The least upper bound is its supremum. \par\noindent\textbf{Applications:} upper bounds control errors, sizes, optimization objectives, and convergence.

\term{Upper Semi-continuous Function}
An upper semi-continuous function cannot jump upward in the limit: nearby values have limit superior at most the value at the point. \par\noindent\textbf{Applications:} upper semi-continuity helps prove existence of maxima and stable optimization solutions.

\term{Upper Triangular}
An upper-triangular matrix has zeros below its main diagonal. Its eigenvalues are the diagonal entries. \par\noindent\textbf{Applications:} triangular form simplifies eigenvalue calculation, linear systems, and matrix exponentials.

\term{Urysohn Lemma}
The Urysohn lemma states that disjoint closed sets in a normal space can be separated by a continuous function taking value $0$ on one and $1$ on the other. \par\noindent\textbf{Applications:} it constructs partitions of unity and supports metrization and extension theorems.

\term{Urysohn Metrization Theorem}
The Urysohn metrization theorem states that every second-countable regular space admits a compatible metric. \par\noindent\textbf{Applications:} it converts abstract topological questions into metric and analytical ones.

\term{Use}
Use is the act of applying a function, rule, or definition to given arguments. \par\noindent\textbf{Applications:} precise use of rules supports symbolic calculation, programming, and formal proofs.

\term{U-statistic}
A U-statistic is an unbiased estimator formed by averaging a kernel over subsets of a sample. \par\noindent\textbf{Applications:} U-statistics estimate population quantities such as moments, correlations, and interaction effects.

\term{Unitary Group}
The unitary group $U(n)$ consists of all complex $n\times n$ unitary matrices. It is a compact Lie group. \par\noindent\textbf{Applications:} it describes quantum symmetries, rotations of complex state spaces, and harmonic analysis.

\term{Unique}
Unique means existing in exactly one form or instance. \par\noindent\textbf{Applications:} uniqueness identifies a determined solution, canonical object, or unambiguous representation.

\term{Unique Factorization Domain (UFD)}
A unique factorization domain is an integral domain in which every nonzero nonunit factors uniquely into irreducible, equivalently prime, elements up to order and units. \par\noindent\textbf{Applications:} UFDs support polynomial factorization and arithmetic arguments.

\term{Unital Ring}
A unital ring has a multiplicative identity $1$ satisfying $1a=a1=a$ for every element $a$. \par\noindent\textbf{Applications:} unital rings provide the standard setting for modules, algebras, matrices, and algebraic geometry.

\term{Unitary Representation}
A unitary representation maps a group to unitary operators on a Hilbert space. \par\noindent\textbf{Applications:} it represents symmetries in quantum mechanics, harmonic analysis, and signal theory.

\term{Unimodular Group}
A locally compact group is unimodular when its left and right Haar measures agree. \par\noindent\textbf{Applications:} unimodularity simplifies integration, harmonic analysis, and representation theory on groups.

\term{Universal Covering Space}
A universal covering space is a simply connected space mapping locally like a given space through a covering map. \par\noindent\textbf{Applications:} it reveals fundamental groups and supports classification of covering spaces and manifolds.

\term{Ultrafilter}
An ultrafilter is a maximal collection of subsets closed under finite intersections and supersets, containing exactly one of each set and its complement. \par\noindent\textbf{Applications:} ultrafilters support compactness proofs, limits, and nonstandard analysis.

\term{Ultraproduct}
An ultraproduct combines a family of structures and identifies sequences that agree on a chosen ultrafilter-large set. \par\noindent\textbf{Applications:} ultraproducts transfer first-order properties and support model theory and functional analysis.

\term{Ultrametric Space}
An ultrametric space satisfies $d(x,z)\leq\max\{d(x,y),d(y,z)\}$, a stronger triangle inequality than for ordinary metrics. \par\noindent\textbf{Applications:} ultrametrics model $p$-adic numbers, hierarchical clustering, and phylogenetic trees.

\term{Unary Operation}
A unary operation takes one input and returns one output, such as negation or complementation. \par\noindent\textbf{Applications:} unary operations define algebraic structures, logic, syntax, and programming expressions.

\term{Unconditional Convergence}
A series converges unconditionally when every rearrangement converges to the same sum. In many important spaces this is equivalent to absolute convergence. \par\noindent\textbf{Applications:} unconditional convergence supports rearranging infinite expansions and stable functional analysis.

\term{Unbounded Operator}
An unbounded operator is a linear operator that has no finite global norm bound on its domain. \par\noindent\textbf{Applications:} unbounded operators model derivatives, momentum, and energy in PDEs and quantum mechanics.

\term{Unit Disc}
The unit disc is $\{z\in\mathbb C:|z|\leq1\}$; its interior is $|z|<1$. \par\noindent\textbf{Applications:} it is a standard domain for complex analysis, conformal maps, and power series.

\term{Unitary Transformation}
A unitary transformation preserves inner products and therefore lengths and angles in a complex vector space. \par\noindent\textbf{Applications:} unitary transformations represent quantum evolution, rotations, and Fourier changes of basis.

\term{Roots of Unity}
The $n$th roots of unity are the complex numbers satisfying $z^n=1$, namely $e^{2\pi ik/n}$ for $k=0,\ldots,n-1$. \par\noindent\textbf{Applications:} they support Fourier analysis, cyclic symmetry, polynomial factorization, and coding.

\term{Univalent Function}
A univalent function is a holomorphic function that is one-to-one on its domain. \par\noindent\textbf{Applications:} univalent functions control conformal mappings and geometric function theory.

\chapter*{V}
\label{chap:v}

\term{Valuation}
A function $v$ on a field $K$ that assigns a value in an ordered abelian group, satisfying $v(xy) = v(x) + v(y)$ and $v(x+y) \geq \min(v(x), v(y))$. Examples include the $p$-adic valuation on $\mathbb{Q}$.

\term{Value (Logic)}
An assignment of truth values (e.g., True or False) to the variables of a logical formula. In multi-valued logic, a formula can take more than two possible values.

\term{Vandermonde Matrix}
A matrix $V$ where each row is a geometric progression: $V_{ij} = x_j^{i-1}$. Its determinant is the product of all differences $(x_j - x_i)$, and it is invertible iff all $x_i$ are distinct.

\term{Van Kampen's Theorem}
A result in algebraic topology that expresses the fundamental group of a space $X = U \cup V$ as the pushout of the fundamental groups of $U$ and $V$ over their intersection $U \cap V$.

\term{Variation}
A measure of the total vertical oscillation of a function. A function has bounded variation on $[a,b]$ if the supremum of $\sum |f(x_{i+1}) - f(x_i)|$ over all partitions is finite.

\term{Variational Formulation}
The weak formulation of a PDE as an equation in a Hilbert space, such as $\int \nabla u \cdot \nabla v = \int f v$, which is the basis of the finite-element method.

\term{Variational Inequality}
The inequality $F(x^*) \cdot (x - x^*) \geq 0$ for all $x$ in the feasible set, characterising equilibria in optimization and game theory.

\term{Variational Method}
A technique for finding the extrema of a functional (such as an integral) by considering small perturbations of the candidate function. It is the basis for the Euler--Lagrange equations.

\term{Vector}
An element of a vector space. In $\mathbb{R}^n$, a vector is an ordered tuple of $n$ real numbers, representing a quantity with both magnitude and direction.

\term{Vector Field}
An assignment of a vector to each point of a manifold or region in $\mathbb{R}^n$. In physics, vector fields represent quantities like electric fields or fluid velocity.

\term{Vector Space}
A set $V$ equipped with addition and scalar multiplication that satisfies eight axioms (associativity, commutativity, distributive laws, etc.). It is the fundamental structure of linear algebra.

\term{Vector Bundle}
A topological construction where each point of a base space $B$ is associated with a vector space $V$ (the fiber), such that the local structure is a product $U \times V$.

\term{Vectorization}
The operation of converting a matrix into a column vector by stacking its columns one on top of another. It is used to transform matrix equations into vector equations.

\term{Velocity}
The rate of change of position with respect to time, the derivative of the position vector.

\term{Verification}
The process of formally checking that a proof or computation is correct.

\term{Vertex (Graph Theory)}
A fundamental unit of a graph, also called a node. A graph $G=(V,E)$ consists of a set of vertices $V$ and a set of edges $E$ connecting them.

\term{Vertex Cover}
A set of vertices in a graph such that every edge of the graph is incident to at least one vertex in the set. Finding the minimum vertex cover is an NP-hard problem.

\term{Vertex Degree}
The number of edges incident to a vertex in a graph. In a directed graph, this is split into in-degree and out-degree.

\term{Vertex Connectivity}
The minimum number of vertices whose removal disconnects the graph. A graph is $k$-connected if it remains connected after removing any $k-1$ vertices.

\term{Vertex (Polytope)}
A 0-dimensional face of a polytope. A vertex is a point that cannot be expressed as a convex combination of any other points in the polytope.

\term{Valued Field}
A field equipped with a valuation. The study of valued fields is central to number theory and the construction of completions like the $p$-adics.

\term{Variety (Algebraic)}
The set of common zeros of a collection of polynomials in $\mathbb{A}^n$ or $\mathbb{P}^n$. Varieties are the primary objects of study in algebraic geometry.

\term{Varifold}
A generalization of a submanifold that allows for the study of surfaces with singularities, using the tools of geometric measure theory.

\term{Vect}
In category theory, the category whose objects are vector spaces over a fixed field and whose morphisms are linear maps.

\term{Vertical}
At right angles to the horizontal; a vertical line has undefined slope.

\term{Violation}
An instance where a rule, inequality, or constraint fails to hold for a particular case.

\term{Visualisation}
The representation of mathematical objects or data by geometric figures to aid understanding.

\term{Voronoi Diagram}
A partition of a plane into regions based on distance to a specific set of seed points. For each seed, the Voronoi cell consists of all points closer to that seed than to any other.

\term{Voronoi Cell}
The convex polytope consisting of all points in space closer to a given seed $p_i$ than to any other seed in a set.

\term{V-polytope}
A polytope defined as the convex hull of a finite set of vertices (V-representation), as opposed to an H-polytope defined by a set of linear inequalities.

\term{Variance}
The second central moment of a random variable, $\operatorname{Var}(X) = \mathbb{E}[(X - \mu)^2]$. It measures the dispersion of a distribution around its mean.

\term{V-statistic}
A type of la Place-style statistic that generalizes the variance and covariance to higher-order interactions in a sample.

\term{V-fold Cross Validation}
A technique in machine learning where a dataset is split into $V$ folds; the model is trained on $V-1$ folds and tested on the remaining one, repeating this process $V$ times.

\term{Vanishing Point}
In projective geometry, the point where parallel lines in 3D space appear to converge when projected onto a 2D plane.

\term{Vanishing Theorem}
A result in complex geometry or topology stating that certain cohomology groups of a sheaf or manifold are zero, often as a consequence of positivity conditions (e.g., Kodaira vanishing).

\term{V-norm}
A specific type of norm used in the study of certain operator algebras or sequence spaces, often related to the $L^p$ norms.

\term{von Neumann Algebra}
A C*-algebra that is also closed in the weak operator topology; equivalently, the algebra of operators commuting with the commutant of a Hilbert space representation.

\chapter*{W}
\label{chap:w}

\term{Waelbroeck Space}
A Waelbroeck space is a topological vector space designed to support a useful theory of bounded operators, often without local convexity. \par\noindent\textbf{Applications:} such spaces occur in generalized-function and distribution theory.

\term{Walk (Graph Theory)}
A walk in a graph is an alternating sequence of vertices and incident edges. Unlike a path, it may repeat vertices or edges. \par\noindent\textbf{Applications:} walks model routes, random processes, search procedures, and network communication.

\term{Wallis Product}
The Wallis product is the convergent infinite product $\frac21\cdot\frac23\cdot\frac43\cdot\frac45\cdots=\frac\pi2$. \par\noindent\textbf{Applications:} it provides a classical product for $\pi$ and illustrates convergence of infinite products.

\term{Walsh--Hadamard Transform}
The Walsh--Hadamard transform expands data in piecewise-constant Walsh functions using only additions and subtractions. \par\noindent\textbf{Applications:} it supports fast signal processing, digital communication, coding, and quantum algorithms.

\term{W-algebra}
A W-algebra is an algebra of operators encoding generalized symmetries, often in conformal field theory. \par\noindent\textbf{Applications:} W-algebras support the study of higher-spin fields, integrable systems, and representation theory.

\term{Ward Identity}
A Ward identity is a relation among correlation functions that follows from a symmetry of a quantum theory. \par\noindent\textbf{Applications:} Ward identities constrain quantum calculations and express conservation laws related to Noether's theorem.

\term{Waring's Problem}
Waring's problem asks whether every positive integer is a sum of a bounded number of $k$th powers for each fixed $k$. \par\noindent\textbf{Applications:} it is a central problem in additive number theory and the study of representations of integers.

\term{Wave Equation}
The wave equation models propagation, commonly as $u_{tt}=c^2\nabla^2u$. \par\noindent\textbf{Applications:} it describes sound, vibrations, elastic waves, electromagnetic waves, and acoustics.

\term{Wavelength}
Wavelength is the distance between successive points with the same phase, such as neighboring crests. \par\noindent\textbf{Applications:} it determines wave frequency, color, resolution, and antenna behavior.

\term{Wavelet Transform}
A wavelet transform decomposes a signal into shifted and scaled copies of a localized wavelet. It records both time and frequency behavior. \par\noindent\textbf{Applications:} wavelets support image compression, denoising, anomaly detection, and signal analysis.

\term{Weak Convergence}
Weak convergence means that every continuous linear functional has the same limit on the sequence as on the candidate limit. \par\noindent\textbf{Applications:} it provides compactness tools for infinite-dimensional optimization and PDEs.

\term{Weak-* Convergence}
Weak-* convergence in a dual space is tested against vectors of the original space rather than all elements of the double dual. \par\noindent\textbf{Applications:} it supports compactness results such as Banach--Alaoglu and convergence of measures.

\term{Weak Derivative}
A weak derivative is defined through integration by parts against smooth test functions, even when the ordinary derivative does not exist. \par\noindent\textbf{Applications:} weak derivatives define Sobolev spaces and allow nonsmooth PDE solutions.

\term{Weak Solution}
A weak solution satisfies a differential equation after integration against suitable test functions, so it may have fewer classical derivatives. \par\noindent\textbf{Applications:} weak solutions are fundamental in PDE theory, finite elements, and fluid mechanics.

\term{Weak Topology}
A weak topology is the coarsest topology making the chosen continuous linear functionals continuous. Its convergence is tested by those functionals. \par\noindent\textbf{Applications:} weak topologies provide compactness and convergence methods in functional analysis.

\term{Weakly Connected Component}
A weakly connected component of a directed graph is a maximal subgraph that becomes connected when edge directions are ignored. \par\noindent\textbf{Applications:} it identifies network regions and communication components.

\term{Weighted Graph}
A weighted graph assigns a numerical weight to each edge, such as a cost, distance, or capacity. \par\noindent\textbf{Applications:} weighted graphs model transport, communication, scheduling, and logistics.

\term{Weighted Mean}
A weighted mean is an average in which observations contribute according to weights: $\sum_iw_ix_i/\sum_iw_i$. \par\noindent\textbf{Applications:} weighted means combine measurements with different reliability and support survey analysis.

\term{Weil Conjecture}
The Weil conjectures are four theorems describing point counts and zeta functions of varieties over finite fields. \par\noindent\textbf{Applications:} they connect algebraic geometry, number theory, and cohomology.

\term{Weil Pairing}
The Weil pairing is a bilinear, alternating pairing from elliptic-curve torsion points to roots of unity. \par\noindent\textbf{Applications:} it reveals elliptic-curve structure and supports pairing-based cryptography.

\term{Weil Transform}
The Weil transform is an integral transform used to relate functions or representations in harmonic analysis. \par\noindent\textbf{Applications:} it appears in automorphic forms, representation theory, and number theory.

\term{Weierstrass Function}
The Weierstrass function is a function continuous everywhere but differentiable nowhere. \par\noindent\textbf{Applications:} it demonstrates the difference between continuity and differentiability and provides a key example in real analysis.

\term{Weierstrass Product Theorem}
The Weierstrass product theorem represents an entire function through its zeros and an exponential factor, with convergence factors when needed. \par\noindent\textbf{Applications:} it constructs entire functions and studies their zeros in complex analysis.

\term{Weierstrass Theorem (Approximation)}
The Weierstrass approximation theorem states that every continuous function on a closed interval can be uniformly approximated by polynomials. \par\noindent\textbf{Applications:} it justifies polynomial approximation in numerical analysis, modeling, and computation.

\term{Weight (Representation Theory)}
A weight in representation theory is a linear functional describing how a Cartan subalgebra acts on a vector. The collection of weights helps characterize a representation. \par\noindent\textbf{Applications:} weights classify symmetries in Lie theory and quantum mechanics.

\term{Weight (Graph Theory)}
The weight of an edge is its assigned numerical value; the weight of a subgraph is often the sum of its edge weights. \par\noindent\textbf{Applications:} weights represent distance, cost, capacity, or reliability in network optimization.

\term{Well-Defined}
A definition is well-defined when its result does not depend on choices made during the construction. \par\noindent\textbf{Applications:} well-definedness is essential for quotient objects, coordinate definitions, and mathematical functions.

\term{Well-founded Relation}
A relation is well-founded when every nonempty subset has a minimal element under the relation. \par\noindent\textbf{Applications:} well-foundedness guarantees termination of recursive definitions and supports induction arguments.

\term{Well-Ordered}
A well-order is a total order in which every nonempty subset has a least element. \par\noindent\textbf{Applications:} well-orders support transfinite induction and constructions in set theory.

\term{Well-ordering Principle}
The well-ordering principle states that every set can be equipped with a well-order; for positive integers this is the least-number principle. \par\noindent\textbf{Applications:} it supports induction, ordinal methods, and the axiom-of-choice framework.

\term{Well-Posed Problem}
A problem is well-posed when a solution exists, is unique, and changes continuously with the data. \par\noindent\textbf{Applications:} well-posedness is a basic requirement for reliable models and numerical solutions of PDEs and inverse problems.

\term{Weyl Group}
A Weyl group is generated by reflections associated with the roots of a root system. It describes the root system's symmetries. \par\noindent\textbf{Applications:} Weyl groups classify Lie algebras, reflection groups, and representations.

\term{Weyl Algebra}
A Weyl algebra is generated by polynomial variables and their formal derivatives with the relation $[\partial,x]=1$. \par\noindent\textbf{Applications:} it describes differential operators and supports D-modules, PDEs, and representation theory.

\term{Weyl Character Formula}
The Weyl character formula gives the character of an irreducible representation of a semisimple Lie algebra using its highest weight and Weyl group. \par\noindent\textbf{Applications:} it computes representation dimensions and symmetry data.

\term{Weyl Tensor}
The Weyl tensor is the trace-free part of the Riemann curvature tensor. It describes gravitational distortion not determined by local matter density. \par\noindent\textbf{Applications:} it models tidal forces and gravitational waves in relativity.

\term{Weyl's Law}
Weyl's law gives the leading asymptotic growth of Laplacian eigenvalues in terms of the volume of a domain. \par\noindent\textbf{Applications:} it connects spectral data with geometry and is used in PDEs and mathematical physics.

\term{Weyl's Inequality}
Weyl's inequalities bound the eigenvalues of a sum of Hermitian matrices using the eigenvalues of the summands. \par\noindent\textbf{Applications:} they support matrix perturbation theory, optimization, and quantum mechanics.

\term{Whitney Embedding Theorem}
The Whitney embedding theorem states that every smooth $n$-manifold can be smoothly embedded in $\mathbb R^{2n}$. \par\noindent\textbf{Applications:} it represents abstract manifolds as concrete Euclidean subspaces for geometry and topology.

\term{Wigner D-matrix}
Wigner D-matrices represent three-dimensional rotations in bases of angular-momentum states. \par\noindent\textbf{Applications:} they are used in quantum mechanics, spherical harmonics, molecular physics, and 3D orientation.

\term{Wigner's Semicircle Law}
Wigner's semicircle law states that the normalized eigenvalue distribution of a large random symmetric matrix approaches a semicircle. \par\noindent\textbf{Applications:} it supports random matrix theory, statistics, wireless systems, and mathematical physics.

\term{Winding Number}
The winding number is the integer counting how many net times a closed curve winds around a point. \par\noindent\textbf{Applications:} it detects zeros and poles, classifies loops, and appears in complex analysis and topology.

\term{Witt Vector}
Witt vectors package sequences of elements of a ring of characteristic $p$ into a ring of characteristic zero. \par\noindent\textbf{Applications:} they are used in $p$-adic arithmetic, algebraic geometry, and cohomology.

\term{Witt Group}
The Witt group classifies quadratic forms up to adding hyperbolic planes. \par\noindent\textbf{Applications:} it studies quadratic forms, algebraic invariants, and arithmetic geometry.

\term{Word}
A word is a finite sequence of symbols from an alphabet. \par\noindent\textbf{Applications:} words are the basic objects of formal languages, string algorithms, coding, and automata theory.

\term{Wronskian}
The Wronskian is the determinant formed from functions and their successive derivatives. A nonzero Wronskian at a point proves linear independence. \par\noindent\textbf{Applications:} it tests solutions of linear ODEs and helps construct fundamental solution sets.

\term{Well-Ordering Theorem}
The well-ordering theorem states that every set admits a well-order, and it is equivalent to the axiom of choice and Zorn's lemma. \par\noindent\textbf{Applications:} it supports transfinite induction and existence proofs in set theory and algebra.

\chapter{X}
\label{chap:x}

\term{X-ray Transform}
An integral transform that maps a function $f$ on $\mathbb{R}^n$ to the set of its integrals over all lines. It is a special case of the Radon transform for $n=2$ and is fundamental in medical imaging and seismic tomography.

\term{XOR (Exclusive OR)}
A logical operation that returns true if and only if exactly one of its two inputs is true. In binary arithmetic, it is equivalent to addition modulo 2.

\term{x-axis}
The horizontal coordinate axis in the Cartesian plane, usually measuring the first coordinate.

\term{x-coordinate}
The first coordinate of a point, measuring its signed distance from the y-axis.

\term{Chi-squared ($\chi^2$) Distribution}
A continuous probability distribution that arises when summing the squares of $k$ independent standard normal random variables. It is widely used in hypothesis testing to determine "goodness of fit" between observed and expected data.

\term{Chi-squared ($\chi^2$) Test}
A statistical test used to determine if there is a significant difference between the expected frequencies and the observed frequencies in one or more categories of data.


\chapter{Y}
\label{chap:y}

\term{Yang-Mills Theory}
A gauge theory that generalizes electromagnetism to non-abelian Lie groups. It provides the mathematical framework for the strong and weak nuclear forces in the Standard Model of particle physics. The existence of a "mass gap" in Yang-Mills theory is one of the Millennium Prize Problems.

\term{Yang-Baxter Equation}
A consistency equation for a system of three particles undergoing pairwise scattering, $R_{12}R_{13}R_{23} = R_{23}R_{13}R_{12}$. It is central to the study of integrable systems, quantum groups, and knot invariants.

\term{y-axis}
The vertical coordinate axis in the Cartesian plane, usually measuring the second coordinate.

\term{y-coordinate}
The second coordinate of a point, measuring its signed distance from the x-axis.

\term{y-intercept}
The point where a line or curve crosses the y-axis, occurring where $x = 0$.

\term{Young Diagram}
A finite collection of boxes arranged in left-aligned rows, with the row lengths in non-increasing order. They provide a visual representation of integer partitions and index the irreducible representations of the symmetric group $S_n$.

\term{Young Tableau}
A filling of a Young diagram with positive integers. A tableau is "standard" if the entries are strictly increasing across each row and down each column. The number of standard Young tableaux of a given shape can be computed using the Hook Length Formula.

\term{Young's Inequality}
A result in analysis stating that for $a, b \geq 0$ and $1/p + 1/q = 1$ (with $p, q > 1$), $ab \leq \frac{a^p}{p} + \frac{b^q}{q}$. This is a generalization of the AM-GM inequality and is a key step in proving Hölder's inequality.

\term{Yoneda Lemma}
A fundamental result in category theory stating that for any object $A$ in a category $\mathcal{C}$, the natural transformations from the representable functor $\operatorname{Hom}(A, -)$ to any functor $F: \mathcal{C} \to \text{Set}$ are in bijection with the elements of $F(A)$. It implies that an object is completely determined by its relationships to other objects.

\term{Y-Delta Transform}
A method in network theory for simplifying electrical circuits by replacing a "star" (Y) configuration of three resistors with an equivalent "delta" ($\Delta$) configuration, and vice versa.

\chapter*{Z}
\label{chap:z}

\term{Z-score}
A measure of how many standard deviations a data point is from the mean of a distribution: $z = (x - \mu) / \sigma$. It is used to standardize different datasets for comparison and is the basis for the z-test in statistics.

\term{Z-test}
A statistical hypothesis test used to determine whether the means of two populations are different when the population variance is known and the sample size is large.

\term{Z-transform}
A discrete-time version of the Laplace transform, mapping a sequence $\{x[n]\}$ to a complex function $X(z) = \sum_{n=-\infty}^\infty x[n] z^{-n}$. It is widely used in digital signal processing and control theory to analyze linear time-invariant systems.

\term{Zeros of a Function}
The set of points $z$ in the domain of a function $f$ such that $f(z) = 0$. In complex analysis, zeros are classified as isolated or non-isolated, and their order (multiplicity) is a key feature of the function's behavior.

\term{Zeta Function (Riemann)}
The function $\zeta(s) = \sum_{n=1}^\infty n^{-s}$ for $\operatorname{Re}(s) > 1$, analytically continued to the whole complex plane. The Riemann Hypothesis, stating that all non-trivial zeros lie on the critical line $\operatorname{Re}(s) = 1/2$, is one of the most important open problems in mathematics.

\term{Zigzag}
A polygonal path alternating direction sharply, used in graph theory and combinatorics as a pattern of signs or steps.

\term{Zonal Polynomial}
A symmetric homogeneous polynomial appearing in multivariate statistics and representation theory, associated with spherical functions.

\term{Zonohedron}
A convex polyhedron whose faces are centrally symmetric, a Minkowski sum of line segments.

\term{Zorn's Lemma}
A statement in set theory: every partially ordered set in which every chain has an upper bound contains at least one maximal element. It is equivalent to the Axiom of Choice and is used to prove the existence of bases in every vector space and maximal ideals in rings.

\term{Zygmund Class}
A function belongs to a Zygmund class when its second difference satisfies a bound such as $|f(x+h)+f(x-h)-2f(x)|\le C|h|$ for all relevant $x$ and $h$. This describes a form of smoothness close to, but slightly weaker than, Lipschitz continuity. Applications: Zygmund classes help analyze Fourier series, approximation error, and the behavior of functions near singularities.

\term{Zeno's Paradoxes}
A set of philosophical problems, attributed to Zeno of Elea, illustrating the difficulties of conceptualizing infinity, motion, and continuity. The Dichotomy paradox suggests an object cannot reach a destination because it must first reach the halfway point, then the halfway point of the remainder, and so on. Achilles and the Tortoise argues that the swift runner Achilles never overtakes a tortoise given a head start, since each gain leaves a further gap. The Arrow paradox claims a flying arrow occupies a fixed position at every instant, so motion is impossible. The Stadium paradox concerns the relative motion of rows of bodies in the smallest indivisible instant. The first two are resolved by the convergence of infinite geometric series; the last two sharpen the distinction between continuous motion and its discrete mathematical description.

\term{Zappa--Szép Product}
A generalization of the semi-direct product of groups where two groups $G$ and $H$ act on each other. The resulting group $G \bowtie H$ is the set $G \times H$ with a multiplication rule that incorporates both actions.

\term{Zariski Topology}
A topology used in algebraic geometry where the closed sets are the zero-sets of collections of polynomials. In this topology, open sets are very large (dense), making the space non-Hausdorff.

\term{Zariski Closure}
The smallest algebraic variety (the set of zeros of a set of polynomials) containing a given subset of an affine or projective space. It is the closure of the set in the Zariski topology.

\term{Z-module}
A module over the ring of integers $\mathbb{Z}$. Every abelian group is naturally a $\mathbb{Z}$-module, making the study of $\mathbb{Z}$-modules identical to the study of abelian groups.

\term{Zero Divisor}
An element $a \neq 0$ in a ring $R$ such that there exists another non-zero element $b \in R$ with $ab = 0$. An integral domain is defined as a commutative ring with no zero divisors.

\term{Zero-sum Game}
A game in game theory where the sum of the gains and losses of all players is zero. In such games, one player's gain is exactly balanced by the other players' losses.

\term{Zero-one Law}
A phenomenon in mathematical logic and probability where the probability of a first-order sentence being true in a random structure is either 0 or 1 as the size of the structure tends to infinity.

\term{Zero-set}
The set of all points in the domain of a function where the function's value is zero. In topology, zero-sets are used to define the fact that a space is completely regular.

\term{Zero-point Energy}
The lowest possible energy that a quantum mechanical physical system may have. In the context of the harmonic oscillator, this is $E_0 = \frac{1}{2}\hbar\omega$.

\term{Zygmund's Theorem}
A result in harmonic analysis concerning the convergence of Fourier series of functions in the Zygmund class, showing that such functions possess a higher degree of regularity than general continuous functions.

\end{document}


\end{document}
